\documentclass[mathserif,12pt,aspectratio=169]{beamer}
\usepackage{iftex}
\ifPDFTeX
  \usepackage[T2A]{fontenc}
  \usepackage[utf8x]{inputenc}
  \usepackage[scaled=0.95]{PTSans}
\else
  \usepackage{fontspec}
  \setmainfont[BoldFont=cmunbx.otf,ItalicFont=cmunti.otf,Ligatures=TeX]{cmunrm.otf}
  \setsansfont[BoldFont=cmunsx.otf,ItalicFont=cmunsi.otf,Ligatures=TeX]{cmunss.otf}
  \setmonofont{cmuntt.otf}
  \def\emразделенной{\em разделенной}
\fi
\usepackage[russian]{babel}
\usepackage{graphicx,multicol,algorithm,algorithmic,latexsym,amssymb,amsthm,ragged2e}
\geometry{paperwidth=210mm,paperheight=118.125mm}
\renewcommand{\index}[1]{}
\makeatletter
\patchcmd{\beamer@@@@frame}
  {\setkeys{beamerframe}{#2}}
  {\setkeys{beamerframe}{shrink=0,#2}}{}{}
% GA_2026.pdf
\NewCommandCopy{\GAoriginalgraphic}{\includegraphics}
\newcommand{\GAgraphiccrop}[3]{%
  \expandafter\def\csname GAgraphic:#1\endcsname##1{%
    \GAoriginalgraphic[page=#2,viewport=#3,clip,##1]{GA_2026.pdf}}}
\RenewDocumentCommand{\includegraphics}{s O{} m}{%
  \ifcsname GAgraphic:#3\endcsname
    \csname GAgraphic:#3\endcsname{#2}%
  \else
    \IfBooleanTF{#1}{\GAoriginalgraphic*[#2]{#3}}{\GAoriginalgraphic[#2]{#3}}%
  \fi}
\expandafter\def\csname GAgraphic:odm33-29.eps\endcsname#1{%
  \GAoriginalgraphic[#1]{odm33-29-eps-converted-to.pdf}}
\GAgraphiccrop{graph}{61}{58.969 16.96 394.573 157.149}
\GAgraphiccrop{tree0}{65}{45.719 70.241 407.825 228.376}
\GAgraphiccrop{1_2_3_3.jpg}{178}{274.016 23.731 444.099 82.606}
\GAgraphiccrop{1_2_3_4.jpg}{179}{14.173 5.295 439.375 85.54}
\GAgraphiccrop{1_2_3_5.jpg}{180}{28.346 105.926 425.217 172.215}
\GAgraphiccrop{1_2_5_44.jpg}{191}{353.386 201.317 421.417 237.325}
\GAgraphiccrop{1_3_1_tree.png}{225}{32.039 -0.287 156.554 104.598}
\GAgraphiccrop{1_3_1_graph.png}{227}{6.572 60.404 159.948 159.616}
\GAgraphiccrop{2_2_2.jpg}{541}{63.78 3.285 389.766 186.017}
\GAgraphiccrop{2_2_3_1.jpg}{543}{120.472 -8.215 333.074 143.323}
\GAgraphiccrop{2_2_6_3.png}{574}{248.852 130.209 390.584 199.666}
\GAgraphiccrop{2_2_7.jpg}{576}{42.52 16.042 411.017 215.551}

% GA_preambule.tex
\usepackage{ragged2e}			
\justifying			
\renewcommand{\raggedright}{\leftskip=0pt \rightskip=0pt plus 0cm}			
\hyphenpenalty=5 			
\exhyphenpenalty=10000			
\doublehyphendemerits=2			
\finalhyphendemerits=500			
\uchyph=0			
\renewcommand{\algorithmicrequire}{\textbf{Вход:}}			
\renewcommand{\algorithmicensure}{\textbf{Выход:}}			
\renewcommand{\algorithmiccomment}{{\color{gray} //}}			
%Здесь располагаются все макросы для всякой математики и алгоритмов -			
%\let\em\it			
\def\em#1{\it\color{blue}#1}
\def\NB#1{{\bf\color{brown}#1}}			
%\def\eqed{\tag{\qed}}			
\def\Boxed#1{#1}			
\def\odmkey#1#2{{\it\color{red} #1}\index{#2}}			
\let\emptyset\varnothing			
\def\Set#1#2{\left\{ #1 \mid #2 \right\}}			
\let\Equiv\Longleftrightarrow			
\def\Assign{\operatorname{:=}}			
\def\Pointer{\operatorname{\uparrow\!}}			
\def\Is{\operatorname{\colon}}			
\def\Def{\stackrel{\text{Def}}{=}}			
\def\SubsetNE{\varsubsetneq}%\stackrel{\ne}{\subset}}			
\def\SubsetOne{\stackrel{+1}{\subset}}			
\def\dDef{\operatorname{::=}}			
\def\A#1#2{\forall\,#1\ \left(#2\right)}			
%\def\A#1#2{\forall\,#1\ \linebreak[2]\left(#2\right)}			
\def\dA#1#2{\forall\,#1\cdots #2}			
\def\Alg#1#2{\left\langle #1;#2\right\rangle}		\def\Ms#1{{\color{red}\mathbf{\,#1\,}}}	
\def\Bl#1{\mathbf{#1}}			
\def\hy#1{\left( #1 \right)}			
\def\DIVC#1#2{\left\langle #1 , #2\right\rangle}
\def\divC#1#2{\left\langle #1 , #2\right\rangle}			
\def\DIVI#1#2{#1\in #2}			
\def\E#1#2{\exists\,#1\ \left(#2\right)}			
%\def\E#1#2{\exists\,#1\ \linebreak[2]\left(#2\right)}			
\def\To#1{\stackrel{#1}{\longrightarrow}}			
\def\ToC#1{\stackrel{#1}{,}}			
\def\ToS#1{\stackrel{#1}{\sim}}			
\def\ToE#1{\stackrel{#1}{=}}			
\def\Tuple#1#2#3{{#1}_{#2},\dots,{#1}_{#3}}			
\def\dTuple#1#2#3{{#1}_{#2}\dots{#1}_{#3}}			
\let\Impl\Longrightarrow			
\let\Impll\Longleftarrow			
\def\And{\operatorname{\,\&\,}}			
\def\Or{\,\vee\,}			
\def\Wedge{\,\wedge\,}			
\def\BigOr{\bigvee}			
\def\BigAnd{\bigwedge}			
\def\nat{\operatorname{\rm nat}}			
\def\sgn{\operatorname{\rm sgn}}			
\def\func{\operatorname{\rm func}}			
\def\ker{\operatorname{\rm ker}}			
\def\Im{\operatorname{\rm Im}}			
\def\Dom{\operatorname{\rm Dom}}			
\def\Ob{\operatorname{\rm Ob}}			
\def\Mor{\operatorname{\rm Mor}}			
\def\SetCat{\operatorname{\rm Set}}			
\def\OrSet{\operatorname{\rm OrSet}}			
\def\Aut{\operatorname{\rm Aut}}			
\def\End{\operatorname{\rm End}}			
\def\Gr{\operatorname{\rm Gr}}			
\def\Suc{\operatorname{\rm Suc}}			
\def\card{\operatorname{\rm card}}			
\def\pr{\operatorname{\rm pr}}			
\def\MOD{\operatorname{\rm\bf mod}}			
\def\DIV{\operatorname{\rm\bf div}}			
\def\Lrf#1{\lfloor #1 \rfloor}			
\def\lrle#1{\langle #1 \rangle}			
\def\llrle#1{\langle {\text{#1}} \rangle}			
\def\Map#1#2#3{{#1}\colon{#2}\nolinebreak\to\nolinebreak{#3}}			
\def\MapI#1#2#3{{#1}\colon{#2}\in{#3}}			
\def\Rule#1#2#3{\frac{#1}{#2}\,#3}			
\def\Seq#1#2#3{#1 \vdash_{#3} #2}			
\def\Prog#1#2#3{\{#1\}\,#2\,\{#3\}}			
\def\Not{\neg}
\def\Mu{\mathrm{M}}			
\def\demo{\noindent $\triangleright$ \vskip 3 em \par\hfill$\triangleleft$}			
\DeclareMathOperator{\Eq}{Eq}			
\DeclareMathOperator{\Unify}{Unify}			
\DeclareMathOperator{\Antilex}{Antilex}			
\DeclareMathOperator{\Reverse}{Reverse}			
\DeclareMathOperator{\Fano}{Fano}			
\DeclareMathOperator{\Med}{Med}			
\DeclareMathOperator{\Huffman}{Huffman}			
\DeclareMathOperator{\Up}{Up}			
\DeclareMathOperator{\Down}{Down}			
\DeclareMathOperator{\KCC}{KCC}			
\DeclareMathOperator{\NewNode}{NewNode}			
\DeclareMathOperator{\Find}{Find}			
\DeclareMathOperator{\New}{new}			
\DeclareMathOperator{\Dispose}{dispose}			
\DeclareMathOperator{\Delete}{Delete}			
\DeclareMathOperator{\BT}{BT}			
\DeclareMathOperator{\Last}{last}			
\DeclareMathOperator{\Selectmax}{Selectmax}			
\DeclareMathOperator{\Sort}{Sort}			
\DeclareMathOperator{\FD}{FD}			
\DeclareMathOperator{\Length}{Length}			
\DeclareMathOperator{\Append}{Append}			
\DeclareMathOperator{\Eval}{Eval}			
\DeclareMathOperator{\TraverseTree}{TraverseTree}			
%Списки, etc			
\def\labelenumi{\theenumi}			
\def\theenumi{\arabic{enumi}.}			
\def\labelenumii{\theenumii}			
\def\theenumii{\arabic{enumii})}			
\def\p@enumii{\theenumi}			
\def\labelenumiii{\theenumiii.}			
\def\theenumiii{\roman{enumiii}}			
\def\p@enumiii{\theenumi(\theenumii)}			
\def\labelenumiv{\theenumiv.}			
\def\theenumiv{\Alph{enumiv}}			
\def\p@enumiv{\p@enumiii\theenumiii}			
%нумерация			
%\numberwithin{equation}{chapter}			
%\numberwithin{algorithm}{chapter}			
%\numberwithin{figure}{chapter}			
%\renewcommand\thesection       {\thechapter.\arabic{section}}			
%\floatname{algorithm}{Алгоритм}			
%\renewcommand{\algorithmicrequire}{\textbf{Вход:}}			
%\renewcommand{\algorithmicensure}{\textbf{Выход:}}			
\let\cal\mathcal			
\let\Bbb\mathbb			
\let\leq\leqslant			
\let\le\leqslant			
\let\geq\geqslant			
\let\ge\geqslant			
\def\rmb#1{#1 \linebreak #1}			
\def\rmbb#1{#1 \\ #1}			
%% М. Лебедев, конец августа 2010			
\def\sign{\operatorname{\rm sign}}			
%% from class odm2007			
%\newtheorem*{theorem}{ТЕОРЕМА}			
%\newtheorem{vartheorem}{ТЕОРЕМА}			
%\newtheorem*{corollary}{\begin{corollary}}			
%\newtheorem{varcorollary}{\begin{corollary}}			
%\newtheorem*{lemma}{ЛЕММА}			
%\newtheorem{varlemma}{ЛЕММА}			
%\newtheorem*{ground}{Обоснование}			
%\def\endground{\qed\endtrivlist\@endpefalse }			
%\newtheorem*{example}{Пример}			
%\def\endexample{\qed\endtrivlist\@endpefalse }			
\def\clip#1{\bgroup\par\addvspace{7pt plus 3pt minus 3pt}%			
	\leavevmode\lower 3pt\hbox{\sffamily\bfseries #1} \hrulefill\par\nopagebreak\small} %\hrulefill\par\nopagebreak}			
\def\endclip{\par\nopagebreak\addvspace{3pt}\hrule width130mm\egroup%			
	\addvspace{12pt plus 3pt minus 3pt}}			
\def\remark{\par{\bf \color{brown} Замечание.}}			
\def\endremark{\par}			
\def\digress{\begin{small}
\par{\bf \color{brown} Отступление.}}			
	\def\enddigress{\end{small}\par}			
%\def\example{\par\addvspace{9pt plus 1pt minus 3pt}{\sffamily\bfseries Пример}\par\nopagebreak}			
%\def\endexample{\par\addvspace{9pt plus 1pt minus 3pt}}			
\def\example{\par{\bf \color{brown} Пример.}}			 
\def\endexample{\par}			
\def\examples{\par{\bf \color{brown} Примеры.}}			
\def\endexamples{\par}			
\def\prototheorem#1{\par{\bf #1}\bgroup\it}			
\def\endprototheorem{\par\egroup}			
\def\theorem{\prototheorem{\color{brown} Теорема.}}			
\def\endtheorem{\endprototheorem}			
\def\corollary{\prototheorem{\color{brown} Следствие.}}			
\def\endcorollary{\endprototheorem}			
\def\lemma{\prototheorem{\color{brown} Лемма.}}			
\def\endlemma{\endprototheorem}			
%\def\proof{\par\addvspace{12pt plus 3pt minus 3pt}{\sffamily\bfseries\small %Д{\scriptsize\uppercase{оказательство}}}\par\nopagebreak}			
%\def\proof{\par\addvspace{12pt plus 3pt minus 3pt}{\sffamily\bfseries\small %Д{\scriptsize\uppercase{оказательство}}}\quad\nopagebreak}			
\def\proof{\par{\bf\color{brown} Доказательство.}}			
\def\endproof{\hfill\qed\par}			
%\def\proofsection#1{\par{\small $\mathsf{\left[\:\text{{\sf#1}}\:\right]}$}\;}			
\def\proofsection#1{\par{\color{brown}\sffamily\small  [\;}{\sffamily\small\color{brown} #1}{\sffamily\small\color{brown}\;]}\;}			
\def\eqed{%			
	\gdef\tagform@##1{\maketag@@@{\ignorespaces##1\unskip}}%			
	\tag{\color{brown}\qed}%			
}			
\def\ground{\par{\bf\color{brown} Обоснование.}}			
\def\endground{\hfill\color{brown}\qed\par}			
% algorithms
%\newcounter{algorithm}[chapter]
\def\ext@algorithm{loa}
\def\fnum@algorithm{\algorithmname\ \thealgorithm}
\def\algorithmname{Алгоритм}
\def\thealgorithm{\thechapter.\arabic{algorithm}.}


\def\algorithm{%
	\def\@captionfont{\normalfont\sffamily\small}%
	\captionindent=-10\fill%
	\def\caption{\refstepcounter{algorithm}\@dblarg{\@caption{algorithm}}}%
	\def\@makecaption##1##2{\addvspace{12pt plus 3pt minus 2pt}\normalfont\sffamily\small{\bfseries ##1.} ##2}
	\small
}
\def\endalgorithm{\relax}


% citations
\def\@cite#1#2{{%
		\m@th\upshape\mdseries[{#1}{\if@tempswa, #2\fi}]}}
		
	% Стиль презентации			
	\usetheme{default}			
	\setbeamertemplate{navigation symbols}{}			
	\setbeamertemplate{footline}[page number]			
	\setbeamersize{text margin left = 6pt}			
	\setbeamersize{text margin right = 6pt}			
			
%	\usepackage{epstopdf}			
	%\renewcommand{\algorithmiccomment}{ //}			
	\def\Reduce{\stackrel{*}{\leadsto}}			
	\DeclareMathOperator{\Clear}{Clear}			
	\DeclareMathOperator{\Conf}{Conf}			
	\DeclareMathOperator{\lcm}{lcm}			
	\def\Mod#1#2#3{#1\equiv #2 \,\left(\bmod\,#3\right)}			
	%\newtheorem{theorem}{ТЕОРЕМА}			
%	\def\digress{\par{\bf Отступление.}}			
%	\def\enddigress{\par}			
			
	\renewcommand{\algorithmiccomment}{ //}			
	\DeclareMathOperator{\GI}{Grey2Int}			
			
	\setbeamertemplate{frametitle continuation}{}			
	\def\TT{\text{T}}			
	\def\FF{\text{F}}			
	\DeclareMathOperator{\firstNeqSubTerms}{firstNeqSubTerms}			
				
	\DeclareMathOperator{\Partition}{Partition}			
	\DeclareMathOperator{\Quicksort}{Quicksort}			
			
	\DeclareMathOperator{\ChangeArray}{ChangeArray}			
	%\newtheorem{corollary}{\begin{corollary}}			
			
	\def\divI#1#2{#1 \in #2}			
	\DeclareMathOperator{\Ordering}{Ordering}			
	\setbeamertemplate{frametitle}			
	{\color{blue}\large\insertframetitle\par\vskip-6pt}			
	%%вставка для оглавления со страницами			
	%\makeatletter			
	%\long\def\beamer@section[#1]#2{%			
	%	\beamer@savemode%		
	%	\mode<all>%		
	%	\ifbeamer@inlecture		
	%	\refstepcounter{section}%		
	%	\beamer@ifempty{#2}%		
	%	{\long\def\secname{#1}\long\def\lastsection{#1}}%		
	%	{\global\advance\beamer@tocsectionnumber by 1\relax%		
	%		\long\def\secname{#2}%	
	%		\long\def\lastsection{#1}%	
	%		\addtocontents{toc}{\protect\beamer@sectionintoc{\the\c@section}{#2\hfill\the\c@page}{\the\c@page}{\the\c@part}%	
	%			{\the\beamer@tocsectionnumber}}}%
	%	{\let\\=\relax\xdef\sectionlink{{Navigation\the\c@page}{\noexpand\secname}}}%		
	%	\beamer@tempcount=\c@page\advance\beamer@tempcount by -1%		
	%	\beamer@ifempty{#1}{}{%		
	%		\addtocontents{nav}{\protect\headcommand{\protect\sectionentry{\the\c@section}{#1}{\the\c@page}{\secname}{\the\c@part}}}%	
	%		\addtocontents{nav}{\protect\headcommand{\protect\beamer@sectionpages{\the\beamer@sectionstartpage}{\the\beamer@tempcount}}}%	
	%		\addtocontents{nav}{\protect\headcommand{\protect\beamer@subsectionpages{\the\beamer@subsectionstartpage}{\the\beamer@tempcount}}}%	
	%	}%		
	%	\beamer@sectionstartpage=\c@page%		
	%	\beamer@subsectionstartpage=\c@page%		
	%	\def\insertsection{\expandafter\hyperlink\sectionlink}%		
	%	\def\insertsubsection{}%		
	%	\def\insertsubsubsection{}%		
	%	\def\insertsectionhead{\hyperlink{Navigation\the\c@page}{#1}}%		
	%	\def\insertsubsectionhead{}%		
	%	\def\insertsubsubsectionhead{}%		
	%	\def\lastsubsection{}%		
	%	\Hy@writebookmark{\the\c@section}{\secname}{Outline\the\c@part.\the\c@section}{2}{toc}%		
	%	\hyper@anchorstart{Outline\the\c@part.\the\c@section}\hyper@anchorend%		
	%	\beamer@ifempty{#2}{\beamer@atbeginsections}{\beamer@atbeginsection}%		
	%	\fi%		
	%	\beamer@resumemode}%		
	%			
	%\def\beamer@subsection[#1]#2{%			
	%	\ifbeamer@inlecture%		
	%	\refstepcounter{subsection}%		
	%	\beamer@ifempty{#2}{\long\def\subsecname{#1}\long\def\lastsubsection{#1}}		
	%	{%		
	%		\long\def\subsecname{#2}%	
	%		\long\def\lastsubsection{#1}%	
	%		\addtocontents{toc}{\protect\beamer@subsectionintoc{\the\c@section}{\the\c@subsection}{#2\hfill\the\c@page}{\the\c@page}{\the\c@part}{\the\beamer@tocsectionnumber}}%	
	%	\addtocontents{nav}{%		
	%		\protect\headcommand{\protect\beamer@subsectionentry{\the\c@part}{\the\c@section}{\the\c@subsection}{\the\c@page}{\lastsubsection}}%	
	%		\protect\headcommand{\protect\beamer@subsectionpages{\the\beamer@subsectionstartpage}{\the\beamer@tempcount}}%	
	%	\edef\subsectionlink{{Navigation\the\c@page}{\noexpand\subsecname}}%		
	%	\def\insertsubsection{\expandafter\hyperlink\subsectionlink}%		
	%	\def\insertsubsectionhead{\hyperlink{Navigation\the\c@page}{#1}}%		
	%	\Hy@writebookmark{\the\c@subsection}{#2}{Outline\the\c@part.\the\c@section.\the\c@subsection.\the\c@page}{3}{toc}%		
	%	\hyper@anchorstart{Outline\the\c@part.\the\c@section.\the\c@subsection.\the\c@page}\hyper@anchorend%		
	%	\beamer@ifempty{#2}{\beamer@atbeginsubsections}{\beamer@atbeginsubsection}%		
	%	\beamer@resumemode}		
	%\makeatother			
	%%конец вставки для оглавления			
			
	\def\Path#1#2{\langle {\stackrel{\rightarrow}{#1 , #2}}\rangle}			
	\DeclareMathOperator{\Init}{Init}			
	\DeclareMathOperator{\Relax}{Relax}			
	\DeclareMathOperator{\ExtractMin}{ExtractMin}			
	\DeclareMathOperator{\Aug}{Aug}			
%	{\color{blue}\large\insertframetitle\par\vskip-8pt}			
			
	\renewcommand{\algorithmicelsif}			
	{\algorithmicelse\algorithmicif}			
	\DeclareMathOperator{\Build}{Build}			
	\DeclareMathOperator{\Modify}{Modify}			
	\DeclareMathOperator{\Heapify}{Heapify}			
	\DeclareMathOperator{\MakeHeap}{MakeHeap}			
	\DeclareMathOperator{\IncElt}{IncElt}			
		
		
\def\DMref#1#2{{\color{blue}\underline{\hyperlink{#2}{#1}}}}	

\def\TODO#1{{\color{gray} TODO {#1}}}	

\def\Ok{{\color{green}\textbf{!}\ }}
\def\Qm{{\color{brown}\textbf{?}\ }}
\def\Er{{\color{red}\textbf{--}\ }}
\makeatother
\emergencystretch=3em
\begin{document}
	\title{Грамматики и автоматы}  
	\author{Ф. А. Новиков\\fedornovikov51@gmail.com}
	
	
	
	\date{Санкт-Петербург\\14 ноября 2026 года} 
	% Создание заглавной страницы
	\frame{\titlepage} 
	% Автоматическая генерация содержания
\begin{multicols}{2}[\frametitle{\vspace*{-25pt}\begin{center}\color{blue}\Large Содержание\end{center}\vspace*{-25pt}}\usebeamertemplate{frametitle}]
\tableofcontents
\end{multicols}
% GA_0_2025.tex, 1–1
\section{0. Введение в предмет}

% GA_0_2025.tex, 386–386
\subsection{0.2. Архитектура курса}\label{0.2.}

% GA_0_2025.tex, 473–1065
\subsubsection{0.2.3. Используемые обозначения}\label{0.2.3.}
\begin{frame}
\frametitle{\Qm 0.2.3. Используемые обозначения}
Изложение теории и практики реализации 
языков предметных областей в этом курсе
ведётся на естественном языке 
(предполагается, что читатель понимает русский язык)
с использованием трёх формальных языков:
\begin{itemize}
\item простого текстового языка записи алгоритмов (псевдокод, \DMref{п.~0.2.4}{0.2.4.});
\item графического унифицированного языка моделирования
(подмножество UML, \DMref{п.~0.2.5}{0.2.5.});
\item языка математических формул (раздел \DMref{0.3}{0.3.});
\end{itemize}  


\TODO{ Выделения в тексте: коричневые слова (теоремы, ...), 
эмфатические ударения (синий), определяю.щие вхождения (красный)}


\end{frame}

\subsubsection{0.2.4. Псевдокод}\label{0.2.4.}
\begin{frame}
\frametitle{\Ok 0.2.4. Псевдокод (1/5)}

Изложение материала в этом учебном курсе конструктивное, 
с указанием конкретных {\em алгоритмов} (п.~\DMref{0.3.3}{0.3.3})
  построения объектов и проведения преобразований.
Алгоритмы записываются на \odmkey{псевдокоде}{Псевдокод}. 
Текст на псевдокоде считается записью алгоритма, 
а доказательное рассуждение о тексте алгоритма считается
обоснованием свойств алгоритма. Псевдокод не исполняется машиной: он
предназначен для чтения человеком.

\smallskip
Предполагается, что студенты не испытывают затруднений в понимании
текстов программ на языке высокого уровня. Именно поэтому в
качестве нотации псевдокода используется 
{\em некоторый неспецифицированный\/} алгоритмический язык,
напоминающий распространённые языки программирования высокого уровня,
но не подменяющий их.\\
Основным свойством псевдокода является достаточная выразительность
при не слишком обременительном синтаксисе. 
Включать в запись алгоритма разрешается что угодно,
лишь бы это было понятно по здравому программистскому смыслу.  

\smallskip
Ключевые слова псевдокода выделены полужирным шрифтом, идентификаторы
объeктов записываются курсивом и часто состоят из одной буквы, 
как это принято для математических
объектов, имена алгоритмов записываются прямым шрифтом, 
как названия математических функций,
примечания отделяются двойной косой чертой //.

\end{frame}

\begin{frame}
\frametitle{\Ok 0.2.4. Псевдокод (2/5)}

\vspace{-1pt}
\begin{remark}
Для операции присваивания и других определений
{\em никогда} не используется знак равенства $=$, 
что является заметным отличием нотации 
псевдокода от языков программирования.
Этот отход от программистских традиций, по мнению автора, {\em необходим}
для записи алгоритмов в математическом контексте.  
\end{remark}

Псевдокод используется для записи различных алгоритмических конструкций, в том числе, по меньшей мере, в следующих трёх случаях.
\begin{itemize}
\item {\em Краткое <<однострочное>>  выражение} для 
определения математического объекта.
\begin{example} Абсолютная величина числа: 
$|x|\Def \textbf{ if } x>0 \textbf{ then } x \textbf{ else } -x \textbf{ end if }$.
\end{example}
\item {\em <<Безымянный>> фрагмент } доказательства или построения в заданном контексте. 
\begin{example}
Устранение $S$-правил в грамматике 
$G=\langle\mathrm{M},\Sigma,\Pi, S\rangle$ и построение эквивалентной грамматики
$G^\prime=\langle\mathrm{M}^\prime,\Sigma,\Pi^\prime, S\rangle$ для
грамматик {\em любого} класса.

\vspace{-1pt}
\begin{algorithmic}
\STATE {$\mathrm{M}^\prime\Assign\mathrm{M}+S^\prime$}
\COMMENT{ новый нетерминал}
\FOR{$(\alpha\to\beta)\in\Pi$}
\STATE {$ (\alpha\to\beta)\Assign(\alpha[S^\prime/S]\to\beta[S^\prime/S])$}
\COMMENT{ замена всех вхождений $S$ вхождениями $S^\prime$}
\ENDFOR
\STATE {$\Pi^\prime\Assign\Pi+(S\to S^\prime)$}
\COMMENT{ новое правило для аксиомы}
\end{algorithmic}
\end{example}
\end{itemize}
\end{frame}

\begin{frame}
\frametitle{\Ok 0.2.4. Псевдокод (3/5) \Qm}

\vspace{-2pt}
\begin{itemize}
\item
{\em Именованный алгоритм}, состоящий из {\em заголовка} и {\em тела},
разделённых знаком $\Def$. 
В заголовке указываются: возможный {\em класс} алгоритма, имя, параметры
с типами и возможный тип результата. 
\begin{example}
Проверка правильности скобочной структуры цепочки $\omega$.
\begin{algorithmic}
\STATE {$ \textbf{func } \mathrm{D}\ ( $}
\COMMENT{ класс и имя алгоритма}
\STATE {$\phantom{\textbf{func } \mathrm{D}\ (}
\omega : \textbf{ array } [1..n] \textbf{ of char}$}
\COMMENT{ параметр --- исходная цепочка}
\STATE {$\phantom{\textbf{func } \mathrm{D}\ (}
)\ :  \textbf{ bool} \Def$}
\COMMENT{ тип возвращаемого значения}
\STATE {$d \Assign 0$}
\COMMENT{ локальная переменная --- счётчик скобок}
\FOR{$i \textbf{ from } 1 \textbf{ to } n$}
\STATE {$\textbf{if } \omega[i] = "(" \textbf{ then } d \Assign d+1
\textbf{ end if }$, $\textbf{ if } \omega[i] = ")" \textbf{ then } d \Assign d-1
\textbf{ end if}$}
\STATE {$\textbf{if } d < 0 \textbf{ then return false end if}$}
\ENDFOR
\STATE {$\textbf{return } d = 0$}
\end{algorithmic}
\end{example}
\end{itemize} 

\vspace{-2pt}

\begin{digress}
Псевдокод, хотя и является искусственным языком, 
по своей природе ближе к естественным языкам, 
чем к  языкам программирования. 
Пишущий на псевдокоде вправе использовать {\em любые} слова
и обозначения, лишь бы они оказались понятны читающему.
\end{digress}
\end{frame}

\begin{frame}
\frametitle{\Ok 0.2.4. Псевдокод (4/5) \Qm}
Следующие обозначения и конструкции предполагаются известными и понятными.

\smallskip
Общеупотребительные структуры
управления:

\begin{itemize}
\item
 {\em ветвления} в краткой ({\bf if} ... {\bf then} ... {\bf
end~if}), полной ({\bf if} ... {\bf then} ... {\bf else} ... {\bf
end~if}) и
 расширенной ({\bf if} ... {\bf then} ... 
 {\bf elseif} ... {\bf then} ...
  {\bf else} ... {\bf
end~if}) формах; 
\item {\em циклы} со счётчиком ({\bf for} ... {\bf from} ...
{\bf to} ... {\bf do}
 ... {\bf end~for}), с постусловием
({\bf~repeat}~...~{\bf~until}), с предусловием \linebreak[3]({\bf
while}~...~{\bf~do}~...~{\bfseries end~while}) и по множеству
({\bf for} ... {\bf do} ... {\bf end~for}).
\end{itemize}

\begin{remark}
Допускается использование  {\em охраняемых команд Дейкстры} 
({\bf if} ... {\bf fi} и {\bf do} ... {\bf od}) вместо традиционной структур управления.
\end{remark}

Из структур данных доступны:
\begin{itemize}
\item массивы {\bf array} [...]
{\bf of} ...,  в том числе многомерные и с произвольными 
{\em упорядоченными} индексами; 
\item 
структуры {\bf struct \{ } ... {\bf \}}, в том числе с рекурсивными определениями типов; 
\item
указатели $\uparrow$, причём
указатель может иметь специальное значение
{\bf nil}, не указывающее ни на один объект.
\end{itemize}

\begin{remark}
Номенклатуpа   {\em примитивных типов} включает, но не ограничивается типами 
{\bf bool}, {\bf char}, {\bf int}, {\bf string}.
\end{remark}

\end{frame}

\begin{frame}
\frametitle{\Ok 0.2.4. Псевдокод (5/5) \Qm}
\begin{remark}
Подразумевается строгая {\em статическая  типизация}, даже приведения не
используются, за исключением разыменования, которое
подразумевается везде, где оно нужно по здравому программистскому
смыслу.
В то же время объявления локальных переменных почти всегда
опускаются, если их тип ясен из контекста.
 \end{remark}
 
\smallskip
Особого упоминания заслуживают семь естественных операторов, аналоги
которых в явном виде редко встречаются в <<настоящих>> языках
программирования.

\begin{itemize}
\item Оператор \textbf{skip} ничего не делает. 
\item Оператор \textbf{exit} прекращает {\em выполнение} того алгоритма,
в тексте которого находится. 
\item Оператор \textbf{yield} $x$ {\em возвращает} значение $x$, но при этом 
{\em выполнение} алгоритма {\em не прекращается\/}, а продолжается со следующего оператора. 
\item Оператор \textbf{return} $x$~--- это сокращение 
для пары операторов \textbf{yield} $x$ ; \textbf{exit}. 
\item Оператор \textbf{next for} $x$ означает 
продолжение выполнения {\em со следующего шага цикла} по переменной $x$.
Этот оператор должен находиться в теле цикла по $x$ (на любом
уровне вложенности). 
\item Оператор \textbf{select} $m\!\in\!M$ 
означает выбор {\em произвольного} элемента $m$ из множества $M$. 
\end{itemize}

\end{frame}

% 0.2.4. Псевдокод. А. Архиповский, 25.09.2026.
% Вторая редакция. Основание: замечания ФАН от 18.09.2026 к TODO_слайд_26
% («программа может не выполниться двумя способами: зависает и падает»;
% «не все программы функции, некоторые итераторы»; «мало print, хочется yield»)
% и к TODO_слайды_47_48 («нужно вводить программные конструкции
% для нормальных алгорифмов Маркова; они уже есть частично: subst(alpha,beta,gamma)»).
% Добавлены: исход fail, заголовок iter и оператор yield, операции над словами
% subst и lsubst. Синтаксис в кадре 5/5 приведён в соответствие.

\begin{frame}
\frametitle{\Er 0.2.4. Псевдокод (1/6)}

Алгоритмы в курсе записываются на \odmkey{псевдокоде}{Псевдокод}. Тезис
Чёрча\,--\,Тьюринга принимается без оговорок (п.~\DMref{0.3.3}{0.3.3}), поэтому
запись на псевдокоде считается записью алгоритма, а рассуждение о её тексте~---
обоснованием свойств алгоритма. Псевдокод не исполняется машиной: он
предназначен для человека.

Алгоритм состоит из {\em заголовка} и {\em тела}. В заголовке~--- имя, параметры
с типами и тип результата либо, если алгоритм не функция, исходные данные и
результат. Тело записывается операторами со сдвигом вправо; комментарий
начинается знаком \textbf{//}.

\begin{example}
Проверка правильности скобочной структуры цепочки $\omega$.
\end{example}
\vspace{-8pt}
\begin{algorithmic}
\STATE {$\textbf{func } \mathrm{balance}\,(\omega : \textbf{array } [1..n]
\textbf{ of char}) : \textbf{bool}$}
\STATE {$d \Assign 0$}
\COMMENT{ счётчик скобок}
\FOR{$i \textbf{ from } 1 \textbf{ to } n$}
\STATE {$\textbf{if } \omega[i] = "(" \textbf{ then } d \Assign d+1
\textbf{ end if}$}
\STATE {$\textbf{if } \omega[i] = ")" \textbf{ then } d \Assign d-1
\textbf{ end if}$}
\STATE {$\textbf{if } d < 0 \textbf{ then return false end if}$}
\ENDFOR
\STATE {$\textbf{return } d = 0$}
\end{algorithmic}
\end{frame}

\begin{frame}
\frametitle{\Qm 0.2.4. Псевдокод (2/6)}
%%% Три исхода и три вида алгоритмов

Запуск алгоритма имеет ровно три исхода: алгоритм возвращает значение
(\textbf{return}); алгоритм {\em падает} (\textbf{fail}); алгоритм {\em
зависает}, то есть не заканчивает работу. Различаются и три вида заголовков.

\vspace{-2pt}
\begin{enumerate}
\item \textbf{func}~--- {\em функция}: возвращает значение объявленного типа.
\item \textbf{proc}~--- {\em процедура}: значения не возвращает, изменяет
данные.
\item \textbf{iter}~--- {\em итератор}: выдаёт оператором \textbf{yield}
последовательность значений, вообще говоря бесконечную.
\end{enumerate}
\vspace{-2pt}

Оператор \textbf{fail} прекращает работу без результата; так же ведёт себя
всякая неопределённая операция (вырезка за границей слова, деление на нуль).
Отказ~--- не авария, а штатный исход: на нём основан перебор с возвратами.

\begin{example}
Итератор натурального ряда~--- правильный алгоритм, который не
останавливается никогда:
\end{example}
\vspace{-8pt}
\begin{algorithmic}
\STATE {$\textbf{iter } \mathrm{nat}\,(\,)$}
\STATE {$n \Assign 0$}
\STATE {$\textbf{while true do yield } n \Assign n+1 \textbf{ end while}$}
\end{algorithmic}
\vspace{-4pt}

\begin{remark}
Цикл \textbf{for} $x\in X$ \textbf{do} \dots\ \textbf{end for} пробегает
значения, выданные итератором $X$.
\end{remark}
\end{frame}

\begin{frame}
\frametitle{\Er 0.2.4. Псевдокод (3/6)}
%%% Сводка конструкций

\begin{center}\scriptsize
\begin{tabular}{|l|l|}
\hline
Конструкция & Запись \\
\hline
Заголовок функции & \texttt{func} имя(параметр : тип) : тип \\
Заголовок процедуры & \texttt{proc} имя(параметры) \\
Заголовок итератора & \texttt{iter} имя(параметры) \\
Исходные данные и результат & Вход: \dots\ Выход: \dots \\
Объявление переменной & имя : тип \\
Присваивание & $x \Assign$ выражение \\
Ветвление & \texttt{if} усл. \texttt{then} \dots\ \texttt{else} \dots\
\texttt{end if} \\
Цикл со счётчиком & \texttt{for} $i$ \texttt{from} $a$ \texttt{to} $b$
\texttt{do} \dots\ \texttt{end for} \\
Цикл по множеству или итератору & \texttt{for} $x\in X$ \texttt{do} \dots\
\texttt{end for} \\
Цикл с условием & \texttt{while} усл. \texttt{do} \dots\ \texttt{end while} \\
Возврат значения (конец работы) & \texttt{return} выражение \\
Отказ (конец работы) & \texttt{fail} \\
Выдача очередного значения & \texttt{yield} выражение \\
Печать & \texttt{print} список \\
Работа с очередью & $x\to Q$ (положить), $x\leftarrow Q$ (взять) \\
Создание объекта & \texttt{new} $A$ \\
Пустой оператор & \texttt{skip} \\
Комментарий & \texttt{//} текст до конца строки \\
\hline
Типы & \textbf{bool}, \textbf{char}, $1..n$, $\mathbb{N}_0$, \textbf{array}
$[\dots]$ \textbf{of} $T$, \textbf{queue of} $T$, \textbf{struct},
$\Pointer T$, $A^*$ \\
Операции над словами & $|\alpha|$~--- длина, $\alpha[i]$~--- буква,
$\alpha[i..j]$~--- вырезка, \\
 & $\alpha\beta$~--- сцепление, $\varepsilon$~--- пустое слово,
$\alpha\subset\beta$~--- вхождение \\
Подстановка & $\mathrm{lsubst}\,(\alpha,\beta,\gamma)$,\ \
$\mathrm{subst}\,(\alpha,\beta,\gamma)$ \\
\hline
\end{tabular}
\end{center}
\end{frame}

\begin{frame}
\frametitle{\Qm 0.2.4. Псевдокод (4/6) \NB{Перенести в 1.1.1?}}
%%% Подстановка

Обе подстановки заменяют одно вхождение цепочки $\beta$ в цепочку $\alpha$ на
цепочку $\gamma$; различаются они тем, какое вхождение берётся.

\vspace{-2pt}
\begin{itemize}
\item $\mathrm{lsubst}\,(\alpha,\beta,\gamma)$~--- функция: берётся самое левое
вхождение $\beta$ в $\alpha$; если вхождений нет, функция падает.
\item $\mathrm{subst}\,(\alpha,\beta,\gamma)$~--- итератор: выдаёт результат
замены каждого вхождения $\beta$ в $\alpha$, слева направо, по одному на
вхождение; если вхождений нет, не выдаёт ничего.
\end{itemize}
\vspace{-2pt}

\begin{examples}
$\mathrm{lsubst}\,(abab,\,ab,\,c)=cab$, а итератор $\mathrm{subst}\,(abab,\,ab,\,c)$
выдаёт $cab$ и $abc$. Вызов $\mathrm{lsubst}\,(abab,\,ba,\,\varepsilon)$ даёт
$ab$, а $\mathrm{lsubst}\,(abab,\,c,\,\varepsilon)$ падает.
\end{examples}

Этих двух операций достаточно, чтобы записывать на псевдокоде переписывающие
системы: $\mathrm{lsubst}$ отвечает {\em детерминированному} переписыванию
(нормальный алгорифм Маркова, п.~\DMref{0.3.3}{0.3.3}), а $\mathrm{subst}$~---
{\em недетерминированному} (порождение цепочек грамматикой,
п.~\DMref{1.1.3}{1.1.3}).

\begin{remark}
Обе операции выражаются через вырезку и сцепление: замена вхождения, начинающегося
в позиции $i$, есть $\alpha[1..i-1]\,\gamma\,\alpha[i+|\beta|..|\alpha|]$.
Поэтому в псевдокод они добавлены ради краткости записи, а не ради новых
возможностей.
\end{remark}
\end{frame}

\begin{frame}
\frametitle{\Er 0.2.4. Псевдокод (5/6) \NB{перенести в 3}}
%%% Псевдокод как сквозной пример

Псевдокод~--- формальный язык, нужный курсу с первого алгоритма и известный
читателю заранее. Поэтому он и берётся сквозным примером: каждое новое средство
курса сначала применяется к нему.

\vspace{-2pt}
\begin{center}\scriptsize
\begin{tabular}{|p{5.5cm}|p{5.4cm}|l|}
\hline
Что описывается & Каким средством & Пункт \\
\hline
Лексика: идентификаторы, числа, ключевые слова, комментарии & регулярные
выражения & 1.2.1--1.2.2 \\
Распознавание лексем & языковой граф, детерминизация, конечный автомат &
1.2.3--1.2.4, 2.2 \\
Синтаксис & КС-грамматика, расширенная запись, БНФ & 1.1.8, 1.3.1 \\
Однозначность описания & проблема однозначности КСГ & 1.3.6 \\
Пригодность к разбору & грамматики LL($k$) и LR($k$) & 1.3.7 \\
Разбор & нисходящий, восходящий, станция Дейкстры & 1.4.2--1.4.4 \\
Контекстные условия: описания, типы, области видимости & атрибутные грамматики &
1.5.2, 1.5.4 \\
Семантика & операционная, исчисление Хоара, денотационная & 1.5.5, 1.5.7--1.5.9 \\
Псевдокод как модель вычислимости & тезис Чёрча\,--\,Тьюринга & 0.3.3 \\
\hline
\end{tabular}
\end{center}
\end{frame}

\begin{frame}
\frametitle{\Er 0.2.4. Псевдокод (6/6) \NB{перенести в 3}}
%%% Сквозной пример: первые два шага

Средства, которыми записаны эти два шага, вводятся позже; здесь достаточно
понимать, что $\mid$~--- выбор, $*$~--- повторение, $\{\alpha \# \beta\}$~---
повторение $\alpha$ через $\beta$, $[\alpha]$~--- необязательная часть, а
терминалы взяты в кавычки.

\smallskip
{\em Лексика} (в п.~\DMref{1.2}{1.2.} то же будет записано регулярными
выражениями):
\vspace{-4pt}
\begin{center}\scriptsize
\begin{tabular}{@{}ll@{}}
буква & : "a" $\mid$ \dots\ $\mid$ "z" $\mid$ "A" $\mid$ \dots\ $\mid$ "Z" ; \\
цифра & : "0" $\mid$ \dots\ $\mid$ "9" ; \\
идентификатор & : буква (буква $\mid$ цифра)$^*$ ; \\
число & : цифра цифра$^*$ ; \\
комментарий & : "//" (любой символ, кроме конца строки)$^*$ . \\
\end{tabular}
\end{center}
\vspace{-4pt}

{\em Синтаксис} (в расширенной записи грамматик из п.~\DMref{1.1.8}{1.1.8.}):
\vspace{-4pt}
\begin{center}\scriptsize
\begin{tabular}{@{}ll@{}}
операторы & : \{ оператор \# ";" \} ; \\
оператор & : присваивание $\mid$ ветвление $\mid$ цикл
 $\mid$ "return"\,[\,выражение\,] \\
 & $\mid$ "yield"\,выражение $\mid$ "print"\,список $\mid$ "fail" $\mid$ "skip" ; \\
присваивание & : переменная\,": ="\,выражение ; \\
ветвление & : "if"\,условие\,"then"\,операторы\,[\,"else"\,операторы\,]\,"end if" ; \\
цикл & : "while"\,условие\,"do"\,операторы\,"end while" \\
 & $\mid$ "for"\,переменная\,("from"\,выражение\,"to"\,выражение $\mid$
 "$\in$"\,выражение)\,"do"\,операторы\,"end for" . \\
\end{tabular}
\end{center}
\end{frame}



%%%%%%%%%%%%%%%%%%%%%%%%%%%%%%%%%%%%%%%%%%%%%%%%%%%%%%%%%%%%%%%
%%  Далее следуют слайды подготовленные Алексеем Архиповским,
%%  содержимое которых может быть использовно в треьей главе
%%%%%%%%%%%%%%%%%%%%%%%%%%%%%%%%%%%%%%%%%%%%%%%%%%%%%%%%%%%%%%
%%\begin{frame}
%%\frametitle{0.2.4. Псевдокод (2/4). Сводка конструкций}
%%
%%\begin{center}\scriptsize
%%\begin{tabular}{|l|l|l|}
%%\hline
%%Конструкция & Запись & Пример в курсе \\
%%\hline
%%Заголовок функции & \texttt{func} имя(параметр : тип) : тип & сл. 41, 43 \\
%%Заголовок итератора & \texttt{iter} имя(параметры) & сл. 270 \\
%%Исходные данные и результат & Вход: \dots\ Выход: \dots & сл. 79, 139 \\
%%Объявление переменной & имя : тип & сл. 269, 271 \\
%%Присваивание & $x \Assign$ выражение & везде \\
%%Несколько операторов в строке & оператор; оператор & сл. 43, 79 \\
%%Ветвление & \texttt{if} усл. \texttt{then} \dots\ \texttt{else} \dots\
%%\texttt{end if} & сл. 43, 79, 149 \\
%%Цикл со счётчиком & \texttt{for} $i$ \texttt{from} $a$ \texttt{to} $b$
%%\texttt{do} \dots\ \texttt{end for} & сл. 41, 43, 79 \\
%%Цикл по множеству & \texttt{for} $x\in X$ \texttt{do} \dots\ \texttt{end for}
%%& сл. 79, 139, 271 \\
%%Цикл с условием & \texttt{while} усл. \texttt{do} \dots\ \texttt{end while}
%%& сл. 43, 271 \\
%%Возврат значения & \texttt{return} выражение & сл. 41, 43, 269 \\
%%Печать & \texttt{print} список & сл. 149 \\
%%Выдача очередного элемента & \texttt{yield} $x$ & ДМ, п.~4.5.3 \\
%%Работа с очередью & $x\to Q$ (положить), $x\leftarrow Q$ (взять) & сл. 271 \\
%%Создание объекта & \texttt{new} $A$ & сл. 79 \\
%%Пустой оператор & \texttt{skip} & сл. 309 \\
%%Комментарий & \texttt{//} текст до конца строки & везде \\
%%\hline
%%Типы & \textbf{bool}, \textbf{char}, $1..n$, $\mathbb{N}_0$, \textbf{array}
%%$[\dots]$ \textbf{of} $T$, \textbf{queue of} $T$, \textbf{struct}, $\Pointer T$
%%& сл. 269, 271 \\
%%Операции над словами & $|\alpha|$~--- длина, $\alpha[i]$~--- вырезка,
%%$\alpha\beta$~--- сцепление, $\varepsilon$~--- пустое слово & сл. 41, 43 \\
%%\hline
%%\end{tabular}
%%\end{center}
%%\end{frame}
%%
%%\begin{frame}
%%\frametitle{0.2.4. Псевдокод (3/4). Псевдокод как сквозной пример}
%%
%%Псевдокод~--- формальный язык, нужный курсу с первого алгоритма и известный
%%читателю заранее. Поэтому он и берётся сквозным примером: каждое новое средство
%%курса сначала применяется к нему.
%%
%%\vspace{-2pt}
%%\begin{center}\scriptsize
%%\begin{tabular}{|p{5.3cm}|p{5.4cm}|l|}
%%\hline
%%Что описывается & Каким средством & Пункт \\
%%\hline
%%Лексика: идентификаторы, числа, ключевые слова, комментарии & регулярные
%%выражения & 1.2.1--1.2.2 \\
%%Распознавание лексем & языковой граф, детерминизация, конечный автомат &
%%1.2.3--1.2.4, 2.2 \\
%%Синтаксис & КС-грамматика, расширенная запись, БНФ & 1.1.8, 1.3.1 \\
%%Однозначность описания & проблема однозначности КСГ & 1.3.6 \\
%%Пригодность к разбору & грамматики LL($k$) и LR($k$) & 1.3.7 \\
%%Разбор & нисходящий, восходящий, станция Дейкстры & 1.4.2--1.4.4 \\
%%Контекстные условия: описания, типы, области видимости & атрибутные грамматики &
%%1.5.2, 1.5.4 \\
%%Семантика & операционная, исчисление Хоара, денотационная & 1.5.5, 1.5.7--1.5.9 \\
%%Псевдокод как модель вычислимости & тезис Чёрча\,--\,Тьюринга & 0.3.3 \\
%%\hline
%%\end{tabular}
%%\end{center}
%%\end{frame}
%%
%%\begin{frame}
%%\frametitle{0.2.4. Псевдокод (4/4). Сквозной пример: первые два шага}
%%
%%Средства, которыми записаны эти два шага, вводятся позже; здесь достаточно
%%понимать, что $\mid$~--- выбор, $*$~--- повторение, $\{\alpha \# \beta\}$~---
%%повторение $\alpha$ через $\beta$, $[\alpha]$~--- необязательная часть, а
%%терминалы взяты в кавычки.
%%
%%\smallskip
%%{\em Лексика} (в п.~\DMref{1.2}{1.2.} то же будет записано регулярными
%%выражениями):
%%\vspace{-4pt}
%%\begin{center}\footnotesize
%%\begin{tabular}{@{}ll@{}}
%%буква & : "a" $\mid$ \dots\ $\mid$ "z" $\mid$ "A" $\mid$ \dots\ $\mid$ "Z" ; \\
%%цифра & : "0" $\mid$ \dots\ $\mid$ "9" ; \\
%%идентификатор & : буква (буква $\mid$ цифра)$^*$ ; \\
%%число & : цифра цифра$^*$ ; \\
%%комментарий & : "//" (любой символ, кроме конца строки)$^*$ . \\
%%\end{tabular}
%%\end{center}
%%\vspace{-4pt}
%%
%%{\em Синтаксис} (в расширенной записи грамматик из п.~\DMref{1.1.8}{1.1.8.}):
%%\vspace{-4pt}
%%\begin{center}\footnotesize
%%\begin{tabular}{@{}ll@{}}
%%операторы & : \{ оператор \# ";" \} ; \\
%%оператор & : присваивание $\mid$ ветвление $\mid$ цикл \\
%% & $\mid$ "return"\,выражение $\mid$ "print"\,список $\mid$ "skip" ; \\
%%присваивание & : переменная\,": ="\,выражение ; \\
%%ветвление & : "if"\,условие\,"then"\,операторы\,[\,"else"\,операторы\,]\,"end if" ; \\
%%цикл & : "while"\,условие\,"do"\,операторы\,"end while" \\
%% & $\mid$ "for"\,переменная\,"from"\,выражение\,"to"\,выражение\,"do"\,операторы\,"end for" . \\
%%\end{tabular}
%%\end{center}
%%\end{frame}
%%

\subsubsection{0.2.5. Диаграммы UML}\label{0.2.5.}
\begin{frame}
\frametitle{\Qm 0.2.5. Диаграммы UML}

\TODO{ 16.08.26. Онтологии = диаграммы структуры + 
диаграммы зависимостей + диаграммы вход-выход, \\
диаграммы использования, \\
диаграммы деятельности для процессов обработки данных. ФАН}

\end{frame}




% GA_0_2025.tex, 1074–1074
\subsection{0.3. Математическое введение}\label{0.3}

% GA_0_2025.tex, 1807–2543
\subsubsection{0.3.4. Ориентированные графы и деревья}\label{0.3.4}
\begin{frame}
\frametitle{\Er 0.3.4. Ориентированные графы и деревья (1/2)}
\TODO{Графы, орграфы, корневые деревья. ФАН}

\medskip
\TODO{Нотация нагрузки на вершины и на рёбра. Представления графов и деревьев. }

\medskip
\TODO{ФАН 22.10.26. Далее следует один старый слайд ФАН и 6 слайдов Архиповского.
Все надобно объединить и переделать с сокращениями. \\
Слайды в начале. Идея: орграф = бинарное отношение. Пометка = уникальная, отсюда помеченный = нумерованный, а нагруженный != нумерованный!!! Это было обсуждено на семинаре, но не попало в текст.\\
Слайды в середине. Картинки мелкие, подписи к картинкам не используем. Идея пометка/нагрузка = вздор, убрать... Например, в автоматах это стимул/реакция.     
Таблица мелкая, не годится. Слайды должны читаться с разрешением 800х600=VGA. Сомневаюсь, что про представления надо говорить вообще, тем более, с ошибками (Если взять матрицу смежности для дерева, то откуда О(р)?) \\
Слайды в конце. Картинка с сечением = последовательность пометок, а не пунктирная линия!!!
Пример с записью рядом с картинкой!}

\end{frame}

\begin{frame}
\frametitle{0.3.4. Ориентированные графы и деревья (2/2)}

\odmkey{Ориентированное дерево}{Дерево!ориентированное}
~--- это орграф, в котором ровно один узел не имеет
входящих дуг (это узел называется \odmkey{корнем}{Корень!дерева}),
а все остальные узлы имеют ровно одну входящую дугу и при этом достижимы из корня.
Если исходящие дуги в каждом узле линейно упорядочены,
то дерево называется \odmkey{упорядоченным}{Дерево!упорядоченное}.
Узлы, не имеющие исходящих дуг, называются 
\odmkey{листьями}{Лист!дерева}. Путь от корня к листу здесь 
называется \odmkey{ветвью}{Ветвь!дерева}.
Легко видеть, что в упорядоченном ордереве ветви,
а значит, и листья линейно упорядочены. 
Упорядоченная последовательность листьев называется
\odmkey{кроной}{Крона!дерева}.
Последовательность узлов, для которой
каждая ветвь содержит ровно один узел из этой последовательности
называется \odmkey{сечением}{Сечение!дерева}.

\smallskip
Здесь чаще используются \odmkey{помеченные}{Дерево!помеченное}
упорядоченные ордеревья. В качестве пометок используются символы 
из заданных алфавитов. Пометки не обязательно различны.

\smallskip
\TODO{про изображение деревьв, пример}

\smallskip
Диаграммы деревьев наглядны, но занимают много места.
Здесь для записи помеченных упорядоченных ордеревьев
используется следующая текстовая запись:
в круглых скобках записывается пометка корня,
а затем последовательность поддеревьев.       

\end{frame}

\begin{frame}
\frametitle{0.3.4. Ориентированные графы и деревья (1/6)}

Принимаются следующие соглашения об изображении: узлы~--- кружки, дуги~--- линии
со стрелкой от начала к концу; пометка узла пишется внутри кружка, пометка
дуги~--- рядом с дугой; если узел или дуга несёт и пометку, и
\odmkey{нагрузку}{Нагрузка}, они записываются через наклонную черту (пометка\,/\,нагрузка), а
несколько нагрузок перечисляются в фиксированном порядке через ту же черту.

\begin{center}
\includegraphics[width=0.76\textwidth]{graph}
\end{center}
\end{frame}

\begin{frame}
\frametitle{0.3.4. Ориентированные графы и деревья (2/6)}

\begin{remark}
Нагрузка не обязана быть числом. В этом курсе дуги чаще всего нагружены буквами
или цепочками: в языковом графе (п.~\DMref{1.2.3}{1.2.3.}) дуга нагружена буквой,
в дереве порождения (п.~\DMref{1.1.6}{1.1.6.})~--- применённым правилом, в графе
переходов автомата Мили (п.~\DMref{2.1.3}{2.1.3.})~--- парой <<вход\,/\,выход>>, а
у автомата Мура выход нагружает узел, а не дугу.
\end{remark}
\vspace{-4pt}

Орграф с $p$ узлами и $q$ дугами представляется в программах одним из четырёх
способов; через $n(p,q)$ обозначен объём памяти.

\vspace{-2pt}
\begin{center}\scriptsize
\begin{tabular}{|l|p{6.6cm}|l|l|}
\hline
Представление & Запись & $n(p,q)$ & Пример \\
\hline
Матрица смежности & \texttt{M : array} $[1..p,1..p]$ \texttt{of} $0..1$;
$M[i,j]=1$, если есть дуга из $i$ в $j$ & $O(p^2)$ & 16 разрядов \\
Матрица инциденций & \texttt{H : array} $[1..p,1..q]$ \texttt{of} $-1..1$;
$-1$~--- начало дуги, $1$~--- конец & $O(pq)$ & 20 разрядов \\
Списки смежности & $\Gamma$ \texttt{: array} $[1..p]$ \texttt{of} $\Pointer N$,
$N=$ \texttt{struct} $\{v:1..p;\ n:\Pointer N\}$ & $O(p+q)$ & 9 элементов \\
Массив дуг & \texttt{E : array} $[1..q]$ \texttt{of struct} $\{b,e:1..p\}$ &
$O(2q)$ & 10 чисел \\
\hline
\end{tabular}
\end{center}

\end{frame}

\begin{frame}
\frametitle{0.3.4. Ориентированные графы и деревья (3/6)}

\begin{example}
Для орграфа с рис.~1а ($p=4$, $q=5$):
\end{example}
\vspace{-14pt}
$$M=\begin{pmatrix}0&1&0&0\\0&0&1&1\\0&0&0&0\\1&0&1&0\end{pmatrix}\!,\qquad
H=\begin{pmatrix}-1&0&0&1&0\\1&-1&-1&0&0\\0&1&0&0&1\\0&0&1&-1&-1\end{pmatrix}\!;$$
\vspace{-12pt}
$$\Gamma[1]=(2),\ \Gamma[2]=(3,4),\ \Gamma[3]=\varnothing,\ \Gamma[4]=(1,3);
\ \ E=\big((1,2),(2,3),(2,4),(4,1),(4,3)\big).$$
\vspace{-18pt}

Нагрузка хранится отдельно: матрицей \texttt{W : array} $[1..p,1..p]$
\texttt{of} $\mathbb{R}$, где на месте отсутствующей дуги стоит $\infty$, либо
дополнительным полем в элементе списка смежности или массива дуг.

\begin{remark}
Для дерева $q=p-1$, поэтому все представления графов, применённые к дереву, дают
$O(p)$.
\end{remark}
\end{frame}

\begin{frame}
\frametitle{0.3.4. Ориентированные графы и деревья (4/6)}

\vspace{-6pt}
\odmkey{Ориентированное дерево}{Дерево!ориентированное}
~--- это орграф, в котором ровно один узел не имеет
входящих дуг (это узел называется \odmkey{корнем}{Корень!дерева}),
а все остальные узлы имеют ровно одну входящую дугу и при этом достижимы из корня.
Если исходящие дуги в каждом узле линейно упорядочены,
то дерево называется \odmkey{упорядоченным}{Дерево!упорядоченное}.
Узлы, не имеющие исходящих дуг, называются
\odmkey{листьями}{Лист!дерева}. Путь от корня к листу здесь
называется \odmkey{ветвью}{Ветвь!дерева}.
Легко видеть, что в упорядоченном ордереве ветви,
а значит, и листья линейно упорядочены.
Упорядоченная последовательность листьев называется
\odmkey{кроной}{Крона!дерева}.
Последовательность узлов, для которой
каждая ветвь содержит ровно один узел из этой последовательности
называется \odmkey{сечением}{Сечение!дерева}.

Здесь чаще используются \odmkey{помеченные}{Дерево!помеченное}
упорядоченные ордеревья. В качестве пометок используются символы
из заданных алфавитов. Пометки не обязательно различны.

При изображении деревьев принято дополнительное соглашение: корень наверху, все
дуги ориентированы сверху вниз, поэтому стрелки не изображаются; поддеревья
каждого узла располагаются слева направо в порядке их следования.

\begin{remark}
Диаграммы свободного, ориентированного и упорядоченного деревьев графически
неразличимы, поэтому при каждом рисунке нужно оговаривать, дерево какого класса
изображено; здесь всюду~--- упорядоченные ордеревья.
\end{remark}
\vspace{-4pt}
\end{frame}

\begin{frame}
\frametitle{0.3.4. Ориентированные графы и деревья (5/6)}

\begin{center}
\includegraphics[width=0.82\textwidth]{tree0}
\end{center}
\vspace{-8pt}

\begin{example}
Помеченное упорядоченное ордерево: корень помечен $a$, листья~--- $d$, $e$, $f$,
$g$, $i$, крона~--- последовательность листьев слева направо. Пунктиром показано
сечение $\{d,e,c\}$: каждая из пяти ветвей $a$--$b$--$d$, $a$--$b$--$e$,
$a$--$c$--$f$, $a$--$c$--$g$ и $a$--$c$--$h$--$i$ содержит ровно один узел этого
множества.
\end{example}
\end{frame}

\begin{frame}
\frametitle{0.3.4. Ориентированные графы и деревья (6/6)}

Диаграммы деревьев наглядны, но занимают много места. Здесь для записи помеченных
упорядоченных ордеревьев используется следующая текстовая запись: в круглых
скобках записывается пометка корня, а затем последовательность поддеревьев.

\begin{example}
Дерево с предыдущего слайда записывается так:
\vspace{-8pt}
$$\texttt{( a ( b (d) (e) ) ( c (f) (g) ( h (i) ) ) )}$$
\vspace{-20pt}

Если оставить одни скобки, получится {\em битовая шкала} из $2p$ разрядов
(открывающая скобка~--- 1, закрывающая~--- 0), и по ней дерево восстанавливается
однозначно:
\vspace{-8pt}
$$\texttt{1 1 1 0 1 0 0 1 1 0 1 0 1 1 0 0 0 0}$$
\end{example}
\vspace{-16pt}

\begin{center}\scriptsize
\begin{tabular}{|l|l|l|}
\hline
Представление дерева & Запись & Объём \\
\hline
Скобочная запись & \texttt{( a ( b (d) (e) ) \dots\ )} & $O(p)$ символов \\
Битовая шкала & \texttt{C : array} $[1..2p]$ \texttt{of} $0..1$ & $2p$ разрядов \\
Списки смежности & $\Gamma$ \texttt{: array} $[1..p]$ \texttt{of} $\Pointer N$ &
$O(p)$ \\
Первый сын и следующий брат & $N=$ \texttt{struct} $\{v;\ s,b:\Pointer N\}$ &
$O(p)$ \\
Массив предков & \texttt{up : array} $[1..p]$ \texttt{of} $0..p$ & $O(p)$ \\
\hline
\end{tabular}
\end{center}
\end{frame}

% 0.3.4. Ориентированные графы и деревья. А. Архиповский, 25.09.2026.
% Вторая редакция. Основание: замечания ФАН от 17.09.2026 к TODO_слайды_27-28.
% Что изменено против первой редакции:
%   1) исправлено главное: сечение проходит через УЗЛЫ, а не через рёбра;
%      сечение — это не разрез, поэтому корень и крона суть сечения.
%      Рисунок переделан: сечение показано выделенными узлами, а не линией,
%      пересекающей дуги, и приведены три сечения одного дерева;
%   2) добавлены кадры, которых требуют пометки ФАН на слайдах 29–30:
%      графы, орграфы, корневые деревья; нотация нагрузки на узлы и на дуги;
%      представления графов и деревьев; изображение деревьев с примером;
%   3) различие пометки и нагрузки вынесено в отдельный кадр.
% Рисунки: fig/graph.drawio, fig/tree0.drawio (исходники в drawio).

\begin{frame}
\frametitle{0.3.4. Ориентированные графы и деревья (1/8)}
%%% Орграф: определения

\odmkey{Ориентированный граф}{Граф!ориентированный} (сокращённо {\em орграф})~---
это пара $D=\langle V,E\rangle$, где $V$~--- конечное множество
\odmkey{узлов}{Узел графа}, а $E\subset V\times V$~--- множество
\odmkey{дуг}{Дуга графа}. Для дуги $e=(u,v)$ узел $u$ называется её
{\em началом}, узел $v$~--- {\em концом}; дуга вида $(u,u)$ называется
{\em петлёй}.

\odmkey{Путь}{Путь в графе}~--- это последовательность дуг, в которой конец
каждой дуги есть начало следующей; {\em длина пути}~--- число дуг в нём.
Путь, начало которого совпадает с концом, называется
\odmkey{контуром}{Контур}. Узел $v$ {\em достижим} из узла $u$, если существует
путь из $u$ в $v$; путь нулевой длины считается допустимым, поэтому всякий узел
достижим из самого себя.

Число дуг, входящих в узел, называется {\em полустепенью захода}, число
исходящих~--- {\em полустепенью исхода}.

\begin{remark}
Если ориентация дуг несущественна, говорят о {\em свободном} графе, а вместо
дуг~--- о {\em рёбрах}, то есть о неупорядоченных парах узлов. В этом курсе
свободные графы почти не встречаются: и языковой граф, и граф переходов
автомата, и деревья разбора ориентированы.
\end{remark}
\end{frame}

\begin{frame}
\frametitle{0.3.4. Ориентированные графы и деревья (2/8)}
%%% Пометка и нагрузка. Соглашения об изображении

\odmkey{Пометка}{Пометка} и \odmkey{нагрузка}{Нагрузка}~--- разные вещи.
Пометка~--- это символ заданного алфавита, отвечающий на вопрос <<что это
такое>>; пометки разных узлов могут совпадать. Нагрузка~--- это дополнительные
данные произвольной природы, отвечающие на вопрос <<что про это известно>>.

Принимаются соглашения об изображении: узлы~--- кружки, дуги~--- линии со
стрелкой от начала к концу; пометка узла пишется внутри кружка, пометка дуги~---
рядом с дугой; пометка и нагрузка записываются через наклонную черту
(пометка\,/\,нагрузка), а несколько нагрузок~--- в фиксированном порядке через
ту же черту.

\begin{center}
\includegraphics[width=0.63\textwidth]{0_3_4_2}
\end{center}
\end{frame}

\begin{frame}
\frametitle{0.3.4. Ориентированные графы и деревья (3/8)}
%%% Нагрузка бывает не числом. Представления орграфа

\begin{remark}
Нагрузка не обязана быть числом. В этом курсе дуги чаще всего нагружены буквами
или цепочками: в языковом графе (п.~\DMref{1.2.3}{1.2.3.}) дуга нагружена буквой,
в дереве порождения (п.~\DMref{1.1.6}{1.1.6.})~--- применённым правилом, в графе
переходов автомата Мили (п.~\DMref{2.1.3}{2.1.3.})~--- парой <<вход\,/\,выход>>, а
у автомата Мура выход нагружает узел, а не дугу.
\end{remark}

Орграф с $p$ узлами и $q$ дугами представляется в программах одним из четырёх
способов; через $n(p,q)$ обозначен объём памяти, в последнем столбце~--- его
значение для орграфа с рисунка а).

\vspace{-2pt}
\begin{center}\scriptsize
\begin{tabular}{|l|p{6.6cm}|l|l|}
\hline
Представление & Запись & $n(p,q)$ & Пример \\
\hline
Матрица смежности & \texttt{M : array} $[1..p,1..p]$ \texttt{of} $0..1$;
$M[i,j]=1$, если есть дуга из $i$ в $j$ & $O(p^2)$ & 16 разрядов \\
Матрица инциденций & \texttt{H : array} $[1..p,1..q]$ \texttt{of} $-1..1$;
$-1$~--- начало дуги, $1$~--- конец & $O(pq)$ & 20 разрядов \\
Списки смежности & $\Gamma$ \texttt{: array} $[1..p]$ \texttt{of} $\Pointer N$,
$N=$ \texttt{struct} $\{v:1..p;\ n:\Pointer N\}$ & $O(p+q)$ & 9 элементов \\
Массив дуг & \texttt{E : array} $[1..q]$ \texttt{of struct} $\{b,e:1..p\}$ &
$O(2q)$ & 10 чисел \\
\hline
\end{tabular}
\end{center}
\end{frame}

\begin{frame}
\frametitle{0.3.4. Ориентированные графы и деревья (4/8)}
%%% Пример представлений

\begin{example}
Для орграфа с рисунка а) ($p=4$, $q=5$):
\end{example}
\vspace{-14pt}
$$M=\begin{pmatrix}0&1&0&0\\0&0&1&1\\0&0&0&0\\1&0&1&0\end{pmatrix}\!,\qquad
H=\begin{pmatrix}-1&0&0&1&0\\1&-1&-1&0&0\\0&1&0&0&1\\0&0&1&-1&-1\end{pmatrix}\!;$$
\vspace{-12pt}
$$\Gamma[1]=(2),\ \Gamma[2]=(3,4),\ \Gamma[3]=\varnothing,\ \Gamma[4]=(1,3);
\ \ E=\big((1,2),(2,3),(2,4),(4,1),(4,3)\big).$$
\vspace{-18pt}

Нагрузка хранится отдельно: матрицей \texttt{W : array} $[1..p,1..p]$
\texttt{of} $\mathbb{R}$, где на месте отсутствующей дуги стоит $\infty$, либо
дополнительным полем в элементе списка смежности или массива дуг. Пометка
хранится так же, но в массиве, занумерованном узлами.

\begin{remark}
Выбор представления определяется тем, какие вопросы задаются графу: матрица
смежности отвечает на вопрос <<есть ли дуга из $i$ в $j$>> за $O(1)$, а списки
смежности перебирают всех соседей узла за время, пропорциональное их числу.
\end{remark}
\end{frame}



\begin{frame}
\frametitle{0.3.4. Ориентированные графы и деревья (5/8)}
%%% Ориентированное дерево

\vspace{-4pt}
\odmkey{Ориентированное дерево}{Дерево!ориентированное}
~--- это орграф, в котором ровно один узел не имеет
входящих дуг (этот узел называется \odmkey{корнем}{Корень!дерева}),
а все остальные узлы имеют ровно одну входящую дугу и при этом достижимы из корня.
Если исходящие дуги в каждом узле линейно упорядочены,
то дерево называется \odmkey{упорядоченным}{Дерево!упорядоченное}.
Узлы, не имеющие исходящих дуг, называются
\odmkey{листьями}{Лист!дерева}. Путь от корня к листу здесь
называется \odmkey{ветвью}{Ветвь!дерева}.
Легко видеть, что в упорядоченном ордереве ветви,
а значит, и листья линейно упорядочены.
Упорядоченная последовательность листьев называется
\odmkey{кроной}{Крона!дерева}.

Здесь чаще используются \odmkey{помеченные}{Дерево!помеченное}
упорядоченные ордеревья. В качестве пометок используются символы
из заданных алфавитов. Пометки не обязательно различны.

При изображении деревьев принято дополнительное соглашение: корень наверху, все
дуги ориентированы сверху вниз, поэтому стрелки не изображаются; поддеревья
каждого узла располагаются слева направо в порядке их следования.

\begin{remark}
Диаграммы свободного, ориентированного и упорядоченного деревьев графически
неразличимы, поэтому при каждом рисунке нужно оговаривать, дерево какого класса
изображено; здесь всюду~--- упорядоченные ордеревья.
\end{remark}
\end{frame}

\begin{frame}
\frametitle{0.3.4. Ориентированные графы и деревья (6/8)}
%%% Сечение: это множество узлов, а не разрез

Множество узлов, для которого каждая ветвь содержит ровно один узел из этого
множества, называется \odmkey{сечением}{Сечение!дерева}. В упорядоченном
ордереве узлы сечения линейно упорядочены.

\begin{center}
\includegraphics[width=0.84\textwidth]{0_3_4_6}
\end{center}
\vspace{-6pt}

\begin{examples}
Множество $\{a,c\}$ сечением не является: ветвь $a$--$c$--$f$ содержит два его
узла; $\{b\}$~--- тоже: та же ветвь не содержит ни одного.
\end{examples}
\vspace{-2pt}

\begin{remark}
Сечение~--- это {\em не} разрез: разрез есть множество дуг, а сечение~---
множество узлов. Линией, пересекающей дуги, сечение не изображается: иначе ни
корень, ни крона сечениями бы не были, хотя по определению они сечения.
\end{remark}
\end{frame}


\begin{frame}
\frametitle{0.3.4. Ориентированные графы и деревья (7/8)}
%%% Текстовая запись дерева

Диаграммы деревьев наглядны, но занимают много места. Здесь для записи помеченных
упорядоченных ордеревьев используется следующая текстовая запись: в круглых
скобках записывается пометка корня, а затем последовательность поддеревьев.

\begin{example}
Дерево с предыдущего кадра записывается так:
\vspace{-8pt}
$$\texttt{( a ( b (d) (e) ) ( c (f) (g) ( h (i) ) ) )}$$
\vspace{-20pt}

Если оставить одни скобки, получится {\em битовая шкала} из $2p$ разрядов
(открывающая скобка~--- 1, закрывающая~--- 0), и по ней дерево восстанавливается
однозначно:
\vspace{-8pt}
$$\texttt{1 1 1 0 1 0 0 1 1 0 1 0 1 1 0 0 0 0}$$
\end{example}
\vspace{-14pt}

\begin{remark}
Крона читается из скобочной записи как последовательность пометок тех узлов, за
которыми сразу идёт закрывающая скобка: $d$, $e$, $f$, $g$, $i$. Всякое сечение
читается так же легко, если известно, на какой глубине оно проходит,~--- но
сечение, вообще говоря, на одной глубине не лежит (см. $\{d,e,c\}$ на рисунке).
\end{remark}
\end{frame}

\begin{frame}
\frametitle{0.3.4. Ориентированные графы и деревья (8/8)}
%%% Представления дерева

Дерево~--- частный случай орграфа, поэтому годятся все четыре представления из
кадра 3/8; но для дерева $q=p-1$, и все они дают $O(p)$. Кроме того, у дерева
есть представления, которых у орграфа нет.

\vspace{-2pt}
\begin{center}\scriptsize
\begin{tabular}{|l|l|l|}
\hline
Представление дерева & Запись & Объём \\
\hline
Скобочная запись & \texttt{( a ( b (d) (e) ) \dots\ )} & $O(p)$ символов \\
Битовая шкала & \texttt{C : array} $[1..2p]$ \texttt{of} $0..1$ & $2p$ разрядов \\
Списки смежности & $\Gamma$ \texttt{: array} $[1..p]$ \texttt{of} $\Pointer N$ &
$O(p)$ \\
Первый сын и следующий брат & $N=$ \texttt{struct} $\{v;\ s,b:\Pointer N\}$ &
$O(p)$ \\
Массив предков & \texttt{up : array} $[1..p]$ \texttt{of} $0..p$ & $O(p)$ \\
\hline
\end{tabular}
\end{center}
\vspace{-2pt}

\begin{remark}
Массив предков задаёт ордерево, но не задаёт порядок сыновей, поэтому для
упорядоченного ордерева его недостаточно. Представление <<первый сын и следующий
брат>> порядок сохраняет и при этом расходует на узел ровно два указателя, сколько
бы сыновей у узла ни было; оно и применяется дальше при построении деревьев
разбора (п.~\DMref{1.4.2}{1.4.2.}).
\end{remark}
\end{frame}



\subsubsection{0.3.5. Решётки и наименьшая неподвижная точка}\label{0.3.5.}

\begin{frame}

\frametitle{{0.3.5. Решётки и наименьшая неподвижная точка (1/1)}}

\odmkey{Решётка}{Решётка} — это множество M с двумя бинарными операциями, $\cap$ и $\cup$, такими, что выполнены следующие условия:
\begin{enumerate}
\item ассоциативность: $(a \cap b) \cap c = a \cap (b \cap c), (a \cup b) \cup c = a \cup (b \cup c)$ ;
\item коммутативность: $a \cap b = b \cap a, a \cup b = b \cup a$ ;
\item поглощение: $(a \cap b) \cup a = a, (a \cup b) \cap a = a$ ;
\item идемпотентность: $a \cap a = a, a \cup a = a $.
\end{enumerate}

\begin{example}
Булеан любого множества M образует решётку $\langle 2^M; \cap, \cup\rangle$ относительно пересечения и объединения.
\end{example}

\TODO{Другая формулировка теоремы, привязанная к решёткам}

\begin{theorem} {\NB{[о наименьшей неподвижной точке функции]}}
Если множество $X$ частично упорядочено, линейно полно, 
имеет наименьший элемент $0 = \min X$ и функция $\Map{f }{X}{X}$ непрерывна и тотальна, $\Dom f = X$, то элемент 
$a = \sup \{f^n(0) \mid n \ge 0\}$ является наименьшей неподвижной точкой функции $f$.
\end{theorem}

\end{frame}

\begin{frame}
\frametitle{0.3.5. Решётка и наименьшая неподвижная точка (1/3)}

\odmkey{Решётка}{Решётка}~--- это множество $L$ с двумя бинарными операциями
$\cap$ и $\cup$, для которых выполнены
\vspace{-4pt}
\begin{enumerate}
\item ассоциативность: $(a\cap b)\cap c=a\cap(b\cap c)$,
$(a\cup b)\cup c=a\cup(b\cup c)$;
\item коммутативность: $a\cap b=b\cap a$, $a\cup b=b\cup a$;
\item поглощение: $(a\cap b)\cup a=a$, $(a\cup b)\cap a=a$.
\end{enumerate}
\vspace{-4pt}

\begin{remark}
Идемпотентность $a\cap a=a$, $a\cup a=a$~--- следствие аксиом 2 и 3, а не
самостоятельная аксиома. Из (3) при $b\Assign a$ получаем $(a\cap a)\cup a=a$;
подставив в (3) $b\Assign a\cap a$, имеем $(a\cup(a\cap a))\cap a=a$, то есть
$a\cap a=a$. Двойственно доказывается $a\cup a=a$.
\end{remark}

{\em Индуцированный порядок.} Положим $a\le b$, если $a\cap b=a$. Тогда $\le$~---
нестрогий частичный порядок, причём $a\cap b=\inf\{a,b\}$, $a\cup b=\sup\{a,b\}$,
и $a\cap b=a$ равносильно $a\cup b=b$. Обратно, если в частично упорядоченном
множестве всякая пара элементов имеет точные грани, то операции
$a\cap b\Assign\inf\{a,b\}$ и $a\cup b\Assign\sup\{a,b\}$ удовлетворяют аксиомам
1--3. Тем самым алгебраическое и порядковое определения решётки эквивалентны.
\end{frame}

\begin{frame}
\frametitle{0.3.5. Решётка и наименьшая неподвижная точка (2/3)}

Решётка $L$ называется \odmkey{полной}{Решётка!полная}, если точные грани есть не
только у пар, но у всякого подмножества $S\subset L$, включая пустое; они
обозначаются $\bigwedge S$ и $\bigvee S$. В полной решётке есть наименьший
элемент $0=\bigvee\emptyset$ и наибольший $1=\bigwedge\emptyset$.

\begin{example}
Булеан $\langle 2^{\mathrm{M}};\subset\rangle$: $\bigwedge S=\bigcap S$,
$\bigvee S=\bigcup S$, $0=\emptyset$, $1=\mathrm{M}$.
\end{example}

Функция $f:L\to L$ {\em монотонна}, если $x\le y$ влечёт $f(x)\le f(y)$; элемент
$x$~--- \odmkey{неподвижная точка}{Точка!неподвижная}, если $f(x)=x$.

\begin{theorem} \NB{[Кнастера\,--\,Тарского]}
Пусть $L$~--- полная решётка и $f:L\to L$ монотонна. Тогда множество неподвижных
точек непусто и само является полной решёткой, а наименьшая и наибольшая из них
суть $\mu f=\bigwedge\Set{x\in L}{f(x)\le x}$ и
$\nu f=\bigvee\Set{x\in L}{x\le f(x)}$.
\end{theorem}

\begin{proof}
Положим $P\Assign\Set{x\in L}{f(x)\le x}$, $m\Assign\bigwedge P$; множество $P$
непусто, так как $1\in P$. Для $x\in P$ имеем $m\le x$, поэтому
$f(m)\le f(x)\le x$; значит $f(m)$~--- нижняя граница $P$, откуда $f(m)\le m$. По
монотонности $f(f(m))\le f(m)$, то есть $f(m)\in P$, а потому $m\le f(m)$. Итак,
$f(m)=m$. Всякая неподвижная точка $x$ удовлетворяет $f(x)\le x$, то есть лежит
в $P$, и потому $m\le x$.
\end{proof}
\end{frame}

\begin{frame}
\frametitle{0.3.5. Решётка и наименьшая неподвижная точка (3/3)}

{\em Цепь}~--- линейно упорядоченное подмножество $L$. Функция $f$
\odmkey{непрерывна}{Функция!непрерывная}, если для всякой возрастающей цепи
$x_0\le x_1\le x_2\le\dots$ выполнено
$f\left(\bigvee_n x_n\right)=\bigvee_n f(x_n)$.

\begin{remark}
Непрерывность влечёт монотонность: при $x\le y$ цепь $x\le y$ даёт
$f(y)=f(x\cup y)=f(x)\cup f(y)$, то есть $f(x)\le f(y)$. Обратное неверно.
\end{remark}

\begin{theorem} \NB{[Клини]}
Если $L$~--- полная решётка с наименьшим элементом $0$ и функция $f$ непрерывна,
то $\mu f=\bigvee\Set{f^n(0)}{n\geqslant 0}$, где $f^0(0)=0$,
$f^{n+1}(0)=f(f^n(0))$.
\end{theorem}
\vspace{-4pt}

\begin{proof}
Из $0\le f(0)$ и монотонности индукцией по $n$ получаем $f^n(0)\le f^{n+1}(0)$,
то есть $f^n(0)$~--- возрастающая цепь. Положим
$a\Assign\bigvee\Set{f^n(0)}{n\geqslant 0}$; тогда
\vspace{-6pt}
$$f(a)=\bigvee\Set{f^{n+1}(0)}{n\geqslant 0}=a,$$
\vspace{-18pt}

поскольку добавление члена $f^0(0)=0$ не меняет точной верхней грани. Пусть
$b$~--- произвольная неподвижная точка; из $f^n(0)\le b$ следует
$f^{n+1}(0)\le f(b)=b$, откуда $a\le b$.
\end{proof}
\vspace{-4pt}
\begin{corollary}
Если решётка $L$ конечна, то всякая монотонная функция на ней непрерывна, а цепь
$0\le f(0)\le f^2(0)\le\dots$ стабилизируется не более чем за $|L|$ шагов.
\end{corollary}

\end{frame}

\begin{frame}
\frametitle{0.3.5. Решётка и наименьшая неподвижная точка (послесловие)}

\TODO{ 18.10.26. ФАН. Предыдущий вариант прислан Алексеем Архиповским 18.09.26.
Этот вариант не вполне удовлетворителен по трём причинам.}

\medskip
\TODO{ 18.10.26. ФАН. Первая причина = плохая математика. Кому мешала аксиоматическая идемпотентность?  Чем отличается супремум от точной верхней грани? Квантор $\mu$ засунут в текст молча. Все это невнятно, увы ... .}

\medskip
\TODO{ 18.10.26. ФАН. Вторая причина = плохое форматирование. Нет дат жизни учёных. Половина имён упомянута для демонстрации <<взезнайста>>, а не для дела. В отечественной 
традиции при переносе формулы знак дублируется на второй строке и т.д. }


\medskip
\TODO{ 18.10.26. ФАН. Третья причина = плохая структура. Где это будет использовано
(решение уравнений с регулярными коэффициентами)? 
Почему там подойдёт такая формулировка? А если использовано не будет, 
то и параграф этот не нужен.}


\end{frame}

% 0.3.5. Решётка и наименьшая неподвижная точка. А. Архиповский, 25.09.2026.
% Протокол ФАН от 18.09.2026, TODO_слайд_29: «Вся эта математика наукообразная,
% красивая, но в теории грамматик и автоматов не очень нужная. Была зацеплена
% только ради решения уравнений с регулярными коэффициентами. Степень
% актуальности не очень большая, предлагаю пока отложить этот материал, но не
% выбрасывать.»
% Кадры сохранены без изменений и в текущую редакцию учебника не предлагаются.

\begin{frame}
\frametitle{0.3.5. Решётка и наименьшая неподвижная точка (1/3)}

\odmkey{Решётка}{Решётка}~--- это множество $L$ с двумя бинарными операциями
$\cap$ и $\cup$, для которых выполнены
\vspace{-4pt}
\begin{enumerate}
\item ассоциативность: $(a\cap b)\cap c=a\cap(b\cap c)$,
$(a\cup b)\cup c=a\cup(b\cup c)$;
\item коммутативность: $a\cap b=b\cap a$, $a\cup b=b\cup a$;
\item поглощение: $(a\cap b)\cup a=a$, $(a\cup b)\cap a=a$.
\end{enumerate}
\vspace{-4pt}

\begin{remark}
Идемпотентность $a\cap a=a$, $a\cup a=a$~--- следствие аксиом 2 и 3, а не
самостоятельная аксиома. Из (3) при $b\Assign a$ получаем $(a\cap a)\cup a=a$;
подставив в (3) $b\Assign a\cap a$, имеем $(a\cup(a\cap a))\cap a=a$, то есть
$a\cap a=a$. Двойственно доказывается $a\cup a=a$.
\end{remark}

{\em Индуцированный порядок.} Положим $a\le b$, если $a\cap b=a$. Тогда $\le$~---
нестрогий частичный порядок, причём $a\cap b=\inf\{a,b\}$, $a\cup b=\sup\{a,b\}$,
и $a\cap b=a$ равносильно $a\cup b=b$. Обратно, если в частично упорядоченном
множестве всякая пара элементов имеет точные грани, то операции
$a\cap b\Assign\inf\{a,b\}$ и $a\cup b\Assign\sup\{a,b\}$ удовлетворяют аксиомам
1--3. Тем самым алгебраическое и порядковое определения решётки эквивалентны.
\end{frame}

\begin{frame}
\frametitle{0.3.5. Решётка и наименьшая неподвижная точка (2/3)}

Решётка $L$ называется \odmkey{полной}{Решётка!полная}, если точные грани есть не
только у пар, но у всякого подмножества $S\subset L$, включая пустое; они
обозначаются $\bigwedge S$ и $\bigvee S$. В полной решётке есть наименьший
элемент $0=\bigvee\emptyset$ и наибольший $1=\bigwedge\emptyset$.

\begin{example}
Булеан $\langle 2^{\mathrm{M}};\subset\rangle$: $\bigwedge S=\bigcap S$,
$\bigvee S=\bigcup S$, $0=\emptyset$, $1=\mathrm{M}$.
\end{example}

Функция $f:L\to L$ {\em монотонна}, если $x\le y$ влечёт $f(x)\le f(y)$; элемент
$x$~--- \odmkey{неподвижная точка}{Точка!неподвижная}, если $f(x)=x$.

\begin{theorem} \NB{[Кнастера\,--\,Тарского]}
Пусть $L$~--- полная решётка и $f:L\to L$ монотонна. Тогда множество неподвижных
точек непусто и само является полной решёткой, а наименьшая и наибольшая из них
суть $\mu f=\bigwedge\Set{x\in L}{f(x)\le x}$ и
$\nu f=\bigvee\Set{x\in L}{x\le f(x)}$.
\end{theorem}

\begin{proof}
Положим $P\Assign\Set{x\in L}{f(x)\le x}$, $m\Assign\bigwedge P$; множество $P$
непусто, так как $1\in P$. Для $x\in P$ имеем $m\le x$, поэтому
$f(m)\le f(x)\le x$; значит $f(m)$~--- нижняя граница $P$, откуда $f(m)\le m$. По
монотонности $f(f(m))\le f(m)$, то есть $f(m)\in P$, а потому $m\le f(m)$. Итак,
$f(m)=m$. Всякая неподвижная точка $x$ удовлетворяет $f(x)\le x$, то есть лежит
в $P$, и потому $m\le x$.
\end{proof}
\end{frame}

\begin{frame}
\frametitle{0.3.5. Решётка и наименьшая неподвижная точка (3/3)}

{\em Цепь}~--- линейно упорядоченное подмножество $L$. Функция $f$
\odmkey{непрерывна}{Функция!непрерывная}, если для всякой возрастающей цепи
$x_0\le x_1\le x_2\le\dots$ выполнено
$f\left(\bigvee_n x_n\right)=\bigvee_n f(x_n)$.

\begin{remark}
Непрерывность влечёт монотонность: при $x\le y$ цепь $x\le y$ даёт
$f(y)=f(x\cup y)=f(x)\cup f(y)$, то есть $f(x)\le f(y)$. Обратное неверно.
\end{remark}

\begin{theorem} \NB{[Клини]}
Если $L$~--- полная решётка с наименьшим элементом $0$ и функция $f$ непрерывна,
то $\mu f=\bigvee\Set{f^n(0)}{n\geqslant 0}$, где $f^0(0)=0$,
$f^{n+1}(0)=f(f^n(0))$.
\end{theorem}
\vspace{-4pt}

\begin{proof}
Из $0\le f(0)$ и монотонности индукцией по $n$ получаем $f^n(0)\le f^{n+1}(0)$,
то есть $f^n(0)$~--- возрастающая цепь. Положим
$a\Assign\bigvee\Set{f^n(0)}{n\geqslant 0}$; тогда
\vspace{-6pt}
$$f(a)=\bigvee\Set{f^{n+1}(0)}{n\geqslant 0}=a,$$
\vspace{-18pt}

поскольку добавление члена $f^0(0)=0$ не меняет точной верхней грани. Пусть
$b$~--- произвольная неподвижная точка; из $f^n(0)\le b$ следует
$f^{n+1}(0)\le f(b)=b$, откуда $a\le b$.
\end{proof}
\vspace{-4pt}
\begin{corollary}
Если решётка $L$ конечна, то всякая монотонная функция на ней непрерывна, а цепь
$0\le f(0)\le f^2(0)\le\dots$ стабилизируется не более чем за $|L|$ шагов.
\end{corollary}
\end{frame}


% GA_1_2025.tex, 1–695
\section{1. Языки и грамматики}\label{1.}
\begin{frame}
\frametitle{1. Языки и грамматики }

Начинаем мы с самого заслуженного и 
до сих пор самого популярного средства описания синтаксиса формальных языков, имеющих текстовое представление,~---
 \odmkey{теории формальных
грамматик}{Теория!формальных!грамматик}.


\medskip
Для особенно простых, но часто применяемых в программировании
формальных языков, которые называются регулярными,
имеется средство описания, которое называется
\odmkey{регулярные выражения}{Выражение!регулярное}.

\medskip
В этой главе приводятся основные определения, факты и алгоритмы из
теории формальных грамматик, излагается техника регулярных выражений
и рассматривается применение описанного аппарата для работы
с языками программирования, в частности,
с языками предметных областей.

\end{frame}


\subsection{1.1. Порождающие грамматики}
\label{1.1.}
\begin{frame}
\frametitle{1.1. Порождающие грамматики (1/3)}
Первым рассматривается простейший случай, когда система знаков~--- это
конечный алфавит символов, а способ
комбинирования~--- это составление линейных последовательностей
символов, то есть {\em текстов}.

\bigskip
\begin{digress}
Исторически идеи порождающих грамматик возникли 
применительно к текстам, но в настоящее время
развились до применения к изображениям, звукам, жестам и многим 
другим символам различной природы. 
Некоторые примеры рассматриваются в главе 3.  
\end{digress}
\end{frame}

\begin{frame}
\frametitle{1.1. Порождающие грамматики (2/3)}

\begin{tabular}{p{10mm}|p{40mm}|p{95mm}}
  № & Название параграфа
  & Ключевые термины и обозначения \\
  \hline
 \text{\DMref{1.1.1.}{1.1.1.}} &  Алфавит, слово и язык
  & {  
  алфавит, буква, символ, слово, цепочка, строка, пустое слово $\varepsilon$,
  вхождение, вырезка, длина слова, конкатенация, сцепление, степень, префикс, суффикс,
  пересечение слов, язык, лексикографический порядок }
   \\
\text{\DMref{1.1.2.}{1.1.2.}} & Формальные грамматики 
  & 
{ порождающая грамматика,
  нетерминал, переменная, терминал, константа,
  правило вывода, продукция, аксиома грамматики,
  вывод, непосредственный вывод, бесповторный вывод, 
  порождаемый язык $\frak{L}$,  допуск слова,
  эквивалентность грамматик  
}
  \\
\text{\DMref{1.1.3.}{1.1.3.}}  & Перечислимость языков &
\\
\text{\DMref{1.1.4.}{1.1.4.}}  & Классификация грамматик и языков &
{Ноам Хомский, грамматика: с фразовой структу\-рой, непосредственных составляющих,
 контекстно-зависи\-мая, контекстно-свободная,
 регулярная, $\frak{T}_i$
 } 
\end{tabular}

\end{frame}


\begin{frame}
\frametitle{1.1. Порождающие грамматики (3/3)}

\begin{tabular}{p{10mm}|p{40mm}|p{95mm}}
  № & Название параграфа
  & Ключевые термины и обозначения \\
  \hline
\text{\DMref{1.1.5.}{1.1.5.}} & Неукорачивающие грамматики &
{разрешимый, полуразрешимый, легко разрешимый язык, распознающий алгоритм}\\   
\text{\DMref{1.1.6.}{1.1.6.}} & Контекстно-свободные грамматики &
{левосторонний вывод, правосторонний вывод, дерево разбора, дерево вывода,
дерево порождения, неоднозначная грамматика}\\   
\text{\DMref{1.1.7.}{1.1.7.}} & Линейные, регулярные и автоматные грамматики &
{правило: начальное, заключительное, леволинейное, праволинейное, 
порождающие и недостижимые символы, цепное правило, 
$\varepsilon$-правило, $S$-правило, приведённая грамматика }\\   
\text{\DMref{1.1.8.}{1.1.8.}} & Языки записи грамматик &
{ метасимвол, нормальная форма Бэкуса--Наура,
металингвистическая переменная, металингвистическая формула }\\   
\text{\DMref{1.1.9.}{1.1.9.}} & Альтернативные средства задания синтаксиса &
{Никлаус Вирт, диаграмма, тривиальная диаграмма, элементарная диаграмма, 
язык разметки   }\\   

\end{tabular}

\end{frame}


\subsubsection{1.1.1. Алфавит, слово и язык}\label{1.1.1}
\begin{frame}
\frametitle{1.1.1. Алфавит, слово и язык (1/9)}
Пусть задано конечное множество $A =\{\Tuple{a}{1}{n}\}$, которое
называется \odmkey{алфавитом}{Алфавит}. Элементы алфавита называются
\odmkey{буквами}{Буква} или \odmkey{символами}{Символ}.
Обычно отдельные символы обозначаются
начальными строчными буквами латинского алфавита.
Последовательность букв называется
\odmkey{словом}{Слово} (в данном алфавите).
Заметим, что в программировании слово часто называют
\odmkey{строкой}{Строка}, а в теории формальных языков~--- \odmkey{цепочкой}{Цепочка}.
Буквы в слове упорядочены (слева направо),
и могут быть перенумерованы.
Одна буква может входить
в слово несколько раз, каждый
отдельный экземпляр буквы в слове называется
\odmkey{вхождением}{Вхождение} буквы в слово.
Вхождения могут быть идентифицированы,
например, номером буквы в слове 
 относительно начала слова.
Обычно слова обозначают
строчными буквами греческого алфавита,
а номера записывают в квадратных скобках.
\begin{example}
Пусть $A=\{a,b\}$. Тогда 
$\alpha=aabb$~--- слово, 
а $\alpha[3]=b$~--- буква.
\end{example}

\begin{remark}
Вхождения можно идентифицировать не номером,
а смещением~--- это вопрос вкуса, который регулируется предварительным соглашением.
\end{remark}

Количество вхождений букв в слово называется \odmkey{длиной}{Длина!слова} слова:
$|\alpha| = |a_1\dots a_k| = k$.
\begin{example} $|$Ура-а-а!$| = 8$.
\end{example}

\odmkey{Пустое}{Слово!пустое} слово обозначается
$\varepsilon$,
$|\varepsilon|= 0$.

\end{frame}


\begin{frame}
\frametitle{1.1.1. Алфавит, слово и язык (2/9)}
Можно рассматривать не только вхождения
отдельных букв в слово, но и 
вхождения любых частей слова как отдельные слова.
Операция взятия части слова называется 
\odmkey{вырезкой}{Операция вырезки}.
Вырезку естественно также обозначить 
квадратными скобками. 
\begin{example}
самовар$[3..6]= $ мова.
\end{example}

Таким образом, определено отношение 
между словами <<быть частью>>.
Это отношение рефлексивно, транзитивно и антисимметрично,
а потому является нестрогим частичным порядком,
и обозначается здесь $\subset$.
\begin{example}
луг $\subset$ слуга.
\end{example}

\begin{remark} Очевидно, что 
$\alpha\subset\beta\Impl|\alpha|\le|\beta|$,
но $\Not(|\alpha|\le|\beta|\Impl\alpha\subset\beta)$!
\end{remark} 
\begin{example}
гул $\not\subset$ слуга.
\end{example}

Для слов определена операция 
\odmkey{конкатенации}{Операция!конкатенации}, 
или \odmkey{сцепления}{Операция!сцепления}.
Обычно операция конкатенации никак не обозначается, а сцепляемые слова 
просто записываются подряд. 

\begin{example}
Конкатенация слов <<пара>> и <<план>> образует слово <<параплан>>.
\end{example}

Сцепление можно применять многократно,
таким образом определяется \odmkey{степень}{Степень!слова} слова. Положим
$\alpha^1 \Def \alpha; \alpha^k\Def\alpha\alpha^{k-1}; k>1$. Можно 
определить нулевую степень слова:
$\alpha^0 \Def \varepsilon$.

\begin{example}
$a^4=aaaa; (ab)^3= ababab$.
\end{example}

\end{frame}


\begin{frame}
\frametitle{1.1.1. Алфавит, слово и язык (3/9)}

Одни и те же буквы могут входить в разные слова.
Это порождает отношения между словами,
связанные с наличием или отсутствием общих частей.

Если $\alpha=\alpha_1\alpha_2$, то $\alpha_1$ называется
\odmkey{началом}{Начало слова}, или \odmkey{префиксом}{Префикс},
слова $\alpha$, а $\alpha_2$~---
\odmkey{окончанием}{Окончание слова}, или
\odmkey{постфиксом}{Постфикс}, слова $\alpha$.
Если при этом $\alpha_1 \neq \varepsilon\And \alpha_2\neq\varepsilon$, то
$\alpha_1$ (соответственно, $\alpha_2$)
называется \odmkey{собственным началом}{Начало слова!собственное}
(соответственно, \odmkey{собственным окончанием}{Окончание
слова!собственное})
слова $\alpha$.

\begin{example}
Слова <<д>> и <<до>> являются собственными префиксами, а слова <<ом>> и <<м>>
являются собственными постфиксами слова <<дом>>.
\end{example}

\begin{remark}
Каждое слово является своим постфиксом и префиксом (несобственным!).
\end{remark}

\odmkey{Пересечением}{Пересечение!слов} слов $\alpha$ и $\beta$ называется
наибольшее слово $\gamma$, являющееся одновременно постфиксом слова $\alpha$ и
префиксом слова $\beta$. Для любых слов $\alpha$ и $\beta$ такое 
слово $\gamma$ существует и единственно, а по нему
однозначно определяются слова $\delta$ и $\zeta$ из равенств
$\alpha=\delta\gamma$, $\beta=\gamma\zeta$.
Если слова $\alpha$ и $\beta$ не имеют общего постфикса и префикса,
то полагают, что $\gamma\Assign\varepsilon$, то есть пересечением
является пустое слово. Таким образом,
пересечение слов является всюду определённой 
некоммутативной (!) бинарной операцией.


\begin{example}
Пусть $\alpha=abb$, $\beta=bbc$. Тогда можно
взять $\gamma\Assign b$, $\delta\Assign ab$, $\zeta\Assign bc$.
Однако наибольшее пересечение даёт единственная тройка
$\gamma\Assign bb$, $\delta\Assign a$, $\zeta\Assign c$.
\end{example}

\end{frame}

\begin{frame}
\frametitle{1.1.1. Алфавит, слово и язык (4/9)}
Множество слов в алфавите
${A}$ обозначается~${A^*}$. 
По определению $\varepsilon \not\in A$, но $\varepsilon\in A^*$.
Если пустое слово исключается из рассмотрения, 
то множество слов в алфавите $A$ 
обозначается~$A^+$.

\smallskip
Произвольное множество $L$ слов в алфавите $A$ называется 
\odmkey{языком}{Язык (language)} в этом алфавите.

\begin{examples}
\begin{enumerate}
\item $A\subset A^*$ --- язык, состоящий из
однобуквенных слов;
\item $\{0,1\}^*$ --- язык двоичных последовательностей;
\item $A^*\subset A^*$ --- универсальный
язык всех слов.
\end{enumerate}
\end{examples}
 
К любым языкам применимы операции над множествами: объединение, пересечение,
разность и симметрическая разность.
Поскольку алфавит языка известен, то определена и
операция дополнения.

\begin{digress}
Обычные операции наивной теории множеств
играют сравнительно небольшую роль в 
теории языков, за исключением
операции объединения.
\end{digress}

\smallskip
Пусть есть два языка над одним алфавитом:
$L_1 \subset A^*, L_2\subset A^*$.  
Тогда операция сцепления слов индуцирует операцию 
сцепления языков: $L_1 L_2 
\Def \Set{\alpha\beta}{\alpha\in L_1 \And
\beta\in L_2}.$
\end{frame}

\begin{frame}
\frametitle{1.1.1. Алфавит, слово и язык (5/9)}
Сцепление языков можно применять многократно,
таким образом определяется \odmkey{степень}{Степень!слова} языка. Пусть $L\subset A^*$. Положим
$L^1 \Def L; L^k\Def L L^{k-1}; k>1$. Можно 
определить нулевую степень языка:
$L^0 \Def \{\varepsilon\}$.

Но сам алфавит $A$~--- это язык,
и значит определены степени этого языка $A^k$
для любого натурального числа $k$. 
Из этого наблюдения и определений
немедленно следует очевидная, 
но важная теорема.

\begin{theorem}
$A^k=\Set{\alpha\in A^*}{|\alpha|=k}$.
\end{theorem}

\begin{remark}
Другими словами, $A^k$~--- это множество слов длины $k$ в алфавите $A$.
\end{remark}

\begin{corollary}
$|A^k|=|A|^k =n^k.$
\end{corollary} 

Рассмотрим $\cal{A} = \{ A^0, A^1, A^2, \dots\}$~--- семейство классов эквивалентности
по отношению равенства длин слов. Тогда
$$A^* = \bigcup_{k=0}^\infty A^k, \qquad
A^+ = \bigcup_{k=1}^\infty A^k,$$
что вполне согласуется с предшествующими определениями.
\end{frame}

\begin{frame}
\frametitle{1.1.1. Алфавит, слово и язык (6/9)}
Множество $A^*$ бесконечно, но счётно, то есть 
его можно занумеровать. Рассмотрим некоторые нумерации, которые особенно часто используются
в других построениях.

Поскольку алфавит $A=\{\Tuple{a}{1}{n}\}$ конечен,
его буквы можно перенумеровать (то есть 
{\em линейно упорядочить}) и представить алфавит словом 
$\alpha\!\beta_A\Def a_1\dots a_{n}$,
то есть конкатенацией всех букв в порядке
возрастания номеров.


\begin{digress}
Слово <<алфавит>> происходит от прочтения
первых букв греческого алфавита:
альфа, бета, ... Аналогичное происхождение
имеет русское слово <<азбука>>: аз, буки, веди, ... \\Обозначение $\alpha\!\beta_A$ допустимо читать <<алфавит $A$>>.
\end{digress}

\begin{example}
Современная кириллица содержит 33 буквы,
и для этих букв узаконен следующий порядок:
$\alpha\!\beta_{\cal{K}}\Assign$ абвгдеёжзийклмнопрстуфхцчшщъыьэюя. 
Индекс $\cal{K}$ символизирует кириллицу.
\end{example}

Зафиксируем биекцию
$A=\{\Tuple{a}{1}{n}\} \sim 1..n$ следующим образом.

\vspace{-16pt}
\begin{columns}[t]
\begin{column}{10cm}
\begin{algorithmic}
\STATE \textbf{func $n_A(a : A) : 1..n$}
\FOR {\textbf{$i$ from $1$ to $n$}}
\STATE{\textbf{if $\alpha\beta_A[i]=a$ then return 
$i$ end if}}
\ENDFOR
\end{algorithmic}
\end{column}    
\begin{column}{5cm}
\begin{algorithmic}
\STATE \textbf{func $a_A(i : 1..n) : A$}
\STATE \textbf{return $\alpha\!\beta_A[i]$}
\end{algorithmic}
\end{column}
\end{columns}    


\end{frame}

\begin{frame}
\frametitle{1.1.1. Алфавит, слово и язык (7/9)}

Нумерация букв индуцирует нумерацию слов.
Пусть $\alpha \in A^k$. Положим 
$$n_A(\alpha)\Assign n_A(\alpha[k])
+ n \cdot n_A(\alpha[k-1]) +\dots
+n^{k-1}\cdot n_A(\alpha[1]) = 
\sum_{i=1}^{k} n_A(\alpha[i]) n^{k-i}. $$
%Тогда номером слова оказывается число,
%записанное в $n$-ричной системе счисления,
%причём младшие цифры записываются слева,
%а старшие --- справа ({\em little-endian}).
%

\vspace{-6pt}
\begin{example}
В этой нумерации для слов в кириллице имеем:
 
\begin{tabular}{l|l|l}
\hline
$n_{\cal{K}}(\text{а})=1$ & $n_{\cal{K}}(\text{аа})=1\cdot 33+1=~~34$ & $n_{\cal{K}}(\text{ааa})=1\cdot 33^2 +1\cdot 33 +1=~1123$  \\
$n_{\cal{K}}(\text{б})=2$ & $n_{\cal{K}}(\text{аб})=1\cdot 33+2=~~35$ & $n_{\cal{K}}(\text{ааб})=1\cdot 33^2+1\cdot 33+2=~1124$  \\
$\dots$ & $\dots$ & $\dots$ \\
$n_{\cal{K}}(\text{я})=33$ & $n_{\cal{K}}(\text{яя})=33\cdot 33+33=1122$ & $n_{\cal{K}}(\text{яяя})=33\cdot 33^2+ 33\cdot 33+33=37059$ \\
\hline
\end{tabular}
\end{example}

\medskip
\begin{theorem}
$\A{A}{A^*\sim \Bbb{N}_0}$.
\end{theorem}

\begin{proof}
Представленная нумерация реализует требуемую 
биекцию. Номер слова
легко вычисляется известным алгоритмом позиционной системы счисления, представленным
на следующем слайде слева.
Обратный алгоритм определения слова по номеру чуть сложнее, поскольку особым образом требуется
обрабатывать числа, кратные основанию $n$. Алгоритм представлен на следующем слайде справа.
Для завершения доказательства остаётся положить $n_A(\varepsilon)\Assign 0$.  
\end{proof}
 
\end{frame}

\begin{frame}
\frametitle{1.1.1. Алфавит, слово и язык (8/9)}

\vspace{-12pt}
\begin{columns}[t]
\begin{column}{5cm}
\begin{algorithmic}
\STATE{\textbf{func $n_A(\alpha : A^*) : 
\Bbb{N}_0$}}
\STATE{$i\Assign 0$}
\FOR {\textbf{$j$ from $1$ to $|\alpha|$ }}
\STATE{$i\Assign n\cdot i +n_A(\alpha[j])$}
\ENDFOR
\STATE{\textbf{return $i$}}
\end{algorithmic}
\end{column}    
\begin{column}{10cm}
\begin{algorithmic}
\STATE \textbf{func $a_A(i : \Bbb{N}_0) : A^*$}
\STATE{$\alpha\Assign \varepsilon$}
\WHILE {$i>0$}
\STATE\textbf{if $i\MOD n =0 $ then $a\Assign a_A(n); i\Assign i-n$}
\STATE\textbf{ else $a\Assign a_A(i\MOD n)$ end if}
\STATE{$\alpha\Assign a \alpha$; $i\Assign i\DIV n$}
\ENDWHILE
\STATE{\textbf{return $\alpha$}}
\end{algorithmic}
\end{column}
\end{columns}    

\medskip
\begin{remark}
Предложенная нумерация обладает рядом приятных свойств. \\Во-первых, нумерация монотонна по длине слова, то есть более длинные слова имеют номера больше, чем более короткие. 
\\Во-вторых, нумерация локально взаимно-однозначна:\\
$A\sim 1..n$, $A^2\sim (n+1)..(n^2+n)$ и вообще
$\A{k\in \Bbb{N}}{A^k \sim \frac{n^{k}-1}{n-1}.. n\frac{n^k-1}{n-1}}
$.\\
В-третьих, нумерация определяет естественный
линейный порядок на множестве слов:\\
$\alpha<\beta \Assign n_A(\alpha)<n_A(\beta)$,
что позволяет перечислить все слова в порядке возрастания.  
\end{remark}


\end{frame}

\begin{frame}
\frametitle{1.1.1. Алфавит, слово и язык (9/9)}
Часто на множестве слов используют 
\odmkey{лексикографический порядок}{Порядок!лексикографический} (или \odmkey{алфавитный}{Порядок!алфавитный}):\\
$\alpha<\beta\Def \beta=\alpha\gamma\And\gamma\neq\varepsilon 
\Or \E{j\le|\alpha|}{\alpha[j]<
\beta[j]\And\A{i \in 1..(j-1)}{\alpha[i]=
\beta[i]}}.$ 
Лексикогра\-фи\-ческий порядок линеен, так же как и введённый естественный порядок.  

\begin{remark}
Лексикографический порядок не позволяет построить генератор, который перечисляет
{\em все} слова в возрастающем порядке. 
\end{remark}

\begin{example}
В естественном порядке 
$ 1 < 2 <11 <12< 21 $, 
а в лексикографическом порядке 
$1<11<12<2<21$. 
\end{example}

\begin{digress}
Нумерация $n_A()$ очень похожа на 
обычную позиционную
систему счисления с основанием $n$. 
Отличие в том, что в позиционных системах
цифры берутся из диапазона $0..(n-1)$,
а в предложенной системе цифры берутся из
диапазона $1..n$. Наличие цифры $0$ с её 
особыми арифметическими свойствами влечёт два важных следствия. С одной стороны,
алгоритмы 
всех арифметических операций в обычной 
системе счисления
оказываются проще, чем алгоритмы 
в системе без нуля,
поэтому в компьютерах при вычислениях всегда применяются позиционные системы с нулём.
С другой стороны, из-за нуля
теряется однозначность соответствия
между словами, составленными из цифр, 
и натуральными числами. Проблему
создают \odmkey{ведущие нули}{Нуль!ведущий}~---
различные строки $0001$ и $1$ в позиционной системе соответствуют одному и тому же
числу <<один>>.
\end{digress}

\end{frame}

\subsubsection{1.1.2. Формальные грамматики}\label{1.1.2.}
\begin{frame}
\frametitle{\Ok 1.1.2. Формальные грамматики (1/4)}

Если язык содержит только конечное множество
слов, то их можно задать просто перечислением.
Если же язык бесконечен, что чаще встречается,
то необходимо использовать специальные средства
для построения конечного описания бесконечного множества.
Возникновение и развитие самого известного из таких средств 
связывают с Ноамом Хомским\footnote{Avram Noam Chomsky, р. 1928}.
  
\smallskip 
\odmkey{Формальная порождающая грамматика}{Грамматика!формальная}
\odmkey{}{Грамматика!порождающая} (далее просто грамматика)
$G=\langle \mathrm{M}, \Sigma, \Pi, S\rangle$~--- это совокупность конечного
алфавита $\mathrm{M}$ \odmkey{нетерминальных}{Символ!нетерминальный}
символов (то есть вспомогательных символов, не входящих в алфавит, над которым определяется язык), конечного алфавита $\Sigma$
\odmkey{терминальных}{Символ!терминальный} символов (то есть  символов алфавита, над которым определяется язык, 
$\mathrm{M}\cap \Sigma = \emptyset$), конечного множества \odmkey{правил вывода}{Правило!вывода}
(правила вывода называют также \odmkey{правилами подстановки}{Правило!подстановки}
или \odmkey{продукциями}{Продукция})
$\Pi=\Set{\alpha\rightarrow\beta}{\alpha,\beta\in (\mathrm{M}\cup \Sigma)^*}$ и
\odmkey{основного символа}{Символ!основной} (или \odmkey{аксиомы
грамматики}{Аксиома!грамматики}) $S \in \mathrm{M}$.


\begin{example}
Совокупность множеств 
$\mathrm{M}\Assign \{N\}$,
$\Sigma\Assign\{a,b\}$, \\
$\Pi\Assign\{N\to NN, N\to \varepsilon, N\to a,
N\to b\}$  образует грамматику  
$G_0=\langle \mathrm{M},\Sigma, \Pi, N\rangle$.
\end{example}

\end{frame}

\begin{frame}
\frametitle{\Ok 1.1.2. Формальные грамматики (2/4)}

\begin{remark}
По традиции нетерминальные символы (или просто
\odmkey{нетерминалы}{Нетерминал}, 
иногда их также называют \odmkey{переменными}{Переменная}) обозначают прописными латинскими
буквами, терминальные символы (или просто
\odmkey{терминалы}{Терминал}) обозначают строчными латинскими
буквами, а цепочки терминальных и нетерминальных символов~---
строчными греческими буквами.
Было бы неудобно {\em множества} символов, то есть алфавиты, 
обозначать символами этих же алфавитов.
Поэтому здесь приняты обозначения алфавитов и множества правил прописными
греческими буквами, по первым буквам греческих слов
$\mathrm{M}\varepsilon\tau\alpha\beta\lambda\eta\tau\eta$ (Переменная),
$\Sigma\tau\alpha\theta o\rho o\varsigma$ (Постоянная),
$\Pi\alpha\rho\alpha\gamma\omega\gamma\iota\kappa o\varsigma
\ \nu o\mu o\varsigma$ (Производящее правило).
\end{remark}

\medskip
Поскольку все множества конечны, 
запись грамматики также конечна,
однако конечная грамматика может порождать
бесконечный язык на основе 
формального \odmkey{вывода}{Вывод}.

\medskip   
Говорят, что из цепочки $\gamma$ \odmkey{непосредственно
выводима}{Выводимость!непосредственная} цепочка $\delta$ в
грамматике $G$ (обозначение $\gamma\Rightarrow_G \delta$), если
$\gamma=\gamma_1\alpha\gamma_2$, $\delta=\gamma_1\beta\gamma_2$ и
$\alpha\rightarrow\beta\in \Pi$. Говорят, что из цепочки $\gamma$
\odmkey{выводима}{Выводимость} цепочка $\delta$ в грамматике $G$,
если существует \odmkey{вывод}{Вывод}, т.е. последовательность цепочек, в которой каждая следующая выводима из предыдущей по одному из правил вывода: 
$$\E{\Tuple{\sigma}{0}{n}}
{\A{i\in 0..(n-1)}{\sigma_{i}\Rightarrow_G\sigma_{i+1} \And \sigma_0 = \gamma
\And \sigma_n = \delta}}.$$ 
\end{frame}

\begin{frame}
\frametitle{\Ok 1.1.2. Формальные грамматики (3/4)}

Число $n$ (количество цепочек в
последовательности, не считая первой) называется
\odmkey{длиной}{Длина!вывода} вывода. Выводимость $\delta$ из
$\gamma$ обозначают 
$\gamma\Rightarrow^n_G\delta$, если нужно
указать длину вывода, или $\gamma\Rightarrow^*_G\delta$, если длина
вывода не указывается. Если грамматика подразумевается, то индекс
$G$ опускают.
Отметим, что $\Rightarrow_G$ является
отношением на множестве цепочек.
В таком случае отношение
 $\Rightarrow_{G}^{*}$ является рефлексивно-транзитивным замыканием 
 отношения $\Rightarrow_{G}$ на множестве цепочек.

\smallskip
\begin{example}
Последовательность формул $NN, NNN,  NNNN, aNNN, abNN, abbN, abba $ является выводом в грамматике $G_0$.
Другими словами, $NN\Rightarrow^6_{G_0} abba$.
\end{example}

\smallskip
\begin{remark}
Таким образом, грамматика~--- 
это {\em формальная теория}, или же
{\em исчисление}, в том строгом значении,
которое имеют эти термины в математической логике.
Теоремами этой теории являются терминальные цепочки, выводимые из аксиомы грамматики.   
\end{remark}

\smallskip
\begin{example}
В грамматике $G_0$ предыдущего примера
все слова из множества $\{a,b\}^*$, включая пустое слово, выводимы
из начального символа грамматики.
Другими словами, грамматика $G_0$~---
это такая малоинтересная теория,
в которой все формулы суть теоремы.   
\end{example}
\end{frame}

\begin{frame}
\frametitle{\Ok 1.1.2. Формальные грамматики (4/4)}
Вывод называется \odmkey{бесповторным}{Вывод!бесповторный}, если
все цепочки в нем различны.

\begin{lemma}
Если существует вывод слова $\alpha$ в грамматике $G$, то существует бесповторный
вывод слова $\alpha$ в грамматике $G$.
\end{lemma}
\begin{proof}
Пусть
$\sigma_0=S,\sigma_1,\dots,\sigma_i,\dots,\sigma_j,\dots\sigma_n=\alpha$~---
вывод цепочки $\alpha$ в $G$, причём $\sigma_i = \sigma_j$. Очевидно, что
$\sigma_0=S,\sigma_1,\dots,\sigma_i,\sigma_{j+1},\dots\sigma_n=\alpha$
также будет выводом цепочки $\alpha$ в $G$.
\end{proof}

\smallskip
\odmkey{Языком}{Язык!порождаемый грамматикой}, порождаемым
грамматикой $G$ (обозначение $\frak{L}(G)$), называется множество всех
цепочек в терминальном алфавите, выводимых из начального символа грамматики:
$$
\frak{L}(G)\Def \Set{\alpha\in \Sigma^*}{S\Rightarrow_G^*\alpha}.
$$
\begin{remark}
В традициях терминологии логических исчислений говорят, что слова языка $\frak{L}(G)$ \odmkey{допускаются}{Слово!допустимое} грамматикой $G$.
\end{remark}

Один и тот же язык может быть порождён
разными грамматиками.
Говорят, что две грамматики $G$ и $G'$
\odmkey{эквивалентны}{Эквивалентность!грамматик} (обозначение
$G\sim G'$), если $\frak{L}(G)=\frak{L}(G')$.

\begin{example}
Пусть $\mathrm{M}\Assign \{N\}$,
$\Sigma\Assign\{a,b\}$,
${\Pi}'\Assign\{N\to \varepsilon, N\to aN,
N\to bN\}$, тогда грамматика  
${G_0}'=\langle \mathrm{M}, \Sigma, {\Pi}', N\rangle$ эквивалентна $G_0$,
то есть $\frak{L}(G_0)= \frak{L}({G_0}')$.
\end{example}
\end{frame}


% GA_1_2025.tex, 928–2269
\subsubsection{1.1.4. Классификация грамматик и языков}
\label{1.1.4.}


\begin{frame}
\frametitle{1.1.4. Классификация грамматик и языков (1/6)}

\begin{enumerate}
\item Грамматика \odmkey{типа 0}{Грамматика!типа 0}
(или грамматика \odmkey{с фразовой структурой}{Грамматика!с
фразовой структурой})~--- это грамматика общего вида без
ограничений на правила вывода. Класс таких грамматик обозначим
$\frak{T}_0$.
\item Грамматика \odmkey{типа 1}{Грамматика!типа 1}
(или грамматика \odmkey{непосредственных
составляющих}{Грамматика!непосредственных составляющих} или
\odmkey{контекст\-но-зависимая}{Грамматика!контекстно-зависимая})
~--- это грамматика, все правила которой имеют вид $\alpha
A\beta\rightarrow\alpha\gamma\beta$, где $\alpha,\beta, \gamma\in (\mathrm{M}\cup
\Sigma)^*$, $|\gamma|\geq 1$, $A\in \mathrm{M}$. Такое правило описывает замену
нетерминального символа $A$ на цепочку $\gamma$ в определенном
\odmkey{контексте}{Контекст!правила} ($\alpha,\beta$), который не
меняется при применении правила. Класс таких грамматик обозначим
$\frak{T}_1$.
\item Грамматика \odmkey{типа 2}{Грамматика!типа 2}
(или
\odmkey{контекстно-свободная}{Грамматика!контекстно-свободная}
грамматика)~---  это грамматика, все правила которой имеют вид
$A\rightarrow\alpha$, где $\alpha\in (\mathrm{M}\cup \Sigma)^*$, $A\in \mathrm{M}$. Такое
правило описывает замену нетерминального символа $A$ на цепочку
$\alpha$ вне зависимости от контекста. Класс таких грамматик
обозначим $\frak{T}_2$.
\item Грамматика \odmkey{типа 3}{Грамматика!типа 3}
(или \odmkey{регулярная}{Грамматика!регулярная} грамматика) ~---
это грамматика, все правила которой имеют вид либо
$A\rightarrow\alpha$, либо $A\rightarrow\alpha B$ (или
$A\rightarrow B\alpha$), где $\alpha\in \Sigma^*$,\\ $A,B\in \mathrm{M}$.
В правой части правила регулярной грамматики может быть не более
одного нетерминала.
Класс таких грамматик обозначим
$\frak{T}_3$.
\end{enumerate}
\end{frame}

\begin{frame}
\frametitle{1.1.4. Классификация грамматик и языков (2/6)}
Ясно, что $\A{i\in 1..3}{G\in \frak{T}_i\Impl G\in \frak{T}_{i-1}}$, поскольку всякое правило
грамматики класса $\frak{T}_i$ 
является частным случаем 
правила грамматики класса $\frak{T}_{i-1}$.

\begin{remark}
Исключением являются правила с пустой правой частью,
но их можно устранить, см. \DMref{п.~1.3.2}{1.3.2.}.
\end{remark}
\medskip
Один и тот же язык может быть описан разными грамматиками, в том
числе грамматиками разного типа. Принято относить язык к
{\em наибольшему} типу из типов описывающих его грамматик. То есть если
для языка $L$ есть две описывающие его грамматики $G'$ и $G''$,
$L=\frak{L}(G')\And L=\frak{L}(G'')\And G'\in \frak{T}_i\And G''\in \frak{T}_j\And i<j$, то
язык $L$ относится к типу $\frak{T}_j$.
\begin{example}
Рассмотрим грамматику
$G=\langle\{S,N,D\},\{0,1,2,3,4,5,6,7,8,9,-,+\},
\linebreak
\{S\rightarrow
N, S\rightarrow
+N, S\rightarrow
-N, N\rightarrow D, N\rightarrow ND, D\rightarrow 0, D\rightarrow 1,\ldots,D\rightarrow 9 \}, S\rangle$, описывающую множество десятичных целых чисел со знаком.
Эта грамматика относится к типу 2 (правило $N\rightarrow ND$
недопустимо в грамматике типа 3). Однако для того же самого языка
можно предложить грамматику
$G'=\langle\{S,N\},\{0,1,2,3,4,5,6,7,8,9,-,+\},
\linebreak
\{S\rightarrow N, 
S\rightarrow +N, 
S\rightarrow -N, 
N\rightarrow 0,
N\rightarrow 1,
\ldots,
N\rightarrow 9,
N\rightarrow 0N,
N\rightarrow 1N,
\ldots,\linebreak
N\rightarrow 9N
\}, S \rangle$,
которая является грамматикой типа 3. Таким образом, язык целых
десятичных чисел со знаком является языком типа 3.
\end{example}

\end{frame}

\begin{frame}
\frametitle{1.1.4. Классификация грамматик и языков (3/6)}

Вопрос об определении типа конкретной грамматики решается очень просто~---
рассмотрением вида правил этой грамматики. К сожалению, вопрос об определении
типа языка оказывается алгоритмически неразрешимым.

\begin{theorem}
Следующая проблема является алгоритмически неразрешимой: дано $L=\frak{L}(G)$,
$G\in\frak{T}_i$, определить, существует ли $G'$, такая что
$L=\frak{L}(G')\And G'\in\frak{T}_j\And i<j$.
\end{theorem}

\begin{proof}
Разбирается основной случай $i=0$: нельзя определить, относится ли язык
$\frak{L}(G)$ к типу $j$, $j\in 1..3$. Доказательство ведётся в обозначениях
наивной теории алгоритмов (п.~\DMref{0.3.3}{0.3.3}) и опирается на четыре факта.
\TODO{ вынести в леммы}


\begin{enumerate}
\item Множество грамматик типа $j$ разрешимо: тип виден по виду правил.\TODO{= было}
\item Грамматика типа 0 и перечисляющий алгоритм~--- два способа записи одного и
того же, и оба перехода алгоритмичны (п.~\DMref{1.1.3}{1.1.3.}). \TODO{= бред}
\item Всякий язык типа 1, 2 или 3 разрешим (п.~\DMref{1.1.5}{1.1.5.}; для типов
2 и 3 сначала устраняются $\varepsilon$-правила). \TODO{= доказать явно в 1.1.6 и 1.1.7.}
\item Существует перечислимое неразрешимое множество $N\subset\mathbb{N}$ и вычислимая
функция $d_1$ с областью определения $N$ \TODO{= какая это функция?}.
\TODO{=достать из 0.3.3. и в явном виде доказать.}

\end{enumerate}
\end{proof}
\end{frame}

\begin{frame}
\frametitle{1.1.4. Классификация грамматик и языков (4/6)}

\begin{proof} (Продолжение)
\proofsection{ Шаг 1. Неразрешимый язык.} Пусть $\Sigma\Assign\{0,1\}$ и
$W\Assign\Set{\text{двоичная запись числа } n}{n\in N}\subset\Sigma^*.$
Язык $W$ перечислим как образ перечислимого множества при вычислимой функции.
Язык $W$ неразрешим: разрешающий его алгоритм разрешал бы и множество $N$,
потому что двоичная запись и её обращение вычислимы. Значит, по факту 3 язык $W$
не является языком типа $\frak{T}_1$, а тем более типов $\frak{T}_2$ и $\frak{T}_3$.

\proofsection{ Шаг 2. Грамматика по номеру.} Рассмотрим алгоритм с одним параметром:
\vspace{-4pt}
\begin{algorithmic}
\STATE {$\textbf{proc } \mathrm{C}\,(n:\mathbb{N})$}
\STATE {$\;\; x\Assign d_1(n)$}
\COMMENT{ при $n\not\in N$ не кончается}
\STATE {$\;\; \textbf{for } a\in W \textbf{ do yield } a \textbf{ end for}$}
\COMMENT{ перечислить $W$}
\STATE {$\textbf{end proc}$}
\end{algorithmic}
\vspace{-4pt}

Остаточная программа $\mathrm{C}_n$ (ДМ, п.~4.5.6 \TODO{Надобно перетаскивать нужное, ссылаться вовне нельзя. Можно без доказательств.})~--- это алгоритм без
аргументов, перечисляющий множество $W$ при $n\in N$ и множество $\emptyset$
при $n\not\in N$. По факту 2 построим по ней грамматику $G_n\in\frak{T}_0$,
порождающую то же множество. Переход \TODO{пешеходный?} $n\mapsto G_n$ вычислим.
\end{proof}
\end{frame}

\begin{frame}
\frametitle{1.1.4. Классификация грамматик и языков (5/6)}

\begin{proof} (Окончание)
\proofsection{Шаг 3. Сведение.} Пусть, вопреки доказываемому, существует алгоритм
$\mathrm{R}$, который по грамматике $G\in\frak{T}_0$ определяет, существует ли
грамматика $G'\in\frak{T}_1$ с $\frak{L}(G')=\frak{L}(G)$. Тогда множество $N$
разрешает алгоритм
$\mathrm{N}(n:\Bbb{N}) : \textbf{bool} \Def
\textbf{return } \Not\mathrm{R}(G_n) $
Действительно, если $n\in N$, то $\frak{L}(G_n)=W$~--- язык неразрешимый, а
потому не язык типа $\frak{T}_1$, и $\mathrm{R}$ даёт $\textbf{false}$. 
Если же $n\not\in N$, то
$\frak{L}(G_n)=\emptyset$, а пустой язык порождается грамматикой
$\langle\{S\},\Sigma,\{S\to S\},S\rangle$, которая относится сразу ко всем
четырём типам, и $\mathrm{R}$ даёт $\textbf{true}$. Значит, множество $N$ разрешимо, что
противоречит факту 4. Для $j=2$ и $j=3$ рассуждение дословно то же: язык $W$ не
относится и к этим типам.
\end{proof}

\begin{remark}
Доказанное~--- частный случай теоремы Райса: \TODO{сноска ученого} всякое нетривиальное свойство
перечислимого множества неразрешимо по программе, задающей это множество. Здесь
свойство~--- принадлежность языка типу $j$, программа~--- грамматика типа 0, а
нетривиальность в том, что пустой язык этим свойством обладает, а $W$~--- нет.
\TODO{Замечание наиважнейшее, но сформулировано на уровне ИИ, то есть безграмотно. Выправить!}
\end{remark}
\end{frame}


%\begin{frame}
%\frametitle{1.1.4. Классификация грамматик и языков (3/4)}
%
%Вопрос об определении типа конкретной грамматики решается очень
%просто~--- рассмотрением вида правил этой грамматики. К сожалению,
%вопрос об определении типа языка оказывается алгоритмически
%неразрешимым.
%\begin{theorem} 
%Следующая проблема является алгоритмически неразрешимой: дано
%$L=\frak{L}(G)$, $G\in \frak{T}_i$, определить, существует ли $G'$,
%такая что $L=\frak{L}(G') \And
%G'\in \frak{T}_j\And i<j $.
%\end{theorem}
%\begin{proof} \TODO{Теорема Райса-Хиггинсона. 
%Доказательство обязательно в обозначениях раздела 4.5 из Дискретной математики для программистов. Целесообразно добавить атом про проблему останова во Введении. 
%Студенты, которые слушали курс Пастора по теории алгоритмов, могут помочь.}
%\end{proof}
%\end{frame}

\begin{frame}
\frametitle{1.1.4. Классификация грамматик и языков (6/6)}

\begin{digress}
Чем выше тип грамматики, 
порождающий язык, тем лучше ---
язык оказывается проще,
его легче анализировать,
описывать и применять.
Поэтому авторы новых
формальных языков стараются
использовать такие грамматические конструкции,
которые выражаются правилами
регулярных грамматик,
в крайнем случае контекстно-свободных грамматик.
Однако иногда приходится иметь дело
с унаследованными языками, 
не только естественными,
но и искусственными, в частности,
с ранними языками программирования.
В таких случаях ради совместимости
приходится жертвовать удобством.  
\end{digress}

\begin{example}
В языке Algol-60 была 
впервые официально использована теория 
формальных грамматик для описания стандарта языка.
При этом принятая  грамматика не позволяла 
определить, к которому $\textbf{ if }$ 
относится $\textbf{ else }$ 
в условном операторе \\ 
$\textbf{ if }
B_1 \textbf{ then if } B_2 \textbf{ then }
S_1 \textbf{ else } S_2$.\\ 
Язык оказался существенно 
синтаксически неоднозначным.
\end{example}

\TODO{Замечание про <<висячий>> \textbf{else} очень нужное, возможно не здесь, а алгоритмических свойствах грамматик. Можно развить эту тему, сказав про \textbf{end if}. }

\end{frame}


\subsubsection{1.1.5. Неукорачивающие грамматики}
\label{1.1.5.}
\begin{frame}
\frametitle{\Ok 1.1.5. Неукорачивающие грамматики (1/5)}

Грамматика $G=\langle \mathrm{M}, \Sigma, \Pi, S\rangle$ называется
\odmkey{неукорачивающей}{Грамматика!неукорачивающая}, если
длина левой части любого правила не превосходит
длины правой части:
$\A{r=\alpha\rightarrow\beta\in \Pi}{|\alpha|\leq|\beta|}$. При
любом выводе в такой грамматике длины цепочек могут разве что
возрастать. Поэтому, в частности, пустое слово $\varepsilon$ не
выводимо в такой грамматике.
\begin{example}
Контекстно-зависимые грамматики являются неукорачивающими.
\end{example}

Язык $L\subset \Sigma^*$ называется
\odmkey{разрешимым}{Язык!разрешимый}
(по аналогии с понятием {\em разрешимое множество}), если существует
алгоритм, который для произвольной цепочки $\alpha\in \Sigma^*$ за
конечное число шагов даёт ответ на вопрос, входит ли эта цепочка в
язык, т.~е. $\alpha\in L$ или $\alpha\not\in L$. Если при этом число шагов алгоритма определяется по длине цепочки
$\alpha$ и может быть указано заранее, то язык называется
\odmkey{легко разрешимым}{Язык!легко разрешимый}.
Если же алгоритм даёт ответ, когда \\$\alpha \in L$, и может не давать ответа (алгоритм не заканчивает работу),
когда $\alpha\not\in L$, то язык называется
\odmkey{полуразрешимым}{Язык!полуразрешимый}.
Алгоритм, разрешающий язык, называют  \odmkey{распознающим}{Язык!распознаваемый}. 

\begin{theorem}
Язык $\frak{L}(G)$, порождаемый неукорачивающей грамматикой
 $G=\langle \mathrm{M}, \Sigma, \Pi, S\rangle$, \\легко разрешим.
\end{theorem}
\begin{proof}
Если $\alpha\in \frak{L}(G)$, то $S\Rightarrow^*\alpha$, существует бесповторный вывод $\Tuple{\sigma}{0}{k}$,
где $\sigma_0 = S$ и
$\sigma_k=\alpha$, причём 
 $\A{i\in 0..k}{|\sigma_i|\leq|\alpha|}$. Обозначим $n\Assign|\alpha|$.
%Если же \\ $\alpha\not\in \frak{L}(G)$, то $\Not
%S\Rightarrow^*\alpha$.
\end{proof}
\end{frame}

\begin{frame}
\frametitle{\Ok 1.1.5. Неукорачивающие грамматики (2/5)}

\begin{proof} (Окончание)
Число $r(n)$  различных цепочек длиной не более
$n$:

\vspace{-12pt}
$$
r(n) = \sum_{i=1}^{n} \left|\mathrm{M}\cup \Sigma\right|^i
=\sum_{i=1}^n (|\mathrm{M}|+|\Sigma|)^i=
(|\mathrm{M}|+|\Sigma|)\frac{(|\mathrm{M}|+|\Sigma|)^n -1}{|\mathrm{M}|+|\Sigma| -1}.
$$

\vspace{-4pt}
Далее, длина бесповторной последовательности, составленной из
таких цепочек, также ограничена (например, величиной $r(n)$). Общее
количество $R(n)$ всех бесповторных последовательностей цепочек длиной не
более $n$ также ограничено:

\vspace{-16pt}
 $$R(n)=\sum_{i=0}^{r(n)} A(r(n),i) =
 \sum_{i=0}^{r(n)} \frac{r(n)!}{(r(n)-i)!}=
  r(n)!\sum_{i=0}^{r(n)}\frac{1}{i!}
  \leq r(n)!\sum_{i=0}^\infty\frac{1}{i!}
  =r(n)!\cdot e.$$ 

\vspace{-4pt}
Тем
более ограничено количество таких бесповторных
последовательностей, которые начинаются с $S$ и заканчиваются
$\alpha$.
Отсюда немедленно вытекает схема алгоритма распознавания. Нужно
перебрать все бесповторные последовательности цепочек длиной не
более $n$, которые начинаются с $S$ и заканчиваются $\alpha$, и для
каждой последовательности проверить, не является ли она допустимым
выводом $\alpha$ в $G$. 
%(то есть проверить, что каждый следующий
%элемент последовательности непосредственно выводим из
%предыдущего). 
Если такой вывод найдётся, то можно прекратить
дальнейший перебор и сделать вывод, что $\alpha\in \frak{L}(G)$. В
противном случае $\alpha\not\in \frak{L}(G)$.
\end{proof}
\end{frame}

\begin{frame}
\frametitle{\Ok 1.1.5. Неукорачивающие грамматики (3/5)}
\begin{remark}
Число шагов перебора,
о котором идёт речь в доказательстве,
очень велико.
Это не случайно.
Можно показать, что в общем случае не существует эффективного
алгоритма распознавания языка, порождённого неукорачивающей
грамматикой. Другими словами, хотя число шагов работы алгоритма
распознавания можно предсказать заранее, но это число
может быть столь велико, что алгоритм практически не применим.
\end{remark}


\begin{digress}
Грамматически можно
описать любое перечислимое множество, в том числе и неразрешимое.
Поскольку грамматики типа $\frak{T}_1$ описывают распознаваемые, то есть
разрешимые языки, ясно что описание нераспознаваемых языков
является уделом грамматик типа $\frak{T}_0$. Хотя грамматики типа $\frak{T}_1$
удобны для описания языков (контекст очень полезен)
и порождают распознаваемые языки,
но алгоритмы распознавания имеют слишком
большую трудоёмкость. Поэтому грамматики
типа $\frak{T}_0$ и $\frak{T}_1$ {\em никогда} не применяются для описания
искусственных языков программирования.
В информационных технологиях
применяются грамматики типа $\frak{T}_2$ и $\frak{T}_3$.
Лингвисты иногда применяют грамматики
типа $\frak{T}_0$ и $\frak{T}_1$ для описания естественных языков,
но это только потому, что они не имеют права
изменять язык в угоду удобству описания,
а программисты могут и должны это делать.
Современная тенденция в языкотворчестве
ведет к исключению из формальных языков конструкций,
которые невозможно описать грамматиками типа $\frak{T}_2$ и $\frak{T}_3$.
\end{digress}

\end{frame}

\begin{frame}
\frametitle{\Ok 1.1.5. Неукорачивающие грамматики (4/5)}

\begin{remark}
Любую грамматику можно доопределить так, 
чтобы все символы, входящие в исходные правила, были нетерминалами.
Для этого достаточно
для каждого терминального символа
$a$ ввести новый нетерминальный символ
$A$ и правило $A\to a$. После этого
нужно все вхождения терминальных
символов
в старые правила заменить 
новыми нетерминальными символами.
\end{remark}



\begin{theorem}
Для всякой неукорачивающей грамматики существует эквивалентная ей \\контекстно-зависимая грамматика.
\end{theorem}

\begin{proof}
Построим правила для контекстно-зависимой
грамматики $G_2$ на основе правил
неукорачивающей грамматики $G_1$.
Каждое правило 
$$X_1 X_2 \ldots X_n \to Y_1 Y_2 \ldots Y_{n+m} \in \Pi,\quad n > 0, m \geq 0$$
заменим на ряд из $2n$ следующих правил:
$X_1 X_2 \dots X_n \rightarrow Z_1 X_2 \dots X_n$,
\\
$Z_1 X_2 \dots X_n \rightarrow Z_1 Z_2 \dots X_n$,
$\ \dots$,
$\ Z_1 Z_2 \dots X_n \rightarrow Z_1 Z_2 \dots Z_n$,
и \\
$Z_1 Z_2 \dots Z_n \rightarrow Y_1 Z_2 \dots Z_n$, $\ Y_1 Z_2 \dots Z_n \rightarrow Y_1 Y_2 \dots Z_n$,
$\ \dots$,\\
$Y_1 Y_2 \dots Z_n \rightarrow Y_1 Y_2 \dots Y_n \dots Y_{n+m}$,
где $Z_1, Z_2, \dots, Z_n$~--- новые нетерминалы, 
для каждого исходного правила они свои.
\end{proof}
\end{frame}

\begin{frame}
\frametitle{\Ok 1.1.5. Неукорачивающие грамматики (5/5)}
\begin{proof} (продолжение)
Правила с новыми нетерминалами могут применяться только в одной указанной последовательности. 
Другими словами, если будет применено правило 
$X_1 X_2 \dots X_n \rightarrow Z_1 X_2 \dots X_n$, то все последующие правила тоже будут применены, потому что иначе нетерминалы $Z_1, Z_2, ... Z_n$ будут присутствовать в выведенном слове, чего быть не должно.
Построенная грамматика эквивалентна исходной, потому что в результате применения новых правил 
цепочка $X_1 X_2 ... X_n$ перейдет в $Y_1 Y_2 ... Y_{n+m}$. 
Новая грамматика является контекстно-зависимой грамматикой по построению.
\end{proof}

\begin{remark}
Неукорачивающие грамматики,
очевидно, не могут порождать
языки, содержащие пустое слово $\varepsilon$.
В большинстве случаев это обстоятельство
не является существенным.
Если, однако, требуется,
чтобы пустое слово входило
в язык, то достаточно
в неукорачивающей 
грамматике разрешить иметь
одно исключительное правило:
$S\to\varepsilon$. 
\end{remark}


\medskip
Из доказанной теоремы и предшествующих наблюдений
следует важный вывод:
классы языков, порождаемых неукорачивающими грамматиками и грамматиками непосредственных составляющих, совпадают. 


\begin{remark}
Для краткости оборот <<контекстно-свободная грамматика>> далее заменяется аббревиатурой КСГ. 
\end{remark}
\end{frame}

\subsubsection{1.1.6. Контекстно-свободные грамматики}\label{1.1.6.}
\begin{frame}
\frametitle{\Ok 1.1.6. Контекстно-свободные грамматики (1/6)}

Пусть задана КСГ $G=\langle \mathrm{M}, \Sigma, \Pi, S\rangle$
и вывод $S\Rightarrow_G^*\omega \in \Sigma^*$
в этой грамматике. Поскольку все правила
имеют вид $A\rightarrow \alpha$, то в выводе
$S=\sigma_0,\sigma_1,\dots,\sigma_n=\omega$
для каждой пары $\sigma_i,\sigma_{i+1}$
соседних цепочек
следующая получается из предыдущей заменой
{\em одного} нетерминального символа.
Вообще говоря, возможно заменить
любой нетерминальный символ,
поэтому последовательность цепочек в выводе 
отнюдь не является однозначно определённой.
Вывод называется
\odmkey{левосторонним}{Вывод!левосторонний}
(\odmkey{правосторонним}{Вывод!правосторонний}),
если на каждом шаге заменяется самое левое
(соответственно, самое правое)
вхождение нетерминального символа в цепочку.

\smallskip
Вывод в КСГ
удобно представить в виде
\odmkey{дерева разбора}{Дерево!разбора}
(или \odmkey{дерева вывода}{Дерево!вывода})~---
упорядоченного ордерева,
корнем которого является аксиома грамматики,
нелистовыми узлами являются нетерминалы,
а листьями~--- терминальные символы цепочки $\omega$. 
Если в выводе было применено
правило  $A\rightarrow \alpha$, то в дереве разбора проводятся
дуги из узла $A$ ко всем узлам, помеченным буквами
цепочки $\alpha$, причем эти узлы располагаются
в том же порядке, что и буквы в цепочке $\alpha$.
Если же было применено правило 
$A\rightarrow \varepsilon$, то можно считать,
что проводится дуга от узла $A$ к листу, помеченному $\varepsilon$, 
или же, не ограничивая общности, можно считать, 
что в дереве разбора в этом случае никаких дуг не проводится. 
\end{frame}


\begin{frame}
\frametitle{\Ok 1.1.6. Контекстно-свободные грамматики (2/6)}

\begin{columns}[t]
\begin{column}{10.1cm}
Дерево вывода компактно
представляет множество эквивалентных выводов,
в которых применены одни и те же правила, но в разном порядке.
\begin{example}
Рассмотрим КСГ $G_1$ с правилами
$E\rightarrow E+T$,
$E\rightarrow T$,
$T\rightarrow T*F$,
$T\rightarrow F$,
$F\rightarrow a$~--- упрощенную грамматику арифметических выражений.
Дерево разбора справа в компактной форме представляет 10 (десять!) возможных выводов
цепочки $a+a$.
\end{example}

Ясно, что по дереву вывода можно однозначно
восстановить левосторонний (равно как и правосторонний)
вывод в данной грамматике.

\end{column}
\begin{column}{5cm}
\vspace{-20pt}
\begin{figure}[h]
\begin{center}
\includegraphics[height=50mm]{1_1_6_2}
\end{center}
\end{figure}
\end{column}
\end{columns}

\vspace{4pt}
Ещё одной наглядной формой 
представления выводов в КСГ является 
\odmkey{дерево порождения}{Дерево!порождения}.
Дерево порождения строится следующим образом.
Корнем дерева является узел, помеченный аксиомой грамматики 
$S$.
Если в грамматике есть правило $A\to\alpha$,
и в дереве построен узел, 
помеченный цепочкой $\beta A \gamma$,
то добавляется узел, помеченный цепочкой
$\beta\alpha\gamma$, и проводится дуга 
$(\beta A \gamma,\beta\alpha\gamma)$,
помеченная правилом $A\to\alpha$.

\end{frame}

\label{1.1.6.3.}
\begin{frame}
\frametitle{\Ok 1.1.6.Контекстно-свободные грамматики (3/5)}

\begin{example}
На рисунке представлено дерево порождения,
в котором все десять выводов предыдущего примера
показаны в явной форме. 
\end{example}
\begin{figure}[h]
\begin{center} 
\end{center}
\end{figure}

\end{frame}



% редакция Алексей Архиповский 18.09.26 исправлена ФАН 6.10.26
\begin{frame}


\frametitle{\Ok 1.1.6. Контекстно-свободные грамматики (4/6)}
Один и тот же язык может порождаться
разными (но эквивалентными) контекстно-свобод\-ными грамматиками.

\smallskip
\begin{example}
Рассмотрим ещё раз КСГ $G_1$ с правилами
$E\rightarrow E+T$,
$E\rightarrow T$,
$T\rightarrow T*F$,
$T\rightarrow F$,
$F\rightarrow a$ и КСГ
$G_2$ с правилами
$E\rightarrow E+E$,
$E\rightarrow E*E$,
$E\rightarrow a$.
Грамматика $G_2$  задает тот же
упрощенный язык арифметических выражений,
что и грамматика $G_1$, т.~е. $\frak{L}(G_1)=\frak{L}(G_2)$.
Чтобы убедиться в этом, достаточно заметить,
что обе грамматики порождают в точности цепочки вида
$a\oplus a\oplus\dots\oplus a$, где каждое $\oplus$ есть $+$ или $*$.
\end{example}

\begin{remark}
Любой вывод в грамматике $G_1$
легко перестраивается в вывод
в грамматике $G_2$ той же терминальной
цепочки.
Действительно, вместо правила
$E\rightarrow E+T$ надо применить правило
 $E\rightarrow E+E$,
вместо правила $T\rightarrow T*F$
 надо применить правило\\
$E\rightarrow E*E$,
вместо правила $F\rightarrow a$ надо применить правило
 $E\rightarrow a$,
а применения правил
$E\rightarrow T$ и
$T\rightarrow F$ надо просто исключить из вывода.
Другими словами, из вывода в грамматике $G_1$
можно получить (более короткий!) вывод в грамматике $G_2$ 
заменой и вычёркиванием правил.
Обратное соответствие правил так не строится,
поскольку из более короткого вывода более длинный вывод заменой правил получить
невозможно.   
\end{remark}

\end{frame}

% редакция Алексей Архиповский 18.09.26 исправлена ФАН 6.10.26
\begin{frame}
\frametitle{\Ok 1.1.6. Контекстно-свободные грамматики (5/6)}
Представленные грамматики $G_1$ и $G_2$
задают один и тот же язык,
но грамматика $G_2$ более компактна, то есть содержит меньше правил и меньше нетерминалов,
чем грамматика $G_1$.
Возникает вопрос: на основании каких свойств
следует предпочесть одну КСГ другой?
Ответ на поставленный вопрос зависит
от целей построения грамматического описания языка.
Например, для {\em публикации}
языка предпочтительна
более компактная грамматика $G_2$.
Однако с точки зрения 
использования языка для {\em программирования }
 грамматика $G_2$ обладает двумя существенными
недостатками.
Во-первых,
КСГ $G_2$ \odmkey{неоднозначна}{Грамматика!неоднозначная}~---
на рисунке представлены два {\em различных} дерева разбора
для выражения $a+a*a$ в этой грамматике (см. \DMref{п.~1.3.6}{1.3.6.6}).

\begin{figure}[h]
\begin{center}
\includegraphics[width=110mm]{1_1_6_5.jpg}
\end{center}
\end{figure}

\end{frame}

% редакция ФАН 6.10.26
\begin{frame}
\frametitle{\Ok 1.1.6.Контекстно-свободные грамматики (6/6)}


Во-вторых, грамматика $G_2$ не учитывает семантику языка
(приоритет арифметических операций).
Если использовать дерево разбора
для вычислений
(что и делается в большинстве реальных компиляторов),
то дерево справа
даст неверный результат.

\smallskip
Конструирование наиболее подходящей грамматики для заданного 
языка является {\em навыком}, который приходит с опытом. 
Всё зависит от {\em целей} составления грамматики.

\begin{example}
Рассмотрим КСГ $G_3$ с правилами
$E\rightarrow a+E$,
$E\rightarrow a*E$,
$E\rightarrow a$.
Грамматика $G_3$  задает тот же
упрощенный язык арифметических выражений,
что и грамматики $G_1$ и $G_2$. При этом 
грамматика $G_3$ столь же компактна,  
как грамматика $G_2$ и однозначна, как грамматика $G_1$.


\begin{columns}[t]
\begin{column}{8.2cm}
На рисунке справа приведено дерево разбора выражения
$a+a*a$ в грамматике $G_3$.
Эта грамматика является праволинейной
(\DMref{п.~1.1.7}{1.1.7.})
и 
обеспечивает наиболее эффективный 
синтаксический анализ \TODO{ref}.
\end{column}
\begin{column}{7cm}
\vspace{-20pt}
\begin{figure}[h]
\begin{center}
\includegraphics[width=70mm]{1_1_6_6.jpg}
\end{center}
\end{figure}
\end{column}
\end{columns}

\vspace{4pt}
На первый взгляд, такая грамматика практически неадекватна, но именно так
работали электронные калькуляторы, распространённые в прошлом веке. 
\end{example} 

\end{frame}




\subsubsection{1.1.7. Линейные, регулярные и автоматные грамматики}\label{1.1.7.}
\begin{frame}
\frametitle{\Ok 1.1.7. Линейные, регулярные и автоматные грамматики (1/5)}

Пусть задана порождающая грамматика $G=\langle \mathrm{M}, \Sigma, \Pi, S\rangle$.
Правило грамматики называется
\odmkey{начальным}{Правило!начальное}, если 
в левой части стоит аксиома грамматики $S$.
Правило грамматики называется
\odmkey{заключительным}{Правило!заключительное}, если оно имеет
вид $A\rightarrow\alpha$, где $A\in \mathrm{M}$, $\alpha\in \Sigma^*$. 
Правило
называется \odmkey{линейным}{Правило!линейное}
(соответственно,
\odmkey{леволинейным}{Правило!леволинейное} или
 \odmkey{праволинейным}{Правило!праволинейное}),
если оно имеет вид $A\rightarrow \alpha B\beta$ (соответственно, $A\rightarrow B\beta$ или
$A\rightarrow \alpha B$), где $A, B \in \mathrm{M}$ и 
$\alpha, \beta\in \Sigma^*$.

\smallskip
Порождающая грамматика, все правила которой заключительные или
линейные (соответственно, леволинейные или праволинейные), называется
\odmkey{линейной}{Грамматика!линейная} (соответственно, 
\odmkey{леволинейной}{Грамматика!леволинейная} или
\odmkey{праволинейной}{Грамматика!праволинейная}). 
Очевидно,
что праволинейные и леволинейные грамматики
являются линейными, а линейные грамматики
являются контекстно-свободными.

\smallskip
Язык, порождаемый линейной (соответственно,
леволинейной или праволинейной) грамматикой,
называется    
\odmkey{линейным}{Язык!линейный} 
(соответственно, 
\odmkey{леволинейным}{Язык!леволинейный} или
\odmkey{праволинейным}{Язык!праволинейный}). 
\begin{example}
Язык $\Set{a^n b^n}{n\in \Bbb{N}}$ 
является линейным языком, потому что порождается 
линейной грамматикой
$\langle \{S, A\}, \{a,b\},
\{S\rightarrow aAb, A\rightarrow 
\varepsilon, A\rightarrow aAb\},S\rangle$.
\end{example}

\begin{remark}
Любое дерево вывода в линейной грамматике
имеет один нетерминальный <<ствол>>,
по бокам от которого висят терминальные 
<<листья>>, причём в леволинейных деревьях листья только справа, а в праволинейных деревьях листья только слева.
\end{remark}

\end{frame}


\begin{frame}
\frametitle{\Ok 1.1.7. Линейные, регулярные и автоматные грамматики (2/5)}
\begin{theorem}
Для любой праволинейной грамматики 
существует эквивалентная леволинейная грамматика
и наоборот.
\end{theorem}
\begin{proof}
Пусть задана праволинейная грамматика. 
Построим для неё эквивалентную 
леволинейную грамматику: 
для каждого начального правила 
$S\rightarrow \alpha A$ введём
заключительное правило 
$A\rightarrow \alpha$,
для каждого праволинейного правила 
$A\rightarrow \alpha B$ введём
леволинейное правило 
$B\rightarrow A\alpha$,
для каждого заключительного правила 
$A\rightarrow \alpha$ введём
начальное правило 
$S\rightarrow A\alpha $.
Теперь любой вывод 
в исходной грамматике 
преобразуется в вывод в построенной грамматике:
достаточно применять новые правила начиная с конца исходного вывода.
\end{proof}

Таким образом, праволинейные и леволинейные грамматики порождают один и тот же класс языков.
Поэтому определение в \DMref{п.~1.1.4}{1.1.4.}, 
гласящее, что леволинейные и
праволинейные грамматики называются
\odmkey{регулярными}{Грамматика!регулярная},
оказывается вполне корректным. 
Для определённости далее рассматриваются леволинейные грамматики. 
Регулярная грамматика
называется \odmkey{автоматной}{Грамматика!автоматная}, если все
терминальные цепочки $\alpha$ в правилах состоят ровно из одного
символа. Другими словами, леволинейная автоматная грамматика
содержит правила вида $A\rightarrow a$ и $A\rightarrow B a$, где
$A, B \in \mathrm{M}$, $a\in \Sigma$.

\end{frame}


\begin{frame}
\frametitle{\Ok 1.1.7. Линейные, регулярные и автоматные грамматики (3/5)}

Терминальные цепочки в правилах регулярной грамматики могут быть
пустыми, и значит, язык, порождаемый регулярной грамматикой $G$,
может содержать пустое слово $\varepsilon$.  В то же время язык,
порождаемый автоматной грамматикой, очевидно, не содержит пустого
слова $\varepsilon$. 
Таким образом, формально эти классы языков различны, но на самом деле они почти эквивалентны.

\smallskip
Чтобы установить эту связь,
необходимо обозначить некоторые специальные свойства
грамматик, которые понадобятся также и в дальнейшем.    

Нетерминальный символ называется
\odmkey{непорождающим}{Символ!непорождающий},
если из него невозможно вывести терминальную цепочку.
Нетерминальный или терминальный  символ называется
\odmkey{недостижимым}{Символ!недостижимый},
если его невозможно вывести из аксиомы.
Если правило имеет вид $A\rightarrow B$, то оно называется
\odmkey{цепным}{правило!цепное}. 
Если в правиле $A\rightarrow
\alpha$ цепочка $\alpha$ пустая,
$\alpha=\varepsilon$, то правило называется
\odmkey{$\varepsilon$-правилом}{Правило!$\varepsilon$}.
Если в одном из начальных правил $S\rightarrow \alpha$
аксиома входит в правую часть,
то правило называется 
\odmkey{$S$-правилом}{Правило!$S$}.

Недостижимые и непорождающие символы, цепные правила,
$\varepsilon$-правила и $S$-правила грамматик
создают различные помехи в рассуждениях и применении алгоритмов, но при этом все такие особенности грамматик являются {\em устранимыми} (см. \DMref{п.~1.3.2}{1.3.2.}).

\end{frame}


\begin{frame}
\frametitle{\Ok 1.1.7. Линейные, регулярные и автоматные грамматики (4/5)}
Грамматики, в которых нет непорождающих 
и недостижимых символов, а также нет 
цепных правил, 
$\varepsilon$-правил и
$S$-правил называются \odmkey{приведёнными}{Грамматика!приведённая}.

\smallskip
На практике языки с пустыми словами, как правило, не используются,
а грамматики составляются квалифицированными авторами и
не содержат непорождающих 
и недостижимых символов, а также неподходящих правил.
Поэтому мы докажем утверждение 
об эквивалентности регулярных и автоматных грамматик
в существенно более
слабой, но практически более полезной форме, а именно в форме
алгоритма преобразования {\em приведённой}  грамматики к нужному виду.

\begin{theorem}
Для любой приведённой регулярной грамматики существует эквивалентная автоматная
грамматика.
\end{theorem}

\begin{proof}
Пусть задана леволинейная регулярная грамматика $G=\langle \mathrm{M}, \Sigma,
\Pi, S\rangle$, где каждое правило $r\in \Pi$ имеет вид $A\rightarrow
\alpha$ или $A\rightarrow B\alpha$, причем $A, B \in \mathrm{M}$,\\
$\alpha\in \Sigma^*$ и $|\alpha|\geq 1$. Построим автоматную грамматику
$G'=\langle \mathrm{M}', \Sigma, \Pi', S\rangle$  в соответствии с алгоритмом на следующем слайде.
 Очевидно, что всякому выводу в грамматике $G$ однозначно
 соответствует вывод в грамматике $G'$. Таким образом, грамматики
 эквивалентны.
\end{proof}
\end{frame}

\begin{frame}
\frametitle{\Ok 1.1.7. Линейные, регулярные и автоматные грамматики (5/5)}


\begin{algorithmic}
\REQUIRE Регулярная леволинейная грамматика
$G=\langle \mathrm{M}, \Sigma, \Pi, S\rangle$
%, не содержащая цепных правил и $\varepsilon$-правил.
\ENSURE
Автоматная леволинейная грамматика
$G'=\langle \mathrm{M}', \Sigma, \Pi',S\rangle$.
\STATE{$\mathrm{M}'\Assign \mathrm{M}; \Pi'\Assign\emptyset$}
 \COMMENT { $r= A \rightarrow \alpha$ или $r= A\rightarrow B\alpha$ }
   \FOR {\textbf{$r\in \Pi$}}
  \STATE {\textbf{if $|\alpha|=1$ then 
  $\Pi'\Assign \Pi'\cup r;$ next for $r$ end if}}
    \COMMENT{ переносим без изменений }
      \STATE {\textbf{$n\Assign|\alpha|-1$}}
    \COMMENT{ количество новых нетерминалов }
    \STATE{\textbf{$\mathrm{M}'\Assign \mathrm{M}'\cup$ new $A_n$ }}
    \COMMENT{ добавляем нетерминал $A_n$ }
    \STATE {\textbf{$\Pi'\Assign \Pi'\cup(A\rightarrow A_n\alpha [n+1])$}}
    \COMMENT{ добавляем правило
     $A\rightarrow A_na_{n+1}$ }
    \FOR{\textbf{$i$ from $n-1$ to 1}}
       \STATE {\textbf{$\mathrm{M}'\Assign \mathrm{M}'\cup$ new $A_i$ }}
       \COMMENT{ добавляем нетерминал $A_i$ }
       \STATE{\textbf{$\Pi'\Assign \Pi'\cup (A_{i+1}\rightarrow A_i \alpha [i+1])$}}
       \COMMENT{ добавляем правило $A_{i+1}\rightarrow A_i a_{i+1}$ }
    \ENDFOR
    \STATE{\textbf{if $r=A\rightarrow \alpha$ then $\Pi'\Assign \Pi'\cup(A_1\rightarrow \alpha[1])$}}
      \COMMENT{ заключительное правило }
     \STATE {\textbf{else $\Pi'\Assign \Pi'\cup(A_1\rightarrow B\alpha[1])$ end if}}
      \COMMENT{ незаключительное правило }
 \ENDFOR
\end{algorithmic}

\end{frame}

\subsubsection{1.1.8. Языки записи грамматик}\label{1.1.8.}
\begin{frame}
\frametitle{\Ok 1.1.8. Языки записи грамматик (1/9)}
Запись правил грамматики в общей форме довольно многословна, не всегда легко читается и
существенно использует рекурсию. 

\begin{example}
Набор правил
$N\to D$, $N\to ND$, $D\to 0$,
$D\to 1$, $D\to 2$, $D\to 3$, 
$D\to 4$, $D\to 5$, 
$D\to 6$, $D\to 7$, $D\to 8$, $D\to 9$
по существу означает, что целое без знака~--- это
последовательность из одной или более цифр, что не всегда очевидно
без дополнительных пояснений.
\end{example}

\smallskip
Были предложены многочисленные усовершенствования записи грамматик, в особенности контекстно-свободных и 
автоматных грамматик,
которые сводятся к введению дополнительных
\odmkey{метасимволов}{Метасимвол}. Дополнительные метасимволы не
меняют выразительной силы  грамматик, но делают запись грамматики более
удобной для восприятия человеком. Обычно рассматриваются следующие
метасимволы.

\begin{enumerate}
\item Вертикальная черта $\Ms{|}$~--- разделение альтернатив. Правило
$$
A\rightarrow \alpha \Ms{|}\beta
$$
эквивалентно паре правил
$$
A\rightarrow \alpha  \qquad A\rightarrow\beta
$$
\end{enumerate}
\end{frame}

\begin{frame}
\frametitle{\Ok 1.1.8. Языки записи грамматик (2/9)}
\begin{enumerate}

\item[2.] Круглые скобки $\Ms{(\ )}$~--- введение неявного нетерминала.
Круглые скобки могут содержать несколько альтернатив, разделенных
метасимволом $\Ms{|}$. Правило
$$
A\rightarrow \alpha\Ms{(}\beta\Ms{)}\gamma
$$
эквивалентно паре правил
$$
A\rightarrow \alpha B\gamma \qquad B\rightarrow\beta
$$
где $B$~--- это новый нетерминал.
\item[3.] Квадратные скобки $\Ms{[\ ]}$~--- необязательная конструкция.
Цепочка, заключенная в квадратные скобки, может присутствовать или
отсутствовать. Правило
$$
A\rightarrow \alpha\Ms{[}\beta\Ms{]}\gamma
$$
эквивалентно паре правил
$$
A\rightarrow \alpha\gamma \qquad A\rightarrow\alpha\beta\gamma.
$$
\end{enumerate}
\end{frame}

\begin{frame}
\frametitle{\Ok 1.1.8. Языки записи грамматик (3/9)}
\begin{enumerate}
\item[4.] Фигурные скобки $\Ms{\{\ \}}$~--- повторение конструкции.
Цепочка, заключенная в фигурные скобки, может повторяться
несколько раз. Правило

\begin{center}
$
A\rightarrow \alpha\Ms{\{}\beta\Ms{\}}\gamma
$
\end{center}

эквивалентно тройке правил

\begin{center}
$
A\rightarrow \alpha B\gamma \qquad B\rightarrow\beta  \qquad B\rightarrow B \beta
$
\end{center}

где $B$~--- это новый нетерминал.

\item[5.] Диез $\Ms{\#}$ в фигурных скобках~--- итерирование одной конструкции через другую конструкцию (\odmkey{итерация Цейтина}{Итерация!Цейтина}
\footnote{Цейтин Григорий Самуилович, 1936--2022}).
Этот метасимвол является обобщением предыдущего:
$\Ms{\{}\alpha\Ms{\}}=\Ms{\{}\alpha\Ms{\#}\varepsilon\Ms{\}}$. Правило

\begin{center}
$
A\rightarrow \alpha\Ms{\{}\beta\Ms{\#}\gamma\Ms{\}}\delta
$
\end{center}

эквивалентно тройке правил

\begin{center}
$
A\rightarrow \alpha B\delta \qquad B\rightarrow\beta 
\qquad B\rightarrow B
\gamma\beta
$
\end{center}

где $B$~--- это новый нетерминал.
\end{enumerate}

\vspace{4pt}

\end{frame}

\begin{frame}
\frametitle{\Ok 1.1.8. Языки записи грамматик (4/9)}
\begin{enumerate}
\item[6.] Точка с запятой $\Ms{;}$~--- разделитель правил в множестве правил.
Строчка
$$
A\rightarrow \alpha \Ms{;} B\rightarrow\beta
$$
эквивалентна паре правил
$$
A\rightarrow \alpha \qquad B\rightarrow\beta 
$$
\item[7.] Точка $\Ms{.}$~--- признак конца множества правил.
\item[8.] Двоеточие $\Ms{:}$~--- замена метасимвола
$\rightarrow$, которого нет на стандартной клавиатуре.
\end{enumerate}

\smallskip
\begin{remark}
Обычно считают, что метасимволы 
$\Ms{[\ ]}$, $\Ms{\{\ \}}$ и $\Ms{\#}$ также вводят
неявный нетерминал, а также что символ разделения альтернатив имеет приоритет перед всеми скобками, так что писать дополнительные круглые скобки
не обязательно.
\end{remark}

\smallskip
\begin{example}
В приведённых обозначениях грамматику целых чисел с возможным знаком можно
описать следующим образом:
$$N \to \Ms{[} + \Ms{|} - \Ms{]}
\Ms{(}  1 \Ms{|} 2 \Ms{|} 3 \Ms{|}
 4 \Ms{|} 5 \Ms{|} 6 \Ms{|} 7 \Ms{|} 8 
 \Ms{|} 9 \Ms{)}
\Ms{[} \Ms{\{} 0 \Ms{|} 1 \Ms{|} 2 \Ms{|} 3 \Ms{|}
 4 \Ms{|} 5 \Ms{|} 6 \Ms{|} 7 \Ms{|} 8 
 \Ms{|} 9 \Ms{\}} \Ms{]} ~\Ms{|}~ 0 ~\Ms{.}$$
Здесь не порождаются ведущие нули, 
а ноль знака не имеет. 
В такой записи грамматики отсутствует явная рекурсия и не используются
вспомогательные нетерминалы.
\end{example}
\end{frame}

\begin{frame}
\frametitle{\Ok 1.1.8. Языки записи грамматик (5/9)}

\vspace{-2pt}
Иногда метасимволы $\Ms{\{\ \}}$ определяют так, чтобы конструкция
повторялась 0 или более раз, т.е. могла отсутствовать. В этом
случае правило
$$
A\rightarrow \alpha\Ms{\{}\beta\Ms{\}}\gamma
$$
считается эквивалентным трём правилам
$$
A\rightarrow \alpha B\gamma 
\qquad B\rightarrow\varepsilon 
\qquad B\rightarrow B \beta
$$
где $B$~--- это новый нетерминал, а $\varepsilon$~--- пустая
цепочка.

\begin{remark} Символы 
$\Ms{:\ |\ (\ )\ [\ ]\ \{\ \}\ \#\ ;\ .}$
присутствуют на клавиатуре и ими 
удобно пользоваться при описании 
грамматики. Однако эти же символы 
часто используются как терминальные 
символы языка, 
и тогда встаёт вопрос,
как отличить терминальный символ
от метасимвола? Здесь этот вопрос
решён {\em визуально}, 
за счёт использования цвета.
Если система программирования
не допускает различения символов по цвету,
то приходится придумывать 
дополнительные средства, например, заключать
терминальные символы в кавычки $"\ "$.
\end{remark}

Нетрудно убедиться, что все эти расширения (а также некоторые
другие, применяемые на практике) не выводят за пределы
контекстно-свободных грамматик типа $\frak{T}_2$.

\end{frame}

\begin{frame}
\frametitle{\Ok 1.1.8. Языки записи грамматик (6/9)}

Таким образом, сосуществуют и используются 
{\em несколько языков} для записи порождающих
формальных грамматик. Такая множественность 
языков и диалектов является скорее правилом,
нежели исключением. Какой вариант использовать~---
дело вкуса. 

\smallskip
Очевидно, что использование метасимволов
сокращает запись грамматики и делает её более наглядной
(при том условии, однако, что читатель грамматики
точно знает значения всех метасимволов).
С другой стороны, использование метасимволов
существенно усложняет, например, доказательство
эквивалентности грамматик, поскольку
приходится рассматривать в доказательстве 
каждый метасимвол отдельно. 
 
\begin{example} 
На следующем слайде для сравнения приведены четыре варианта
записи грамматики двух языков:
строгого языка  записи грамматик с использованием $\rightarrow$ и
расширенного языка с использованием метасимволов.
Для наглядности в этих записях использованы
одинаковые нетерминальные символы:
$G$~--- собственно грамматика,
$N$~--- алфавит нетерминальных символов целевого языка,
$T$~--- алфавит терминальных символов целевого языка,
$R$~--- правило грамматики,
$S$~--- цепочка символов.
Все терминальные символы заключены в кавычки $"\ "$.
Правила строгого языка считаются разделенными
некоторым непечатаемым символом (конец строки, пробел, ...),
который здесь обозначен $\sqcup$.
\end{example}

\end{frame}

\begin{frame}
\frametitle{\Ok 1.1.8. Языки записи грамматик (7/9)}
\begin{columns}[t]
\begin{column}{7.0cm}
Грамматика строгого языка, записанная на строгом языке

$G \rightarrow R$

$G \rightarrow R\sqcup G$

$R \rightarrow N\ "\!\rightarrow\!"\ S$

$S \rightarrow N $

$S \rightarrow T $

$S \rightarrow N S$

$S \rightarrow T S$

\end{column} 
\begin{column}{8.0cm}
Грамматика расширенного языка, записанная на строгом языке

\begin{columns}[t]
\begin{column}{3.7cm}
$G \rightarrow R\ \ "."$

$G \rightarrow R\ ";"\ G$

$R  \rightarrow N\ "\!:\!"\ S$

$S \rightarrow N $

$S \rightarrow T $

$S \rightarrow N S$

$S \rightarrow T S$
\end{column}
\begin{column}{4.3cm}
$S \rightarrow S\  "\!\mid\!"\  S $

$S \rightarrow "\{"\ S\  "\#"\  S\  "\}"$

$S \rightarrow "("\ S\ ")" $

$S \rightarrow "["\ S\ "]" $

$S \rightarrow "\{"\ S\ "\}" $
\end{column}
\end{columns}
\end{column} 
\end{columns}

\medskip
\begin{columns}[t]
\begin{column}{7.0cm}
Грамматика строгого языка, записанная на расширенном языке

$G\ :\ \{\ R\ \#\ \sqcup\ \}\ ; $

$R\ :\ N\ "\!\rightarrow\!"\ S\ ;$

$S\ :\ \{\ N \mid T \}\ . $
\end{column} 
\begin{column}{8.0cm}
Грамматика расширенного языка, записанная на расширенном языке

$G\ :\ \{\ R\ \#\ ";" \}\ "."\ ; $

$R\ :\ N\ \ "\!:\!"\ S\ ;$

$S\ :\ \{\ N \mid T \}\ \ \mid \ \ 
"(" S\ \#\ "\!\mid\!"\ ")" \ \ \mid \ \ \hfill\linebreak 
"["\ S\ \#\ "\!\mid\!"\ "]" \ \ \mid \ \ 
"\{" S \ [\ "\#"\ S\ ]\ "\}" \ 
. $
\end{column} 
\end{columns}
\end{frame}


\begin{frame}
\frametitle{\Ok 1.1.8. Языки записи грамматик (8/9)}
Почти одновременно с появлением {\em теории} формальных
грамматик для
практических целей были предложены средства {\em записи}
синтаксиса языка.
Исторически самым заметным явилось описание языка программирования
Алгол-60 с помощью \odmkey{нормальных форм
Бэкуса-Наура}{Нормальная форма!Бэкуса-Наура}
 (сокращенно БНФ).
Описание языка c помощью нормальных форм 
Бэкуса\footnote{John Warner Backus (1924--2007)}--
Наура\footnote{Peter Naur (1928--2016)} 
состоит из набора 
\odmkey{основных символов}{Символ!основной} языка, 
набора названий конструкций (или
\odmkey{металингвистических
переменных}{Переменная!металингвистическая}) и набора
\odmkey{металингвистических формул}{Формула!металингвистическая}.
Каждая металингвистическая формула состоит
из левой части, в которой
стоит металингвистическая переменная, \odmkey{металингвистической
связки}{Связка!металингвистическая} $::=$ (читается <<это есть>>)
и правой части, в которой указывается один или несколько вариантов
построения конструкции в левой части. Каждый вариант~--- это
последовательность металингвистических переменных и основных
символов. Варианты разделяются символом $|$ (читается <<или>>).
Металингвистические переменные заключаются в угловые скобки, чтобы
не путать их с основными символами. Считается, что метасимволы $<,
>$, $|$ и $::=$ не входят в число основных символов.

\end{frame}

\begin{frame}
\frametitle{\Ok 1.1.8. Языки записи грамматик (9/9)}

\vspace{-1pt}
\begin{example}
Язык целых десятичных чисел со знаком в обозначениях Бэкуса-Наура: \\[-2pt]
<Целое со знаком> $::=$ 0 | <Натуральное> | <Знак> <Натуральное> \\[-2pt]
<Знак> $::=$ $+$ | $-$ \\[-2pt]
<Натуральное> $::=$ <Цифра не ноль>  |  <Натуральное> <Цифра> \\[-2pt]
<Цифра> ::= 0 | <Цифра не ноль> \\[-2pt]
<Цифра не ноль> ::=  1 | 2 | 3 | 4 | 5 | 6 | 7 | 8 | 9
\end{example}

\begin{digress}
В первом варианте БНФ нетерминалы заключали в скобки < >, 
а терминалы не выделяли. Сейчас, наоборот,
нетерминалы не выделяют, а терминалы заключают в кавычки $"\ "$.   
\end{digress}

Совершенно ясно, что нормальные формы Бэкуса-Наура~--- это один из
способов записать контекстно-свободную грамматику, 
то есть грамматику типа $\frak{T}_2$.  Действительно, основные
символы~--- это алфавит терминальных символов $\Sigma$,
металингвистические переменные~--- это алфавит нетерминальных
символов $\mathrm{M}$, 
металингвистические формулы~--- это множество правил
вывода $\Pi$, в которых вместо символа $\to$
использована связка $::=$, а аксиома грамматики обычно ясна из контекста.

\begin{digress}
Форма Бэкуса--Наура
была предназначена, по замыслу авторов,
прежде всего
для точного и в то же время
{\em удобочитаемого}
описания синтаксиса
языков программирования,
поэтому используются развёрнутые названия конструкций.
\end{digress}
\end{frame}


% GA_1_2025.tex, 2880–4853
\subsection{1.2. Регулярные выражения}
\begin{frame}
\frametitle{\Ok 1.2. Регулярные выражения (1/3)}

Порождающие грамматики~--- не единственное средство определения
языков. Языки являются множествами, поэтому над ними определены
все операции над множествами 
(объединение, пересечение, разность и 
симметрическая разность). Кроме того, для
элементов языков~--- слов~--- определена операция конкатенации,
что позволяет определить для языков и другие операции. 
Таким образом,
существует возможность определить язык
{\em алгебраически}, как выражение,
построенное из более простых элементов.

\medskip
Основополагающий вклад в развитие теории регулярных выражений 
внес Клини.\footnote{Stephen Cole Kleene (1909--1994)}

\medskip
Удивительный факт: далёкие как от исчислений (формальные грамматики),
так и от алгебраических структур (регулярные выражения) формализмы,
в частности, ориентированные {\em графы}, при использовании для определения 
языков приводят к одним и тем же результатам. Обсуждение этого наблюдения
начинается в этом разделе, а продолжается и заканчивается в разделе
\DMref{2.2}{2.2}.     

\end{frame}

\begin{frame}
\frametitle{\Ok 1.2. Регулярные выражения (2/3)}

\begin{tabular}[t]{p{10mm}|p{40mm}|p{95mm}}
  № & Название параграфа
  & Ключевые термины и обозначения \\
  \hline
\text{\DMref{1.2.1.}{1.2.1.}} &  Регулярные операции
  & { регулярные операции, 
  объединение, альтернация, сложение, умножение, конкатенация, сцепление,
  итерация, замыкание Клини}
   \\
\text{\DMref{1.2.2.}{1.2.2.}} & Алгебра регулярных выражений 
  & 
{ элементарный язык,
  регулярный язык, регулярное выражение, 
  алгебра регулярных множеств $\cal{R}$, алгебра регулярных выражений $\frak{R}$,
  эквивалентность регулярных выражений}
  \\
\text{\DMref{1.2.3.}{1.2.3.}}  & Языковые графы &
{граф языка, источник, полный источник, эквивалентные источники,
  двухполюсник, начальные и заключительные узлы, значение пути, размер источника,
  замкнутое множество узлов } 
  \\ 
\text{\DMref{1.2.4.}{1.2.4.}}  & Детерминированные и недетерминированные источники &
{ теорема детерминизации } 
\end{tabular}

\end{frame}


\begin{frame}
\frametitle{\Ok 1.2. Регулярные выражения (3/3)}

\begin{tabular}{p{10mm}|p{45mm}|p{90mm}}
  № & Название параграфа
  & Ключевые термины и обозначения \\
  \hline
\text{\DMref{1.2.5.}{1.2.5.}} & Расширение языка регулярных выражений &
{операции повторения, итерация Цейтина, символы шаблона, символьные классы}\\   
\text{\DMref{1.2.6.}{1.2.6.}} & Программирование в регулярных выражениях &
{POSIX, PERL, SNOBOL4, REFAL }\\   
\text{\DMref{1.2.7.}{1.2.7.}} & Уравнения с регулярными коэффициентами &
{ }\\   
\text{\DMref{1.2.8.}{1.2.8.}} & Характеристика класса регулярных языков &
{лемма о разрастании, язык Дика  }   
\end{tabular}

\end{frame}


\subsubsection{1.2.1. Регулярные операции}
\label{1.2.1.}
\begin{frame}
\frametitle{\Ok 1.2.1. Регулярные операции (1/3)}

Пусть $L_1, L_2\subset A^*$. Рассмотрим следующие три операции над языками $L_1, L_2$.

\odmkey{Объединение}{Объединение!языков}
(\odmkey{альтернация}{Операция!альтернации},
\odmkey{сложение}{Сложение!языков}) обозначается знаками $\cup, ~|~, +$:
$$
L_1\cup L_2 = L_1 ~|~ L_2 = L_1 + L_2\Def \Set{\alpha \in A^*}{\alpha \in L_1 \Or \alpha\in L_2}.$$

\begin{remark}
В определении указано,
что языки определены 
над одним алфавитом, \\$L_1, L_2\subset A^*$,
но это сделано только ради удобства 
обозначений и не ограничивает общности.
Можно считать, что $L_1\subset A_1^*$, 
$L_2\subset A_2^*$, и тогда $L_1\cup L_2\subset (A_1\cup A_2)^*$.
\end{remark}

\odmkey{Умножение}{Умножение!языков}
(\odmkey{конкатенация}{Операция!конкатенации},
\odmkey{сцепление}{Сцепление!языков}) обозначается знаком $\cdot$ или
опускается:
$$
L_1 L_2= L_1\cdot L_2\Def \Set{\alpha\in A^*}
{\alpha=\alpha_1\alpha_2 \And \alpha_1\in L_1 \And \alpha_2\in
L_2}.
$$

\begin{remark}
Можно сцеплять язык сам с собой:
$L, LL, LLL, \dots$,
что естественно обозначить с
помощью показателя степени,
то есть $L^1\Assign L, L^2\Assign LL,
L^3\Assign LLL$ и так далее.
При этом по определению 
считают, что $L^0\Assign \{\varepsilon\}$.  
\end{remark}

\odmkey{Итерация}{Итерация!языков} (\odmkey{замыкание Клини}{Замыкание!языков}) обозначается знаком $*$:

\vspace{-12pt}
$$
L_1^*\Def \varepsilon\cup L_1\cup L_1 L_1\cup\dots =
\bigcup_{i=0}^\infty L_1^i
.$$

\end{frame}

\begin{frame}
\frametitle{\Ok 1.2.1. Регулярные операции (2/3)}

Введённые операции называются
\odmkey{регулярными}{Операция!регулярная}. Эти операции имеют
особое значение для теории языков и автоматов.

\medskip
Поскольку результатом регулярной 
операции над языками является язык,
из этих операций с помощью суперпозиции 
можно строить
сколь угодно сложные  выражения,
которые называются \odmkey{регулярными}{Выражение!регулярное}. 

\medskip
Для регулярных операций обычно устанавливают следующий приоритет:
$*$, $\cdot$, $|$, что позволяет опускать лишние скобки при записи
выражений.

\begin{remark}
Операция альтернации является
привычной операцией объединения из
наивной теории множеств,
но операции конкатенации и итерации
не являются 
и не выражаются операциями наивной теории множеств. 
\end{remark}

\begin{examples}
Пусть язык 
$L_1 =\{a,b,c,\dots,x,y,z\}$~--- это латинский алфавит, то есть множество однобуквенных слов,
а язык 
$L_2=\{1,2,3,\dots\}$~--- это натуральные
числа в десятичной записи.
Тогда $L_1 L_2$~--- это множество 
идентификаторов, которые начинаются
с буквы и продолжаются цифрами,
$L_1^*$~--- это множество всех 
слов, составленных из букв латинского алфавита,
включая пустое слово $\varepsilon$,
а $L_2^* = L_2$.
\end{examples}

\end{frame}


\begin{frame}
\frametitle{\Ok 1.2.1. Регулярные операции (3/3)}
\begin{theorem}
Пусть $L_1, L_2, L_3\subset A^*$. Тогда регулярные операции
обладают следующими свойствами.
\begin{enumerate}
  \item Альтернация коммутативна:\\[-2pt]
  $L_1~|~ L_2 = L_2~|~ L_1$;\\[-1pt]
  \item Альтернация ассоциативна:\\[-2pt]
  $L_1~|~ (L_2~|~ L_3) = (L_1~|~ L_2)~|~L_3$;\\[-1pt]
  \item Конкатенация дистрибутивна относительно альтернации слева:\\[-2pt]
  $L_1(L_2~|~ L_3) = L_1 L_2~|~ L_1 L_3$;\\[-1pt]
  \item Конкатенация дистрибутивна относительно альтернации справа:\\[-2pt]
$(L_1~|~ L_2)L_3 = L_1 L_3~|~ L_2 L_3$;\\[-1pt]
  \item Конкатенация ассоциативна:\\[-2pt]
  $(L_1 L_2)L_3 = L_1(L_2 L_3)$;\\[-1pt]
    \item Итерация идемпотентна:\\[-2pt]
  $(L_1^*)^* = L_1^*$;\\[-1pt]
      \item Альтернация идемпотентна:\\[-2pt]
  $L_1~|~L_1 = L_1$.\\[-1pt]
\end{enumerate}
\end{theorem}
\begin{proof}
Очевидно из определений.
\end{proof}
\end{frame}

\subsubsection{1.2.2. Алгебра регулярных выражений}\label{1.2.2.}
\begin{frame}
\frametitle{\Ok 1.2.2. Алгебра регулярных выражений (1/8)}

Пусть задан конечный алфавит $A=\{\Tuple{a}{1}{k}\}$. Язык $L_{a_i}\Assign\{a_i\}$,
состоящий из одного однобуквенного слова, называется
\odmkey{элементарным}{Язык!элементарный} и обозначается $a_i$.
Пустое слово $\varepsilon$ тоже считается элементарным языком.  

\begin{digress}
Используемые обозначения многозначны.
Поясним это утверждение примером.\\
Пусть $A=\{a,b,c,\dots,x,y,z\}$. 
Тогда символ $a$ именует одновременно три различных объекта:
\begin{enumerate}
\item $a$~--- это буква латинского алфавита, $a\in A$;\\[-2pt]
\item $a$~--- это однобуквенное слово, $a\in A^*$;\\[-2pt]
\item $a$~--- это элементарный язык, $a\subset A^*$.\\[-2pt]
\end{enumerate}
Аналогично символ $\varepsilon$ именует одновременно три различных объекта:
\begin{enumerate}
\item $\varepsilon$~--- это символ {\em греческого}
алфавита, $\varepsilon\not\in A$;\\[-2pt]
\item $\varepsilon$~--- это обозначение пустого слова, $\varepsilon\in A^*$;\\[-2pt]
\item $\varepsilon$~--- это элементарный язык, $\varepsilon\subset A^*$.\\[-2pt]
\end{enumerate}
Такие обозначения удобны, 
потому что лаконичны, и не создают 
путаницы, потому что 
тип объекта всегда восстанавливается из контекста. В частности,
аргументы регулярных операций~--- это 
всегда языки, но не слова и не буквы.
\end{digress} 
\end{frame}

\begin{frame}
\frametitle{\Ok 1.2.2. Алгебра регулярных выражений (2/8)}


Язык (то есть множество слов в алфавите $A$) называется \odmkey{регулярным}{Язык!регулярный}, если он
может быть построен из элементарных языков (букв алфавита $A$ и пустого слова 
$\varepsilon$)
с помощью суперпозиции регулярных операций. Семейство регулярных языков является подмножеством
множества всех языков над алфавитом $A$.
Множество регулярных языков обозначим $\cal{R}$:
$\cal{R} \subset 2^{A^*}.$
Алгебра $\langle \cal{R}; ~|~, ~\cdot~, ~{}^*\rangle$ называется \odmkey{алгеброй
регулярных множеств}{алгебра!регулярных!множеств}. Эта алгебра
конечнопорождённая; ее образующими являются
элементарные языки.

\smallskip
Каждый регулярный язык может быть описан {\em формулой}, которая называется \odmkey{регулярным
выражением}{Выражение!регулярное}. Операндами регулярных выражений являются обозначения 
элементарных языков (то есть буквы алфавита $A$), а операциями~--- регулярные операции 
$|~\cdot~{}^*$. 
Множество регулярных выражений обозначим $\frak{R}$.
Алгебра $\langle \frak{R}; ~|~, ~\cdot~, ~*\rangle$ называется \odmkey{алгеброй
регулярных выражений}{алгебра!регулярных!выражений}.  Эта алгебра
также конечнопорождённая; ее образующими являются
буквы алфавита $A$ и обозначение $\varepsilon$.

\smallskip
Каждое регулярное выражение порождает единственный регулярный язык.
Один и тот же язык может быть порождён бесконечным числом
различных регулярных выражений.
\begin{example}
Различные регулярные выражения 
$\varepsilon,\ \varepsilon\varepsilon,\ \varepsilon~|~\varepsilon,\ \varepsilon^*,\ \dots$
порождают один и тот же регулярный язык $\{\varepsilon\}$. 
\end{example}

\end{frame}

\begin{frame}
\frametitle{\Ok 1.2.2. Алгебра регулярных выражений (3/8)}
Соответствие регулярных выражений и регулярных языков является 
{\em гомоморфизмом}, который уважает все регулярные операции.
Говорят, что регулярные выражения
\odmkey{эквивалентны}{Эквивалентность!регулярных выражений}, если они задают
один и тот же язык. 
Эквивалентность регулярных 
выражений обозначают, как обычно, знаком $=$.
Таким образом, по теореме о гомоморфизме
(\DMref{п.~0.3.2}{0.3.2.}) алгебра классов эквивалентности
регулярных выражений {\em изоморфна} алгебре регулярных множеств.
На этом обстоятельстве основана общепринятая в математике практика:
манипулировать конечными дескрипторами (в данном случае регулярными выражениями)
бесконечных объектов (в данном случае, регулярных языков).

\begin{examples}
\begin{enumerate}
  \item Обозначение $A^*$~--- это правильное регулярное выражение:
  $
A^*=(a_1~|~\dots~|~ a_k)^*.
$
  \item Любое слово регулярно: оно построено из букв
  конкатенацией.
  \item Любой конечный язык регулярен:
$
L=\{\alpha_1,\dots, \alpha_k\},\quad
  L=\alpha_1~|~\dots~|~\alpha_k.
$
  \item  Язык $(aba)^*$  регулярен.
\end{enumerate}
\end{examples}

В предыдущем параграфе
установлены семь важных эквивалентностей,
с помощью которых можно установить многие другие
полезные свойства операций алгебры регулярных
выражений.  

\end{frame}

\begin{frame}
\frametitle{\Ok 1.2.2. Алгебра регулярных выражений (4/8)}
Пустое множество $\emptyset\subset A^*$ 
является языком, этот язык считается регулярным по определению, $\emptyset\in \cal{R}, \emptyset\ne\varepsilon$. 
Напомним, что $A^*= \varepsilon +A + AA+ \dots =
\bigcup_{i=0}^\infty A^i$ и введём дополнительное обозначение:
$A^+ \Def A + AA+ \dots = \bigcup_{i=1}^\infty A^i$.
  
Пользуясь свойствами регулярных операций и определениями, 
нетрудно установить, что в алгебре 
регулярных выражений
выполняются следующие соотношения:

\vspace{-16pt}
\begin{columns}[t]
\begin{column}{7cm}
\begin{enumerate}
\item[8.] $A+\emptyset = A$;
\item[9.] $A\emptyset = \emptyset A=\emptyset$;
\item[10.] $A^* +A =A^*$;
\item[11.] $\emptyset^*=\varepsilon$;
\item[12.] $\varepsilon^*=\varepsilon$;
\item[13.] $A^* + \varepsilon=A^*$;
\end{enumerate}
\end{column}
\begin{column}{7cm}
\begin{enumerate}
\item[14.] $A\varepsilon = \varepsilon A = A$;
\item[15.] $A^+=AA^*$;
\item[16.] $A^+ + A =A^+$;
\item[17.] $\emptyset^+ = \emptyset$;
\item[18.] $\varepsilon^+=\varepsilon$;
\item[19.] $A^+ + \varepsilon=A^*$.
\end{enumerate}
\end{column}
\end{columns} 

\vspace{2pt}
\begin{remark}
Алгебра регулярных выражений является моноидом относительно конкатенации и языка $\varepsilon$,
является коммутативным моноидом
относительно альтернации и пустого языка
$\emptyset$, и является некоммутативным
полукольцом с единицей относительно 
указанных операций и элементов.  
\end{remark}
\end{frame}

\begin{frame}
\frametitle{\Ok 1.2.2. Алгебра регулярных выражений (5/8)}

Пользуясь указанными простыми свойствами, в алгебре регулярных выражений
можно доказывать различные равенства.

\begin{example}
 Доказать, что
$(b+aa^*b)+(b+aa^*b)(a+ba^*b)^*(a+ba^*b) = a^*b(a+ba^*b)^*$.
Доказательство:
$(b+aa^*b)+(b+aa^*b)(a+ba^*b)^*(a+ba^*b) =$\\
$=(b+aa^*b)(\varepsilon+(a+ba^*b)^*(a+ba^*b)) =
(b+aa^*b)(\varepsilon+(a+ba^*b)^+) = $\\
$=(b+aa^*b)(a+ba^*b)^* = 
(\varepsilon+aa^*)b(a+ba^*b)^* = (\varepsilon+a^+)b(a+ba^*b)^* = a^*b(a+ba^*b)^*$.
\end{example}

Имеют место некоторые менее очевидные свойства.

\begin{lemma}
$\A{L\subset A^*}
   {\A{\alpha\in L^*, \alpha\neq\varepsilon}
      {\E{\Tuple{\alpha}{1}{k}\in L}
         {\alpha=\alpha_1\dots\alpha_k}
      }
   }$.
\end{lemma}
Другими словами, если слово принадлежит итерации языка, то это слово всегда можно
разложить в конкатенацию слов, принадлежащих языку и обратно.  

\begin{theorem}
$(A+B)^*=(A^*+B^*)^* =(A^*B^*)^*$.
\end{theorem}
\begin{proof}
Обозначим $L_1\Assign (A+B)^*, L_2\Assign 
(A^*+B^*)^*, L_3=(A^*B^*)^*$.
\proofsection{$L_1\subset L_2$}
$A+B\subset A^* + B^* \Impl (A+B)^*\subset 
(A^* + B^*)^*$.
\proofsection{$L_2\subset L_1$}
Если $\alpha\in (A^*+B^*)^*$, то по лемме 
$\E{\Tuple{\alpha}{1}{k}\in A^*+B^*}{\alpha=
\alpha_1\dots\alpha_k}$.
Но по той же лемме
$\A{\alpha_i}{\E{\Tuple{\alpha}{i_1}{i_j}\in A+
B}{\alpha_i=\alpha_{i_1}\dots\alpha_{i_j}}}$.
Значит $\alpha\in (A+B)^*$.
\end{proof}
\end{frame}

\begin{frame}
\frametitle{\Ok 1.2.2. Алгебра регулярных выражений (6/8)}
\begin{proof} (Окончание)
\proofsection{$L_1\subset L_3$}
Если $\alpha\in (A+B)^*$, то по лемме 
$\E{\Tuple{\alpha}{1}{k}\in A+B}{\alpha=
\alpha_1\dots\alpha_k}$.
Но если $\alpha_i\in A+B$, 
то $\alpha_i\in A^*B^*$ 
и значит $\alpha\in (A^*B^*)^*$.
\proofsection{$L_3\subset L_1$}
Если $\alpha\in (A^*B^*)^*$, то по лемме 
$\E{\Tuple{\alpha}{1}{k}\in A^*B^*}{\alpha=
\alpha_1\dots\alpha_k}$.
Далее, если $\alpha_i \in A^*B^*$,
то $\E{\beta_i\in A^{k_1},
\gamma_i \in B^{k_2}}{\alpha_i = \beta_i \gamma_i}$.
Но слова $\beta_i$ и $\gamma_i$
также можно разложить по лемме:
$\E{\Tuple{\beta}{i_1}{i_{k_1}}\!\in\!A}{\beta_i=\beta_{i_1}\dots\beta_{i_{k_1}}}$,
$\E{\Tuple{\gamma}{i_1}{i_{k_2}}\!\in\!B}{\gamma_i=\gamma_{i_1}\dots\gamma_{i_{k_2}}}$,
 и значит слово $\alpha_i$
 разложено в конкатенацию слов,
 каждое из которых принадлежит $A+B$,
 откуда $\alpha\in (A+B)^*$. 
\end{proof}

\begin{corollary} {\NB{[1]}}
$(AB)^*A = A(BA)^*$.
\end{corollary}

\begin{proof} Раскрасим буквы в доказываемом равенстве, чтобы было видно, какому языку принадлежит какое слово:
$(A\textcolor{blue}{B})^*\textcolor{red}{A}=A(BA)^*$.
Если $\alpha\in(A\textcolor{blue}{B})^*\textcolor{red}{A}$,
то  $\alpha=\alpha_1\textcolor{red}{\alpha_2}$, 
причём $\alpha_1\in(A\textcolor{blue}{B})^*\And \alpha_2\in \textcolor{red}{A}$. 
Далее, если $\alpha_1\in(A\textcolor{blue}{B})^*$, 
то по лемме\\ 
$\E{\Tuple{\alpha_{1}}{1}{k}\in A\textcolor{blue}{B}}{\alpha_{1}=\alpha_{1_1}\ldots\alpha_{1_k}}$, 
причём
 $\A{i}{\alpha_{1_i}=\gamma_i\textcolor{blue}{\beta_i}\And
\gamma_i\in A \And \textcolor{blue}{\beta_i}\in \textcolor{blue}{B}}$.\\ 
Таким образом, $\alpha=\gamma_1\textcolor{blue}{\beta_1}\ldots\gamma_k\textcolor{blue}{\beta_k}\textcolor{red}{\alpha_2}\And\forall\gamma,\textcolor{red}{\alpha}\in A,\forall\textcolor{blue}{\beta} \in \textcolor{blue}{B}$. По ассоциативности конкатенации: $\alpha=\gamma_1(\textcolor{blue}{\beta_1}\gamma_2\ldots\textcolor{blue}{\beta_{k-1}}\gamma_k\textcolor{blue}{\beta_k}\textcolor{red}{\alpha_2})$. Значит $\alpha\in A(BA)^*$. Мы доказали, что \\$(AB)^*A\subset A(BA)^*$. Вложение $A(BA)^*\subset (AB)^*A$ доказывается аналогично.
\end{proof}
\end{frame}

\begin{frame}
\frametitle{\Ok 1.2.2. Алгебра регулярных выражений (7/8)}

\begin{corollary} {\NB{[2]}}
$(A^* B)^* =(A+B)^* B + \varepsilon$.
\end{corollary}
\begin{proof}
 Пусть $\alpha=\varepsilon$. 
 Тогда $\alpha\in (A^*B)^* \And \alpha\in (A+B)^*B+\varepsilon$ и можно сразу предположить, что
 $\alpha\neq\varepsilon$.
\proofsection{$(A^*\textcolor{blue}{B})^*\subset (A+B)^*B+\varepsilon$} Если $\alpha\in(A^*\textcolor{blue}{B})^*$, то по лемме \\
$\E{\alpha_1,\ldots,\alpha_k\in A^*\textcolor{blue}{B}}
{\alpha=\alpha_1\ldots\alpha_k}$,
причём $\alpha_i=\gamma_i\textcolor{blue}{\beta_i}
\And\gamma_i\in A^{j_i}\And\textcolor{blue}{\beta_i}\in \textcolor{blue}{B}$. По лемме 
$\E{\gamma_{i_1},\ldots,\gamma_{i_{j_i}}\in A}
{\gamma_i=\gamma_{i_1}\ldots\gamma_{i_{j_i}}}$. 
В итоге $\alpha=\gamma_{1_1}\ldots\gamma_{1_{j_1}}\textcolor{blue}{\beta_1}\ldots\gamma_{k_1}\ldots\gamma_{k_{j_k}}\textcolor{blue}{\beta_k}$, где\\ 
$\gamma_s\in A, \textcolor{blue}{\beta_r} \in \textcolor{blue}{B}$. Расставим скобки и получим $\alpha=(\gamma_{1_1}\ldots\gamma_{1_{j_1}}\textcolor{blue}{\beta_1}\ldots\gamma_{k_1}\ldots\gamma_{k_{j_k}})\textcolor{blue}{\beta_k}$, 
\\$\alpha\in(A+B)^*B$ или, расширяя множество, $\alpha\in(A+B)^*B+\varepsilon$.
 \proofsection{$(\textcolor{red}{A}+\textcolor{blue}{B})^*B+\varepsilon\subset(A^*B)^*$}. Пусть $\alpha\in(\textcolor{red}{A}+\textcolor{blue}{B})^*B+\varepsilon$ и $\alpha\neq\varepsilon$, тогда $\alpha\in(\textcolor{red}{A}+\textcolor{blue}{B})^*B$. По определению конкатенации $\alpha=\textcolor{magenta}{\gamma}\beta$, где $\gamma\in (\textcolor{red}{A}+\textcolor{blue}{B})^j,\beta\in B$. 
 Слово $\textcolor{magenta}{\gamma}$ раскладываем по лемме: 
 $\E{\textcolor{magenta}{\gamma_1},\ldots,\textcolor{magenta}{\gamma_j}\in (\textcolor{red}{A}+\textcolor{blue}{B})}{\textcolor{magenta}{\gamma}=\textcolor{magenta}{\gamma_1}\ldots\textcolor{magenta}{\gamma_j}}$, $\alpha=\textcolor{magenta}{\gamma_1}\ldots\textcolor{magenta}{\gamma_j}\beta$. Получили $\alpha\in(A^*B)^*$.
\end{proof}

\medskip 
\begin{corollary} {\NB{[3]}}
$(AB^*)^*=A(A+B)^*+\varepsilon$.
\end{corollary}
\begin{proof}
Упражнение.
\end{proof}

\medskip
Доказанные эквивалентности позволяют упрощать регулярные выражения. 
\end{frame}

\begin{frame}
\frametitle{\Ok 1.2.2. Алгебра регулярных выражений (8/8)}

\begin{example}
Доказать, что $(ABA(BA)^*B+AB)(A(BA)^*B)^*=A(BA)^*B$. \\
 Для начала проведем несколько преобразований\\ $(ABA(BA)^*B+AB)(A(BA)^*B)^*\overset{14}{=}(ABA(BA)^*B+A\varepsilon B)(A(BA)^*B)^*\overset{3,4}{=}$\\ $\overset{3,4}{=}(A(BA(BA)^*+\varepsilon) B)(A(BA)^*B)^*\overset{15}{=}(A((BA)^++\varepsilon) B)(A(BA)^*B)^*\overset{19}{=}$ \\
 $\overset{19}{=}(A(BA)^* B)(A(BA)^*B)^*\overset{15}{=}(A(BA)^*B)^+$\\
 Теперь осталось доказать, что $(A(BA)^*B)^+=A(BA)^*B$\\
 \proofsection{$A(BA)^*B\subset (A(BA)^*B)^+$} По определению введенной  операции $L^+$.\\
 \proofsection{$A(\textcolor{red}{B}\textcolor{green}{A})^*\textcolor{blue}{B})^+\subset A(BA)^*B$} Заметим, 
 что если пустое слово $\varepsilon$ не содержится ни в языке $A$, ни в языке $B$, то оно не содержится и в выражениях.
  Поэтому можно доказывать не $(A(\textcolor{red}{B}\textcolor{green}{A})^*\textcolor{blue}{B})^+\subset A(BA)^*B$, 
 а $(A(\textcolor{red}{B}\textcolor{green}{A})^*\textcolor{blue}{B})^*\subset A(BA)^*B+\varepsilon$.
 Пусть $\alpha\in(A(\textcolor{red}{B}\textcolor{green}{A})^*\textcolor{blue}{B})^*$. 
По лемме 
$\E{\alpha_1,\ldots,\alpha_k\!\in\!A(\textcolor{red}{B}\textcolor{green}{A})^*\textcolor{blue}{B}}
{\alpha\!=\!\alpha_1\ldots\alpha_k}$, причём
$\alpha_i\!=\!\gamma_i\theta_i\textcolor{blue}{\beta_i}$, 
$\gamma_i\!\in\! A$, 
$\theta_i\!\in\!(\textcolor{red}{B}\textcolor{green}{A})^{l_i}$,
$\textcolor{blue}{\beta_i}\!\in\!\textcolor{blue}{B}$. 
Далее $\E{\theta_{i_1},\ldots,\theta_{i_{l_i}}\!\in\!(\textcolor{red}{B}\textcolor{green}{A})}{\theta_i\!=\!\theta_{i_1}\ldots\theta_{i_{l_i}}}$. 
Наконец, $\theta_{i_j}\!=\!\textcolor{red}{\lambda_{ij}}\textcolor{green}{\mu_{ij}}$, 
где $\textcolor{red}{\lambda_{ij}}\!\in\!\textcolor{red}{B}$, 
$\textcolor{green}{\mu_{ij}}\!\in\!\textcolor{green}{A}$.
Имеем $\alpha\!=\!\gamma_1\textcolor{red}{\lambda_{11}}\textcolor{green}{\mu_{11}}\ldots\textcolor{red}{\lambda_{1l_1}}\textcolor{green}{\mu_{1l_1}}\textcolor{blue}{\beta_1}\gamma_2\ldots\textcolor{blue}{\beta_{k-1}}\gamma_k\textcolor{red}{\lambda_{k1}}\textcolor{green}{\mu_{k1}}\ldots\textcolor{red}{\lambda_{kl_k}}\textcolor{green}{\mu_{kl_k}}\textcolor{blue}{\beta_k}$, где $\gamma,\textcolor{green}{\mu}\!\in\! A$, 
$\textcolor{red}{\lambda},\textcolor{blue}{\beta}\in B$. Получили $\alpha\in A(BA)^*B$.
 \end{example}
 Особенность этого примера: мы построили такой язык $L$, что $L^+=L$. 
\end{frame}

\subsubsection{1.2.3. Языковые графы }
\label{1.2.3.}
\begin{frame}
\frametitle{\Ok 1.2.3. Языковые графы (1/5)}
Для описания языка можно использовать граф. 
\odmkey{Граф языка над
алфавитом $A$}{Граф!языка} (или \odmkey{источник}{Источник})~---
это ориентированный (мульти)граф, на дугах которого написаны буквы (или пустое слово
$\varepsilon$) и в котором выделены
\odmkey{начальные}{Узел!начальный} и
\odmkey{заключительные}{Узел!заключительный} узлы.

\smallskip
Далее
используются следующие обозначения. Если $s$ и $t$~--- узлы
источника $H$, то $\langle s, t\rangle$ обозначает путь из $s$ в
$t$, а $\alpha= V(\langle s, t\rangle)$ означает слово, полученное
конкатенацией букв, написанных на дугах пути $\langle s,
t\rangle$. Такое слово $\alpha \in A^*$ будем называть
\odmkey{значением}{Значение!пути} пути. Очевидно, что $|V(\langle
s, t \rangle)|\leq |\langle s, t\rangle|$  и $|V(\langle s, t
\rangle)| = |\langle s, t \rangle|$ тогда и только тогда, когда
$\langle s, t \rangle$  не содержит дуг, нагруженных~$\varepsilon$.

\smallskip
Говорят, что источник {\em представляет} некоторый язык, а именно,
язык, состоящий из слов, которые являются значениями всех путей,
ведущих из начальных узлов в заключительные.
Очевидно, что всякий источник представляет некоторый язык (но не
всякий язык может быть представлен графом).

\begin{remark}
Множества начальных и заключительных узлов могут пересекаться и
даже совпадать. На диаграммах начальные узлы помечают входной стрелкой
<<из ниоткуда>>, а заключительные узлы выделяют жирной границей. Мультидугу обычно изображают одной линией и указывают
сразу множество букв через запятую.
\end{remark}
\end{frame}

\begin{frame}
\frametitle{\Ok 1.2.3. Языковые графы (2/5)}

\begin{columns}[t]
\begin{column}[T]{8.7cm}
\begin{example}
Справа представлена диаграмма источника
десятичных записей натуральных чисел.
\end{example}
\end{column}
\begin{column}[T]{6.5cm}
\includegraphics[width=6.5cm]{1_2_3_2.jpg}
\end{column}
\end{columns}

\smallskip
Наряду с путями, которые ведут из начальных
узлов в заключительные узлы 
и представляют <<правильные>> слова языка,
в источнике могут быть пути,
которые не ведут из начальных узлов 
в заключительные, и тем самым представляют
<<неправильные>>  слова, то есть 
не принадлежащие языку.
Источник, в котором существуют пути для 
{\em всех} слов в алфавите, естественно назвать
\odmkey{полным}{Источник!полный}. 
Если два источника представляют один язык, то их
называют \odmkey{эквивалентными}{Источник!эквивалентный}.

\begin{theorem}
Для любого источника существует 
эквивалентный полный источник.
\end{theorem}

\begin{proof}
Если в  графе из любого узла исходят дуги, помеченные всеми буквами алфавита,
то источник, очевидно, полный.
Любой источник можно превратить в полный
следующим образом.
Нужно добавить ещё один узел <<поглотитель>> и
назначить его не заключительным. 
Для каждого из старых узлов,
в которых не все буквы алфавита использованы
на исходящих дугах,
нужно провести новые дуги в узел поглотитель и пометить их недостающими
буквами.
В самом этом узле поглотителе нужно добавить петли для всех букв.
Источник станет полным, а представляемый язык не изменится.      
\end{proof}

\end{frame}


\begin{frame}
\frametitle{\Ok 1.2.3. Языковые графы (3/5)}

\odmkey{Размером}{Размер!источника}
(\odmkey{величиной}{Величина!источника}) источника $H$ называется
количество узлов в нем (обозначение $|H|$).

\medskip
Подмножество $X\subset H$ узлов источника называется
\odmkey{замкнутым}{Подмножество!замкнутое}, если
$$
x\in X\And x
\stackrel{\varepsilon}\longrightarrow y \Impl y\in X
$$
(то~есть
в замкнутое подмножество входят все узлы, в которые можно попасть
<<даром>>~--- по слову $\varepsilon$).

\begin{remark}
Если в источнике нет пустых дуг, то все подмножества узлов
замкнутые.
\end{remark}

\medskip
Языковой граф (источник) называется
\odmkey{двухполюсным}{Граф!двухполюсный} или
\odmkey{двухполюсником}{Граф!двухполюсник}, если он имеет один
начальный и один заключительный узел. Обозначим их $s$ и $t$
соответственно.

\begin{theorem}
Для любого источника существует эквивалентный ему двухполюсник.
\end{theorem}

\begin{columns}[t]
\begin{column}[T]{9cm}
\begin{proof}
Пусть $H$~--- заданный источник, $S$~--- множество начальных
узлов, $T$~--- множество заключительных узлов. Построим
эквивалентный двухполюсный источник так, как показано на
рисунке.
\end{proof}
\end{column}
\begin{column}[T]{6cm}

\vspace{-12pt}
\begin{figure}[h]
\begin{center}
\includegraphics[width=6cm]{1_2_3_3.jpg}
\end{center}
\end{figure}
\end{column}
\end{columns}

\end{frame}

\begin{frame}
\frametitle{\Ok  1.2.3. Языковые графы (4/5)}
Введенное понятие двухполюсника исчерпывающим образом
характеризует множество языков, описываемых регулярными
выражениями.

\begin{theorem}
Для любого регулярного выражения $E$ существует двухполюсник $H$,
представляющий $E$.
\end{theorem}
\begin{proof}
Построим $H$ индукцией по структуре $E$.

Буквы: для каждой буквы  $a$, входящей в выражение $E$, построим
источник, как показано на рисунке~{\em а}.
Альтернация: построим параллельное соединение источников,  как
показано на рисунке~{\em б}.
Конкатенация: построим последовательное соединение источников,  как
показано на рисунке~{\em в}.
Итерация: построим зацикливание источников,  как показано на
рисунке~{\em г}.
\end{proof}

\begin{figure}[h]
\begin{center}
\includegraphics[width=15cm]{1_2_3_4.jpg}
\end{center}
\end{figure}

\end{frame}

\begin{frame}
\frametitle{\Ok 1.2.3. Языковые графы (5/5)}
\begin{remark}
Пустые дуги нужны для наглядности построения и простоты
доказательства. При проведении реального построения графа по
регулярному выражению, как правило, можно обойтись склейкой узлов,
как показано на рисунке. Склеенные узлы и добавленные дуги выделены красным.
\end{remark}

\begin{figure}[h]
\begin{center}
\includegraphics[width=140mm]{1_2_3_5.jpg}
\end{center}
\end{figure}

\begin{example}
На рисунке показан двухполюсник,
представляющий общепринятый язык записи 
вещественных чисел с фиксированной запятой с возможным знаком.

%\vspace{-8pt}
\begin{figure}[h]
\begin{center}
\includegraphics[width=150mm]{1_2_3_5_.jpg}
\end{center}
\end{figure}

\end{example}

\end{frame}

% подготовил Буслама Анис 20.04.2026
% Вставить после слайда 1.2.3. (1/5) и перенумеровать последующие слайды
%
% Пример тривиальный, убран 24.08.26 ФАН
%\begin{frame}
%\frametitle{1.2.3. Языковые графы (6/6)}
%
%
%\begin{example}
%Пусть $A=\{a,b\}$ и $H$~--- двухполюсник, представленный на рисунке.
%Тогда $H$ представляет язык
%\[
%L(H)=\Set{a^n b}{n\ge 0}.
%\]
%\end{example}
%\begin{figure}[h]
%\begin{center}
%\includegraphics[width=10cm]{1.2.3_1.png}
%\end{center}
%\end{figure}
%\begin{remark}
%Здесь $s$~--- начальный узел, $t$~--- заключительный узел.\\
%Путь $\langle s,s\rangle$, состоящий из $n$ петель с меткой $a$, имеет
%значение $a^n$, а путь $\langle s,t\rangle$, в котором после этих петель
%проходит дуга с меткой $b$, имеет значение $a^n b$.\\
%Тем самым источник представляет все слова, в которых после произвольного
%числа букв $a$ стоит ровно одна буква $b$.
%\end{remark}
%
%\end{frame}
%

\subsubsection{1.2.4. Детерминированные и недетерминированные источники}\label{1.2.4.}
\begin{frame}
\frametitle{1.2.4. Детерминированные и недетерминированные источники (1/7)}

Источник (граф языка) называется
\odmkey{детерминированным}{источник!детерминированный}, если он
содержит одну начальную вершину, не содержит пустых дуг и
все дуги, исходящие из одного узла, помечены разными буквами.
В противном случае,
то есть если нарушено хотя бы одно из
трёх условий, источник называется
\odmkey{недетерминированным}{источник!недетерминированный}.


\begin{example}
На рисунке слева показан детерминированный источник,
представляющий общепринятый язык записи 
вещественных чисел с фиксированной запятой,
а справа показан эквивалентный недетереминированный источник,
в котором нарушены все три условия детерминированности.


\begin{center}
\includegraphics[width=13cm]{1_2_4_1.jpg}
\end{center}


\end{example}

\end{frame}


\begin{frame}[shrink=22]
\frametitle{1.2.4. Детерминированные и недетерминированные источники (2/7)}

\begin{theorem} {\NB{[Детерминизации]}}
Для любого  источника $H$ с $n$ узлами существует эквивалентный
ему детерминированный источник $H'$, имеющий не более $2^n$ узлов.
\end{theorem}
\begin{proof}
Пусть $\bar{H}$~--- множество узлов исходного источника $H$,
причем $S\subset\bar{H}$~--- множество начальных узлов,
а $T\subset\bar{H}$~--- множество заключительных узлов.
Построим источник $H'$. Рассмотрим все замкнутые подмножества
$\bar{H}$ (включая пустое). Обозначим их $\bar{H}'\Assign 2^{[\bar{H}]}$
и будем считать их узлами источника $H'$.  Таких узлов не более
$2^n$, $|H'|\leq 2^n$. Положим $S'\in \bar{H}'$~--- наименьшее замкнутое подмножество
$\bar{H}$, содержащее все начальные узлы $H$,
$S'\Assign[S]$. Все элементы $\bar{H}'$,
содержащие хотя бы один заключительный узел $\bar{H}$, объявим
заключительными и обозначим 
$T'\Assign\Set{q'\in\bar{H}'}{\E{q\in\bar{H}}{q\in q'\And q\in T}}$.

Построим теперь дуги источника $H'$. Рассмотрим $q_i'\in \bar{H}'$
($q_i'\subset \bar{H}$) и символ $a$.
Положим $q_j'\Assign \Set{q_j\in \bar{H}}{
\E{q_i\in q_i'}{\E{ \langle q_i, q_j\rangle}{ V(\langle q_i,
q_j\rangle)=a}}}$~--- множество узлов, в которые ведут пути из
$q_i'$ со значением $a$. Если $q_j'$ не замкнуто, то положим
$q_j'\Assign [q_j']$. Если нет путей со значением $a$, исходящих из
узлов множества $q_i'$,  то положим $q_j'\Assign \emptyset$, $\emptyset\in \bar{H}'$.
Множество $q_j'$ является узлом источника $H'$, $q_j'\in\bar{H}'$.
Проведем дугу $q_i'\stackrel{a}\longrightarrow q_j'$.
Таким
образом, каждому $q_i'$ и каждому $a$ соответствует ровно одна
дуга, так что источник $H'$ детерминированный по построению.
\end{proof}
\end{frame}


\begin{frame}
\frametitle{1.2.4. Детерминированные и недетерминированные источники (3/7)}
\begin{proof} (Продолжение)
Осталось показать, что $H\sim H'$. Для этого сначала покажем, что
$
\E{\langle S', q_j'\rangle}
{V(\langle S', q_j'\rangle) = \alpha}
\Equiv
\A{q_j\in q_j'}
{\E{q_i\in S}
{V(\langle q_i, q_j \rangle)=\alpha}}.
$
Индукция по длине $\alpha$. База: если
$\alpha=\varepsilon$, то по построению $q_j' = S'$, так~как в $H'$
нет пустых дуг. Если $\alpha=a$,
то, база индукции верна по построению $H'$.
Индукционный переход: пусть
доказываемое утверждение верно для $|\alpha|\leq k$. Рассмотрим
$\alpha a$.

Пусть $\E{\langle S', q_j'\rangle}
{V(\langle S', q_j'\rangle)= \alpha a}$.
Обозначим
$q_\ell'$~--- предпоследний узел на этом пути, так что
$V(\langle S', q_\ell'\rangle)= \alpha$ и
$q_\ell'\stackrel{a}\longrightarrow q_j'$.
По индукционному предположению в источнике $H$
$\A{q_\ell\in q_\ell'}{\E{q_i\in S}{V(\langle q_i, q_\ell \rangle)=\alpha}}$.
По построению источника $H'$, если
$V(\langle q_\ell', q_j'\rangle) =a$, то
$\A{q_j\in q_j'}{\E{q_\ell\in q_\ell'}{V(\langle q_\ell, q_j\rangle)=a}}$ и, таким
образом, в источнике $H$ имеем путь $\langle q_i,\dots, q_\ell,
q_j\rangle$, причём $V(\langle q_i, q_\ell\rangle)=\alpha$ и
$V(\langle q_\ell, q_j\rangle)=a$.
Пусть теперь \\$\A{q_j\in q_j'}
{\E{q_i\in S}
{V(\langle q_i, q_j \rangle)=\alpha a}}$. Рассмотрим
множество $q_\ell'$ узлов, в которые можно попасть из $q_i$ по путям
со значением $\alpha$. Аналогичным образом имеем путь
$\langle S',\dots, q_\ell', q_j'\rangle$, причём
$V(\langle S',q_\ell'\rangle)=\alpha$ и $V(\langle q_\ell', q_j'\rangle)=a$.
Из доказанного свойства и определения заключительных узлов следует, что в источнике $H'$ путь $\alpha$  из начального узла $S'$ в какой-либо
заключительный узел существует тогда и только тогда, когда в источнике $H$ существует
путь $\alpha$ из некоторого начального узла в какой-либо
заключительный узел.
\end{proof}
\end{frame}


\begin{frame}
\frametitle{1.2.4. Детерминированные и недетерминированные источники (4/7)}
\begin{remark}
Если интерпретировать источник как алгоритм
порождения слов языка, то теорема детерминизации показывает,
что в некоторых случаях по недетерминированному алгоритму
можно построить функционально эквивалентный детерминированный.
\end{remark}

\begin{example}
В доказательстве теоремы используется построение, при котором
количество узлов экспоненциально возрастает. Однако это не
означает, что детерминированный источник всегда имеет больше
узлов, чем недетерминированный. На рисунке приведён
соответствующий пример.

\vspace{-4pt}
\begin{figure}[h]
\begin{center}
\includegraphics[width=110mm]{1_2_5_4.jpg}
\end{center}
\end{figure}
\end{example}


\end{frame}


\begin{frame}[shrink=22]
\frametitle{1.2.4. Детерминированные и недетерминированные источники (5/7)}

\begin{remark}
Детерминированный источник~--- это граф автомата без выходов
(см. \DMref{п.~2.2.1}{2.2.1.}).
Недетерминированный источник~--- это недетерминированный конечный автомат,
то~есть автомат, в котором нарушено условие
автоматной детерминированности (см. 
\DMref{п.~2.1.4}{2.1.4.}).
Если интерпретировать источник как алгоритм
порождения слов языка, то недетерминированный источник
на каждом
шаге  <<выбирает>> одно из ребер, по которому можно
перейти. Теорема детерминизации показывает,
что с точки зрения порождения языков, возможности 
этих двух классов автоматов совпадают,
хотя эти классы автоматов различны с точки зрения
программной реализации.
\end{remark}

\bigskip
\TODO{ Вопрос о программной реализации детерминированных и детерминированных 
источников обозначен, но не закончен, это открытый TODO. Нужно либо подготовить хороший атом с примером, либо перенести этот слайд в следующий семестр.}

\begin{remark}
Свойства источника, как графа, связаны со свойствами языка, как множества слов.
Например, источник конечного языка всегда можно представить в виде
леса корневых деревьев, то есть бесконтурного орграфа. 
Источник бесконечного языка, напротив, обязательно содержит хотя бы один контур.  
\end{remark}

\end{frame}


\begin{frame}
\frametitle{1.2.4. Детерминированные и недетерминированные источники (6/7)}

\begin{example}
Пусть $H$~--- детерминированный источник с единственным начальным узлом $s$
и множеством заключительных узлов $T$.
Тогда при чтении слова $\alpha=a_1a_2\dots a_n$ достаточно хранить только
текущий узел:
\[
q_0\Assign s,\qquad q_i\Assign \delta(q_{i-1},a_i),\quad 1\le i\le n.
\]
Слово $\alpha$ допускается тогда и только тогда, когда $q_n\in T$.
\end{example}

\begin{remark}
Для недетерминированного источника вместо одного текущего узла приходится
хранить множество активных узлов $X_i\subset\bar H$.
После чтения очередной буквы строится множество всех достижимых узлов
и, при необходимости, его $\varepsilon$-замыкание.
Таким образом, программная реализация детерминированного источника сводится
к одному текущему состоянию, а реализация недетерминированного~--- к
множеству активных состояний.
\end{remark}

\begin{example}
Пусть после чтения префикса $a_1a_2\dots a_i$ активны вершины множества $X_i$.
Тогда для недетерминированного источника
\[
X_0=\varepsilon(S),\qquad
X_{i+1}=\varepsilon\bigl(\mathrm{P}_{<q,p>_{i+1}}(X_i)\bigr),
\]
\end{example}
\end{frame}


\begin{frame}
\frametitle{1.2.4. Детерминированные и недетерминированные источники (7/7)}


где
\[
\mathrm{P}_{<q,p>}(X)=\Set{p}{\E{q\in X}{}}.
\]
Слово допускается, если $X_n\cap T\neq\emptyset$.

\medskip
\begin{algorithmic}
\REQUIRE входное слово $\alpha=a_1a_2\dots a_n$, источник $H=\langle \bar H,S,T,\delta\rangle$
\STATE $X\Assign\varepsilon\text{-closure}(S)$
\FOR{\textbf{$i$ from $1$ to $n$}}
  \STATE $Y\Assign\emptyset$
  \FOR{\textbf{$q\in X$}}
    \STATE $Y\Assign Y\cup\Set{p\in \bar H}{\langle q,p\rangle\in H\ \And\ V(\langle q,p\rangle)=a_i}$
  \ENDFOR
  \STATE $X\Assign\varepsilon\text{-closure}(Y)$
\ENDFOR
\STATE {\textbf{return}} $X\cap T\neq\emptyset$
\end{algorithmic}

\TODO{Два последних слайда = примерный план. ФАН 24.08.26}
\end{frame}


\subsubsection{1.2.5. Расширения языка регулярных выражений}
\label{1.2.5.}

\begin{frame}
\frametitle{1.2.5. Расширения языка регулярных выражений (1/4)}

Трёх введённых регулярных операций
(альтернация, конкатенация и замыкание Клини) {\em достаточно} 
для описания класса регулярных языков.
Однако набор регулярных операций может быть {\em расширен} многими способами,
например, введением дополнительных операций, которые не меняют
выразительной мощности регулярных выражений, но делают их более
{\em удобными}.

\smallskip
Рассмотрим операцию дополнения:
$$\overline{L}\Def
\Set{\alpha\in A^*}{\alpha\not\in L}.$$

\begin{theorem}
Если $L$~--- регулярный язык, то
его дополнение $\overline{L}$ также регулярно.
\end{theorem}
\begin{proof}
Если в полном
источнике языка $L$ заключительные узлы
объявить незаключительными и наоборот, то полученный граф будет
источником языка $\overline{L}$.
\end{proof}

\smallskip
Имея регулярную операцию дополнения, легко
вывести, что операция пересечения также регулярна:
$$
L_1\cap L_2 =
\overline{\overline{L}_1~|~\overline{L}_2}.
$$

\begin{remark}
Алгебра $\langle \cal{R}; \cup, \cap, \overline{\phantom{L}}\rangle$ является булевой алгеброй. 
\end{remark}

\end{frame}


\begin{frame}
\frametitle{1.2.5. Расширения языка регулярных выражений (2/4)}

Помимо введения дополнительных операций над регулярными языками, 
можно вводить сокращённые обозначения для регулярных выражений.

\begin{example}
\odmkey{Символы шаблона}{Символ!шаблона}, которые используются в операционных
системах для задания маски имени файла, имеют следующий смысл:
знак <<?>>, который обозначает один произвольный символ, это сокращение
для регулярного выражения $a_1 | \dots | a_n$,
а знак <<*>>, соответственно, сокращение для
$(a_1 | \dots | a_n)^*$, где
$\Tuple{a}{1}{n}$~--- все символы, которые разрешается использовать
в именах файлов.
\end{example}

\begin{remark}
Важно понимать, что конкретное значение
символов шаблона зависит от контекста,
потому что в разных операционных системах и даже
в разных версиях одной системы наборы допустимых символов различны.
\end{remark}

\smallskip
Пусть далее в этом параграфе $A=\{\Tuple{a}{1}{n}\}$~--- некоторый фиксированный конечный алфавит, причём буквы 
алфавита считаются перенумерованными, а $P$ и $Q$~--- произвольные регулярные 
выражения над этим алфавитом. Между алгеброй регулярных языков и алгеброй 
регулярных выражений установлен изоморфизм, и это означает, что к регулярным выражениям можно применять регулярные операции, получая новые регулярные выражения.        

\end{frame}


\begin{frame}
\frametitle{1.2.5. Расширения языка регулярных выражений (3/4)}
В частности, в соответствии с \DMref{п.~1.2.1}{1.2.1.} и \DMref{п.~1.2.2}{1.2.2.}\\
$ 
PQ \Def\Set{\alpha\beta}{\alpha\in P \And \beta\in Q}, \
P^1\Def P, \ 
P^k\Def PP^{k-1},\  
P^*\Def\bigcup_{i=0}^{\infty}P^i,\ P^+\Def PP^* .$

Введем дополнительное обозначение для опциональной, необязательной конструкции, которое  выражается через базовые регулярные операции:
$$
P? \Def \varepsilon \mid P.
$$

\begin{example}
$(ab)?c = c|abc.
$
\end{example}

\begin{remark}
Операции $P^*$, $P^+$ и $P?$ можно рассматривать как
\odmkey{операции повторения}{Операция!повторения}:\\
$P^*$ --- нуль или более повторений,
$P^+$ --- одно или более повторений,
$P?$ --- нуль или одно повторение $P$.
\end{remark}

Удобно ввести \odmkey{итерацию с разделителем}{Итерация!с разделителем}
(обобщённую итерацию, итерацию Цейтина\footnote{Григорий Самуилович Цейтин (1936--2022)}).
Пусть $P$ и $Q$~--- регулярные выражения:
$$
P \operatorname{\#} Q \Def P(QP)^*.
$$

\end{frame}


\begin{frame}
\frametitle{1.2.5. Расширения языка регулярных выражений (4/4)}
%\begin{example}
%$A\#B = A~|~ABA~|~ABABA~|~ABABABA~|~\dots$
%\end{example}

\vspace{-12pt}
\begin{columns}
\begin{column}{11cm}
\begin{digress}
Выразительную силу и пользу итерации Цейтина легко оценить,
взглянув на источник выражения $P \operatorname{\#} Q$.  
\end{digress}
\end{column}
\begin{column}{4cm}
\begin{center}
\includegraphics[width=2.4cm]{1_2_5_44.jpg}
\end{center}
\end{column}
\end{columns}

\smallskip
Часто используют \odmkey{символьные классы}{Класс!символьный}, то есть
подмножества символов заданного алфавита $A$. 
Подмножество можно задать перечислением элементов или же диапазоном, 
поскольку алфавит линейно упорядочен. 

\begin{example}
$[abcde] = [a-e] = (a|b|c|d|e)$. 
\end{example}

\smallskip
Символьному классу 
можно дать имя и многократно использовать это имя в других регулярных выражениях 

\begin{example} Символьный класс <<цифра>>:
$\backslash d = [0-9] = [0123456789] =  (0|1|2|3|4|5|6|7|8|9)$. 
\end{example}

\smallskip
Используют в том числе и инвертированные символьные классы:
$
[\char`\^ S] = [~\overline{S}~]
$

\begin{remark}
Операция дополнения определена только
относительно фиксированного алфавита.
Следует быть осторожными при использовании 
инвертированных символьных классов, поскольку результат
зависит от ряда трудно учитываемых факторов:
версии операционной системы, 
действующей в данный момент схемы кодирования и т.д.
\end{remark}

\begin{example}
Если $A=\{0,1,2,3\}$, то ${[\char`\^{}01]} = {2|3}$.
\end{example}

\end{frame}


%\begin{frame}
%\frametitle{1.2.5. Расширения языка регулярных выражений (5/5)}
%Ещё два популярных варианта сокращения:
%
%\smallskip
%\begin{itemize}
%\item \odmkey{точка}{Точка (.)} \texttt{.}~--- В одних программах точка является сокращенным обозначением символьного класса, пораждающего элементарный язык, совпадающий с любым элементом некоторого базового языка, а в других– с любым элементом некоторого базового языка, кроме символа новой строки.
%\item \odmkey{сокращённые классы}{Класс!сокращённый} вида
%\texttt{\textbackslash d}, \texttt{\textbackslash w}, \texttt{\textbackslash s}
%(и отрицания \texttt{\textbackslash D}, \texttt{\textbackslash W}, \texttt{\textbackslash S})
%обычно понимаются как сокращения для некоторых символьных классов.
%\end{itemize}
%\end{frame}

\subsubsection{1.2.6. Программирование в регулярных выражениях}
\label{1.2.6.}

\begin{frame}
\frametitle{1.2.6. Программирование в регулярных выражениях (1/5)}
Регулярные выражения являются элегантным и эффективным средством
решения многих задач обработки текстов и встроены в большинство
унаследованных и современных систем программирования.

\bigskip 
\TODO{Посмотреть POSIX, PERL, SNOBOL4, REFAL. Понять основные идеи. \\
Эффективный алгоритм сопоставления с образцом = парсинг с присваиваниями. \\
Обусловленные (условные, с условиями) регулярные выражения. \\
Обязательно хороший пример. Я люблю такой 08.01.2026 => 2026, Jan, 8 \\ 
но можно взять любой, только понятный до конца и без ошибок в коде.\\
Если это пример вдобавок будет выдавать :) в ответ на дату 31.06.2026, то честь и слава. }

\end{frame}

\begin{frame}
\frametitle{1.2.6. Программирование в регулярных выражениях (1/5)}
\TODO{ Замечания ФАН 20 января. \\
В целом хорошо = 4 * 4 = 16. \\
Слить слайды 1 и 2.\\
Расшифровки всех аббревиатур. \\
Извлечение значений для переменных = присваивание переменным значений, это важнейшее место сказано коряво. \\
Backtracing = поиск с возвратами. Есть великолепные графики, показывающие сравнительную производительность юниксоидных регулярных выражений и питона, из которых хорошо видно, что питон = отстой. Это неплохо бы дать в отступлении. \\
В примерах надо объяснять (кратко) все буквы и имена методов}




Регулярные выражения являются элегантным и эффективным средством
решения многих задач обработки текстов и встроены в большинство
унаследованных и современных систем программирования.
\end{frame}


% Текст слайда подготовил Грушин Андрей 13.01.2026
\begin{frame}
\frametitle{1.2.6. Программирование в регулярных выражениях (2/5)}
Исторически и практически регулярные выражения развивались в разных традициях.

\smallskip
\begin{itemize}
\item \odmkey{POSIX}{POSIX} (BRE/ERE) --- стандарт для многих UNIX-утилит
(\texttt{grep}, \texttt{sed}, \texttt{awk}):
акцент на переносимости и предсказуемости.

\smallskip
\item \odmkey{PERL/PCRE}{PERL@PERL/PCRE} --- <<расширенные>> регулярные выражения: захваты, подстановки, дополнительные проверки контекста, обратные ссылки и др.

\smallskip
\item \odmkey{SNOBOL4}{SNOBOL4} --- сопоставление с образцом является центральной
идеей языка: образцы можно комбинировать, успешное сопоставление
естественно связывает переменные с найденными фрагментами.

\smallskip
\item \odmkey{REFAL}{REFAL} --- сопоставление и переписывание выражений по правилам: акцент не на поиске в строке, а на систематических преобразованиях структур.
\end{itemize}
\end{frame}


% Текст слайда подготовил Грушин Андрей 13.01.2026
\begin{frame}
\frametitle{1.2.6. Программирование в регулярных выражениях (3/5)}
В программировании важная задача при работе с регулярными выражениями ---
\odmkey{сопоставление с образцом}{Сопоставление с образцом}:
проверка соответствия строки или её фрагмента некоторому образцу,
и извлечение значений для переменных.

\smallskip
Существует два подхода:
\begin{itemize}
\item Автоматный подход (POSIX): регулярное выражение
преобразуется в автомат, сопоставление выполняется проходом по слову.

\item Backtracking (Perl): перебираются
альтернативы и разбиения, фиксируется первый успешный разбор.
\end{itemize}

\begin{remark}
Если выражение допускает несколько разборов одной строки, то результат
определяется правилами сопоставления в выбранной системе.
\end{remark}

\begin{example}
Выражение \texttt{(a|aa)}, строка \texttt{'aa'}.\\
PERL (backtracking): совпадение \texttt{'a'}.
POSIX (leftmost-longest): совпадение \texttt{'aa'}.
\end{example}
\end{frame}


% Текст слайда подготовил Грушин Андрей 13.01.2026
\begin{frame}
\frametitle{1.2.6. Программирование в регулярных выражениях (4/5)}
В некоторых системах регулярных выражений встречаются
\odmkey{условные регулярные выражения}{Условные регулярные выражения} ---
конструкции, позволяющие выбирать один из образцов в зависимости от условия:
\[
\texttt{(?(cond)\;yes\;|\;no)}.
\]

\begin{remark}
Условные конструкции не входят в POSIX-стандарт.
Они подчёркивают, что на практике регулярные выражения нередко дополняются
средствами управления ходом сопоставления.
\end{remark}

\begin{example}
Требуется распознавать время в одном из форматов: \texttt{HH} или \texttt{HH:MM},
при условии: минуты обязательны, если присутствует двоеточие:
\[
\texttt{(\textbackslash d\textbackslash d)(:)?(?(2)\textbackslash d\textbackslash d)}.
\]
Группа 2 соответствует символу \texttt{':'}.
Условная часть \texttt{(?(2)\textbackslash d\textbackslash d)} требует две цифры минут, если двоеточие было сопоставлено.

\smallskip
Строки \texttt{'09'} и \texttt{'09:30'} подходят, а \texttt{'09:'} и \texttt{'0930'} --- нет.
\end{example}
\end{frame}


% Текст слайда подготовил Грушин Андрей 14.01.2026
\begin{frame}
\frametitle{1.2.6. Программирование в регулярных выражениях (5/5)}
\begin{example}
Строка задаёт дату в формате
\texttt{'DD.MM.YYYY'}. Требуется преобразовать её к виду \texttt{'YYYY, Mon, D'},
при некорректной дате вывести \texttt{':)'}.
\begin{algorithmic}
\STATE \textbf{Input:} $s$

\STATE $R_{\mathrm{valid}} \Assign$ \texttt{\^{}(?:}
\ \texttt{31\textbackslash.(?:01|03|05|07|08|10|12)\textbackslash.\textbackslash d\{4\}}
\STATE \quad \texttt{|30\textbackslash.(?:01|03|04|05|06|07|08|09|10|11|12)\textbackslash.\textbackslash d\{4\}}
\STATE \quad \texttt{|29\textbackslash.(?:01|03|04|05|06|07|08|09|10|11|12)\textbackslash.\textbackslash d\{4\}}
\STATE \quad \texttt{|(?:0[1-9]|1\textbackslash d|2[0-8])\textbackslash.(?:0[1-9]|1[0-2])\textbackslash.\textbackslash d\{4\}}\texttt{)\$}

\STATE $R_{\mathrm{extract}} \Assign$
\texttt{\^{}0?(\textbackslash d+)\textbackslash.0?(\textbackslash d+)\textbackslash.(\textbackslash d\{4\})\$}

\STATE $Mon[1..12] \Assign
[\texttt{Jan},\texttt{Feb},\texttt{Mar},\texttt{Apr},\texttt{May},\texttt{Jun},\texttt{Jul},\texttt{Aug},\texttt{Sep},\texttt{Oct},\texttt{Nov},\texttt{Dec}]$

\IF{$match(R_{\mathrm{valid}}, s)$}
    \STATE $(d, m_{\mathrm{str}}, y)\Assign capture(R_{\mathrm{extract}}, s)$
    \STATE $m\Assign toInt(m_{\mathrm{str}})$
    \STATE \textbf{print} $y\texttt{, }Mon[m]\texttt{, }d$
\ELSE
    \STATE \textbf{print} \texttt{:)}
\ENDIF
\end{algorithmic}
\end{example}
\end{frame}


\subsubsection{1.2.7. Уравнения с регулярными коэффициентами}\label{1.2.7.}

% Текст слайда подготовил Зеленов Никита 10.04.2026
\begin{frame}
\frametitle{1.2.7. Уравнения с регулярными коэффициентами (1/7)}
В алгебре регулярных выражений можно 
составлять и решать {\em уравнения}, пользуясь 
свойствами регулярных операций, установленными
в \DMref{п.~1.2.2}{1.2.2.}.
Коэффициентами и неизвестными таких
уравнений являются регулярные множества 
слов, поэтому их называют
\odmkey{уравнениями с  регулярными коэффициентами}{Уравнение!с регулярными коэффициентами}.

\begin{example}
Уравнение $X=aX+b$ имеет решение $X=a^*b$.\\
Действительно, $X=a^*b=(\varepsilon+aa^*)b
=\varepsilon b +aa^*b = b+aX=aX+b$.
\end{example}

\begin{digress}
Регулярные операции обладают многими 
удобными свойствами, но они не имеют
обратных операций, 
а потому <<школьные>> приёмы
решения алгебраических уравнений,
наподобие правила 
<<перенести неизвестные в левую часть>>
оказываются неприменимыми.
\end{digress}

В уравнении коэффициентами могут 
быть не только отдельные символы или конкретные слова,
как в предыдущем примере, 
но и любые регулярные выражения.
В этом случае решение ищется в форме
регулярного выражения над коэффициентами
уравнения.

\begin{example}
Уравнение $X=EX+B$ имеет решение $X=E^*B$
при любых регулярных коэффициентах $E$ и $B$.
\end{example} 

Однако даже для линейного уравнения решение
может не быть единственным.
\end{frame}

\begin{frame}
\frametitle{1.2.7. Уравнения с регулярными коэффициентами (2/7)}
\begin{example}
Если в уравнении $X=EX+B$
окажется, что $\varepsilon\in E$
(т.~е. $E^+=E^*$), то
выражение $X=E^*(B+C)$ 
является решением для произвольного $C$.  
Действительно, $EX +B= EE^*(B+C) +B =
E^+(B+C) +B= (E^+B+B)+E^+C =
(E^++\varepsilon)B + E^+C$,
и учитывая, что $E^+=E^++\varepsilon=E^*$  в данном случае, имеем\\ 
$(E^++\varepsilon)B + E^+C = E^*B+E^*C=E^*(B+C)=X$.
\end{example}

\begin{remark}
Причина неоднозначности: при $\varepsilon\in E$ уравнение $X=EX+B$
имеет целое семейство неподвижных точек (решений), отличающихся добавкой $E^*C$.
Ниже будет выделено \emph{наименьшее} (по включению) решение.
\end{remark}

\begin{theorem}[Лемма Ардена \footnote{Дэвид Арден, лемма Ардена(1960 год)}]
Рассмотрим уравнение $X=EX+B$ над регулярными множествами.
\begin{enumerate}
\item Множество $E^*B$ является решением.
\item $E^*B$ --- \emph{наименьшее} решение по включению.
\item Если $\varepsilon\notin E$, то решение единственно и равно $E^*B$.
\end{enumerate}
\end{theorem}
\end{frame}
\begin{frame}
\frametitle{1.2.7. Уравнения с регулярными коэффициентами (3/7)}


\begin{proof}
(1) уже проверено в примере: $E(E^*B)+B=E^+B+B=(E^++\varepsilon)B=E^*B$.

(2) Пусть $X$ --- любое решение, т.~е. $X=EX+B$.
Тогда $X\supseteq B$.
Подставляя рекуррентно, получаем
$X=EX+B\supseteq EB + B$,
далее $X\supseteq E(EB+B)+B=E^2B+EB+B$, и вообще
$X\supseteq B+EB+\cdots+E^nB$ для всех $n\ge 0$.
Переходя к объединению по $n$, имеем $X\supseteq \bigcup_{n\ge 0}E^nB=E^*B$.

(3) Пусть $\varepsilon\notin E$ и $X$ --- решение.
Покажем $X\subseteq E^*B$.
Возьмём слово $\omega\in X$ минимальной длины среди слов из $X\setminus E^*B$ (если такие есть).
Из равенства $X=EX+B$ следует: $\omega\in B$ или $\omega=uv$, где $u\in E$, $v\in X$.
Если $\omega\in B$, то $\omega\in E^0B\subseteq E^*B$ --- противоречие.
Если $\omega=uv$ с $u\in E$, то поскольку $\varepsilon\notin E$, имеем $|u|\ge 1$, значит $|v|<|\omega|$.
По минимальности $\omega$ получаем $v\in E^*B$, следовательно $\omega=uv\in EE^*B=E^+B\subseteq E^*B$ --- противоречие.
Итак, $X\subseteq E^*B$, вместе с (2) получаем $X=E^*B$.
\end{proof}

\begin{digress}
Связь с теоремой о наименьшей неподвижной точке (ср. Ахо--Ульман, т.~1, стр.~126 и далее):
уравнение $X=EX+B$ можно понимать как задачу найти неподвижную точку функции
$F(X)=EX+B$ на решётке регулярных языков.
\end{digress}
\end{frame}

\begin{frame}
\frametitle{1.2.7. Уравнения с регулярными коэффициентами (4/7)}


\begin{theorem}
Пусть $\mathcal{P}(\Sigma^*)$ упорядочено по включению.
Функция $F(X)=EX+B$ монотонна, и её наименьшая неподвижная точка равна
\[
\mu F=\bigcup_{n\ge 0}F^n(\varnothing)=\bigcup_{n\ge 0}E^nB=E^*B.
\]
\end{theorem}

\begin{proof}
Монотонность: если $X\subseteq Y$, то $EX\subseteq EY$, значит $F(X)\subseteq F(Y)$.

Обозначим $X_0=\varnothing$, $X_{n+1}=F(X_n)$.
Тогда $X_1=B$, $X_2=EB+B$, $X_3=E^2B+EB+B$ и по индукции
$X_n=\bigcup_{i=0}^{n-1}E^iB$.
Следовательно, $\bigcup_{n\ge 0}X_n=\bigcup_{i\ge 0}E^iB=E^*B$.

Проверим неподвижность:
$F(E^*B)=E(E^*B)+B=E^+B+B=E^*B$.
Пусть $X$ --- любая неподвижная точка $F$, т.~е. $X=F(X)$.
Тогда $X$ является решением уравнения $X=EX+B$, а значит по лемме Ардена
$E^*B\subseteq X$.

Итак, $E^*B$ --- наименьшая неподвижная точка.
\end{proof}
\end{frame}

\begin{frame}
\frametitle{1.2.7. Уравнения с регулярными коэффициентами (5/7)}
\odmkey{Система линейных уравнений с рег. коэффициентами}{Уравнение!система с регулярными коэффициентами}
--- это набор уравнений вида
\[
\begin{cases}
X_1 = E_{11}X_1 + E_{12}X_2 + \cdots + E_{1n}X_n + B_1,\\
\ \ \dots\\
X_n = E_{n1}X_1 + E_{n2}X_2 + \cdots + E_{nn}X_n + B_n,
\end{cases}
\]
где все $E_{ij}$ и $B_i$ --- регулярные выражения (языки), а $X_i$ --- неизвестные языки.


\begin{remark}
Удобно писать в векторной форме: $X=EX+B$, где
$X=(X_1,\dots,X_n)^T$, $B=(B_1,\dots,B_n)^T$,
а $E=(E_{ij})$ --- матрица регулярных коэффициентов.
\end{remark}

\begin{theorem}
Система $X=EX+B$ имеет наименьшее (по покомпонентному включению) решение,
равное наименьшей неподвижной точке монотонной функции $F(X)=EX+B$.
Это решение получается итерациями Клини:
$X^{(0)}=\varnothing$, $X^{(k+1)}=F(X^{(k)})$, и
$X^{(\infty)}=\bigcup_{k\ge 0}X^{(k)}$.
\end{theorem}

\begin{proof}
Функция $F$ монотонна покомпонентно (как и в одномерном случае).
На полной решётке $\mathcal{P}(\Sigma^*)^n$ существует наименьшая неподвижная точка,
равная объединению итераций от $\varnothing$.
Это объединение даёт вектор языков, который удовлетворяет $X=EX+B$ и минимален среди решений.
\end{proof}
\end{frame}

\begin{frame}
\frametitle{1.2.7. Уравнения с регулярными коэффициентами (6/7)}
\begin{example}
Рассмотрим язык
\[
L=1(0\mid 1)^*,
\]
то есть все двоичные записи натуральных чисел без ведущих нулей (кроме $0$).
Построим систему уравнений по автомату:

\smallskip

начальное состояние $q_0$ --- ещё ничего не прочитано;\\
состояние $q_1$ --- прочитана начальная $1$, далее допускаются любые $0/1$;\\
переходы: из $q_0$ по $1$ в $q_1$, из $q_1$ по $0$ и $1$ в $q_1$;\\
принимающее состояние: только $q_1$.


Обозначим $X_0$ --- язык слов, ведущих из $q_0$ в принимающее состояние,
$X_1$ --- язык слов, ведущих из $q_1$ в принимающее состояние.
Тогда:
\[
\begin{cases}
X_0 = 1X_1,\\
X_1 = 0X_1 + 1X_1 + \varepsilon.
\end{cases}
\]
\end{example}
\end{frame}

\begin{frame}
\frametitle{1.2.7. Уравнения с регулярными коэффициентами (7/7)}

Во втором уравнении сгруппируем:
$X_1=(0+1)X_1+\varepsilon$.
Так как $\varepsilon\notin (0+1)$, по лемме Ардена имеем
\[
X_1=(0+1)^*\varepsilon=(0+1)^*=(0\mid 1)^*.
\]
Подставляя в первое уравнение, получаем
\[
X_0=1X_1=1(0\mid 1)^*.
\]


\begin{remark}
Если нужно включить число $0$, то язык записей натуральных чисел будет
$L_0=\{0\}+1(0\mid 1)^*$.
\end{remark}

\end{frame}

% Текст слайда подготовил Зеленов Никита 05.05.2026
\begin{frame}
\frametitle{1.2.7. Уравнения с регулярными коэффициентами (1/6)}
В алгебре регулярных выражений можно 
составлять и решать {\em уравнения}, пользуясь 
свойствами регулярных операций, установленными
в \DMref{п.~1.2.2}{1.2.2.}.
Коэффициентами и неизвестными таких
уравнений являются регулярные множества 
слов, поэтому их называют
\odmkey{уравнениями с  регулярными коэффициентами}{Уравнение!с регулярными коэффициентами}.

\begin{example}
Уравнение $X=aX+b$ имеет решение $X=a^*b$.\\
Действительно, $X=a^*b=(\varepsilon+aa^*)b
=\varepsilon b +aa^*b = b+aX=aX+b$.
\end{example}

\begin{digress}
Регулярные операции обладают многими 
удобными свойствами, но они не имеют
обратных операций, 
а потому <<школьные>> приёмы
решения алгебраических уравнений,
наподобие правила 
<<перенести неизвестные в левую часть>>
оказываются неприменимыми.
\end{digress}

В уравнении коэффициентами могут 
быть не только отдельные символы или конкретные слова,
как в предыдущем примере, 
но и любые регулярные выражения.
В этом случае решение ищется в форме
регулярного выражения над коэффициентами
уравнения.

\begin{example}
Уравнение $X=EX+B$ имеет решение $X=E^*B$
при любых регулярных коэффициентах $E$ и $B$.
\end{example} 

Однако даже для линейного уравнения решение
может не быть единственным.
\end{frame}



\begin{frame}
\frametitle{1.2.7. Уравнения с регулярными коэффициентами (2/6)}
\begin{example}
Рассмотрим $X=EX+B$, где $E$ и $B$ --- регулярные множества слов
над алфавитом $A$.
Это уравнение равносильно $X=F(X)$, где $F(X)=EX+B$.

Если $\varepsilon\in E$ (т.~е. $E^+=E^*$), то
выражение $X=E^*(B+C)$ 
является решением для произвольного множества $C$.  
Действительно, $EX +B= EE^*(B+C) +B =
E^+(B+C) +B= (E^+B+B)+E^+C =
(E^++\varepsilon)B + E^+C$,
и учитывая, что $E^+=E^++\varepsilon=E^*$  в данном случае, имеем 
$(E^++\varepsilon)B + E^+C = E^*B+E^*C=E^*(B+C)=X$.
\end{example}

\begin{remark}
Таким образом, уравнение $X=EX+B$ можно рассматривать как поиск
\emph{неподвижной точки} $F(X)=EX+B$, то есть множества $X$, для которого
$F(X)=X$.

При $\varepsilon\in E$ уравнение имеет целое семейство решений,
отличающихся добавкой $E^*C$.
\end{remark}

\begin{theorem} {\NB{[Лемма Ардена\footnote{Дэвид Арден, ХХ в.}]} }
Рассмотрим уравнение $X=EX+B$ над регулярными множествами.
\begin{enumerate}
\item Множество $E^*B$ является решением.
\item $E^*B$ --- \emph{наименьшее} решение по включению.
\item Если $\varepsilon\notin E$, то решение единственно и равно $E^*B$.
\end{enumerate}
\end{theorem}
\end{frame}



\begin{frame}
\frametitle{1.2.7. Уравнения с регулярными коэффициентами (3/6)}


\begin{proof}
\proofsection{Утв. 1}
Это уже проверено в примере: $E(E^*B)+B=E^+B+B=(E^++\varepsilon)B=E^*B$.

\proofsection{Утв. 2} 
Пусть $X$ --- любое решение, т.~е. $X=EX+B$.
Тогда $X\supseteq B$.
Подставляя рекуррентно, получаем
$X=EX+B\supseteq EB + B$,
далее $X\supseteq E(EB+B)+B=E^2B+EB+B$, и вообще
$X\supseteq B+EB+\cdots+E^nB$ для всех $n\ge 0$.
Переходя к объединению по $n$, имеем $X\supseteq \bigcup_{n\ge 0}E^nB=E^*B$.

\proofsection{Утв. 3} 
Пусть $\varepsilon\notin E$ и $X$ --- решение.
Покажем $X\subseteq E^*B$.
Возьмём слово $\omega\in X$ минимальной длины среди слов из $X\setminus E^*B$ (если такие есть).
Из равенства $X=EX+B$ следует: $\omega\in B$ или $\omega=\alpha\beta$, где $\alpha\in E$, $\beta\in X$.
Если $\omega\in B$, то $\omega\in E^0B\subseteq E^*B$ --- противоречие.
Если $\omega=\alpha\beta$ с $\alpha\in E$, то поскольку $\varepsilon\notin E$, имеем $|\alpha|\ge 1$, значит $|\beta|<|\omega|$.
По минимальности $\omega$ получаем $\beta\in E^*B$, следовательно $\omega=\alpha\beta\in EE^*B=E^+B\subseteq E^*B$ --- противоречие.
Итак, $X\subseteq E^*B$, вместе с (2) получаем $X=E^*B$.
\end{proof}

\begin{digress}
Связь с теоремой о наименьшей неподвижной точке:
уравнение $X=EX+B$ можно понимать как задачу найти неподвижную точку функции
$F(X)=EX+B$ на решётке регулярных языков (см. \DMref{п.~0.3.5}{0.3.5.}).
\end{digress}
\end{frame}

\begin{frame}
\frametitle{1.2.7. Уравнения с регулярными коэффициентами (4/6)}


\begin{theorem}
Пусть $\mathcal{P}(A^*)$ упорядочено по включению.
Функция $F(X)=EX+B$ монотонна, и её наименьшая неподвижная точка (см.\DMref{п.~0.3.5}{0.3.5.}) равна
\[
\mu F=\bigcup_{n\ge 0}F^n(\varnothing)=\bigcup_{n\ge 0}E^nB=E^*B.
\]
\end{theorem}

\begin{proof}
Монотонность: если $X\subseteq Y$, то $EX\subseteq EY$, значит $F(X)\subseteq F(Y)$.

Обозначим $X_0=\varnothing$, $X_{n+1}=F(X_n)$.
Тогда $X_1=B$, $X_2=EB+B$, $X_3=E^2B+EB+B$ и по индукции
$X_n=\bigcup_{i=0}^{n-1}E^iB$.
Следовательно, $\bigcup_{n\ge 0}X_n=\bigcup_{i\ge 0}E^iB=E^*B$.

Проверим неподвижность:
$F(E^*B)=E(E^*B)+B=E^+B+B=E^*B$.
Пусть $X$ --- любая неподвижная точка $F$, т.~е. $X=F(X)$.
Тогда $X$ является решением уравнения $X=EX+B$, а значит по лемме Ардена
$E^*B\subseteq X$.

Итак, $E^*B$ --- наименьшая неподвижная точка.
\end{proof}
\end{frame}


\begin{frame}
\frametitle{1.2.7. Уравнения с регулярными коэффициентами (5/6)}
\odmkey{Система линейных уравнений с рег. коэффициентами}{Уравнение!система с регулярными коэффициентами}
--- это набор уравнений вида
\[
\begin{cases}
X_1 = E_{11}X_1 + E_{12}X_2 + \cdots + E_{1n}X_n + B_1,\\
\ \ \dots\\
X_n = E_{n1}X_1 + E_{n2}X_2 + \cdots + E_{nn}X_n + B_n,
\end{cases}
\]
где все $E_{ij}$ и $B_i$ --- регулярные выражения (языки), а $X_i$ --- неизвестные языки.

\begin{remark}
Удобно писать в векторной форме: $\overline{X}=\overline{E}$ $\overline{X}+\overline{B}$, где
$\overline{X}=(X_1,\dots,X_n)^T$, $\overline{B}=(B_1,\dots,B_n)^T$,
а $\overline{E}=(E_{ij})$ --- матрица регулярных коэффициентов.
\end{remark}

\begin{theorem}
Система $\overline{X}=\overline{E}$ $\overline{X}+\overline{B}$ имеет наименьшее решение,
равное наименьшей неподвижной точке монотонной функции $F(\overline{X})=\overline{E}$ $\overline{X}+\overline{B}$.
Это решение получается итерациями Клини:
$\overline{X}^{(0)}=\varnothing$, $\overline{X}^{(k+1)}=F(\overline{X}^{(k)})$, и
$\overline{X}^{(\infty)}=\bigcup_{k\ge 0}\overline{X}^{(k)}$.
\end{theorem}

\begin{proof}
Функция $F$ монотонна покомпонентно (как и в одномерном случае).
На полной решётке $\mathcal{P}(A^*)^n$ существует наименьшая неподвижная точка,
равная объединению итераций от $\varnothing$.
Это объединение даёт вектор языков, который удовлетворяет $\overline{X}=\overline{E}$ $\overline{X}+\overline{B}$ и минимален среди решений.
\end{proof}
\end{frame}







\begin{frame}
\frametitle{1.2.7. Уравнения с регулярными коэффициентами (6/6)}

\begin{example}
Рассмотрим язык $L=1(0\mid 1)^*$,
то есть все двоичные записи натуральных чисел без ведущих нулей (кроме $0$).

Построим систему уравнений по автомату (см. \DMref{п.~1.2.4}{1.2.4.}):

$Q=\{q_0,q_1\},  q_0 \text{ --- начальное состояние}, F=\{q_1\}$.

\vspace{-0.8em}

\begin{columns}[t]
\begin{column}{0.5\textwidth}
\[
\begin{array}{c|cc}
 & 0 & 1\\ \hline
q_0 & \varnothing & q_1\\
q_1 & q_1 & q_1
\end{array}
\]
\end{column}

\begin{column}{0.5\textwidth}
\vspace{-1em}
\begin{center}
\includegraphics[width=5.5cm]{1_2_7_1.png}
\end{center}
\end{column}
\end{columns}

\vspace{-1em}
Обозначим $X_0$ --- язык слов, ведущих из $q_0$ в состояние $q_1$,
$X_1$ --- язык слов, ведущих из $q_1$ в состояние $q_1$. Тогда:

\vspace{-0.9em}

\[
\begin{cases}
X_0 = 1X_1,\\
X_1 = 0X_1 + 1X_1 + \varepsilon
\end{cases}
\quad
\Rightarrow
\quad
X_1=(0+1)X_1+\varepsilon
\]

Так как $\varepsilon\notin (0+1)$, по лемме Ардена:
\[
X_1=(0+1)^*=(0\mid 1)^*, \quad
X_0=1X_1=1(0\mid 1)^*.
\]
\end{example}
\end{frame}


% Текст слайда подготовил Грушин Андрей 05.01.2026
% поправил ФАН 09.01.2026
\subsubsection{1.2.8. Характеристика класса регулярных языков}\label{1.2.8.}
\begin{frame}
\frametitle{1.2.8. Характеристика класса регулярных языков (1/3) }
Следующее наблюдение помогает очертить класс регулярных языков. 

\begin{lemma} {\NB{[о разрастании]}}
Пусть $L$~--- регулярный язык над алфавитом $\Sigma$.
Тогда существует такое $n\in \Bbb{N}$, что для любого слова $\omega\in L$
длины $|\omega|\ge n$ найдутся слова $\alpha,\beta,\gamma\in\Sigma^*$, 
для которых верно:
$\quad
\alpha\beta\gamma=\omega,\quad \beta\ne\varepsilon,\quad |\alpha\beta|\le n,\quad
\A{i\ge 0}{\alpha \beta^i \gamma\in L} 
$.
\end{lemma}


\begin{proof}
Так как язык $L$ регулярен, существует полный детерминированный источник
$H$ языка $L$ (см.~\DMref{п.~1.2.3}{1.2.3.}, \DMref{п.~1.2.4}{1.2.4.}). Обозначим через $n$ число узлов
источника $H$.

Рассмотрим слово $\omega\in L$, $|\omega|\ge n$, и зафиксируем путь в $H$,
порождающий слово $\omega$.
Пусть $q_k$~--- узел, в котором оказывается источник после порождения первых $k$
символов слова $\omega$ (в частности, $q_0$~--- начальный узел).
Среди узлов $q_0,q_1,\dots,q_n$ имеется совпадение:
$q_k=q_m$ для некоторых $0\le k < m \le n$ (принцип Дирихле).
Положим
$\alpha$~--- префикс слова $\omega$ длины $k$,
$\beta$~--- подслово слова $\omega$, прочитанное на участке пути от $q_k$ до $q_m$,
и $\gamma$~--- оставшийся суффикс, так что $\omega=\alpha\beta\gamma$.
Тогда $|\beta|=m-k>0$ и $|\alpha\beta|= m \le n$.

Поскольку участок пути $q_k\To{\beta}q_m$ является циклом ($q_k=q_m$),
его можно пройти $i$ раз, после чего продолжить путь по $\gamma$.
Следовательно, для любого $i\ge 0$ существует путь со значением $\alpha\beta^i\gamma$
из начального узла в заключительный, то есть $\alpha\beta^i\gamma\in L$.
\end{proof}

\end{frame}

% Текст слайда подготовил Грушин Андрей 06.01.2026
% ФАН 10.01.2026
\begin{frame}
\frametitle{1.2.8. Характеристика класса регулярных языков (2/3) }

\begin{remark}
Лемма о разрастании применима только к {\em бесконечным} языкам,
в которых длина слов не ограничена сверху.
Все конечные языки регулярны.   
\end{remark}

\begin{example}
Язык $D\Assign\Set{a^nb^n}{n\ge 1}$ не является регулярным языком.
Действительно, рассмотрим слово $\omega=a^nb^n \in D$ при некотором 
достаточно большом $n$
и допустим, что оно разбито на три части $\omega=\alpha\beta\gamma$ так,
что $|\beta|>0$. Рассмотрим слово $\alpha\beta\beta\gamma$.
Если $\beta\in a^+$ или $\beta\in b^+$, то
$\alpha\beta\beta\gamma\not\in D$, потому что нарушается баланс $a$ и $b$.
Если же $\beta\in a^+b^+$, то
$\alpha\beta\beta\gamma\not\in D$, потому что возникает чередование $ba$.
\end{example}

\begin{digress}
Рассмотренный пример является частным случаем 
\odmkey{языка Дика}{Язык!Дика}\footnote{Walther Franz Anton von Dyck (1856--1934)},
то есть языка правильно построенных скобочных выражений в математике.
Правильно построенные последовательности открывающих и закрывающих скобок
являются контекстно-свободным языком, поскольку порождаются грамматикой
$S\to ()\mid(S)\mid SS$, но не являются регулярным языком, то есть не могут быть 
описаны и/или распознаны никаким регулярным выражением.
В этом проявляется ограниченность регулярных выражений: такие выражения
<<не умеют считать>>, точнее говоря, умеют считать только до числа $n$~--- количества
узлов в графе-источнике языка.
\end{digress}
\end{frame}

% Текст слайда подготовил Грушин Андрей 06.01.2026
% ФАН 10.01.2026
\begin{frame}
\frametitle{1.2.8. Характеристика класса регулярных языков (3/3) }

\begin{theorem}
Язык регулярных выражений $\frak{R}$ не является регулярным языком.
\end{theorem}
\begin{proof}
От противного.
Пусть $\Sigma$  
содержит некоторый исходный алфавит 
плюс используемые метасимволы регулярных выражений $(~)~|~*$. По лемме о разрастании существует число $n$,
такое что если $|\omega|\ge n \And \omega\in \frak{R}$, то существует
такое слово $\beta$, $|\beta|>0$, что $\omega=\alpha\beta\gamma$
и $\A{i}{\alpha\beta^i\gamma\in\frak{R}}$, причём $|\alpha\beta|\le n$.
Рассмотрим слово
\[
\omega=\underbrace{( \dots ( \dots (}_{n~\text{символов}~(}\ \cdots\ 
\underbrace{)\dots )\dots)}_{n~\text{символов}~)}.
\]
Это корректное регулярное выражение, содержащее $n$ пар вложенных скобок.
Не имеет значения, что обозначено многоточием $\cdots$. 
Так как первые $n$ символов слова $\omega$ --- это открывающие скобки,
то слово $\beta$ состоит только из открывающих скобок, 
и не содержит закрывающих скобок. 
Но тогда при разрастании слова $\beta$ нарушается баланс скобок и значит
  $\alpha\beta^i\gamma\not\in\frak{R}$ для $i>1$.
  Противоречие, число $n$ не существует. 
\end{proof}

\begin{remark}
Если вложенность скобок ограничена константой, то язык математических выражений
со скобками (язык Дика) оказывается регулярным.
\end{remark}

\end{frame}




% GA_1_2025.tex, 4889–8205
\subsection{1.3. Контекстно-свободные грамматики}\label{1.3.}
\begin{frame}
\frametitle{1.3. Контекстно-свободные грамматики}

Для определения синтаксиса реальных
языков программирования и 
языков предметной области 
чаще всего используются
контекстно-свободные грамматики, (напомним
также введённое обозначение КСГ, \DMref{п.~1.1.5}{1.1.5.}). 
Это обстоятельство объясняется тем,
что выразительной силы регулярных языков иногда оказывается
недостаточно для достижения 
практических целей, а для 
контекстно-зависимых грамматик 
неизвестно достаточно эффективных алгоритмов 
синтаксического анализа.
Таким образом,
КСГ образуют <<золотую середину>>:
эти грамматики позволяют, 
с одной стороны,
формально описать <<почти весь>> 
реальный язык программирования,
и, с другой стороны,
позволяют <<достаточно эффективно>>
реализовать все необходимые 
алгоритмы работы с реальным языком программирования.

\medskip  
Поэтому изучение свойств КСГ 
и алгоритмов, связанных с этим
классом грамматик, необходимо 
для всех последующих построений. 
В этом разделе изложена {\em теория}
 контекстно\-свободных грамматик, 
которая используется в прикладных 
построениях других разделов. 
\end{frame}

\subsubsection{1.3.1. Выводимость в контекстно-свободных грамматиках}\label{1.3.1.}
\begin{frame}
\frametitle{1.3.1. Выводимость в контекстно-свободных грамматиках (1/3)}

Выводимость в КСГ является частным случаем общего
понятия выводимости, рассматриваемого в \DMref{п.~1.1.2}{1.1.2.}. Однако, поскольку форма правил вывода (продукций)
в КСГ ограничена (левая часть правила~--- это всегда одиночный
нетерминал), выводимость в КСГ обладает 
рядом дополнительных полезных свойств и допускает наглядное 
представление в виде {\em упорядоченного корневого дерева},
которое называется  \odmkey{дерево разбора}{Дерево!разбора} или  \odmkey{дерево вывода}{Дерево!вывода} 
(см. \DMref{п.~1.1.5}{1.1.5.}).


\begin{remark}
Для грамматик более общего вида, контекстно-зависимых
грамматик и грамматик с фразовой структурой (\DMref{п.~1.1.4}{1.1.4.}), 
также можно предложить графоподобные представления
выводов в этих грамматиках.
Однако такие представления оказываются
принадлежами более сложным классам объектов
(в частности, гиперорграфам), 
диаграммы которых не столь наглядны, как привычные
диаграммы корневых деревьев.
В то же время, хотя представление вывода деревом разбора
безусловно возможно в случае линейных грамматик
(поскольку в таких грамматиках левая 
часть каждого правила~---  
это одиночный нетерминал, \DMref{п.~1.1.6}{1.1.6.}),
но представление вывода в виде корневого дерева
не добавляет наглядности, поскольку диаграмма 
такого дерева всегда имеет форму 
<<ствола>>, состоящего из одной линии, к которой
подвешены терминалы.  
\end{remark}

\end{frame}

\begin{frame}
\frametitle{1.3.1. Выводимость в контекстно-свободных грамматиках (2/3)}

Напомним, что в корневом ориентированном упорядоченном дереве
\odmkey{сечением}{Сечение!дерева}
называется множество узлов,
которое содержит ровно по одному излу из каждого пути из корня
в лист. 
В частности, сам корень, а также крона, являются сечениями.


\odmkey{Дерево разбора}{Дерево!разбора}, 
или \odmkey{дерево вывода}{Дерево вывода}
в КСГ $G=\langle N, T, R, S\rangle$~---
это ориентированное упорядоченное дерево,
все узлы которого помечены 
символами из $N\cup T \cup \{\varepsilon\}$
таким образом, что если в дереве
какой-то узел помечен символом $A$,
а потомки этого узла помечены символами,
образующими слово 
$\alpha$,
то в множестве правил грамматики $R$
существует правило $A\to\alpha$.


\begin{theorem}
Если $G=\langle N, T, R, S\rangle$~--- контекстно-свободная грамматика, то 
$S\Rightarrow^*\alpha$ тогда и только тогда,
когда существует дерево вывода с корнем $s$ и кроной $\alpha\in T^*$.
\end{theorem}

\begin{proof}
 \TODO{Эссе 1.3. часть 1. Доказательство. По выводу можно построить дерево, а по дереву можно построить вывод. 
}
\end{proof}

\begin{remark} 
Теорема утверждает существование дерева вывода,
но не утверждает, что такое дерево единственно
(см. \DMref{п.~1.3.6}{1.3.6.}).
\end{remark}
\TODO{09.09.26 ФАН. Далее наслоение трёх слоёв слайдов, подготовленных в разное время. Требуется слить в один хорошо отформатированный параграф на 7--9 слайдов.}
\end{frame}

\begin{frame}
\frametitle{1.3.1. Выводимость в контекстно-свободных грамматиках (3/3)}
\TODO{Эссе 1.3. часть 2. Отступление. Аналогия с логическими исчислениями. 
\\Секвенциальная форма = сечение = узел в дереве порождения, см. 0.3.4. 
}

\TODO{Эссе 1.3. часть 3. Отступление. Аналогия с орграфами.  Правило $A\to \alpha B \beta$ порождает дугу $A\to B$. Отношение достижимости соответствует выводимости. }

Таким образом, вопрос о выводимости в КСГ
сводится к вопросу о построении дерева разбора. 

\medskip 
Понятие выводимости в КСГ полезно обобщить, введя выводимость     
из одного символа и из цепочки символов, как терминальных, так и нетерминальных.
Действительно, пусть
задана грамматика $G=\langle \mathrm{M}, \Sigma, \Pi, S \rangle$ и
произвольная цепочка $\alpha \in (\mathrm{M}\cup\Sigma)^*$.
Положим $G' \Assign \langle \mathrm{M}+S', \Sigma, \Pi+(S'\to\alpha), S' \rangle$,
то есть добавим в грамматику новую аксиому $S'$
и новое правило $S' \to\alpha$. Легко видеть, что
$\A{\omega\in\Sigma^*}{S\Rightarrow^*_G \omega \Equiv S'\Rightarrow^*_{G'} \omega}$. 
Это наблюдение позволяет считать, что
$ \alpha\Rightarrow^*_G \omega \Def S'\Rightarrow^*_{G'} \omega$.
Другими словами, любой символ или 
цепочка символов $\alpha$ в грамматике $G$ также порождают некоторый язык, а именно множество выводимых слов, поэтому правомерно использовать обозначение
$\frak{L}(\alpha)$:
$$\frak{L}(\alpha)\Def\Set{\omega\in\Sigma^*}{\alpha\Rightarrow^*_G \omega}.$$
       
\end{frame}


% СЛАЙД 1: ДОКАЗАТЕЛЬСТВО ТЕОРЕМЫ
% Текст слайда подготовил Швачко Никита 5030102/20202
% Дата 11.01.2026
\begin{frame}
\frametitle{1.3.1. Выводимость в КСГ: доказательство (4/7)}

\begin{proof}
\proofsection{Необходимость: по выводу строим дерево.}
Пусть $S = \alpha_0 \Rightarrow \alpha_1 \Rightarrow \ldots \Rightarrow \alpha_n = \alpha$~--- вывод длины $n$.
Построим последовательность деревьев $D_0, D_1, \ldots, D_n$, где крона $D_i$ есть $\alpha_i$.

Положим $D_0$~--- дерево из одной вершины, помеченной $S$.

Пусть $\alpha_i = \beta A \gamma$ и применяется правило $A \to X_1 \ldots X_k$.
Тогда $D_{i+1}$ получается из $D_i$ добавлением к листу, помеченному $A$, 
потомков $X_1, \ldots, X_k$.
Крона $D_{i+1}$ будет $\alpha_{i+1} = \beta X_1 \ldots X_k \gamma$.

Дерево $D_n$~--- искомое: его корень~--- $S$, его крона~--- $\alpha$.

\proofsection{Достаточность: по дереву строим вывод.}
Рассмотрим последовательность сечений $C_0, C_1, \ldots, C_m$, где $C_0 = \{S\}$ (корень), 
а $C_m$~--- крона дерева.
Каждое $C_{i+1}$ получается заменой нетерминала $A \in C_i$ его потомками в дереве.
Последовательность крон сечений образует вывод $S \Rightarrow^* \alpha$.
\end{proof}

\end{frame}



% СЛАЙД 2: ПРИМЕР 1 — Дерево разбора арифметического выражения
% Текст слайда подготовил Швачко Никита 5030102/20202
% Дата 11.01.2026
\begin{frame}
\frametitle{1.3.1. Пример: дерево разбора выражения (5/7)}

\begin{example}
Рассмотрим КСГ арифметических выражений 
$G = \langle N, T, R, E \rangle$, где $N = \{E, T, F\}$, $T = \{a, +, *, (, )\}$, с правилами:
\begin{align*}
(1)\ E &\to E + T, & (2)\ E &\to T, \\
(3)\ T &\to T * F, & (4)\ T &\to F, \\
(5)\ F &\to (E), & (6)\ F &\to a.
\end{align*}

\odmkey{Левосторонний вывод}{Вывод!левосторонний} цепочки $a + a * a$:
\[
E \Rightarrow E + T \Rightarrow T + T \Rightarrow F + T \Rightarrow a + T \Rightarrow a + T * F \Rightarrow a + F * F \Rightarrow a + a * F \Rightarrow a + a * a.
\]

Левый разбор: $\pi = (1, 2, 4, 6, 3, 4, 6, 6)$.

Дерево разбора отражает приоритет: $a + (a * a)$, т.е. умножение выполняется раньше.
\end{example}

\end{frame}



% СЛАЙД 3: ОТСТУПЛЕНИЕ — Аналогия с логикой и орграфами
% Текст слайда подготовил Швачко Никита 5030102/20202
% Дата 11.01.2026
\begin{frame}
\frametitle{1.3.1. Аналогии: логика и орграфы (6/7)}

\begin{digress}[Аналогия с логическими исчислениями]
В секвенциальном исчислении доказательство 
представляется деревом, где каждый узел~--- секвенция (утверждение), 
а рёбра соответствуют правилам вывода.

Аналогия: \emph{секвенциальная форма соответствует \odmkey{сечению дерева разбора}{Сечение!дерева}} 
(см. \DMref{п.~0.3.4}{0.3.4.}).

В обоих случаях корень~--- исходное утверждение (аксиома $S$), 
а листья~--- элементарные факты (терминалы).
\end{digress}

\begin{digress}[Аналогия с ориентированными графами]
Построим граф достижимости нетерминалов: 
для каждого правила $A \to \alpha B \beta$, где $B \in \mathrm{M}$, 
добавим дугу $A \to B$.

Тогда \emph{отношение достижимости} в этом графе соответствует выводимости:
если $B$ достижим из $A$ в графе, то существует $\omega$, такое что 
$A \Rightarrow^* \omega$ и $\omega$ содержит $B$.
\end{digress}

\bigskip
\TODO{Хорошая работа. \\
Семь удовлетворительных атомов. \\
Ошибки форматирования, нет картинок. }

\end{frame}



% СЛАЙД 4: ПРИМЕР 2 — Граф достижимости нетерминалов
% Текст слайда подготовил Швачко Никита 5030102/20202
% Дата 11.01.2026
\begin{frame}
\frametitle{1.3.1. Пример: граф достижимости нетерминалов (7/7)}

\begin{example}
Для грамматики арифметических выражений с правилами
\[
E \to E + T \mid T, \quad T \to T * F \mid F, \quad F \to (E) \mid a
\]
построим граф достижимости нетерминалов.

\NB{Дуги графа:}
\begin{itemize}
\item Из $E \to E + T$: дуги $E \to E$, $E \to T$
\item Из $E \to T$: дуга $E \to T$
\item Из $T \to T * F$: дуги $T \to T$, $T \to F$
\item Из $T \to F$: дуга $T \to F$
\item Из $F \to (E)$: дуга $F \to E$
\end{itemize}

Получаем: $E \leftrightarrow T \leftrightarrow F$, причём $F \to E$.
Все нетерминалы взаимно достижимы (граф сильно связен).
\end{example}

\begin{remark}
Взаимная достижимость означает: из любого нетерминала можно вывести цепочку, 
содержащую любой другой нетерминал.
\end{remark}

\end{frame}


% СЛАЙД 1: ДОКАЗАТЕЛЬСТВО ТЕОРЕМЫ
% Текст слайда подготовил Швачко Никита 5030102/20202
% Дата 11.04.2026
\begin{frame}
\frametitle{1.3.1. Выводимость в КСГ: доказательство (4/7)}

\begin{proof}
\proofsection{Необходимость.}
Пусть
$
S = \alpha_0 \Rightarrow \alpha_1 \Rightarrow \cdots \Rightarrow \alpha_n = \alpha.
$
\\
Построим по индукции деревья $D_0, D_1, \dots, D_n$ так,
чтобы крона дерева $D_i$ совпадала с цепочкой $\alpha_i$.

База индукции: $D_0$ состоит из одной вершины, помеченной $S$.

Индукционный шаг: пусть
$
\alpha_i = \beta A \gamma \Rightarrow \beta X_1 \dots X_k \gamma = \alpha_{i+1}
$
\\
по правилу $A \to X_1 \dots X_k$.
Тогда дерево $D_{i+1}$ получается из $D_i$ добавлением к листу,
помеченному $A$, потомков $X_1, \dots, X_k$.
Очевидно, что крона дерева $D_{i+1}$ есть $\alpha_{i+1}$.

Следовательно, $D_n$ есть дерево вывода с корнем $S$ и кроной $\alpha$.

\proofsection{Достаточность.}
Пусть дано дерево вывода с корнем $S$ и кроной $\alpha$.
Рассмотрим последовательность сечений $C_0, C_1, \dots, C_m$,
где $C_0=\{S\}$, а $C_m$ есть крона дерева.

Переход от $C_i$ к $C_{i+1}$ состоит в замене одного нетерминала
его потомками.
По определению дерева вывода такая замена соответствует применению
некоторого правила грамматики.
Поэтому последовательность крон сечений задаёт вывод
$
S \Rightarrow^* \alpha.
$
\end{proof}

\end{frame}



% СЛАЙД 2: ПРИМЕР 1 — Дерево разбора арифметического выражения
% Текст слайда подготовил Швачко Никита 5030102/20202
% Дата 11.04.2026
\begin{frame}
    \frametitle{1.3.1. Пример: дерево разбора выражения (5/7)}
    
    \begin{example}
    Рассмотрим КСГ арифметических выражений:
    \begin{align*}
    (1)\ E &\to E + T, & (2)\ E &\to T, \\
    (3)\ T &\to T * F, & (4)\ T &\to F, \\
    (5)\ F &\to (E),   & (6)\ F &\to a.
    \end{align*}
    
    Рассмотрим \odmkey{левосторонний вывод}{Вывод!левосторонний} цепочки
    $a + a * a$:
    \[
    E \Rightarrow E + T \Rightarrow T + T \Rightarrow F + T \Rightarrow a + T
    \Rightarrow a + T * F \Rightarrow a + F * F \Rightarrow a + a * F
    \Rightarrow a + a * a.
    \]
    
    \smallskip
{\small     Левый разбор, то есть последовательность номеров применённых правил:
   $\pi = (1, 2, 4, 6, 3, 4, 6, 6)$.}
    
    \smallskip
    \begin{columns}[T,totalwidth=\textwidth]
    \begin{column}{0.40\textwidth}
    \vspace*{-6mm}
    \begin{center}
    \includegraphics[height=37mm]{1_3_1_tree.png}
    \end{center}
    \end{column}
    \begin{column}{0.57\textwidth}
    \NB{Что видно по дереву:}
    \begin{itemize}
    \item корень помечен аксиомой $E$;
    \item крона равна слову $a + a * a$;
    \item выражение интерпретируется как $a + (a * a)$, то есть умножение
    имеет более высокий приоритет.
    \end{itemize}
    \end{column}
    \end{columns}
    \end{example}
    
    \end{frame}



% СЛАЙД 3: ОТСТУПЛЕНИЕ — Аналогия с логикой и орграфами
% Текст слайда подготовил Швачко Никита 5030102/20202
% Дата 11.04.2026
\begin{frame}
\frametitle{1.3.1. Аналогии: логика и орграфы (6/7)}
\begin{columns}[T,totalwidth=\textwidth]
\begin{column}{0.49\textwidth}
\begin{digress}[Логические исчисления]
Доказательство в секвенциальном исчислении тоже имеет вид дерева:
\begin{itemize}
\item корень~--- исходное утверждение;
\item внутренние узлы~--- результаты применения правил;
\item листья~--- элементарные факты.
\end{itemize}

Поэтому \odmkey{сечение дерева разбора}{Сечение!дерева}
играет ту же роль, что и текущая секвенциальная форма доказательства
(см. \DMref{п.~1.1.5}{1.1.5.}).
\end{digress}
\end{column}
\begin{column}{0.49\textwidth}
\begin{digress}[Ориентированные графы]
Построим граф достижимости нетерминалов:
для каждого правила $A \to \alpha B \beta$
проведём дугу $A \to B$.

\smallskip
Если в этом графе $B$ достижим из $A$, то существует вывод
\[
A \Rightarrow^* \omega,
\]
в котором цепочка $\omega$ содержит символ $B$.

Следовательно, достижимость в графе отражает возможность появления
одного нетерминала в выводе из другого.
\end{digress}
\end{column}
\end{columns}

\end{frame}



% СЛАЙД 4: ПРИМЕР 2 — Граф достижимости нетерминалов
% Текст слайда подготовил Швачко Никита 5030102/20202
% Дата 11.04.2026
\begin{frame}
\frametitle{1.3.1. Пример: граф достижимости нетерминалов (7/7)}
\begin{example}
Для грамматики арифметических выражений с правилами
\[
E \to E + T \mid T, \qquad
T \to T * F \mid F, \qquad
F \to (E) \mid a
\]
построим граф достижимости нетерминалов.
\begin{columns}[T,totalwidth=\textwidth]
\begin{column}{0.35\textwidth}
\begin{center}
\includegraphics[height=35mm]{1_3_1_graph.png}
\end{center}
\end{column}
\begin{column}{0.53\textwidth}
    \NB{Дуги графа:}
    \begin{itemize}
        \item из $E \to E + T$ получаем $E \to E$ и $E \to T$;
        \item правило $E \to T$ новых дуг не добавляет;
        \item из $T \to T * F$ получаем $T \to T$ и $T \to F$;
        \item правило $T \to F$ новых дуг не добавляет;
        \item из $F \to (E)$ получаем $F \to E$;
        \item правило $F \to a$ новых дуг не даёт.
    \end{itemize}
    
Следовательно, все три нетерминала лежат в одной сильно связной компоненте.
\end{column}
\end{columns}
\end{example}


\begin{remark}
Взаимная достижимость означает: из любого нетерминала можно вывести цепочку,
содержащую любой другой нетерминал.
\end{remark}

\end{frame}

\begin{frame}
\frametitle{1.3.1. Выводимость в контекстно-свободных грамматиках (1/6)}

Выводимость в КСГ является частным случаем общего
понятия выводимости (\DMref{п.~1.1.2}{1.1.2.}). 
Однако, поскольку форма правил вывода (продукций)
в КСГ ограничена (левая часть правила~--- это всегда одиночный
нетерминал), выводимость в КСГ обладает 
рядом дополнительных полезных свойств и допускает наглядное 
представление в виде {\em упорядоченного корневого дерева},
которое называется  \odmkey{дерево разбора}{Дерево!разбора} 
или  \odmkey{дерево вывода}{Дерево!вывода} 
(см. \DMref{п.~1.1.5}{1.1.5.}).


\begin{remark}
Для грамматик более общего вида, контекстно-зависимых
грамматик и грамматик с фразовой структурой (\DMref{п.~1.1.3}{1.1.3.}), 
известны попытки предложить графоподобные представления
выводов в этих грамматиках.
Однако такие представления оказываются
принадлежащими более сложным классам объектов
(в частности, гиперорграфам), 
диаграммы которых не столь наглядны, как привычные
диаграммы корневых деревьев.
В то же время, хотя представление вывода деревом разбора
безусловно возможно в случае линейных грамматик
(поскольку в таких грамматиках левая 
часть каждого правила~---  
это одиночный нетерминал, \DMref{п.~1.1.6}{1.1.6.}),
но представление вывода в виде корневого дерева
не добавляет наглядности, поскольку диаграмма 
такого дерева всегда имеет форму 
<<ствола>>, состоящего из одной линии, к которой
подвешены терминалы. Поэтому понятие <<дерево разбора>> присуще исключительно
контекстно-свободным грамматикам.   
\end{remark}

\end{frame}

\begin{frame}
\frametitle{1.3.1. Выводимость в контекстно-свободных грамматиках (2/6)}

Напомним (\DMref{п. 0.3.4}{0.3.4.}), что в корневом ориентированном упорядоченном дереве
\odmkey{сечением}{Сечение!дерева}
называется множество узлов,
которое содержит ровно по одному узлу из каждого пути из корня
в лист. В частности, сам корень, а также крона дерева, являются сечениями.


\odmkey{Дерево разбора}{Дерево!разбора}, 
или \odmkey{дерево вывода}{Дерево вывода}
в КСГ $G=\langle \mathrm{M}, \Sigma, \Pi, S\rangle$~---
это ориентированное упорядоченное дерево,
все узлы которого помечены 
символами из множества $\mathrm{M}\cup \Sigma \cup \{\varepsilon\}$
таким образом, что если в дереве
какой-то узел помечен символом $A$,
а потомки этого узла помечены символами,
образующими слово 
$\alpha$,
то в множестве правил грамматики $\Pi$
существует правило $A\to\alpha$.


\begin{theorem}
Если $G=\langle \mathrm{M}, \Sigma, \Pi, S\rangle$~--- 
контекстно-свободная грамматика, то 
$S\Rightarrow^*\omega$ тогда и только тогда,
когда существует дерево вывода $G$ с корнем $S$ и кроной $\omega\in \Sigma^*$.
\end{theorem}

\begin{proof} Используем строковую скобочную запись деревьев (\DMref{п.~0.3.4}{0.3.4.}).
\proofsection{Необходимость}
Пусть
$
S = \alpha_0 \Rightarrow \alpha_1 \Rightarrow \cdots \Rightarrow \alpha_n = \omega.
$
Построим  деревья $D_0, D_1, \dots, D_n$ так,
чтобы крона дерева $D_i$ совпадала с цепочкой $\alpha_i$.
База: $D_0\Assign (S)$.
Шаг построения: пусть
$
\alpha_i = \beta A \gamma \Rightarrow \beta X_1 \dots X_k \gamma = \alpha_{i+1}
$
по правилу $A \to X_1 \dots X_k$.
Тогда дерево $D_{i+1}$ получается из дерева $D_i$ так: 
$(A)\Assign(A(X_1)\dots(X_n))$
%помеченному $A$, потомков $X_1, \dots, X_k$.
Очевидно, что крона дерева $D_{i+1}$ есть $\alpha_{i+1}$.
Следовательно, $D_n$ есть дерево вывода с корнем $S$ и кроной $\omega$.
\end{proof}

\begin{remark} 
Теорема утверждает существование дерева вывода,
но не утверждает, что такое дерево единственно
(см. \DMref{п.~1.3.6}{1.3.6.}).
\end{remark}
\end{frame}

\begin{frame}
\frametitle{1.3.1. Выводимость в контекстно-свободных грамматиках (3/6)}
\begin{proof} (продолжение)
\proofsection{Достаточность}
Пусть дано дерево вывода с корнем $S$ и кроной $\omega$.
Рассмотрим последовательность сечений $C_0, C_1, \dots, C_m$,
где $C_0=\{S\}$, а $C_m$ есть крона дерева.
Переход от $C_i$ к $C_{i+1}$ состоит в замене одного нетерминала
его потомками.
По определению дерева вывода такая замена соответствует применению
некоторого правила грамматики.
Поэтому последовательность сечений задаёт вывод
$
S \Rightarrow^* \omega.
$
\end{proof}

\begin{remark} 
Таким образом,  преобразования вывода в дерево, а дерева в вывод не 
представляют алгоритмических затруднений.
Теорема утверждает существование дерева вывода,
но не утверждает, что такое дерево единственно
(см. \DMref{п.~1.3.6}{1.3.6.}).
\end{remark}

\smallskip
Понятие выводимости в КСГ полезно обобщить, введя выводимость     
из одного символа или из цепочки символов, как терминальных, так и нетерминальных.
Действительно, пусть
задана грамматика $G=\langle \mathrm{M}, \Sigma, \Pi, S \rangle$ и
произвольная цепочка $\alpha \in (\mathrm{M}\cup\Sigma)^*$.
Положим $G' \Assign \langle \mathrm{M}+S', \Sigma, \Pi+(S'\to\alpha), S' \rangle$,
то есть добавим в грамматику новую аксиому $S'$
и новое правило $S' \to\alpha$. Легко видеть, что
$\A{\omega\in\Sigma^*}{S\Rightarrow^*_G \omega \Equiv S'\Rightarrow^*_{G'} \omega}$. 
Это наблюдение позволяет считать, что
$ \alpha\Rightarrow^*_G \omega \Def S'\Rightarrow^*_{G'} \omega$ 
и
$\frak{L}(\alpha)\Def\Set{\omega\in\Sigma^*}{\alpha\Rightarrow^*_G \omega}
.$
\end{frame}

\begin{frame}
\frametitle{1.3.1. Выводимость в контекстно-свободных грамматиках (4/6)}
\begin{remark}       
Другими словами, {\em любой} символ или 
{\em любая} цепочка символов $\alpha$ в грамматике $G$ также порождают 
некоторый язык $\frak{L}(\alpha)$, который вовсе не обязан 
быть подмножеством языка $\frak{L}(G)$, даже если цепочка $\alpha$ 
входит в выводимую секвенциальную форму:
$$ \Not\left( \omega\in\frak{L}(\alpha) \And S\Rightarrow_G^*\beta\alpha\gamma 
\Impl \omega\in\frak{L}(G)\right) .$$ 
\end{remark}
\begin{example}
Арифметическое выражение, скажем $a+b*c$, допускается синтаксисом
большинства языков программирования, в частности, языка Си,
однако не является синтаксически правильной программой языка Си.  
\end{example}

\smallskip
Дерево разбора некоторой цепочки, принадлежащей заданному языку, вовсе 
не является свойством, имманентно присущим этой цепочке. Дело в том, 
что один и тот же язык может быть задан разными {\em эквивалентными} грамматиками, 
и деревья разбора одной и той же цепочки в разных грамматиках могут отличаться!
\begin{example}
Рассмотрим упрощённую грамматику выражений (см. \DMref{п.~1.1.6}{1.1.6.}) КСГ $G_1$ правилами\\
$(1)E\rightarrow E+T$,
$(2)E\rightarrow T$,
$(3)T\rightarrow T*F$,
$(4)T\rightarrow F$,
$(5)F\rightarrow a$,
$(6)F\rightarrow b$,
$(7)F\rightarrow c$\\
(для удобства правила перенумерованы) и эквивалентную КСГ
$G_2$ с правилами\\
$(1)E\rightarrow E*E$,
$(2)E\rightarrow E+E$,
$(3)E\rightarrow a$,
$(4)E\rightarrow b$,
$(5)E\rightarrow c$.

\end{example}
   

\end{frame}

\begin{frame}[shrink=22]
\frametitle{1.3.1. Выводимость в контекстно-свободных грамматиках (5/6)}
\begin{example} (продолжение)
    Рассмотрим \odmkey{левосторонний вывод}{Вывод!левосторонний} цепочки
    $a + b * c$ в КСГ $G_1$:
\[
    E \stackrel{(1)}{\Rightarrow} 
    E + T \stackrel{(2)}{\Rightarrow} 
    T + T \stackrel{(4)}{\Rightarrow} 
    F + T \stackrel{(5)}{\Rightarrow} 
    a + T \stackrel{(3)}{\Rightarrow} 
    a + T * F \stackrel{(4)}{\Rightarrow} 
    a + F * F \stackrel{(6)}{\Rightarrow} 
    a + b * F \stackrel{(7)}{\Rightarrow} 
    a + b * c,
\]
которому соответствует дерево разбора $(E(T(F(a)))(+)(T(F(b))(*)(F(c))))$,
представленное на рисунке слева. Рассмотрим также 
{\em левосторонний} вывод той же цепочки в КСГ $G_2$:
$\displaystyle{
    E \stackrel{(1)}{\Rightarrow} 
    E * E \stackrel{(2)}{\Rightarrow} 
    E + E * E \stackrel{(3)}{\Rightarrow} 
    a + E * E \stackrel{(4)}{\Rightarrow} 
    a + b * E \stackrel{(5)}{\Rightarrow} 
    a + b * c}
$, которому соответствует дерево разбора $(E(E(a)(+)(b))(*)(E(c)))$,
представленное справа.
\end{example}


\begin{columns}[T,totalwidth=\textwidth]
\begin{column}{0.6\textwidth}

\vspace{-12pt}
\begin{center}
\includegraphics[height=4cm]{1_3_1_1.jpg}
\end{center}
\end{column}
\begin{column}{0.4\textwidth}
Дерево справа ниже, имеет меньше узлов, получено быстрее,
но полученная синтаксическая структура неверно передаёт семантику,
потому что $2+2*2 = 6 \ne 8$.\\
\end{column}
\end{columns}
\end{frame}

\begin{frame}
\frametitle{1.3.1. Выводимость в контекстно-свободных грамматиках (6/6)}

\begin{digress}
Теория КСГ имеет явные аналогии с матлогикой и теорией графов.
В матлогике используются исчисления двух типов.
В исчислениях {\em гильбертовского типа} основное содержание формальной 
теории заключается в аксиомах, а правила вывода немногочисленны, обычно modus ponens;
примером является исчисление предикатов.
В исчислениях {\em генценовского типа}, напротив, основное содержание формальной 
теории заключается в правилах вывода, а аксиомы тривиальны; примером является исчисление секвенций.
Совершенно очевидна аналогия между КСГ и логическими исчислениями генценовского типа,
отсюда и прямое заимствование терминологии (<<секвенциальная форма>>).
Контекстно-свободную грамматику можно представить в виде гиперорграфа специального
вида, так называемого \odmkey{графа И/ИЛИ}{Граф!И/ИЛИ}. В таком графе
есть как обычные дуги из узла в узел, так и гипердуги, которые ведут из одного узла сразу в несколько узлов.
Если пометить узлы графа символами грамматики, то гипердуги представляют правила грамматики (продукции). В таком графе поиск вывода соответствует поиску так 
называемого <<графа решения>>, то есть гиперпути, который начинается в узле
$S$ и заканчивается в терминальных узлах. 
Следует подчеркнуть, что хотя связь между теорией КСГ и другими разделами существует,
эта связь в определённом смысле односторонняя. Матлогика и теория графов суть более общие теории, а потому их методы более универсальны,
но менее эффективны, чем специфические методы теории КСГ.    
\end{digress}
\end{frame}




\subsubsection{1.3.2. Преобразования контекстно-свободных грамматик}\label{1.3.2.}
\begin{frame}
\frametitle{1.3.2. Преобразования контекстно-свободных грамматик (1/9)}


Если $\frak{L}(G_1)=\frak{L}(G_2)$, 
то грамматики 
$G_1$ и $G_2$
называются \odmkey{эквивалентными}{Грамматика!эквивалентная} и записывают это так:
$G_1\sim G_2$.

\begin{remark}
Ясно, эквивалентность грамматик 
рефлексивна, симметрична и транзитивна,
а потому введённая терминология и обозначения
вполне оправданы.
\end{remark}

\medskip
\TODO{Эссе 1.3. Часть 4. Пример двух существенно различных, но эквивалентных
грамматик. Пример с арифметическими выражениями.
}

\medskip
Эквивалентные грамматики
порождают один язык, и с этой точки зрения
неразличимы, 
но эквивалентные грамматики
могут иметь различные свойства,
которые при применении, 
например, алгоритмов синтаксического 
анализа, могут давать различные практические результаты.   

\medskip
\TODO{ Эссе 1.3. Часть 4. Ещё пример с арифметическими выражениями.
Разные результаты вычислений, приоритет операций. 
}

Известно несколько частных приёмов,
которые позволяют придать
контекстно-свободной грамматике
желательные свойства, не меняя языка.
Эти приёмы называются
\odmkey{эквивалентными преобразованиями}{Преобразование!эквивалентное!грамматик}
грамматик.
\end{frame}


\begin{frame}
\frametitle{1.3.2. Преобразования контекстно-свободных грамматик (2/9)}

Первым приёмом является устранение бесполезных символов из грамматики.

\smallskip
Символ $X$ называется 
\odmkey{полезным}{Символ!полезный} в грамматике 
$G=\langle{\Mu,\Sigma,\Pi,S}\rangle$, если существует некоторый 
\odmkey{терминальный}{Вывод!терминальный}
вывод (по правилам грамматики $G$) вида $S \Rightarrow^* \alpha X \beta \Rightarrow^* \omega$, 
где $\omega \in \Sigma^*$.

\begin{remark}
Символ $X$ может быть как терминальным,
так и нетерминальным.
Выводимая цепочка $ \alpha X \beta$ может быть первой или последней в порождении.
\end{remark}    

\smallskip
Символ $X$ называется \odmkey{бесполезным}{Символ!бесполезный}, если он не является полезным.
Бесполезные символы не нужны в грамматике, 
и от них можно и нужно избавиться.
Для выявления бесполезных символов определим свойства, характерные для полезных символов:

\begin{enumerate}
\item Символ $X$ называется  \odmkey{порождающим}{Символ!порождающий}, 
если $\E{\omega \in \Sigma^*}{X \Rightarrow^*\omega}$. 
\item Символ $X$ называется  \odmkey{достижимым}{Символ!достижимый}, если 
$\E{\alpha, \beta\in (\Mu\cup \Sigma)^*}{S \Rightarrow^* \alpha X \beta}$ .
\end{enumerate}

Полезный символ является одновременно порождающим и достижимым. Таким образом, если удалить из грамматики все непорождающие и все недостижимые символы, то останутся  полезные символы и только они.


\end{frame}


\begin{frame}
\frametitle{1.3.2. Преобразования контекстно-свободных грамматик (3/9)}

\begin{example}
Рассмотрим грамматику 
$\langle\{S,A,B\}, \{b\}, \{S \rightarrow A\mid b,  B \rightarrow b\},S\rangle$.
Все символы, кроме $A$, являются порождающими: $b$ порождает себя, $S$ порождает $b$, $B$ порождает $b$.
Все символы, кроме $B$, являются достижимыми.
\end{example}

\smallskip
Зададим алгоритм индуктивного построения множества порождающих символов.
\begin{enumerate}
\item Каждый терминальный символ порождающий.
\item Если все символы в правой части $\alpha$ правила 
 $A\rightarrow \alpha$ порождающие, то символ $A$~--- порождающий.
\end{enumerate}

\smallskip
Зададим аналогичное индуктивное построение множества достижимых символов.
\begin{enumerate}
\item Начальный символ $S$ достижимый.
\item Если символ $A$ достижимый, то все символы в правой части $\alpha$ правила 
 $A\rightarrow \alpha$ достижимые.
\end{enumerate}

Таким образом, построив множество порождающих символов и
множество достижимых символов, и взяв пересечение этих множеств, получаем множество полезных символов данной грамматики. 
\begin{example}
$\langle\{S,A,B\}, \{b\}, \{S \rightarrow A\mid b,  B \rightarrow b\},S\rangle
\sim \langle\{S\}, \{b\}, \{S \rightarrow b\},S\rangle$.
\end{example}

\end{frame}


\begin{frame}
\frametitle{1.3.2. Преобразования контекстно-свободных грамматик (4/9)}

Из определений следует, что  любой {\em терминальный} 
вывод $S \Rightarrow^* \omega$, где $\omega\in \Sigma^*$, 
содержит только {\em полезные} символы, а значит бесполезные символы, хотя и могут присутствовать в
грамматике, но не влияют на язык, определяемый грамматикой. 

\begin{theorem}
Если $G_1$~--- произвольная контекстно-свободная 
грамматика, а грамматика $G_2$ получена из грамматики
$G_1$ удалением всех бесполезных символов и всех содержащих их правил, то $\frak{L}(G_1)=\frak{L}(G_2)$. 
\end{theorem}

\begin{remark}
Возможен вырожденный случай, когда начальный 
символ $S$ является бесполезным
в грамматике $G_1$, то есть непорождающим,
и язык, порождаемый грамматикой $G_1$~--- пустое множество. В таком случае 
грамматика $G_2$ оказывается пустой, $\langle\emptyset,
\emptyset, \emptyset, \quad\rangle$,
но заключение теоремы остаётся справедливым,
поскольку язык, определяемый пустой 
грамматикой, естественно считать пустым.
Поэтому, не ограничивая общности,
можно считать, что начальные символы
в обеих грамматиках совпадают.  
\end{remark}

\begin{proof}
Обозначим $L_1\Assign\frak{L}(G_1), L_2\Assign\frak{L}(G_2)$. 
Пусть $\alpha\in L_2$, т.е. $S\Rightarrow^*_{G_2} \alpha$.
Тогда $S\Rightarrow^*_{G_1} \alpha$, поскольку всякое
правило $G_2$ является правилом $G_1$ по построению $G_2$.
Отсюда $\alpha\in L_1$ и $L_2\subset L_1$.
Пусть $\alpha\in L_1$, т.е. $S\Rightarrow^*_{G_1} \alpha$.
Тогда $S\Rightarrow^*_{G_2} \alpha$, поскольку терминальный вывод не может содержать бесполезных символов.
Отсюда $\alpha\in L_2$ и $L_1\subset L_2$.
\end{proof}

\end{frame}


\begin{frame}
\frametitle{1.3.2. Преобразования контекстно-свободных грамматик (5/9)}

Вторым приёмом построения эквивалентных грамматик 
является \odmkey{устранение}{Устранение!правила} 
некоторых правил за счёт {\em подстановки} 
одних правил в другие правила.

\smallskip
\begin{theorem}
Пусть КСГ $G_1=\langle \Mu, \Sigma, \Pi, S\rangle$,
$A\in \Mu$, $(A\to \alpha_1 B \alpha_2)\in \Pi$,
$B \to \beta_1 |\dots|\beta_k$ и
пусть КСГ $G_2=\langle \Mu, \Sigma, \Pi_2, S\rangle$,
где $\Pi_2=R\setminus \{A\to \alpha_1 B \alpha_2\}
\cup\{A\to \alpha_1\beta_1\alpha_2|\dots|\alpha_1\beta_k\alpha_2\}$. 
Тогда $\frak{L}(G_1)=\frak{L}(G_2)$.
\end{theorem}

\begin{proof}
Обозначим $L_1\Assign\frak{L}(G_1), L_2\Assign\frak{L}(G_2)$. 
Пусть $\omega\in L_2$, т.е. $S\Rightarrow^*_{G_2} \omega$.
Тогда заменим в этом выводе 
каждое применение нового правила вида
$A\to\alpha_1\beta_i\alpha_2$ последовательным
применением двух старых правил $A\to\alpha_1 B\alpha_2$, 
$B\to\beta_i$ и получим вывод 
$S\Rightarrow^*_{G_2} \omega$.
Отсюда $\omega\in L_1$ и $L_2\subset L_1$.
Пусть теперь $\omega\in L_1$, т.~е. $S\Rightarrow^*_{G_1} \omega$.
Тогда заменим в этом выводе 
каждое применение удаленного правила 
$A\to\alpha_1 B \alpha_2$ и соответствующего
(не обязательного соседнего!) 
применения правила $B\to\beta_i$
применением нового правила $A\to\alpha_1 \beta_i\alpha_2$, 
 и получим вывод 
$S\Rightarrow^*_{G_2} \omega$.
Отсюда $\alpha\in L_2$ и $L_1\subset L_2$.
\end{proof}

\begin{corollary} Пусть в КСГ 
$G_1=\langle \Mu, \Sigma, \Pi , S \rangle$
есть некоторый вывод $A\Rightarrow^*_{G_1}\alpha$,
где $A\in \Mu$, $\alpha\in (\Mu\cup \Sigma)^*$ и
пусть КСГ $G_2=\langle \Mu, \Sigma, \Pi_2, S \rangle$,
где $\Pi_2=R\cup\{A\to \alpha\}$. \\
Тогда $\frak{L}(G_1)=\frak{L}(G_2)$.
\end{corollary}

\begin{remark}
В исчислениях любой вывод можно использовать как 
\odmkey{производное}{Правило!производное} правило.
\end{remark}

\end{frame}


\begin{frame}
\frametitle{1.3.2. Преобразования контекстно-свободных грамматик (6/9)}


Во многих случаях желательно избавляться от 
\odmkey{цепных}{Правило!цепное} правил $A\to B$ 
(см. \DMref{п.~1.1.7.}{1.1.7.}).
Цепные правила можно удалять по одному на основании предыдущей теоремы,
а можно удалить все разом, 
пользуясь следующим индуктивным построением.

\begin{enumerate}
\item Для каждого нетерминала $A$ построим
множество 
$\cal{C}(A)\Assign\Set{B\in \Mu}{A\Rightarrow^* B}$\\
нетерминалов, выводимых из нетерминала $A$ 
цепными правилами из заданного множества правил $\Pi_1$ (при этом $A\in \cal{C}(A)$).
\item Для каждого нетерминала $B$ построим
множество $\cal{B}(B)\Assign\Set{A\in \Mu}{B\in \cal{C}(A)}$\\
нетерминалов, из которых выводим нетерминал $B$ 
цепными правилами из заданного множества правил $\Pi_1$ (при этом $B\in \cal{B}(B)$) 
\item Для каждого {\em нецепного}
правила $A\to\alpha$ из заданного множества правил $\Pi_1$ 
добавим в новое множество правил $R_2$
множество правил вида $B\to \alpha$ для каждого нетерминала 
$B\in\cal{B}(A)$.
\end{enumerate}

\begin{example}
Устранение цепных 
правил в грамматике арифметических выражений. 
\begin{columns}
\begin{column}{3cm}
$E\to E+T\mid T$

$T\to T*F\mid F$

$F\to (E)\mid a$
\end{column}
\begin{column}{6cm}
$\cal{C}(E)\Assign \{E, T, F\},\cal{B}(E)\Assign \{E\}$

$\cal{C}(T)\Assign \{T, F\},\cal{B}(E)\Assign \{E, T\}$

$\cal{C}(E)\Assign \{F\},\cal{B}(E)\Assign \{E, T, F\}$

\end{column}
\begin{column}{6cm}
$E\to E+T\mid T*F \mid (E) \mid a $

$T\to T*F\mid (E) \mid a $

$F\to (E)\mid a$
\end{column}
\end{columns}
\end{example}
\end{frame}


\begin{frame}
\frametitle{1.3.2. Преобразования контекстно-свободных грамматик (7/9)}
\begin{theorem}
Пусть КСГ $G_1\!=\!\langle \Mu, \Sigma, \Pi_1, S\rangle$ 
содержит цепные правила, 
а из КСГ \\$G_2\!=\!\langle \Mu, \Sigma, \Pi_2, S\rangle$ 
цепные правила удалены указанным построением. 
Тогда $\frak{L}(G_1)\!=\!\frak{L}(G_2)$.
\end{theorem}

\begin{proof}
Обозначим $L_1\Assign\frak{L}(G_1), L_2\Assign\frak{L}(G_2)$. 
Пусть $\omega\in L_2$, т.е. $S\Rightarrow^*_{G_2} \omega$.
Тогда если в этом выводе применяется правило
$(B\to\alpha)\in R_2$, то это значит, что  
$\E{A\in G_1}{B\Rightarrow_{G_1}^* A \And (A\to\alpha)\in \Pi_1}$ (возможно, что $B=A$). 
Заменим правило из $\Pi_2$ последовательностью 
правил из $\Pi_1$ и получим вывод 
$S\Rightarrow^*_{G_1} \omega$.
Отсюда $\omega\in L_1$ и $L_2\subset L_1$.
Пусть теперь $\omega\in L_1$, т.е. $S\Rightarrow^*_{G_1} \omega$.
Тогда заменим в этом выводе 
каждое применение 
последовательности цепных правил, которая заканчивается 
нецепным правилом
$A\to\dots\to B\to\alpha$ применением
правила 
$A\to\alpha$ и 
 и получим вывод 
$S\Rightarrow^*_{G_2} \omega$.
Отсюда $\omega\in L_2$ и $L_1\subset L_2$.
\end{proof}

\medskip
Если в начальном правиле $S \rightarrow \alpha$ аксиома 
$S$ входит в правую часть, $S\in \alpha$ то правило называется \odmkey{$S$-правилом}{Правило!$S$-правилом}.
Покажем, что случай $S$-правил является частным случаем цепных правил, и к нему можно применить алгоритмы избавления от цепных правил.

Действительно, для устранения $S$-правил 
достаточно ввести новое {\em цепное}
правило \\$S_1\to S$ и объявить $S_1$ 
новой аксиомой грамматики. Множество 
выводимых терминальных цепочек 
при этом, очевидно, не изменится.
\end{frame}


\begin{frame}
\frametitle{1.3.2. Преобразования контекстно-свободных грамматик (8/9)}

К нежелательным часто относят $\varepsilon$-правила,
или \odmkey{укорачивающие}{Правило!укорачивающее} правила вида $A\to\varepsilon$.

\begin{remark}
Алгоритм устранения цепных правил может оказаться
неприменимым в случае наличия $\varepsilon$-правил.
\end{remark}

Устранить укорачивающие правила возможно следующим
построением, во многом следующим идеям предыдущих построений. 
 
\begin{enumerate}
\item Построим
множество 
$\cal{E}\Assign\Set{A\in N}{A\Rightarrow^+ \varepsilon}$
нетерминалов, из которых выводимо
пустое слово.
\item Каждое {\em неукорачивающее} 
правило из множества $R_1$ представим в виде\\ 
$A\to \alpha_0 B_1 \alpha_1 \dots
 \alpha_{k-1}B_k\alpha_k,$
где $\A{i\in 1..k}{B_i\in \cal{E}}$, 
все символы в $\alpha_i$, напротив,  
не входят в $\cal{E}$ и $k\ge 0$.
\item Для каждого неукорачивающего правила,
содержащего $k$ символов из $\cal{E}$,
в множество $R_2$ включаются  $2^k$
правил вида
$A\to \alpha_0 X_1 \alpha_1 \dots \alpha_{k-1}X_k\alpha_k$,
где $X_i$ означает либо ничего, пустое слово, либо
символ $B_i$.
Такие правила соответствуют  
всем возможным подмножествам множества 
$\{\Tuple{B}{1}{k}\}$.
\item Если $S\in\cal{E}$, то добавляется правило
$S\to \varepsilon$ для возможности порождения пустого слова.
\end{enumerate}

\end{frame}


\begin{frame}
\frametitle{1.3.2. Преобразования контекстно-свободных грамматик (9/9)}
 
Третий приём, добавление нового нетерминала, очень прост: 
часть цепочки, составляющей правую  часть  правила  исходной грамматики можно обозначить новым нетерминалом, введя дополнительное правило, определяющее  новый  нетерминал.

\begin{example} Пусть в исходной грамматике $G$ одно из  правил  имеет вид $A \rightarrow abBc$.  Обозначив $bB$ через новый  нетерминал  $C$,  заменим это правило двумя:  $A \rightarrow aCc, C \rightarrow bB$.
\end{example}


Грамматика, в которой нет бесполезных символов,
а также $\varepsilon$-правил
(за исключением разве что одного правила 
$S\to\varepsilon)$, $S$-правил и цепных правил 
называется
\odmkey{приведённой}{Грамматика!приведённая}.

\begin{theorem}
Для любой контекстно-свободной грамматики $G_1$
существует такая приведённая контекстно-свободная грамматика $G_2$,
что  $\frak{L}(G_2)=\frak{L}(G_1)$.
\end{theorem}

\begin{proof}
Существование требуемой грамматики
обеспечивается ранее перечисленными алгоритмами
устранения нежелательных элементов.
Важно соблюдать порядок преобразований. 
\begin{enumerate}
\item Устранить $\varepsilon$-правила.
\item Устранить $S$-правила.
\item Устранить цепные правила.
\item Устранить бесполезные символы. 
\end{enumerate}
\end{proof}


\end{frame}

\subsubsection{1.3.3. Нормальная форма Хомского}\label{1.3.3.}
\begin{frame}
\frametitle{1.3.3. Нормальная форма Хомского}
%
%\TODO{Нормальная форма Хомского}
%
%
%\TODO{Грицаенко / Мельник. }
%
%\bigskip
%\TODO{Материал получен давно, $\approx 15$ атомов.
%Возможно, требуется ревизия.
%{\color{blue} Варисов, Самутичев}} 

\TODO{10.09.26 ФАН. Очень простой и неплохо подготовленный параграф, который необходимо довести до совершенства,
уплотнить и отформатировать слайды так, как принято в этом учебнике. }

Говорят, что КСГ находится в 
\odmkey{нормальной форме Хомского}{Нормальная форма!Хомского}, если все правила этой грамматики
имеют одну из двух следующих форм: 
либо $A\to BC$, либо $A\to a$, \\где $A, B, C \in N, a\in T$. 


\begin{theorem}
Для любой контекстно-свободной грамматики $G_1$
существует такая контекст\-но-свободная грамматика $G_2$
в нормальной форме Хомского,
что  $\frak{L}(G_2)=\frak{L}(G_1)-\varepsilon$.
\end{theorem}
\begin{proof}
\TODO{Эссе 1.3. Часть 5. }
\end{proof}

\begin{example}
\TODO{Эссе 1.3. Часть 5. Упрощенная грамматика выражений в привычной форме\\ и в форме Хомского
}
\end{example}

\medskip
\begin{remark}
Если грамматика имеет нормальную форму Хомского,
то дерево разбора в такой грамматике
является двоичным деревом. Для 
двоичных деревьев известно множество
полезных и эффективных представлений и алгоритмов.
\end{remark}

\end{frame}


%Слайды подготовил И.В. Кан 03.05.2026 на основе материалов Е. Гребенкина от 10.01.2026
\begin{frame}
  \frametitle{1.3.3. Нормальная форма Хомского (1/4)}

  Говорят, что КСГ находится в
  \odmkey{нормальной форме Хомского}{Нормальная форма!Хомского},
  если все правила этой грамматики имеют одну из двух следующих форм:
  \[
  A \to BC \text{ или } A \to a \text{,}
  \]
  где $A, B, C \in \mathrm{M}$, $a \in \Sigma$.

  \begin{theorem}
    Для любой контекстно-свободной грамматики $G_1$ существует такая
    контекст\-но-свободная грамматика $G_2$ в нормальной форме Хомского,
    что $\mathfrak{L}(G_2)=\mathfrak{L}(G_1)-\varepsilon$.
  \end{theorem}

  \medskip
  Грамматика $G_1$ должна быть предварительно
  приведённой (см. \DMref{п.~1.3.2}{1.3.2.}) и не содержать $\varepsilon$-правил.
  Если пустое слово принадлежит языку, то допускается единственное правило
  $S\to\varepsilon$.

\end{frame}

\begin{frame}
  \frametitle{1.3.3. Нормальная форма Хомского (2/4)}

  \begin{proof}
    Пусть $G_1$ — приведённая грамматика, не содержащая $\varepsilon$-правил.
    
    Построим $G_2$ в два этапа.
    \begin{enumerate}
      \item \odmkey{Устранение смешанных правил}{Устранение!смешанных правил}.
        Для каждого терминала $a\in\Sigma$, входящего в правила длины более~$1$,
        введём новый нетерминал $X_a$ и правило $X_a\to a$.
        В исходных правилах заменим все вхождения $a$ на $X_a$.
        Теперь правила имеют вид $A\to a$ или $A\to B_1\ldots B_k$, где $B_i\in \mathrm{M}$.
      \item \odmkey{Бинаризация длинных правил}{Бинаризация!правил}.
        Для каждого правила вида $A\to B_1B_2\ldots B_k$, $k\ge 3$,
        введём новые нетерминалы $D_1,\ldots,D_{k-2}$ и заменим исходное
        правило на цепочку:
        \[
          A \to B_1 D_1,\quad D_1 \to B_2 D_2,\quad \dots,\quad
          D_{k-2} \to B_{k-1} B_k .
        \]
    \end{enumerate}
    Полученная грамматика $G_2$ содержит только правила вида $A\to BC$ и $A\to a$.
    Язык не изменился: каждое преобразование является эквивалентным.
  \end{proof}

\end{frame}


\begin{frame}
\frametitle{1.3.3. Нормальная форма Хомского (3/4)}

\begin{example}
Рассмотрим приведённую грамматику арифметических выражений (без скобок):
\[
E \to E+T \mid T, \quad T \to T*F \mid F, \quad F \to a \mid b \mid c.
\]

\begin{enumerate}
    \item Устраним терминалы в длинных правилах. Вводим $P \to +$ и $M \to *$.
    Правила принимают вид:
    $E \to EPT \mid TMF \mid a \mid b \mid c$.
    \item Разобьём длинные правила.
    $E \to EPT$ заменим на $E \to E Z_1,\; Z_1 \to PT$.
    
    $E \to TMF$ заменим на $E \to T Z_2,\; Z_2 \to MF$.
\end{enumerate}
Результирующая грамматика в НФХ:
\[
\begin{aligned}
E &\to E Z_1 \mid T Z_2 \mid a \mid b \mid c, \\
Z_1 &\to P T, \quad P \to +, \\
Z_2 &\to M F, \quad M \to *.
\end{aligned}
\]
\end{example}
\end{frame}

\begin{frame}
  \frametitle{1.3.3. Нормальная форма Хомского (4/4)}

  \begin{remark}
    Важнейшим свойством грамматик в нормальной форме Хомского является
    структура дерева разбора. Поскольку каждое правило (кроме правил вида $A\to a$) содержит
    ровно два нетерминала в правой части, дерево разбора любой цепочки
    длины $n\ge 2$ является \em{двоичным}.
  \end{remark}

  \begin{corollary}
    Если $G$ — грамматика в нормальной форме Хомского и $S\Rightarrow^*\alpha$,
    где $\alpha\in\Sigma^*$, $|\alpha|=n$, то глубина дерева разбора не менее
    $\log_2 n$.
  \end{corollary}

  \medskip
  Это свойство ограничивает ширину дерева относительно его высоты и
  активно используется при доказательстве свойств контекстно-свободных языков,
  в частности, леммы о разрастании (\DMref{п.~1.3.4}{1.3.4.}).

\end{frame}


% Текст слайда подготовил Гребенкин Егор 
% Дата 10.01.2026
\begin{frame}
\frametitle{1.3.3. Нормальная форма Хомского (2/4)}

\begin{proof}
Пусть $G_1$ — приведённая грамматика (см. \DMref{п.~1.3.2}{1.3.2.}), не содержащая \\$\varepsilon$-правил.
Построим $G_2$ в два этапа.
\begin{enumerate}
    \item \odmkey{Устранение смешанных правил}{Устранение!смешанных правил}. 
    Для каждого терминала $a\in \Sigma$, входящего в правила длины более 1, введём новый нетерминал $X_a$ и правило $X_a \to a$.
    В исходных правилах заменим все вхождения $a$ на $X_a$.
    Теперь правила имеют вид $A \to a$ или $A \to B_1 \dots B_k$, где $B_i \in \mathrm{M}$.
    \item \odmkey{Бинаризация длинных правил}{Бинаризация!правил}.
    Для каждого правила вида $A \to B_1 B_2 \dots B_k$, где $k \ge 3$, введём новые нетерминалы $D_1, \dots, D_{k-2}$ и заменим исходное правило на цепочку:
    $$
    A \to B_1 D_1, \quad D_1 \to B_2 D_2, \quad \dots, \quad D_{k-2} \to B_{k-1} B_k.
    $$
\end{enumerate}
Полученная грамматика $G_2$ содержит только правила вида $A \to BC$ и $A \to a$.
Язык не изменился, так как каждое преобразование является эквивалентным.
\end{proof}
\end{frame}
% Текст слайда подготовил Гребенкин Егор 
% Дата 10.01.2026
\begin{frame}
\frametitle{1.3.3. Нормальная форма Хомского (3/4)}

\begin{example}
Рассмотрим приведённую грамматику арифметических выражений (без скобок):
$E \to E+T \mid T*F \mid a$.
\begin{enumerate}
    \item Устраняем терминалы в длинных правилах. Вводим $P \to +$ и $M \to *$.
    Правила принимают вид:
    $E \to EPT \mid TMF \mid a$.
    \item Разбиваем длинные правила.
    Правило $E \to EPT$ заменяем на $E \to E Z_1, Z_1 \to PT$.
    Правило $E \to TMF$ заменяем на $E \to T Z_2, Z_2 \to MF$.
\end{enumerate}
Результирующая грамматика в НФХ:
$$
\begin{aligned}
E &\to E Z_1 \mid T Z_2 \mid a, \\
Z_1 &\to P T, \quad P \to +, \\
Z_2 &\to M F, \quad M \to *.
\end{aligned}
$$
\end{example}
\end{frame}

% Текст слайда подготовил Гребенкин Егор 
% Дата 10.01.2026
\begin{frame}
\frametitle{1.3.3. Нормальная форма Хомского (4/4)}

\begin{remark}
Важнейшим свойством грамматик в нормальной форме Хомского является структура дерева разбора.
Поскольку каждое правило (кроме правил вида $A\to a$) содержит ровно два нетерминала в правой части, дерево разбора любой цепочки длины $n \ge 2$ является двоичным.
\end{remark}

\begin{corollary}
Если $G$ — грамматика в нормальной форме Хомского, и $S \Rightarrow^* \alpha$, где $\alpha \in \Sigma^*$ и $|\alpha|=n$, то глубина дерева разбора не менее $\log_2 n$.
\end{corollary}

Это свойство ограничивает <<ширину>> дерева относительно его высоты,
 что используется при доказательстве свойств контекстно-свободных языков, в частности, леммы о разрастании.
 
 \bigskip
 \TODO{Необходимо слить два текста в один без повторов и в разумном порядке, распределить по слайдам, ваправить форматы и обозначения.}
\end{frame}


\subsubsection{1.3.4. Лемма Огдена}\label{1.3.4.}
\begin{frame}
\frametitle{1.3.4. Лемма Огдена}

\TODO{10.09.26. ФАН. Изменение в терминологии: эта лемма для регулярных языков называется <<лемма о разрастании>>, а для КС-языков называется <<лемма Огдена>>. Дать в начале замечание о терминологии и вставить сноску про Огдена. Проверить употребление термина везде.}

%\TODO{Материал получен, $\approx 10$ атомов.
%{\color{blue} Василевский, Счастливцев\\Митрофанова / Суриков}} 


\TODO{Эссе 1.3. Часть 6. Лемма о разрастании. (теорема Огдена). \\
Очень важно неформально объяснить, что лемма о разрастании относится к языкам, хотя формулируется и доказывается с помощью грамматик. \\
Можно сначала лемму о разрастании для регулярных языков.\\
Доказательство лучше не прямое а через нормальную форму Хомского. \\
 Обязательно пример: 
как с помощью леммы проверяется, что язык 
контекстно-свободный или нет.}
\end{frame}

% Текст слайда подготовил Гребенкин Егор 
% Дата 10.01.2026
\begin{frame}
\frametitle{1.3.4. Лемма о разрастании (1/4)}

Для контекстно-свободных языков существует аналог леммы о накачке для регулярных языков. Более сильный вариант известен как \odmkey{лемма Огдена}{Лемма!Огдена}.

\begin{theorem}
Для любого контекстно-свободного языка $L$ существует такая константа $n \in \mathbb{N}$, что любую цепочку $z \in L$, в которой отмечено не менее $n$ позиций, можно представить в виде $z = u v w x y$, так что:
\begin{enumerate}
    \item цепочка $vwx$ содержит \textbf{не более} $n$ отмеченных позиций;
    \item цепочка $vx$ содержит \textbf{хотя бы одну} отмеченную позицию;
    \item $\A{i \ge 0}{u v^i w x^i y \in L}$.
\end{enumerate}
\end{theorem}

\begin{remark}
Классическая лемма о разрастании для КС-языков (\odmkey{Bar-Hillel lemma}{Лемма!Бар-Хиллела}) является частным случаем леммы Огдена, где в цепочке отмечены все позиции.
\end{remark}
\end{frame}

% Текст слайда подготовил Гребенкин Егор 
% Дата 10.01.2026
\begin{frame}
\frametitle{1.3.4. Лемма о разрастании (2/4)}

\begin{proof}
Пусть $L=\frak{L}(G)$, где $G=\langle \mathrm{M}, \Sigma, \Pi, S\rangle$ — грамматика в нормальной форме Хомского. Пусть $k=|\mathrm{M}|$. Положим $n = 2^{k+1}$.
Рассмотрим дерево разбора цепочки $z$. Назовём узел дерева \odmkey{ветвящимся}{Узел!ветвящийся}, если оба его поддерева содержат отмеченные позиции.
Построим путь от корня к листу, каждый раз выбирая потомка, содержащего больше отмеченных позиций.
Так как отмеченных позиций $\ge 2^{k+1}$, на этом пути встретится не менее $k+1$ ветвящихся узлов.
\end{proof}
\end{frame}

% Текст слайда подготовил Гребенкин Егор 
% Дата 10.01.2026
\begin{frame}
\frametitle{1.3.4. Лемма о разрастании (3/4)}

\begin{proof} (Окончание)
На выбранном пути находится более $k$ нетерминалов. \\Согласно принципу Дирихле, найдётся нетерминал $A$, который встречается на пути дважды (как $A_1$ и $A_2$, где $A_2$ — потомок $A_1$).
Это позволяет разбить вывод на части:
$$
S \Rightarrow^* u A_1 y, \quad A_1 \Rightarrow^* v A_2 x, \quad A_2 \Rightarrow^* w.
$$
В силу того, что $A_1$ — ветвящийся узел, в поддеревьях, соответствующих $v$ и $x$, содержится как минимум одна отмеченная позиция, следовательно $vx \neq \varepsilon$.
Повторяя вывод $A \Rightarrow^* v A x$ произвольное число раз (или пропуская его), получаем цепочки $u v^i w x^i y$, принадлежащие языку $L$.
\end{proof}
\end{frame}

% Текст слайда подготовил Гребенкин Егор 
% Дата 10.01.2026
\begin{frame}
\frametitle{1.3.4. Лемма о разрастании (4/4)}

\begin{example}
Докажем, что язык $L = \Set{a^m b^m c^m}{m \ge 1}$ не является контекстно-свободным.
Предположим противное. Пусть $n$ — константа леммы Огдена.
Рассмотрим цепочку $z = a^n b^n c^n$. Отметим в ней все позиции с символом $a$.
Разложим $z = u v w x y$.
Согласно лемме, $vx$ обязан содержать отмеченные символы (то есть $a$).
В то же время, $vwx$ содержит не более $n$ отмеченных символов.
Если мы <<накачаем>> цепочку ($i=2$), количество символов $a$ увеличится.
Однако, так как $v$ и $x$ не могут содержать символы $c$ (из-за структуры цепочки и ограничения на длину $vwx$), количество $c$ останется неизменным.
Получим цепочку с разным количеством $a$ и $c$, которая не принадлежит $L$.
Противоречие.
\end{example}

\bigskip
\TODO{Хороший материал, сделано в точности то, что требовалось. \\
Требуются следующие улучшения: \\
Перемещение материала по слайдам с обычной плотностью.\\
Примечание про Огдена и наведение порядка с именованием. \\
Греческие буквы для цепочек и переносы формул. \\
Семь хороших атомов.  }

\end{frame}


%Слайды подготовил И.В. Кан 03.05.2026 на основе материалов Е. Гребенкина от 10.01.2026
\begin{frame}
  \frametitle{1.3.4. Лемма о разрастании (1/5)}

  Для контекстно-свободных языков существует аналог леммы о накачке
  для регулярных языков. Более сильный вариант известен как
  \odmkey{лемма Огдена}{Лемма!Огдена}.

  \begin{theorem}
    Для любого контекстно-свободного языка $L$ существует такая константа
    
    $n \in \mathbb{N}$, что для любой цепочки $\omega\in L$, в которой отмечено
    не менее $n$ позиций, существует разложение
    $\omega = \alpha\beta\gamma\delta\eta$, где:
    \begin{enumerate}
      \item цепочка $\beta\gamma\delta$ содержит \textbf{не более} $n$ отмеченных позиций;
      \item цепочка $\beta\delta$ содержит \textbf{хотя бы одну} отмеченную позицию;
      \item $\forall\,i\ge 0\;\; \alpha\,\beta^{\,i}\,\gamma\,\delta^{\,i}\,\eta \in L$.
    \end{enumerate}
  \end{theorem}

\end{frame}

\begin{frame}
  \frametitle{1.3.4. Лемма о разрастании (2/5)}

  \begin{remark}
    Лемма Огдена (W.~Ogden,~1968) усиливает классическую
    \odmkey{лемму Бар-Хиллела}{Лемма!Бар-Хиллела}
    (Bar-Hillel,~1961), которая получается, если в цепочке
    отмечены \emph{все} позиции.

    \medskip
    \textbf{Суть усиления.}
    В лемме Огдена разрешается отмечать лишь \emph{некоторые} позиции,
    например все вхождения одного конкретного символа.  Благодаря этому
    ``накачка'' ($i\ge 2$) затрагивает только участки, содержащие отмеченные
    позиции, и не меняет остальные части цепочки.  Это даёт более тонкий
    контроль над структурой накачиваемых фрагментов и позволяет доказывать,
    что язык не является контекстно-свободным, в тех случаях, когда
    классическая лемма недостаточно чувствительна.
  \end{remark}

\end{frame}

\begin{frame}
  \frametitle{1.3.4. Лемма о разрастании (3/5)}

  \begin{proof}
    Пусть $L = \mathfrak{L}(G)$, где $G = \langle \mathrm{M},\Sigma,\Pi,S\rangle$ —
    грамматика в нормальной форме Хомского (см. \DMref{п.~1.3.3}{1.3.3.}).
    Обозначим $k = |\mathrm{M}|$ и положим $n = 2^{\,k+1}$.

    Рассмотрим дерево разбора цепочки $\omega$.
    Назовём узел \odmkey{ветвящимся}{Узел!ветвящийся},
    если оба его поддерева содержат хотя бы одну отмеченную позицию.
    Построим путь от корня к листу, каждый раз выбирая потомка,
    содержащего \emph{больше} отмеченных позиций.

    Так как общее число отмеченных позиций не меньше $2^{k+1}$,
    на этом пути окажется не менее $k+1$ ветвящихся узлов.
  \end{proof}

\end{frame}

\begin{frame}
  \frametitle{1.3.4. Лемма о разрастании (4/5)}

  \begin{proof}(Окончание)
    На выбранном пути находится более $k$ нетерминалов.
    
    Согласно \emph{принципу Дирихле}, найдётся нетерминал $A$, который
    встречается на пути дважды — обозначим верхнее вхождение через $A_1$,
    нижнее через $A_2$ (потомок $A_1$).

    Это позволяет разбить вывод на части:
    \[
      S \;\Rightarrow^*\; \alpha A_1 \eta,\qquad
      A_1 \Rightarrow^* \beta A_2 \delta,\qquad
      A_2 \Rightarrow^* \gamma .
    \]

    В силу того, что $A_1$ — ветвящийся узел, в поддеревьях, соответствующих
    $\beta$ и $\delta$, содержится как минимум одна отмеченная позиция,
    следовательно $\beta\delta \neq \varepsilon$.
    Повторяя вывод $A \Rightarrow^* \beta A \delta$ произвольное число раз
    (или пропуская его), получаем цепочки
    $\alpha \beta^{\,i} \gamma \delta^{\,i} \eta$, принадлежащие языку $L$.
  \end{proof}

\end{frame}

\begin{frame}
  \frametitle{1.3.4. Лемма о разрастании (5/5)}

  \begin{example}
    Докажем, что язык
    $L = \{\, a^m b^m c^m \mid m\ge 1 \,\}$
    не является контекстно-свободным.

    Предположим противное.  Пусть $n$ — константа из леммы Огдена.
    Рассмотрим цепочку $\omega = a^n b^n c^n$ и отметим в ней все позиции
    символа $a$ (их ровно $n$).
    Разложим $\omega = \alpha\beta\gamma\delta\eta$ согласно лемме.

    По условию, $\beta\delta$ содержит хотя бы одну отмеченную позицию,
    т.\,е. символ $a$; при этом $\beta\gamma\delta$ содержит не более $n$
    отмеченных позиций.  Следовательно, $\beta$ и $\delta$ состоят
    исключительно из символов $a$ и, возможно, $b$, но не могут содержать
    символов $c$.

    Если ``накачать'' цепочку ($i=2$), количество $a$ возрастает, тогда как
    количество $c$ остаётся неизменным.  Полученная цепочка имеет
    разное число $a$ и $c$ и потому не принадлежит $L$ — противоречие.
  \end{example}

\TODO{10.09.26. ФАН. По существу верно, но форматирование ужасное. Очень плохо с терминологией и именными ссылками.}

\end{frame}


\subsubsection{1.3.5. Нормальная форма Грейбах}\label{1.3.5.}
\begin{frame}
\frametitle{1.3.5. Нормальная форма Грейбах (1/2)}
%\TODO{Нормальная форма Грейбах. }
%
%\bigskip
%\TODO{Грицаенко / Мельник. Очень много.  
%Нужна прополка. Можно взять из Ахо второй вариант}
%
%\bigskip
%\TODO{Материал получен давно, $\approx 5$ атомов.
%Возможно, требуется ревизия.
%{\color{blue} Варисов, Самутичев}} 

Говорят, что КСГ находится в 
\odmkey{нормальной форме Грейбах}{Нормальная форма!Грейбах}, если все правила этой грамматики
имеют следующую форму: 
$A\to a\alpha$, где $A \in N, a\in T, \alpha\in N^*$. 

Построение нормальной формы Грейбах 
основано на устранении левой рекурсии.

Нетерминал $A\in N$ называется 
\odmkey{рекурсивным}{Нетерминал!рекурсивный},
если $A\Rightarrow^+\alpha A \beta$. 
Если при этом $\alpha =\varepsilon$, то
нетерминал называется 
\odmkey{леворекурсивным}{Нетерминал!леворекурсивный},
а если $\beta =\varepsilon$, то
нетерминал называется 
\odmkey{праворекурсивным}{Нетерминал!праворекурсивный}.

\medskip
Грамматика, в которой существуют леворекурсивные
(праворекурсивные) нетерминалы,
называется 
\odmkey{леворекурсивной}{Грамматика!леворекурсивная} (соответственно, \odmkey{праворекурсивной}{Грамматика!праворекурсивная}).


\begin{theorem}
Для любой приведённой контекстно-свободной грамматики $G_1$
существует не леворекурсивная контекстно-свободная грамматика $G_2$,
такая что  $\frak{L}(G_2)=\frak{L}(G_1)$.
\end{theorem}

\begin{proof}
\TODO{Эссе 1.3. Часть 7. \\Алгоритм устранения левой рекурсии
\\желательно связать с уравнениями с регулярными коэффициентами 
}
\end{proof}

\begin{example}
\TODO{Эссе 1.3. Часть 7. \\Упрощенная грамматика выражений в привычной форме\\ и с устранением левой рекурсии }
\end{example}

\end{frame}


\begin{frame}[shrink=22]
\frametitle{1.3.5. Нормальная форма Грейбах (2/2)}
\begin{theorem}
Для любой не леворекурсивной контекстно-свободной грамматики $G_1$
существует контекстно-свободная грамматика $G_2$ в нормальной форме Грейбах,
такая что  $\frak{L}(G_2)=\frak{L}(G_1)$.
\end{theorem}

\begin{proof}
\TODO{Эссе 1.3. Часть 8. \\ Алгоритм приведения к нормальной форме Грейбах 
}
\end{proof}

\begin{example}
\TODO{Эссе 1.3. Часть 8. Упрощенная грамматика выражений в привычной форме\\ и в нормальной форме Грейбах}
\end{example}

\bigskip
\TODO{{\color{red} Замечания ФАН 20 января}\\
Удовлетворительно, большая часть материала пригодна, но необходимо связать с другими частями и исправить огромное количество ошибок. \\
Существенная дырка: а зачем это вообще нужно?\\
Зеленов, долг = 15 баллов
 }

\end{frame}

% Текст слайда подготовил Зеленов Никита 12.01.2026

\begin{frame}
\frametitle{1.3.5. Нормальная форма Грейбах (1/6)}

Говорят, что КСГ находится в 
\odmkey{нормальной форме Грейбах}{Нормальная форма!Грейбах}, если все правила этой грамматики
имеют следующую форму: 
$A\to a\alpha$, где $A \in N, a\in T, \alpha\in N^*$. 

Построение нормальной формы Грейбах 
основано на устранении левой рекурсии.

Нетерминал $A\in N$ называется 
\odmkey{рекурсивным}{Нетерминал!рекурсивный},
если $A\Rightarrow^+\alpha A \beta$. 
Если при этом $\alpha =\varepsilon$, то
нетерминал называется 
\odmkey{леворекурсивным}{Нетерминал!леворекурсивный},
а если $\beta =\varepsilon$, то
нетерминал называется 
\odmkey{праворекурсивным}{Нетерминал!праворекурсивный}.

\medskip
Грамматика, в которой существуют леворекурсивные
(праворекурсивные) нетерминалы,
называется 
\odmkey{леворекурсивной}{Грамматика!леворекурсивная} (соответственно, \odmkey{праворекурсивной}{Грамматика!праворекурсивная}).

\begin{theorem}
Для любой приведённой контекстно-свободной грамматики $G_1$
существует не леворекурсивная контекстно-свободная грамматика $G_2$,
такая что  $\frak{L}(G_2)=\frak{L}(G_1)$.
\end{theorem}
\begin{proof}
Доказательство основано на двух стандартных преобразованиях и
на алгоритме последовательного устранения косвенной и непосредственной левой рекурсии.
\end{proof}

\begin{lemma}[Устранение правила]
Пусть в грамматике есть правило $A\to B\alpha$, а все $B$-правила имеют вид
$B\to \beta_1\mid \dots \mid \beta_k$.
Тогда замена правила $A\to B\alpha$ на правила
$A\to \beta_1\alpha \mid \dots \mid \beta_k\alpha$
не меняет язык грамматики.
\end{lemma}
\end{frame}


\begin{frame}
\frametitle{1.3.5. Нормальная форма Грейбах (2/6)}


\begin{proof}
Правило $A\to B\alpha$ допускает ровно те выводы, где сначала выводится некоторая цепочка из $B$,
а затем приписывается $\alpha$.
Замена на $A\to \beta_i\alpha$ непосредственно воспроизводит все эти случаи. 
\end{proof}

\begin{lemma}[Устранение непосредственной левой рекурсии]
Пусть все $A$-правила имеют вид
$A\to A\alpha_1\mid \dots \mid A\alpha_k \mid \beta_1\mid \dots \mid \beta_m$,
где ни одна $\beta_i$ не начинается с $A$.
Тогда при введении нового нетерминала $A'$ и замене этих правил на
\[
A\to \beta_1A' \mid \dots \mid \beta_mA',\qquad
A'\to \alpha_1A' \mid \dots \mid \alpha_kA' \mid \varepsilon
\]
язык грамматики не меняется.
\end{lemma}

\begin{proof}
В исходной грамматике из $A$ выводятся в точности цепочки
$\beta_i\alpha_{j_1}\alpha_{j_2}\cdots \alpha_{j_\ell}$, где $\ell\ge 0$.
В новой грамматике:
$A\Rightarrow \beta_iA'\Rightarrow \beta_i\alpha_{j_1}A'\Rightarrow \cdots
\Rightarrow \beta_i\alpha_{j_1}\cdots \alpha_{j_\ell}A'\Rightarrow
\beta_i\alpha_{j_1}\cdots \alpha_{j_\ell}$.
Обратно, любое применение $A'$ порождает ровно такую же последовательность $\alpha$-блоков
(или $\varepsilon$), значит множества выводимых цепочек совпадают.
\end{proof}

\smallskip



\end{frame}


\begin{frame}
\frametitle{1.3.5. Нормальная форма Грейбах (3/6)}
\begin{remark}
Для устранения \emph{косвенной} левой рекурсии упорядочим нетерминалы
$N=\{A_1,\dots,A_n\}$ и последовательно для $i=1,\dots,n$ заменяем
каждое правило $A_i\to A_j\alpha$ при $j<i$ по лемме об устранении правила,
после чего устраняем непосредственную левую рекурсию для $A_i$
по предыдущей лемме. Тогда в результате не остаётся правил,
из которых может начаться левый вывод $A_i\Rightarrow^+ A_i\gamma$.
\end{remark}
\begin{digress}
Связь с уравнениями: непосредственная левая рекурсия
$A\to A\alpha_1\mid\dots\mid A\alpha_k\mid \beta_1\mid\dots\mid \beta_m$
соответствует ``линейному'' уравнению языков $X=XE+B$,
где $E=\alpha_1+\dots+\alpha_k$, $B=\beta_1+\dots+\beta_m$.
Преобразование с $A'$ реализует фактор $E^*$: из $A'$ выводятся любые повторения $\alpha$-блоков.
\end{digress}

\begin{example}[Упрощённая грамматика выражений и устранение левой рекурсии]
Пусть $T=\{x,+,\ast,(,)\}$.
Обычная (леворекурсивная) грамматика арифметических выражений:
\[
\begin{array}{l}
E\to E+T \mid T,\\
T\to T\ast F \mid F,\\
F\to (E)\mid x.
\end{array}
\]
\end{example}
\end{frame}


\begin{frame}[shrink=22]
\frametitle{1.3.5. Нормальная форма Грейбах (4/6)}

Устраняя левую рекурсию, получаем эквивалентную грамматику:
\[
\begin{array}{l}
E\to TE' ,\\
E'\to +TE'\mid \varepsilon,\\
T\to FT',\\
T'\to \ast FT'\mid \varepsilon,\\
F\to (E)\mid x.
\end{array}
\]


\begin{remark}
Полученная грамматика удобна для нисходящего разбора:
ни один нетерминал не является леворекурсивным.
\end{remark}

\begin{theorem}
Для любой не леворекурсивной контекстно-свободной грамматики $G_1$
существует контекстно-свободная грамматика $G_2$ в нормальной форме Грейбах,
такая что  $\frak{L}(G_2)=\frak{L}(G_1)$.
\end{theorem}

\begin{digress}
Обычно перед построением НФГ грамматику упрощают (удаляют бесполезные символы,
$\varepsilon$-правила и цепные правила), а затем приводят к форме,
где каждое правило имеет вид $A\to a$ или $A\to BC$ (при $\varepsilon\notin L$).
Эти преобразования сохраняют язык.
\end{digress}

\end{frame}


\begin{frame}
\frametitle{1.3.5. Нормальная форма Грейбах (5/6)}
\begin{proof}
Опишем конструкцию (для случая $\varepsilon\notin L$), которая после предварительного упрощения
строит НФГ из грамматики, у которой правила имеют вид $A\to a$ и $B\to CD$.

Введём новые нетерминалы $(A\backslash B)$ для всех $A,B\in N$ и положим
$N' = N \cup \{(A\backslash B)\mid A,B\in N\}$.
Смысл: $(C\backslash F)$ порождает такие $\gamma\in N^*$, что $F\Rightarrow^* C\gamma$.

Зададим правила:
\[
(A\backslash A)\to \varepsilon\quad (A\in N),
\]
\[
(C\backslash E)\to A(A\backslash D)(B\backslash E)
\quad \text{для всех }(B\to CD)\in R,\ A\in N,\ E\in N,
\]
\[
A\to a \quad \text{для всех }(A\to a)\in R,
\qquad
S\to a(A\backslash S)\ \text{для всех }(A\to a)\in R.
\]
Теперь заменим в правилах вида $(C\backslash E)\to A(A\backslash D)(B\backslash E)$
первый нетерминал $A$ на терминал $a$, используя все варианты правил $A\to a$.
Получим правила
\[
(C\backslash E)\to a(A\backslash D)(B\backslash E),
\]
то есть каждое правило начинается с терминала, а далее идут только нетерминалы:
получена нормальная форма Грейбах.



\end{proof}
\end{frame}


\begin{frame}
\frametitle{1.3.5. Нормальная форма Грейбах (6/6)}
Корректность (вкратце) доказывается индукцией по длине $\gamma\in N^*$:
\begin{enumerate}
\item Если $F\Rightarrow^* C\gamma$ в исходной грамматике, то $(C\backslash E)\Rightarrow^* \gamma(F\backslash E)$ в новой.
\item Для всех $C,F$ верно: $F\Rightarrow^* C\gamma$ тогда и только тогда, когда $(C\backslash F)\Rightarrow^* \gamma$.
\end{enumerate}
Из (2) следует равенство языков, в частности для стартового символа.
\begin{example}[Упрощённая грамматика выражений в НФГ]
Ниже приведён один из возможных вариантов НФГ для языка выражений из примера (3/6)
(используем вспомогательный нетерминал $R\to )$ для закрывающей скобки).
Все правила имеют вид $A\to a\alpha$.
\small
\[
\begin{array}{l}
R\to ),\\
F\to x \mid (ER,\\
M\to \ast FM \mid \ast F,\\
P\to +TP \mid +T,\\[2mm]
T\to x \mid xM \mid (ER \mid (ERM,\\[2mm]
E\to x \mid xM \mid xP \mid xMP \mid (ER \mid (ERM \mid (ERP \mid (ERMP.
\end{array}
\]
\end{example}

\begin{remark}
Грамматика получена из нелеворекурсивной формы подстановкой правил для $T$ в правила для $E$
так, чтобы в каждой продукции первым символом был терминал.
\end{remark}
\end{frame}




\subsubsection{1.3.6. Алгоритмические свойства грамматик}\label{1.3.6.}
\begin{frame}
\frametitle{1.3.6. Алгоритмические свойства грамматик (1/Х)}

\TODO{Повторное Определение термина неоднозначные грамматики и двусмысленные грамматки.  Два развёрнутых примера. Первый пример Зубковой 1.1.5. Второй пример листочки Чепулис. {\color{blue}  Войнова Алёна}}
\end{frame}


\begin{frame}
\frametitle{1.3.6. Алгоритмические свойства грамматик (2/Х)}


\TODO{Замкнутость, детерминированность, неоднозначность. \\
Теоретический параграф про разрешимое и неразрешимое.\\
Является ли проблема однозначности грамматики алгоритмически разрешимой?}

\begin{theorem}
Проблема распознавания однозначности 
контекстно-свободной грамматики алгоритмически неразрешима.
\end{theorem}

\bigskip
\TODO{\color{blue} Параграф повис в воздухе. {\color{blue} Яков Караев, сделан доклад. Нужны хорошие атомы.}}

\medskip
К сожалению, в общем случае
проблема эквивалентности двух
контекстно-свободных грамматик,
равно как и проблема определения
неоднозначности контекстно-свободной
грамматики оказываются алгоритмически неразрешимыми. 

\medskip
\TODO{разобраться и доложить!}

\medskip
\TODO{Проблемы: Неоднозначность. Пустота языка. }

\medskip
\TODO{{\color{red} Замечание ФАН 21 января.} 
Параграф трудный, но очень важный и полезный для самообразования. Помогу с источниками. }

\end{frame}


% Источники к разделу: Карпов Ю.Г. Основы построения трансляторов, 2005г.
% Источники к пункту: определения КСГ, языка цепочки и дерева вывода согласованы с \DMref{п.~1.3.1}{1.3.1.} и \DMref{п.~1.1.5}{1.1.5.}.
% Источники к пункту: замкнутость КС-языков и незамкнутость относительно пересечения см. Хопкрофт Дж., Мотвани Р., Ульман Дж. Введение в теорию автоматов, языков и вычислений. 3-е изд., §7.3.
% Источники к пункту: разрешимые свойства КСГ см. Хопкрофт Дж., Мотвани Р., Ульман Дж. Введение в теорию автоматов, языков и вычислений. 3-е изд., §7.4.
% Источники к пункту: неразрешимость однозначности КСГ через ПСП см. Галлье Ж. Глава 8. Проблема соответствия Поста, §8.2, с. 555; см. также Хопкрофт, Мотвани, Ульман, §9.4.
% Источники к пункту: пример с условным оператором согласован с \DMref{п.~1.4.1}{1.4.1.}; пример с выражениями согласован с \DMref{п.~1.1.6}{1.1.6.}.

\begin{frame}[t]
	\frametitle{1.3.6. Алгоритмические свойства грамматик }
	\setlength{\parskip}{1pt}
	\sloppy
	
	Алгоритмические свойства грамматик относятся к вопросам, для которых требуется
	
	общий метод, применимый ко всем грамматикам данного класса: сохранение класса 
	
	языков при операциях, однозначность синтаксического описания, а также разрешимость задач о грамматике и о цепочке. Понятия алгоритмической разрешимости и полуразрешимости используются в смысле \DMref{п.~1.1.5}{1.1.5.}; язык, порождаемый цепочкой $\omega$, обозначается $\frak{L}(\omega)$ в смысле \DMref{п.~1.3.1}{1.3.1.}
	% Обоснование: терминология разрешимости уже введена в п.~1.1.5; обозначение \frak{L}(\alpha) введено в п.~1.3.1; Хопкрофт, Мотвани, Ульман, 3-е изд., §7.1--7.4.
	
	
	Грамматика $G=\langle \mathrm{M}, \Sigma, \Pi, S \rangle$ называется \odmkey{неоднозначной}{Грамматика!неоднозначная}, если существует цепочка 
	
	$\omega\in\frak{L}(G)$, имеющая два различных дерева разбора в $G$. Термин \odmkey{двусмысленная грамматика}{Грамматика!двусмысленная} употребляется как синоним. Грамматика называется \odmkey{однозначной}{Грамматика!однозначная}, если каждая 
	
	цепочка языка $\frak{L}(G)$ имеет ровно одно дерево разбора.
	% Обоснование: определение следует из определения дерева вывода в п.~1.1.5 и КСГ в п.~1.3.1; Хопкрофт, Мотвани, Ульман, 3-е изд., §5.2.
	
	
	Неоднозначность является свойством грамматики, а не только языка. Эквивалентные грамматики могут порождать один и тот же язык, но задавать различные деревья разбора для одних и тех же цепочек; поэтому равенство $\frak{L}(G_1)=\frak{L}(G_2)$ не означает совпадения синтаксических структур, задаваемых грамматиками.
	% Обоснование: понятие эквивалентности относится к языкам по п.~1.3.2, а дерево разбора является дополнительной структурой вывода по п.~1.1.5.
\end{frame}

\begin{frame}[t]
	\frametitle{1.3.6. Алгоритмические свойства грамматик }
	\setlength{\parskip}{1pt}
	\sloppy
	
	\begin{digress}
		Существенно неоднозначным называется язык, для которого всякая порождающая КСГ неоднозначна. Здесь существенно неоднозначные языки не используются для доказательств; термин отделяет свойство языка от свойства выбранной грамматики.
		% Обоснование: термин приведён для разграничения свойств языка и грамматики; Хопкрофт, Мотвани, Ульман, 3-е изд., §7.4.
	\end{digress}
	
	\begin{example}
		Первый пример связан с арифметическими выражениями. В \DMref{п.~1.1.6}{1.1.6.} сопоставляются разные грамматические описания выражений; компактная КСГ с правилами 
		
		$E\to E+E\mid E*E\mid a$ неоднозначна. Цепочка $a+a*a$ имеет два дерева разбора: в одном корневым является правило $E\to E+E$, и выражение разбивается как $a+(a*a)$; в другом корневым является правило $E\to E*E$, и выражение разбивается как $(a+a)*a$.
		% Обоснование: пример опирается на материал о деревьях вывода и выражениях из п.~1.1.6; Хопкрофт, Мотвани, Ульман, 3-е изд., §5.2.
	\end{example}
	
	
	Оба дерева из предыдущего примера соответствуют корректным выводам одной цепочки, но задают разные синтаксические структуры. Поэтому при построении анализатора недостаточно знать только язык $\frak{L}(G)$: требуется контролировать также структуру 
	
	вывода.
	% Обоснование: связь языка и структуры разбора следует из постановки задачи синтаксического анализа в \DMref{п.~1.4.1}{1.4.1.}.
	
\end{frame}

\begin{frame}[t]
	\frametitle{1.3.6. Алгоритмические свойства грамматик }
	\setlength{\parskip}{1pt}
	\sloppy
	
	
	\begin{example}
		Второй пример связан с условным оператором. Фрагмент грамматики с правилами $S\to \mathrm{if}\ B\ \mathrm{then}\ S$, $S\to \mathrm{if}\ B\ \mathrm{then}\ S\ \mathrm{else}\ S$ и $S\to a$ неоднозначен. Цепочка $\mathrm{if}\ B_1\ \mathrm{then}\ \mathrm{if}\ B_2\ \mathrm{then}\ a\ \mathrm{else}\ a$ имеет два дерева разбора: в одном часть $\mathrm{else}\ a$ относится к внешнему условному оператору, в другом~--- к внутреннему. Эта неоднозначность является классическим случаем ``обособленного else''; в \DMref{п.~1.4.1}{1.4.1.} она связана с требованиями к постановке синтаксического анализа.
		% Обоснование: пример согласован с материалом п.~1.4.1 о неоднозначности условного оператора; Хопкрофт, Мотвани, Ульман, 3-е изд., §5.2.
	\end{example}
	
	\begin{remark}
		Неоднозначность устраняется не объяснением вне грамматики, а изменением формального описания: например, соглашением о связывании $\mathrm{else}$ с ближайшим
		
		незавершённым $\mathrm{if}$ или введением явного маркера конца условного оператора.
		% Обоснование: указанная связь с ближайшим незавершённым if использована в п.~1.4.1;
	\end{remark}
	
	Детерминированность синтаксического анализа означает, что на каждом шаге разбора следующее действие определяется текущей конфигурацией и ограниченным фрагментом входной цепочки. В таком методе анализа не возникает двух конкурирующих деревьев разбора одной цепочки. Подклассы $LL(k)$ и $LR(k)$ вводятся в \DMref{п.~1.3.7}{1.3.7.}; здесь используется только связь между детерминированностью и устранением неоднозначности.
	% Обоснование: определения LL(k) и LR(k) вынесены в п.~1.3.7; Хопкрофт, Мотвани, Ульман, 3-е изд., §4.4 и §5.4.
	
\end{frame}

\begin{frame}[t]
	\frametitle{1.3.6. Алгоритмические свойства грамматик }
	\setlength{\parskip}{1pt}
	\sloppy
	
	
	
	\begin{remark}
		Эти свойства различны. Компактная грамматика может быть неоднозначной; однозначная грамматика может не принадлежать удобному детерминированному
		
		классу; детерминированная грамматика задаёт управляемую процедуру распознавания.
		% Обоснование: различие следует из примеров этого пункта и из последующего выделения классов LL(k), LR(k) в п.~1.3.7: Хопкрофт, Мотвани, Ульман, 3-е изд., §5.4.
	\end{remark}
	
	\begin{theorem}
		Класс контекстно-свободных языков замкнут относительно объединения,
		
		конкатенации и итерации Клини.
		% Обоснование: классический результат о КС-языках: Хопкрофт, Мотвани, Ульман, 3-е изд., §7.3.
	\end{theorem}
	
	\begin{proof}
		\par\vspace{-4pt}
		Пусть $L_i=\frak{L}(G_i)$, где $G_i=\langle \mathrm{M}_i,\Sigma,\Pi_i,S_i\rangle$. После переименования нетерминалов множества $\mathrm{M}_1$ и $\mathrm{M}_2$ не пересекаются. Для $L_1\cup L_2$ добавляется новая аксиома $S$ и правило $S\to S_1\mid S_2$; для $L_1L_2$ добавляется $S\to S_1S_2$; для $L_1^*$ добавляются $S\to S_1S\mid\varepsilon$. Все
		
		новые правила контекстно-свободны, а проверка первого шага вывода даёт
		
		соответственно $L_1\cup L_2$, $L_1L_2$ и $L_1^*$.
	\end{proof}
	
\end{frame}

\begin{frame}[t]
	\frametitle{1.3.6. Алгоритмические свойства грамматик }
	\setlength{\parskip}{1pt}
	\sloppy
	
	\begin{theorem}
		Класс контекстно-свободных языков не замкнут
		
		относительно пересечения и дополнения.
		% Обоснование: классический результат о КС-языках; Хопкрофт, Мотвани, Ульман, 3-е изд., §7.3.
	\end{theorem}
	
	\begin{proof}
		\par\vspace{-4pt}
		Пусть $L_1$ состоит из слов $a^m b^m c^n$, а $L_2$ состоит из слов $a^m b^n c^n$, где $m,n\geq 0$. Оба языка контекстно-свободны: первый синхронно порождает равное число $a$ и $b$, второй синхронно порождает равное число $b$ и $c$. Пересечение этих языков равно $\Set{a^n b^n c^n}{n\geq 0}$, а этот язык не является контекстно-свободным по лемме о разрастании из \DMref{п.~1.3.4}{1.3.4.}. Замкнутости относительно пересечения нет. Если бы класс был замкнут относительно дополнения, то по закону де Моргана и замкнутости относительно объединения он был бы замкнут относительно пересечения; это противоречит доказанному.
	\end{proof}
	
	Алгоритмическая проблема о грамматиках называется разрешимой, если существует
	
	алгоритм, который для каждого допустимого входа останавливается и даёт правильный ответ <<да>> или <<нет>>. Проблема называется неразрешимой, если такого алгоритма не существует.
	% Обоснование: определение согласовано с понятием разрешимости из п.~1.1.5; Сипсер М. Введение в теорию вычислений, 2-е изд., §3.1--4.1.
	
\end{frame}

\begin{frame}[t]
	\frametitle{1.3.6. Алгоритмические свойства грамматик }
	\setlength{\parskip}{1pt}
	\sloppy
	
	\begin{theorem}
		Проблема пустоты языка контекстно-свободной грамматики алгоритмически разрешима.
		% Обоснование: стандартная разрешимая задача для КСГ; Хопкрофт, Мотвани, Ульман, 3-е изд., §7.4.
	\end{theorem}
	
	\begin{proof}
		\par\vspace{-4pt}
		По грамматике $G=\langle \mathrm{M},\Sigma,\Pi,S\rangle$ отмечаются порождающие нетерминалы. Сначала
		
		отмечаются те $A$, для которых есть правило $A\to\omega$ с $\omega\in\Sigma^*$; затем отмечается каждый $A$, для которого есть правило $A\to\alpha$, где все нетерминалы из $\alpha$ уже отмечены. Так как $\mathrm{M}$ конечно, процесс останавливается. Язык $\frak{L}(G)$ непуст тогда и только тогда, когда 
		
		отмечен $S$: в одну сторону обратный просмотр отметок даёт терминальный вывод, в 
		
		другую сторону индукция по высоте дерева разбора отмечает все нетерминальные вершины терминального вывода из $S$.
	\end{proof}
	
	\begin{theorem}
		Проблема принадлежности цепочки языку контекстно-свободной грамматики алгоритмически разрешима.
		% Обоснование: стандартная разрешимая задача для КСГ; Хопкрофт, Мотвани, Ульман, 3-е изд., §7.4; алгоритмы анализа далее относятся к разд.~1.4.
	\end{theorem}
	
\end{frame}

\begin{frame}[t]
	\frametitle{1.3.6. Алгоритмические свойства грамматик }
	\setlength{\parskip}{1pt}
	\sloppy
	
	\begin{proof}
		\par\vspace{-4pt}
		Грамматика переводится в нормальную форму Хомского по \DMref{п.~1.3.3}{1.3.3.}. Для непустой цепочки длины $n$ каждое дерево разбора в такой грамматике имеет конечное число внутренних вершин и ровно $n$ терминальных листьев. Следовательно, существует конечное множество возможных деревьев с заданными листьями. Перебор этих деревьев и проверка согласованности внутренних вершин с правилами грамматики дают конечную процедуру. Более эффективные способы проверки принадлежности относятся к задаче синтаксического анализа из разд.~\DMref{1.4}{1.4.}
	\end{proof}
	
	\begin{remark}
		Разрешимость принадлежности не означает эффективного общего разбора в практическом смысле. Для произвольных КСГ требуются общие алгоритмы, тогда как детерминированные подклассы из \DMref{п.~1.3.7}{1.3.7.} выделяются ради управляемого выбора действия.
		% Обоснование: связь принадлежности и синтаксического анализа проводится в разд.~1.4; Хопкрофт, Мотвани, Ульман, 3-е изд., §7.4.
	\end{remark}
	
	\begin{theorem}
		Проблема распознавания однозначности контекстно-свободной грамматики алгоритмически неразрешима.
		% Обоснование: классическая неразрешимая задача для КСГ; Галлье Ж. Глава 8. Проблема соответствия Поста, §8.2, теорема 8.2, с. 555.
	\end{theorem}
	
\end{frame}

\begin{frame}[t]
	\frametitle{1.3.6. Алгоритмические свойства грамматик }
	\setlength{\parskip}{1pt}
	\sloppy
	
	\begin{proof}
		\par\vspace{-4pt}
		Доказательство основано на сведении проблемы соответствия Поста к проблеме однозначности. Экземпляр ПСП состоит из пар непустых слов $(u_i,v_i)$, $1\leq i\leq k$, над
		
		алфавитом $\Sigma$; требуется определить, существует ли непустая последовательность индексов $i_1,\ldots,i_m$, для которой верхняя и нижняя конкатенации совпадают: 
		
		$u_{i_1}\cdots u_{i_m}=v_{i_1}\cdots v_{i_m}$. Проблема соответствия Поста алгоритмически неразрешима.
		% Обоснование: определение ПСП и её неразрешимость; Галлье Ж. Глава 8, §8.1, с. 551--554.
		
		Каждому экземпляру ПСП сопоставляется грамматика $G$, неоднозначная в точности тогда, когда экземпляр имеет решение. Вводятся новые терминалы $c_1,\ldots,c_k$ и грамматика
		
		$G=\langle\{S,A,B\},\Sigma\cup\{c_1,\ldots,c_k\},\Pi,S\rangle$. Правила имеют вид $S\to A\mid B$, $A\to u_iAc_i\mid u_ic_i$ и $B\to v_iBc_i\mid v_ic_i$ для всех $1\leq i\leq k$.
		% Обоснование: конструкция сведения ПСП к неоднозначности КСГ; Галлье Ж. Глава 8, §8.2, с. 555.
		
		Нетерминал $A$ порождает цепочки с верхней конкатенацией $u_{i_1}\cdots u_{i_m}$ и суффиксом $c_{i_m}\cdots c_{i_1}$; нетерминал $B$ порождает такие же цепочки с нижней конкатенацией $v_{i_1}\cdots v_{i_m}$. Суффикс из символов $c_i$ однозначно восстанавливает последовательность индексов, поскольку все $c_i$ различны.
		% Обоснование: индукция по числу применений правил A и B; Галлье Ж. Глава 8, §8.2, с. 555--556.
	\end{proof}
	
\end{frame}

\begin{frame}[t]
	\frametitle{1.3.6. Алгоритмические свойства грамматик }
	\setlength{\parskip}{1pt}
	\sloppy
	
	\begin{proof}
		\par\vspace{-4pt}
		Если ПСП имеет решение $i_1,\ldots,i_m$, то одна и та же цепочка с общим префиксом $u_{i_1}\cdots u_{i_m}=v_{i_1}\cdots v_{i_m}$ и суффиксом $c_{i_m}\cdots c_{i_1}$ выводится из $S$ двумя способами: через $A$ и через $B$. Поэтому грамматика $G$ неоднозначна.
		
		Обратно, если грамматика $G$ неоднозначна, то неоднозначность не возникает внутри ветви $A$ или внутри ветви $B$: суффикс $c_{i_m}\cdots c_{i_1}$ однозначно задаёт все применённые правила и их порядок. Следовательно, неоднозначная цепочка выводится и через $A$, и через $B$. Для некоторой непустой последовательности индексов выполняется равенство $u_{i_1}\cdots u_{i_m}=v_{i_1}\cdots v_{i_m}$, то есть исходный экземпляр ПСП имеет решение.
		
		Построенная грамматика неоднозначна тогда и только тогда, когда исходный экземпляр ПСП имеет решение. Алгоритм распознавания однозначности КСГ дал бы алгоритм решения ПСП, что невозможно. Следовательно, проблема распознавания однозначности КСГ алгоритмически неразрешима.
		% Обоснование: заключительная часть сведения; Галлье Ж. Глава 8, §8.2, с. 555--556.
	\end{proof}
	
\end{frame}

\begin{frame}[t]
	\frametitle{1.3.6. Алгоритмические свойства грамматик }
	\setlength{\parskip}{1pt}
	\sloppy
	
	\begin{remark}
		Суффикс $c_{i_m}\cdots c_{i_1}$ записан в обратном порядке, чтобы последовательность индексов была восстановима из правой части вывода. Иначе совпадение слов ПСП не отделяется от случайного совпадения разбиений на блоки.
		% Обоснование: различимость символов c_i используется в доказательстве однозначного восстановления индексов; Галлье Ж. Глава 8, §8.2, с. 555.
	\end{remark}
	
	\begin{corollary}
		Не существует общего алгоритма, который по произвольной контекстно-свободной грамматике $G=\langle \mathrm{M},\Sigma,\Pi,S\rangle$ устанавливает, имеет ли каждая цепочка из $\frak{L}(G)$ единственное дерево разбора.
		% Обоснование: непосредственная переформулировка предыдущей теоремы.
	\end{corollary}
	
	\begin{proof}
		\par\vspace{-4pt}
		Такой алгоритм был бы алгоритмом распознавания однозначности КСГ. Предыдущая
		
		теорема устанавливает невозможность его существования, поскольку иначе решалась бы 
		
		алгоритмически неразрешимая проблема соответствия Поста.
	\end{proof}
	
	\begin{remark}
		Неразрешимость проблемы однозначности не противоречит разрешимости проблемы принадлежности. Принадлежность фиксированной цепочки проверяется
		
		конечной процедурой, тогда как однозначность требует общего утверждения обо всех цепочках языка и обо всех их деревьях разбора.
	\end{remark}
	
\end{frame}

\begin{frame}
	\frametitle{1.3.6. Алгоритмические свойства грамматик }
	\setlength{\parskip}{1pt}
	\sloppy
	
	\begin{remark}
		Поэтому практические методы анализа используют проверяемые синтаксические ограничения, например принадлежность грамматики классам $LL(k)$ или $LR(k)$, а не общий алгоритм распознавания однозначности произвольной КСГ.
		% Обоснование: сопоставление разрешимости принадлежности и неразрешимости однозначности следует из теорем этого пункта; Хопкрофт, Мотвани, Ульман, 3-е изд., §7.4; Галлье Ж. Глава 8, §8.2.
	\end{remark}
	
	\begin{digress}
		Проблема эквивалентности двух КСГ формулируется как проверка равенства $\frak{L}(G_1)=\frak{L}(G_2)$ для произвольных грамматик. В этом пункте используется только доказанная неразрешимость однозначности; задача эквивалентности требует отдельного сведения и не подменяет вопрос о деревьях разбора.
		% Обоснование: неразрешимость эквивалентности КСГ относится к тому же кругу задач, но доказывается отдельно; Галлье Ж. Глава 8, §8.2, теорема 8.7, с. 561--562.
	\end{digress}
	
\end{frame}

\subsubsection{1.3.7. Грамматики $LL(k)$ и $LR(k)$ }\label{1.3.7.}

% Текст слайда подготовила Королева Дарья 5030102/20201 
% Дата 24.12.25
% полностью переписал ФАН 25.12.25
\begin{frame}
\frametitle{1.3.7. Грамматики $LL(k)$ и $LR(k)$ (1/9)}

Как показано в \DMref{п. 1.4.4}{1.4.4.}, 
наилучшие известные алгоритмы синтаксического анализа \\
контекстно-свободных языков без ограничений на грамматику  
имеют асимптотическую оценку трудоёмкости 
разбора $ \cal{O}(n^3)$, где $n$~--- длина входной цепочки.  

\smallskip
В этом параграфе рассматриваются такие подклассы КСГ, 
для которых возможно построить эффективные анализаторы, 
расходующие на обработку входной цепочки длины \(n\) время и память 
порядка $\cal{O}(n)$. 

\smallskip
За эту эффективность приходится расплачиваться тем, что ни один из классов грамматик,
для которых можно построить такие эффективные анализаторы,
{\em не} порождает {\em все} контекстно-свободные языки. 
Однако эти ограниченные классы грамматик оказываются 
{\em достаточными} для описания синтаксиса современных языков программирования.

\smallskip
\begin{digress}
По традиции рассматриваемые здесь подклассы КСГ
принято обозначать $LL(k)$ и $LR(k)$. 
В этих обозначениях первая буква означает направление просмотра
входной цепочки ($L$~--- слева направо), 
вторая буква означает ориентацию вывода
($L$~--- левосторонний вывод, $R$~--- правосторонний вывод, 
\DMref{п.~1.1.5}{1.1.5.}),
а $k$~--- число символов, на которые анализатор должен <<заглянуть вперёд>>
для того, чтобы разбор оставался безвозвратным. 
\end{digress}

\end{frame}

% Текст слайда подготовила Королева Дарья 5030102/20201 
% Дата 24.12.25
% ревизия ФАН 28.12.25
% Изменения внесены 06.01.26 Королева Дарья Сергеевна
\begin{frame}
\frametitle{1.3.7. Грамматики $LL(k)$ и $LR(k)$ (2/9)}
Пусть задана однозначная КСГ $G=\langle \mathrm{M}, \Sigma, \Pi, S \rangle$ и её язык 
$L = \frak{L}(G)$. Рассмотрим функцию, по традиции обозначаемую
$\mathrm{FIRST}^G_k$, доставляющую множество всех 
возможных префиксов терминального порождения слова $\alpha\in(\mathrm{M}\cup\Sigma)^*$:      
$$
    \mathrm{FIRST}_k^G(\alpha) \Def
    \Set{\omega\in\Sigma^*}{(|\omega| < k \And \omega\in\frak{L}(\alpha))
    \Or (|\omega| = k \And \E{\beta\in\Sigma^*}{\omega\beta\in\frak{L}(\alpha) })}.
$$
\begin{remark}
Для доказательства теорем необходимо рассматривать 
функцию $\mathrm{FIRST}^G_k$ как определённую на множестве всех
цепочек, но для реализации простых алгоритмов разбора 
достаточно определить эту функцию только для нетерминальных символов:\\
$\Map{\mathrm{FIRST}^G_k}{\mathrm{M}}{2^{\Sigma^*}}$. 
Если индексы $G$ и $k$ ясны из контекста, их опускают.   
\end{remark}

\smallskip
КСГ $G=\langle \mathrm{M}, \Sigma, \Pi, S \rangle$ 
называется \odmkey{$LL(k)$-грамматикой}{Грамматика!$LL(k)$},
 если из существования двух левых выводов
$ S \Rightarrow_l^* \omega A\alpha \Rightarrow_l \omega\beta\alpha \Rightarrow^* \omega
\delta $ и
$ S \Rightarrow_l^* \omega A\alpha \Rightarrow_l \omega\gamma\alpha \Rightarrow^* \omega\lambda $,
для которых \( \mathrm{FIRST}_k(\delta) = \mathrm{FIRST}_k(\lambda) \), вытекает, 
что \( \beta = \gamma \).

\begin{remark} $LL(k)$ условие означает, что при левом разборе
цепочки \( \omega A\alpha \), где цепочка $\omega$ уже разобрана и $A$ первый слева нетерминал, существует {\em не более одного правила} для \( A \), позволяющего продолжить вывод с данными \( k \) терминалами.
\end{remark}
\end{frame}

% Текст слайда подготовила Королева Дарья 5030102/20201 
% Дата 24.12.25
% Изменения внесены 06.01.26 Королева Дарья Сергеевна
% Изменения внесены 15.02.26 Королева Дарья Сергеевна
\begin{frame}
\frametitle{1.3.7. Грамматики $LL(k)$ и $LR(k)$ (3/9)}
\begin{example}
Рассмотрим грамматику \( G_1 \) с правилами 
\(S \rightarrow aAS \mid b, A \rightarrow a \mid bSA \) 
и покажем, что $G_1$ 
является $LL(1)$-грамматикой. 
Чтобы убедиться в этом, надо рассмотреть две цепочки вывода:
 \[
 S \Rightarrow_l^* \omega S\alpha \Rightarrow_l \omega\beta\alpha \Rightarrow^* \omega\delta \quad
 S \Rightarrow_l^* \omega S\alpha \Rightarrow_l \omega\gamma\alpha \Rightarrow^* \omega\lambda ,\]
где \(\mathrm{FIRST}_1(\delta) = \mathrm{FIRST}_1(\lambda)\).
 Eсли \( \delta \) и \( \lambda \) начинаются символом \( a \), то в выводе участвует \( S \to aAS \) и \( \beta = \gamma = aAS \). Eсли же \( \delta \) и \( \lambda \) начинаются \( b \), то применяется \( S \to b \) и \( \beta = \gamma = b \).
% Откуда $\beta=\gamma$.
\end{example}

\begin{theorem}
Приведенная КСГ $G = \langle \mathrm{M}, \Sigma, \Pi, S \rangle$ является $LL(k)$-грамматикой тогда и только тогда, когда для двух различных правил \( A \to \beta \) и \( A \to \gamma \) пересечение\\ \( \mathrm{FIRST}_k(\beta\alpha) \cap \mathrm{FIRST}_k(\gamma\alpha) = \emptyset \) при всех таких \( wA\alpha \), что \( S \Rightarrow_l^* wA\alpha \).
\end{theorem}
\begin{proof}
\proofsection{Необходимость.}
Допустим, что \( \omega, A, \alpha, \beta \) и \( \gamma \) 
удовлетворяют условиям теоремы, а \\
\( \operatorname{FIRST}_k(\beta\alpha) \cap \operatorname{FIRST}_k(\gamma\alpha) \) 
содержит \( \delta \). 
Тогда по определению \(\operatorname{FIRST}\) для некоторых \( \lambda \) и \( \mu \) найдутся выводы 
\(
S \Rightarrow_i^* \omega A\alpha \Rightarrow_i^* \omega \beta\alpha \Rightarrow_i^* \omega \delta \lambda
\)
и 
\(
S \Rightarrow_i^* \omega A \alpha \Rightarrow_i^* \omega\gamma\alpha \Rightarrow_i^* \omega \delta \mu.
\)
Если \( |\delta| < k \), то \( \lambda = \mu = \varepsilon \). Так как \( \beta \neq \gamma \), то \( G \) не \( LL(k) \)-грамматика.
\end{proof}
\begin{remark}
Здесь использован тот факт, что \( \mathrm{M} \) не содержит бесполезных нетерминалов
\end{remark}
\end{frame}

% Текст слайда подготовила Королева Дарья 5030102/20201 
% Дата 24.12.25
% Изменения внесены 06.01.26 Королева Дарья Сергеевна
% Изменения внесены 15.02.26 Королева Дарья Сергеевна
\begin{frame}
\frametitle{1.3.7. Грамматики $LL(k)$ и $LR(k)$ (4/9)}
\begin{proof} (Окончание)
\proofsection{Достаточность.}
Допустим, что \( G \) не $LL(k)$-грамматика. Тогда найдутся такие два вывода 
\(
S \Rightarrow_i^* \omega A\alpha \Rightarrow_l \omega \beta \alpha \Rightarrow_i^* \omega \delta
\)
и
\(
S \Rightarrow_i^* \omega A\alpha \Rightarrow_l \omega \gamma \alpha \Rightarrow_i^* \omega \lambda,
\)
что цепочки \( \delta \) и \( \lambda \) совпадают в первых \( k \) позициях, но \( \beta \neq \gamma \). Поэтому \( A \to \beta \) и \( A \to \gamma \)~---
 различные правила из \( \Pi \) и каждое из множеств \( \operatorname{FIRST}_k(\beta \alpha) \) и \( \operatorname{FIRST}_k(\gamma \alpha) \) содержит цепочку\\ \( \operatorname{FIRST}_k(\delta) \) = \( \operatorname{FIRST}_k(\lambda) \).
\end{proof}


\begin{example} Грамматика $G$, состоящая из двух правил $S \to aS \mid a$, 
не является $LL(1)$-грамматикой, 
так как $\operatorname{FIRST}_1(aS) = \operatorname{FIRST}_1(a) = a$. 
Интуитивно это можно объяснить так: видя при разборе цепочки, начинающейся символом $a$, только этот первый символ, мы не знаем, какое из правил $S \to aS$ или $S \to a$ надо применить к $S$.

С другой стороны, $G$~--- это $LL(2)$-грамматика. В самом деле, если $S \Rightarrow^*_l \omega A\alpha$, то $A = S$ и $\alpha = \varepsilon$. Так как для $S$ даны только два указанных правила, то $\beta = aS$ и $\gamma = a$. Поскольку $\operatorname{FIRST}_2(aS) = aa$ и $\operatorname{FIRST}_2(a) = a$, то по теореме $G$ будет $LL(2)$-грамматикой.
\end{example} 

\begin{digress} Можно определить функцию 
$\operatorname{FOLLOW}_k(\beta)$, доставляющую
множество первых $k$ терминальных символов, которые могут следовать непосредственно за цепочкой $\beta$ в выводе из начального символа:
$
\operatorname{FOLLOW}_k(\beta) \Def \Set{\omega\in \Sigma^*}{ S \Rightarrow^*
 \alpha \beta \gamma \And \omega \in \operatorname{FIRST}_k(\gamma) }
$.
\end{digress}


\end{frame}

% Текст слайда подготовила Королева Дарья 5030102/20201 
% Дата 15.03.26
\begin{frame}
\frametitle{1.3.7. Грамматики $LL(k)$ и $LR(k)$ (5/9)}
Пусть $\langle \mathrm{M}, \Sigma, \Pi, S \rangle$ -- контекстно-свободная грамматика. 
\odmkey{Пополненной грамматикой}{Грамматика!Пополненная}, полученной из $G$, называется грамматика 
$
G' = \langle \mathrm{M} \cup \{ S' \}, \Sigma, \Pi \cup \{ S' \to S \}, S' \rangle,
$
где $S'$ --- новый начальный символ, не принадлежащий $\mathrm{M}$, $S' \to S$ --- нулевое правило грамматики $G'$, 
а остальные правила занумерованы числами $1, 2, \ldots, p$. 
\begin{remark} Это начальное правило добавляется для того, чтобы свёртку, в которой 
используется нулевое правило, можно было интерпретировать как сигнал о том, 
что цепочка допускается.
\end{remark}

Пусть $\langle \mathrm{M}, \Sigma, \Pi, S \rangle$ --- КСГ, 
а $G' = \langle \mathrm{M}', \Sigma, \Pi', S'\rangle $ --- полученная из неё пополненная грамматика. Тогда 
$G$ называется \odmkey{LR($k$)-грамматикой}{Грамматика!$LR(k)$} для $k \ge 0$, 
если из условий\\
$
S' \Rightarrow^*_{r} \alpha A \omega \Rightarrow_{r} \alpha \beta \omega,\quad
S' \Rightarrow^*_{r} \gamma B \delta \Rightarrow_{r} \gamma \beta \lambda, \quad
\operatorname{FIRST}_k(\omega) = \operatorname{FIRST}_k(\lambda)
$
следует, что $\alpha A \lambda = \gamma B \delta$, то есть $\alpha = \gamma$, $A = B$ и $\delta = \lambda$.
\begin{example}
Грамматика с правилами \(S \to Sa \mid a \) не  является  \(LR(0)\)-грамматикой . 
По определению \(LR(0)\)-грамматики: 
\(
S' \Rightarrow_{r}S, \space
S' \Rightarrow_{r}S \Rightarrow_{r}Sa, \space
\operatorname{FIRST}_0(\varepsilon) = \operatorname{FIRST}_0(a) = \varepsilon
\).
В первом выводе нетерминал $A \Assign S'$, во втором --- нетерминал $B \Assign S$. Tак как $A \neq B$, анализатор, видя в стеке символ $S$, не может детерминированно решить: свернуть его в $S'$ или
оставить его в стеке и ждать символ $a$ для применения правила \(S \to a \).
\end{example}
\end{frame}

% Текст слайда подготовила Королева Дарья 5030102/20201 
% Дата 15.03.26
\begin{frame}
\frametitle{1.3.7. Грамматики $LL(k)$ и $LR(k)$ (6/9)}
Рассмотрим некоторые следствия определения 
$LR(k)$-грамматики, необходимые для конструирования
эффективного синтаксического анализатора.
Пусть $\langle \mathrm{M}, \Sigma, \Pi, S \rangle$ --- некоторая КСГ и
$S\Rightarrow^*_r \alpha A \omega \Rightarrow_r \alpha \beta \omega
\Rightarrow_r^*\gamma\omega \in \Sigma^*$~---
правый терминальный вывод.
В этом случае говорят, что цепочку $\alpha \beta \omega$ 
можно \odmkey{свернуть слева}{Свёртка!слева} по правилу
$A\to\beta$, а вхождение цепочки $\beta$ в цепочку 
$\alpha \beta \omega$ называется \odmkey{основой}{Основа},
а любой префикс цепочки $\alpha\beta$ называется 
\odmkey{активным}{Префикс!активный}.
При этом структура
\(
[A \to \beta_1 \cdot \beta_2,\, \gamma]
\)
называется \odmkey{$LR(k)$-ситуацией}{Ситуация!$LR(k)$},  
если $A \to \beta_1 \beta_2$ --- правило из $\Pi$, 
а $\gamma \in \Sigma^{*}$. 
Говорят, что $LR(k)$-ситуация $[A \to \beta_1 \cdot \beta_2,\, \gamma]$ 
\odmkey{допустима для активного префикса}{Ситуация!Допустимая} $\alpha \beta_1$, если существует вывод
\(
S' \Rightarrow_r^* \alpha A \omega \Rightarrow_r \alpha \beta_1 \beta_2 \omega,
\)
такой, что $\gamma = \operatorname{FIRST}_k(\omega)$.    
\begin{example}
    Рассмотрим грамматику \(G\) с правилами 
    \(S \to C \mid D, C \to aC \mid b, D \to aD\mid c\). 
    Ситуация \([C \to a \cdot C, \varepsilon]\) допустима для \(aaa\), 
    поскольку \(S \Rightarrow_r^* aaC \Rightarrow_r aaaC\), 
    т.~е. в этом примере 
    \(\alpha = aa\) и \(\omega = \varepsilon\)
\end{example} 

\smallskip
Определим функцию $\operatorname{EFF}_k^G(\alpha)$ так:
\begin{enumerate}
    \item если $\alpha$ начинается терминалом, то
    \(
        \mathrm{EFF}_k(\alpha) = \mathrm{FIRST}_k(\alpha);
    \)
    \item если $\alpha$ начинается нетерминалом, то
    \(
        \mathrm{EFF}_k(\alpha) = 
        \{\, \omega \mid 
        \omega \in \mathrm{FIRST}_k(\alpha)\) 
         и существует вывод 
        $\alpha \Rightarrow_r^* \beta \Rightarrow_r \omega \delta$,
        где  $\beta \neq A \omega \delta$  для $ A \in \mathrm{M} \,\}$.
\end{enumerate}

\end{frame}

% Текст слайда подготовила Королева Дарья 5030102/20201 
% Дата 15.03.26
\begin{frame}
\frametitle{1.3.7. Грамматики $LL(k)$ и $LR(k)$ (7/9)}
\begin{theorem}
Пусть $\langle \mathrm{M}, \Sigma, \Pi, S \rangle$ --- КС-грамматика.  
Тогда $G$ является $LR(k)$-грамматикой тогда и только тогда, когда для каждой цепочки $u \in \Sigma^*$ выполняется условие:
если $\alpha \beta$ --- активный префикс правовыводимой цепочки в пополненной грамматике $G'$ и если $LR(k)$-ситуация $[A \to \beta \cdot, \gamma]$ допустима для $\alpha \beta$,
то не существует другой $LR(k)$-ситуации 
$[A_1 \to \beta_1 \cdot \beta_2, \delta],$
допустимой для той же $\alpha \beta$, где $\gamma \in \mathrm{EFF}_k(\beta_2 \delta)$.
\end{theorem}
\TODO{с этого места необходимо продолжить переработку}
\begin{proof}
\proofsection{Необходимость.}
Пусть $[A \to \beta \cdot, u]$ и $[A_1 \to \beta_1 \cdot \beta_2, v]$ —
различные ситуации, допустимые для $\alpha \beta$.  
Тогда существуют выводы:
\[
\begin{aligned}
S' &\Rightarrow_r^* \alpha A \omega \Rightarrow_r \alpha \beta \omega,\\
S' &\Rightarrow_r^* \alpha_1 A_1 \delta \Rightarrow_r \alpha_1 \beta_1 \beta_2 \delta,
\end{aligned}
\]
где $\mathrm{FIRST}_k(\omega)=u$, $\mathrm{FIRST}_k(\delta)=v$ и $\alpha \beta = \alpha_1 \beta_1$. Кроме того, $\beta_2 \delta \Rightarrow^*_r u \lambda$ для некоторой цепочки $\lambda$ и в этом (возможно, пустом) выводе терминал левом конце никогда не заменяется пустой цепочкой. 
\end{proof}
\end{frame}

% Текст слайда подготовила Королева Дарья 5030102/20201 
% Дата 15.03.26
\begin{frame}
\frametitle{1.3.7. Грамматики $LL(k)$ и $LR(k)$ (8/9)}
\begin{proof} (Продолжение)
Рассмотрим 3 случая:
\begin{enumerate}
    \item $\beta_2 = \varepsilon$.
    Тогда $u=v$, и выводы имеют вид
    \(
    S' \Rightarrow_r^* \alpha A \omega \Rightarrow_r \alpha \beta \omega,\)
    \(S' \Rightarrow_r^* \alpha A_1 \delta \Rightarrow_r \alpha \beta_1 \delta.
    \)
    Так как ситуации различны, то либо $A \ne A_1$, либо $\beta \ne \beta_1$,
    и нарушается $LR(k)$-условие.
\item $\beta_2 = \mu$, $\mu \in \Sigma^+$.
Тогда \(
S' \Rightarrow_r^* \alpha A \omega \Rightarrow_r \alpha \beta \omega,\)
\(S' \Rightarrow_r^* \alpha A_1 \delta \Rightarrow_r \alpha \beta_1 \mu \delta.
\)
По условию $u \in \mathrm{EFF}_k(\beta_2 v) = \mathrm{EFF}_k(\mu v)$
Так как $\beta_2 = \mu$ состоит только из терминалов:
$\mathrm{EFF}_k(\mu v) = \mathrm{FIRST}_k(\mu v)$
Таким образом  $u \in \mathrm{FIRST}_k(\mu v)$
Противоречие аналогично
\item $\beta_2$ содержит нетерминал.
Пусть $\beta_2 \Rightarrow_r^* u_1 B u_3 \Rightarrow_r u_1 u_2 u_3$, где $u_1 u_2 \ne \varepsilon$.
Тогда:
\[
\begin{aligned}
S' &\Rightarrow_r^* \alpha A \omega \Rightarrow_r \alpha \beta \omega,\\
S' &\Rightarrow_r^* \alpha_1 A_1 \delta \Rightarrow_r \alpha_1 \beta_1 \beta_2 \delta 
     \Rightarrow_r \alpha_1 \beta_1 u_1 B u_3 \delta \Rightarrow_r \alpha_1 \beta_1 u_1 u_2 u_3 \delta.
\end{aligned}
\]
В $LR(k)$-грамматике требуется, чтобы
\(
\alpha A u_1 u_2 u_3 \delta = \alpha_1 \beta_1 u_1 B u_3 \delta,
\) т.е. \\
\(
\alpha A u_1 u_2 = \alpha_1 \beta_1 u_1 B,
\)
\(
A u_1 u_2 = \beta u_1 B,
\)
но это невозможно, так как $u_1 u_2 \ne \varepsilon$.
\end{enumerate}

\vspace{-12pt}
\end{proof}

\begin{remark}
именно здесь используется условие, что \(u \in \mathrm{EFF}_k(\beta_2, v).\)
Если бы в формулировке использовать только множество $\mathrm{FIRST}_k$,
цепочка $u_1 u_2$ могла бы быть пустой, и доказательство бы не прошло.
\end{remark}

\end{frame}

% Текст слайда подготовила Королева Дарья 5030102/20201 
% Дата 15.03.26
\begin{frame}
\frametitle{1.3.7. Грамматики $LL(k)$ и $LR(k)$ (9/9)}
\begin{proof} (Окончание)
\proofsection{Достаточность.}
Пусть $G$ --- не $LR(k$)-грамматика.  
Тогда существуют два вывода:
\[
\begin{aligned}
\quad & S' \Rightarrow_r^* \alpha A \omega \Rightarrow_r \alpha \beta \omega,\\
\quad & S' \Rightarrow^*_{r} \gamma B \delta \Rightarrow_{r} \gamma \beta \lambda= \alpha \beta \lambda,
\end{aligned}
\]
где $\mathrm{FIRST}_k(\omega) = \mathrm{FIRST}_k(\lambda) = u$,
но $\alpha A \lambda \ne \gamma B \delta$.
Выберем $\alpha \beta$ минимальной длины.  
Пусть $\alpha_1 A_1 y_1$ — последняя открытая цепочка в выводе:
\(
S' \Rightarrow_r^* \gamma B \delta \Rightarrow_r^* \alpha_1 A_1 y_1
\Rightarrow_r \alpha_1 \beta_1 \beta_2 y_1 \Rightarrow_r^* \alpha_1 \beta_1 \lambda.
\)
Тогда $\alpha_1 \beta_1 = \alpha \beta$, и
$u = \mathrm{FIRST}_k(\lambda) \in \mathrm{EFF}_k(\beta_2, \lambda)$. 

Ситуация \([A_1 \to \beta_1 \cdot \beta_2, v]\) допустима для \(\alpha \beta\), где \(v = \mathrm{FIRST}_k(y_1)\).
В то же время ситуация \([A \to \beta \cdot , u]\) допустима для \(\alpha \beta\), поэтому осталось доказать, что \(A_1 \to \beta_1 \cdot \beta_2\) отличается от \(A \to \beta \). 

Пусть \(A_1 \to \beta_1 \cdot \beta_2\) совпадает с \(A \to \beta \). Тогда 
\(
S' \Rightarrow_r^*  \alpha_1 A \lambda
\Rightarrow_r \alpha_1 \beta \lambda
\), где \(\alpha_1 \beta = \alpha \beta \), поэтому \(\alpha_1 = \alpha\) и \(\alpha A \lambda = \alpha B \delta \), что противоречит предположению о том, что \(G\) не $LR(k)$-грамматика 
\end{proof}
\end{frame}

% Текст слайда подготовила Королева Дарья 5030102/20201 
% Дата 15.03.26
\begin{frame}{Специальные классы грамматик I (10/10)}
\begin{remark}
     Можно дать неформальное, но наглядное определение \(LR(k)\)-грамматики в терминах деревьев выводов. Грамматика \(G\) будет \(LR(k)\)-грамматикой, если, просмотрев только часть кроны дерева вывода в этой грамматике, расположенную слева от данной внутренней вершины, и часть кроны, выведенную из нее, а также следующие \(k\) терминальных символов, можно установить, какое правило было применено к этой вершине.
\end{remark}

\TODO{ Этому замечанию надобно найти место где-то в другом месте. Слайд лишний. } 
 
\end{frame}

\subsubsection{1.3.8. Разделённые грамматики}\label{1.3.8.}
% Текст слайда подготовила Королева Дарья 5030102/20201 
% Дата 17.02.26
\begin{frame}
\frametitle{1.3.8. Разделённые грамматики (1/5)}
Введение понятия $LL(1)$-грамматики — это один из возможных способов наложить на грамматику такие условия, чтобы шаг предсказания очередного правила вывода можно было сделать детерминированным, используя только текущий входной символ.

\smallskip
Существует обобщающий подход разбора, соответствующий {\em разделённой} грамматике. 
Для разделенной грамматики шаг предсказания тривиален. В связи с этим современное языкотворчество имеет тенденцию к использованию преимущественно разделённых грамматик,
ввиду простоты их реализации и возможности элементарной проверки условий разделенности по записи правил.
В то время как проверка выполнения свойства 
$LL(k)$ требует применения специального алгоритма.
Более того, проблема определения того, существует ли такое число $k$,
что грамматика является $LL(k)$-грамматикой,
является алгоритмически неразрешимой.  
\end{frame}

\begin{frame}
\frametitle{1.3.8. Вставной слайд в разделённые грамматики}
Этот слайд не является частью материала. УДАЛИТЬ после понимания.

\smallskip
Предыдущий и особенно последующий тексты содержат опечатки,
пропуски определений, ошибки в обозначениях и т.~д.
Всё это надобно устранить обычным порядком.
Здесь дополнительно имеются ошибки в ПОСЛЕДОВАТЕЛЬНОСТИ изложения.
Что такое подход разбора? Откуда взялся шаг предсказания?
Как работает МП-преобразователь?  
Изложение построено верно по существу,
но неверно методически.
ПРИЛЕЖНЫМ студентам (для которых пишется этот учебник!)
затруднено изучение материала.
В этом учебнике алгоритмы анализа
изложены в \DMref{1.4.}{1.4.},
автоматы-распознаватели [должны быть] описаны в главе 2.
Нельзя забегать вперёд без предупреждения: студенты тоже люди.

\smallskip
Как решить проблему?
Первый вариант: насытить текст ссылками вперёд,
убедившись, что ссылки ведут туда, куда надо.

\smallskip
Второй вариант: 
изложить свойства разделённых грамматик чисто грамматически, в частности,
ДОКАЗАТЬ первое замечание и все недоказанные утверждения далее,
 а рассказы про разбор убрать.
Второй подход предпочтельнее, немного труднее,
существенно полезнее для автора (надобно разобраться, в чём дело)
и поэтому намного ЦЕННЕЕ в баллах. 
Хотя количество слайдов СОКРАТИТСЯ, но оценка ВОЗРАСТЁТ! 

\end{frame}

% Текст слайда подготовила Королева Дарья 5030102/20201 
% Дата 17.02.26
\begin{frame}
\frametitle{1.3.8. Разделённые грамматики (2/5)}

КСГ $G=\langle \mathrm{M}, \Sigma, \Pi, S \rangle$ называется \odmkey{простой}{Грамматика!простая}
(или \odmkey{разделённой}{Грамматика!разделённая}),
если для любого нетерминала $A \in\mathrm{M}$
правые части всех правил для этого нетерминала
начинаются с различных терминальных символов.
Другими словами, для любой пары $a\in \Sigma$, $A\in \mathrm{M}$
существует не более одного правила вида $A\to a\alpha$. 
\begin{remark}
Очевидно, что простая грамматика не содержит $\varepsilon$-правил
и является \\$LL(1)$-грамматикой.
\end{remark}
Для разделенной грамматики шаг предсказания максимально прост: если \( a \) — текущий входной символ, \( A \) — верхний символ магазина и в грамматике есть правило \( A \to a \alpha \), то в магазине надо заменить \( A \) на \( \alpha \) и сдвинуть входную головку. 

\smallskip
Обобщающий подход разбора: к \( A \)-правилам разделенной грамматики разрешается добавить \( e \)-правило \( A \to e \), а на шаге предсказания для текущих символов \( a \) и \( A \) разрешается применить правило \( A \to e \) (т.е. удалить \( A \) из магазина), если в грамматике нет правила вида \( A \to a \alpha{} \).

\begin{remark}
 Добавление \( e \)-правил к разделенным грамматикам позволяет выделить новые классы грамматик, позволяющие описывать более сложные языковые конструкции, сохраняя при этом эффективность разбора $\cal{O}(n)$
\end{remark}

\end{frame}

% Текст слайда подготовила Королева Дарья 5030102/20201 
% Дата 17.02.26
\begin{frame}
\frametitle{1.3.8. Разделённые грамматики (3/5)}

КСГ $G=\langle \mathrm{M}, \Sigma, \Pi, S \rangle$ в \odmkey{слабой нормальной форме Грейбах}{Грамматика!в слабой нормальной форме Грейбах} \DMref{п. 1.3.5}{1.3.5.}, у которой правые части правил вида \( A \to a\alpha \) и \( A \to b\beta \) начинаются разными символами, т. е. \( a \neq b \), то она называется \odmkey{псевдоразделенной}{Псевдоразделенная!Грамматика}.

Для формализации выбора правила по текущему символу для псевдоразделенной грамматики строится простой МП-распознаватель \( M_G = \langle \Sigma, \mathrm{M}, \delta, S \rangle \), где функция переходов \( \delta \) определяется следующим образом:
\begin{enumerate}
  \item Если в стеке нетерминал \( A \), а на входе символ \( a \), и существует правило \( A \to a\alpha \), то \( A \) заменяется на \( a\alpha \) в стеке без сдвига входной головки. \(\delta(a, A) = (a, 1)\)
  \item Если правила для символа \( a \) нет, но существует \( e \)-правило \( A \to e \), то нетерминал \( A \) удаляется из стека (также без сдвига головки). \(\delta(b, S) = (e, 0)\)
\end{enumerate}
Автомат служит для детерминизации шага предсказания в левом выводе. Он позволяет сопоставить нетерминал на вершине стека с текущим входным символом без преждевременного сдвига входной головки
\end{frame}

% Текст слайда подготовила Королева Дарья 5030102/20201 
% Дата 17.02.26
\begin{frame}
\frametitle{1.3.8. Разделённые грамматики (4/5)}

\begin{example}
Пусть \( G_1 \) — псевдоразделенная грамматика с правилами
\(
S \to aSA \mid e,  A \to a \mid b.
\)

Тогда \( M_{G_1} = (\{a, b\}, \{S, A\}, \delta, S) \), где
\[
\begin{aligned}
\delta(a, S) &= (SA, 1)  \\
\delta(b, S) &= (e, 0) \\
\delta(\$, S) &= (e, 0) \\
\delta(a, A) &= (e, 1) \\
\delta(b, A) &= (e, 1)
\end{aligned}
\]
\end{example}


\begin{remark}
МП-распознаватель ограничивается анализом одного текущего символа. \\  
Обобщением этого подхода является $k$-предсказывающий алгоритм, где решение принимается на основе аванцепочки длиной $k$. Конфигурация алгоритма представляет собой необработанный вход, содержимое магазина и выходную ленту (номера правил). Работа алгоритма координируется управляющей таблицей, которая фактически реализует функцию \( \delta \) МП-распознавателя. 
\end{remark}
\end{frame}

% Текст слайда подготовила Королева Дарья 5030102/20201 
% Дата 17.02.26
\begin{frame}{1.3.8. Разделённые грамматики (5/5)}
Псевдоразделенная грамматика G называется \odmkey{слаборазделенной}{Слаборазделенная! грамматика}, если $\frak{L}(G) = \frak{L}(M_G)$ для соответствующего простого МП-преобразователя $M_G$

\begin{example}
    Грамматика $G$ с правилами \(S \to a S A \mid e, A \to b \mid e\) слаборазделенная

    Порождаемый язык \(\frak{L} (G) = \{ a^n b^n \mid  0 \leq k \leq n \} \)

    Логика простого МП-распознавателя \( M_G \)

\begin{enumerate}
\item \( \delta(S, a) = (SA, 1) \). При чтении \( n \) символов \( a \) в стеке накапливается \( n \) \( A\).
\item \( \delta(S, b) = (e, 0) \). При встрече первого \( b \) нетерминал \( S \) удаляется из стека без сдвига входа.
\item \( \delta(A, b) = (e, 1) \). Символы \( b \) последовательно удаляют \( A \) из стека.
\item \( \delta(A, \$) = (e, 0) \). Избыточные нетерминалы \( A \) (если \( k < n \)) стираются на маркере конца строки.
\end{enumerate}

\( \frak{L}(G) = \frak{L}(M_G) \), грамматика является слаборазделенной.

\begin{remark}
    Если цепочка начинается с \( b \), стек опустеет раньше времени. Если количество \( b \) превысит \( n \), для лишних символов в стеке не останется нетерминалов \( A\).
\end{remark}
\end{example}
\begin{remark}
    Не существует алгоритма, который по произвольной КСГ определяет, является ли она слаборазделенной.
\end{remark}
\end{frame}

% Текст слайда подготовила Королева Дарья 5030102/20201 
% Дата 17.02.26
% Внесены изменения 15.05.26
\begin{frame}
\frametitle{1.3.8. Разделённые грамматики (1/5)}
Введение понятия $ LL(1) $-грамматики представляет собой один из способов наложить на контекстно-свободную грамматику такие ограничения, при которых выбор очередного правила вывода можно выполнить детерминированно, используя лишь один текущий входной символ.

\smallskip

Существует более строгий и в то же время более простой подход, основанный на понятии {\emразделенной} грамматики. 
Для таких грамматик шаг предсказания правила тривиален: он определяется непосредственно первым терминальным символом правой части правила. В связи с этим современное языкотворчество имеет тенденцию к использованию преимущественно разделённых грамматик,
ввиду простоты их реализации и возможности элементарной проверки условий разделенности по записи правил.

\smallskip

В отличие от этого, проверка свойства $ LL(k) $ для заданного $  k  $ требует применения специального алгоритма, а проблема существования такого $ k $, при котором грамматика является $ LL(k) $-грамматикой, алгоритмически неразрешима.


\end{frame}

% Текст слайда подготовила Королева Дарья 5030102/20201 
% Дата 17.02.26
% Внесены изменения 15.05.26
\begin{frame}
\frametitle{1.3.8. Разделённые грамматики (2/5)}

КСГ $G=\langle \mathrm{M}, \Sigma, \Pi, S \rangle$ называется \odmkey{простой}{Грамматика!простая}
(или \odmkey{разделённой}{Грамматика!разделённая}),
если для любого нетерминала $A \in\mathrm{M}$
правые части всех правил для этого нетерминала
начинаются с различных терминальных символов.
Другими словами, для любой пары $a\in \Sigma$, $A\in \mathrm{M}$
существует не более одного правила вида $A\to a\alpha$. 

\begin{example} Грамматика $ G $ с правилами
$  S \to aS \mid bA \mid c  $,
$  A \to dA \mid e  $
является разделённой.
\end{example}

\smallskip

\begin{remark}
Очевидно, что простая грамматика не содержит $\varepsilon$-правил
и является \\$LL(1)$-грамматикой.
\end{remark}

\begin{proof}
Отсутствие $\varepsilon$-правил очевидно: $\varepsilon$-правило начинается не с терминального символа, что нарушило бы условие разделённости при наличии хотя бы одного правила $A \to a\alpha$. 
Покажем, что $G$ является $LL(1)$-грамматикой. Предположим, что существуют два вывода
\(
S \Rightarrow_l^* \omega A\alpha \Rightarrow_l \omega\beta\alpha \Rightarrow^* \omega\delta,
\)
\(
S \Rightarrow_l^* \omega A\alpha \Rightarrow_l \omega\gamma\alpha \Rightarrow^* \omega\lambda,
\)
такие что $\mathrm{FIRST}_1(\delta) = \mathrm{FIRST}_1(\lambda)$, но $\beta \neq \gamma$.
Поскольку $G$ простая и $\beta \neq \gamma$, правые части правил начинаются с различных терминальных символов $\beta = a\beta'$, $\gamma = b\gamma'$ и $a \neq b$. Тогда
\(
\mathrm{FIRST}_1(\beta\alpha) = \{a\},  \mathrm{FIRST}_1(\gamma\alpha) = \{b\}.
\)
Так как первый символ является терминалом, он не изменяется при дальнейшем выводе, следовательно
\(
a \in \mathrm{FIRST}_1(\delta), b \in \mathrm{FIRST}_1(\lambda).
\)
Это противоречит равенству $\mathrm{FIRST}_1(\delta) = \mathrm{FIRST}_1(\lambda)$. Таким образом, $\beta = \gamma$.

Следовательно, $G$ является LL(1)-грамматикой.
\end{proof}

\end{frame}

% Текст слайда подготовила Королева Дарья 5030102/20201 
% Дата 17.02.26
% Внесены изменения 15.05.26
\begin{frame}
\frametitle{1.3.8. Разделённые грамматики (3/5)}

КСГ $G=\langle \mathrm{M}, \Sigma, \Pi, S \rangle$ в {\em слабой нормальной форме Грейбах} (см. \DMref{п.~1.3.5}{1.3.5.}), у которой правые части правил вида \( A \to a\alpha \) и \( A \to b\beta \) начинаются разными символами, т. е. \( a \neq b \), называется \odmkey{псевдоразделенной}{Грамматика!Псевдоразделенная}.

\begin{remark}
    Добавление $\varepsilon$-правил к разделенным грамматикам позволяет выделить новые классы грамматик, которые описывают более сложные языковые конструкции, сохраняя эффективность $\cal{O}(n)$.
\end{remark}

\smallskip

\begin{example}
Грамматика $G$ с правилами $S \to aSA \mid \varepsilon, A \to b \mid \varepsilon$
является псевдоразделённой.
\end{example} 
\begin{proof}
Согласно определению грамматика $  G  $ находится в слабой нормальной форме Грейбах.
Проверим условие псевдоразделённости.
Нетерминал $ S $ имеет только одно непустое правило $ S \to aSA $, которое начинается с символа $ a $.
Нетерминал $ A $ имеет только одно непустое правило $ A \to b $, которое начинается с символа $ b $.
Поскольку у каждого нетерминала все непустые правила начинаются с различных терминальных символов, грамматика $  G  $ является псевдоразделённой.
\end{proof}


\end{frame}

% Текст слайда подготовила Королева Дарья 5030102/20201 
% Дата 15.05.26
\begin{frame}
\frametitle{1.3.8. Разделённые грамматики (4/5)}
КСГ $G=\langle \mathrm{M}, \Sigma, \Pi, S \rangle$ называется 
\odmkey{слаборазделённой}{Грамматика!слаборазделённая}, если она 
является псевдоразделённой
и для любого нетерминала $A$ такого, что $A \to \varepsilon$, 
и для любого правила вида $B \to \beta A \gamma$ (где $\beta$ не содержит 
нетерминал $A$ в начале), первый терминальный символ цепочки $\gamma$ 
не принадлежит множеству $\mathrm{FIRST}(A)$.

\smallskip

\begin{example}
Грамматика $G$ с правилами
\(
S \to aSA \mid \varepsilon,  A \to b \mid \varepsilon
\) слаборазделенная и 
\(
\frak{L}(G) = \{ a^n b^k \mid 0 \le k \le n \}.
\)
\end{example}

\begin{proof} Сначала покажем, что \( \frak{L}(G) = \{ a^n b^k \mid 0 \le k \le n \} \) 

$(\subseteq)$ Каждый шаг применения $S \to aSA$ добавляет один $a$ и один $A$. После завершения $S \to \varepsilon$ остаётся ровно $n$ нетерминалов $A$, каждый из которых порождает не более одного $b$. Следовательно, $w = a^n b^k$ с $k \le n$.

$(\supseteq)$ Для цепочки $a^n b^k$ ($k \le n$) существует левый вывод:
\(
S \Rightarrow^* a^n S A^n \Rightarrow a^n A^n \Rightarrow a^n b^k A^{n-k} \Rightarrow a^n b^k.
\)

\smallskip

Рассмотрим слабую разделенность $ G $. Предположим, что существует $  w \in \frak{L}(G)  $, имеющая два различных левых вывода. Рассмотрим первый шаг, где они расходятся. 

\end{proof}

\end{frame}

% Текст слайда подготовила Королева Дарья 5030102/20201 
% Дата 15.05.26
\begin{frame}{1.3.8. Разделённые грамматики (5/5)}
\begin{proof} (продолжение)

В обоих выводах присутствует сентенциальная форма $  \omega A\alpha  $, и применяются разные правила $  A\to\beta  $ и $  A\to\gamma  $ ($  \beta\neq\gamma $). Поскольку $G$ является псевдоразделённой, возможны только два случая.

\textbf{Случай 1.} Оба правила начинаются с терминалов: $  \beta=a\beta'  $, $  \gamma=b\gamma'  $, $  a\neq b  $. Тогда 

\(\mathrm{FIRST}_1(\beta\alpha)=\{a\} \neq \{b\}=\mathrm{FIRST}_1(\gamma\alpha)  \), поэтому дальнейшие выводы не могут привести к одной цепочке $ w $.

\textbf{Случай 2.} Одно правило — $  \varepsilon  $-правило, без потери общности $ A\to\varepsilon $ и $ A\to a\delta $.
Применение $ A\to a\delta $ даёт первый символ $  a  $, который должен принадлежать $  \mathrm{FIRST}_1(\alpha) $.
Однако по условию псевдоразделённости первый терминальный символ цепочки $  \alpha  $ (следующей за $ A $) не принадлежит $  \mathrm{FIRST}(A) $. Таким образом, $ a\notin\mathrm{FIRST}(A) $.

В обоих случаях получаем противоречие. Следовательно, каждый $ w $ имеет единственный левый вывод.
\end{proof}

\smallskip

\begin{remark}
    Не существует алгоритма, который по произвольной КСГ определяет, является ли она слаборазделенной.
\end{remark}

\end{frame}


\subsubsection{1.3.9. Классификация контекстно-свободных грамматик}\label{1.3.9.}
\begin{frame}
\frametitle{1.3.9. Классификация контекстно-свободных грамматик}
\TODO{Итоговый обзор на тему: в каких случаях какие грамматики целесообразно применять. 
Здесь же про грамматики в регулярной форме.}

\TODO{Общая диаграмма онтология всех классов грамматик. 
Здесь же про класс КС-языков, что 
язык $\{a^n b^n c^n\}$ НС-язык, но не КС-язык,
и при этом $\{a^n b^n\}$ КС-язык, но 
не регулярный язык.}

\bigskip
\TODO{\color{blue} Параграф повис в воздухе. Кто???}
\end{frame}

% Текст слайда подготовил Швачко Никита 5030102/20202
% Дата 11.01.2026
\begin{frame}
\frametitle{1.3.9. Классификация контекстно-свободных грамматик (2/5)}

В курсе рассмотрены различные классы грамматик. 
Выбор подходящего класса зависит от задачи:

\medskip
\begin{tabular}{p{45mm}|p{100mm}}
\textbf{Класс грамматики} & \textbf{Когда применять} \\
\hline
\odmkey{Регулярные грамматики}{Грамматика!регулярная} & 
Лексический анализ, токенизация, поиск по образцу. 
Разбор за $\mathcal{O}(n)$, минимальные ресурсы. \\
\hline
\odmkey{$LL(k)$-грамматики}{Грамматика!$LL(k)$} & 
Ручное написание парсера (рекурсивный спуск). 
Простая реализация, понятная структура. \\
\hline
\odmkey{$LR(k)$-грамматики}{LR(k)!Грамматика} & 
Автоматическая генерация парсера (yacc, bison). 
Мощнее $LL(k)$, но сложнее для понимания. \\
\hline
Общие КСГ & 
Когда важна естественность описания, а не эффективность разбора.
Алгоритмы CYK, Эрли: $\mathcal{O}(n^3)$. \\
\end{tabular}

\end{frame}



% СЛАЙД 2: ИЕРАРХИЯ КЛАССОВ ГРАММАТИК
% Текст слайда подготовил Швачко Никита 5030102/20202
% Дата 11.01.2026
\begin{frame}
\frametitle{1.3.9. Классификация контекстно-свободных грамматик (3/5)}

\odmkey{Иерархия классов грамматик}{Иерархия!грамматик} по Хомскому и их свойства:

\begin{center}
\begin{tabular}{|c|c|c|c|}
\hline
\textbf{Тип} & \textbf{Название} & \textbf{Распознаватель} & \textbf{Сложность} \\
\hline
0 & С фразовой структурой & Машина Тьюринга & Полуразрешимо \\
\hline
1 & Контекстно-зависимые (НС) & ЛО-автомат & $\mathcal{O}(c^n)$ \\
\hline
2 & Контекстно-свободные & МП-автомат & $\mathcal{O}(n^3)$ \\
\hline
3 & Регулярные & Конечный автомат & $\mathcal{O}(n)$ \\
\hline
\end{tabular}
\end{center}

\begin{remark}
Внутри класса КСГ (тип 2) существует своя иерархия подклассов:
\[
\text{Регулярные} \subset LL(1) \subset LL(k) \subset \text{Однозначные КСГ} \subset \text{КСГ}.
\]
При этом для каждого фиксированного $k$: $LL(k) \subset LR(k)$.
\end{remark}

\end{frame}



% СЛАЙД 3: ПРИМЕРЫ ЯЗЫКОВ РАЗНЫХ КЛАССОВ
% Текст слайда подготовил Швачко Никита 5030102/20202
% Дата 11.01.2026
\begin{frame}
\frametitle{1.3.9. Классификация контекстно-свободных грамматик (4/5)}

\NB{Примеры языков, разделяющие классы:}

\begin{itemize}
\item \textbf{Регулярный язык:} $L_1 = \{a^n \mid n \ge 0\} = a^*$~--- распознаётся конечным автоматом.

\item \textbf{КС-язык, но не регулярный:} $L_2 = \{a^n b^n \mid n \ge 0\}$~--- 
требует МП-автомат (нужен стек для подсчёта $a$). 
Доказательство нерегулярности~--- лемма о разрастании.

\item \textbf{НС-язык, но не КС-язык:} $L_3 = \{a^n b^n c^n \mid n \ge 0\}$~--- 
требует линейно-ограниченный автомат. 
Доказательство~--- лемма Огдена (см. п.~1.3.4).

\item \textbf{Рекурсивно перечислимый, но не КЗ:} 
язык $\{w \mid \text{машина Тьюринга } M_w \text{ останавливается}\}$~--- 
проблема остановки неразрешима линейно ограниченным автоматом.
\end{itemize}

\begin{digress}
Каждый следующий класс строго содержит предыдущий. 
Это обосновывает необходимость изучения всей иерархии Хомского.
\end{digress}

\end{frame}



% СЛАЙД 4: ГРАММАТИКИ В РЕГУЛЯРНОЙ ФОРМЕ
% Текст слайда подготовил Швачко Никита 5030102/20202
% Дата 11.01.2026
\begin{frame}
\frametitle{1.3.9. Классификация контекстно-свободных грамматик (5/5)}

\odmkey{Грамматика в регулярной форме}{Грамматика!в регулярной форме}~--- 
это КСГ, в которой правые части правил записаны с использованием 
регулярных операций (объединение, конкатенация, итерация).

\begin{example}
Вместо правил $A \to aA \mid a$ можно записать $A \to a^+$ или $A \to aa^*$.
\end{example}

\NB{Преимущества регулярной формы:}
\begin{itemize}
\item Более компактная запись грамматики
\item Естественное представление повторяющихся конструкций
\item Используется в РБНФ (\DMref{п.~1.1.8}{1.1.8.}) и современных парсер-генераторах
\end{itemize}

\begin{remark}
Грамматика в регулярной форме эквивалентна обычной КСГ~--- 
регулярные операции можно раскрыть, введя дополнительные нетерминалы.
Например, $A \to \alpha^*$ заменяется на $A \to \varepsilon \mid \alpha A$.
\end{remark}

\bigskip
\TODO{Отличный ПЛАН! \\
Теперь каждый пункт раскрыть, обосновать, насытить ссылками и примерами. \\
Интересно, кто такие современные парсер-генераторы?}

\end{frame}


% СЛАЙД 1: КЛАССЫ КСГ И ИХ ОПРЕДЕЛЕНИЯ
% Текст слайда подготовил Швачко Никита 5030102/20202
% Дата 11.04.2026
\begin{frame}
\frametitle{1.3.9. Классификация контекстно-свободных грамматик (1/6)}

\odmkey{Классификация контекстно-свободных грамматик}{Грамматика!классификация}
важна как с точки зрения выразительной силы,
так и с точки зрения возможности эффективного синтаксического анализа.

\medskip
\begin{tabular}{p{38mm}|p{107mm}}
Класс & Анализатор и основная область применения \\ \hline
\odmkey{Регулярные}{Грамматика!регулярная} &
Конечные автоматы; лексический анализ; линейное время
(\DMref{п.~1.2}{1.2.}). \\ \hline

\odmkey{$LL(k)$-грамматики}{LL(k)!Грамматика} &
Нисходящий анализ; рекурсивный спуск, ANTLR
(\DMref{п.~1.3.5}{1.3.5.}, \DMref{п.~1.3.7}{1.3.7.}). \\ \hline

\odmkey{$LR(k)$-грамматики}{LR(k)!Грамматика} &
Восходящий анализ; yacc, bison
(\DMref{п.~1.3.6}{1.3.6.}, \DMref{п.~1.3.7}{1.3.7.}). \\ \hline

КСГ общего вида &
Универсальные алгоритмы разбора; CYK, Эрли
(\DMref{п.~1.3.2}{1.3.2.}, \DMref{п.~1.4.4}{1.4.4.}). \\
\end{tabular}

\begin{remark}
При проектировании языка желательно задавать правила
в простой и единообразной форме:
это облегчает чтение грамматики и помогает построению анализатора.
\end{remark}

\end{frame}


% СЛАЙД 2: ВЛОЖЕНИЯ ВНУТРИ ТИПА 2
% Текст слайда подготовил Швачко Никита 5030102/20202
% Дата 11.04.2026
\begin{frame}
\frametitle{1.3.9. Классификация контекстно-свободных грамматик (2/6)}

\NB{Вложения основных подклассов:}

\[
\text{Регулярные} \subsetneq LL(1) \subseteq LL(k) \subseteq LR(k) \subsetneq \text{КСГ},
\qquad k \ge 1.
\]

\begin{remark}
Для каждого фиксированного $k \ge 1$ справедливо включение
$LL(k)\subseteq LR(k)$
(\DMref{п.~1.3.7}{1.3.7.}).
Тем самым $LR(k)$-грамматики образуют более широкий
детерминированный класс, чем $LL(k)$-грамматики.
\end{remark}

\begin{digress}
Чем уже рассматриваемый класс грамматик,
тем обычно проще построить для него эффективный синтаксический анализатор.
\end{digress}

\end{frame}

% СЛАЙД 3: МЕСТО КСГ В ИЕРАРХИИ ХОМСКОГО
% Текст слайда подготовил Швачко Никита 5030102/20202
% Дата 11.04.2026
\begin{frame}
\frametitle{1.3.9. Классификация контекстно-свободных грамматик (3/6)}
\NB{Место КС-грамматик в иерархии Хомского:}

\[
\text{Тип 0} \supsetneq \text{Тип 1}
\supsetneq \text{Тип 2: контекстно-свободные грамматики}
\supsetneq \text{Тип 3}.
\]
\centerline{\DMref{п.~1.1.4}{1.1.4.}}

\medskip
Следовательно, контекстно-свободные грамматики занимают
промежуточное положение:
они существенно выразительнее регулярных грамматик,
но слабее контекстно-зависимых.

\medskip
\begin{remark}
Внутри типа~2 выделяются важные детерминированные подклассы
$LL(k)$ и $LR(k)$
(\DMref{п.~1.3.7}{1.3.7.}),
а регулярные грамматики оказываются частным случаем
контекстно-свободных.
\end{remark}

\medskip
\begin{digress}
Тем самым классификация КС-грамматик уточняет
общую иерархию Хомского с точки зрения алгоритмов разбора.
\end{digress}

\end{frame}
% СЛАЙД 4: ПРИМЕРЫ ЯЗЫКОВ, РАЗДЕЛЯЮЩИЕ КЛАССЫ
% Текст слайда подготовил Швачко Никита 5030102/20202
% Дата 11.04.2026
\begin{frame}
\frametitle{1.3.9. Классификация контекстно-свободных грамматик (4/6)}

\NB{Примеры языков, разделяющих классы:}

\begin{itemize}
\item \textbf{Регулярный язык:}
$L_1 = a^*$;
он распознаётся конечным автоматом
(\DMref{п.~1.2}{1.2.}).

\item \textbf{Контекстно-свободный, но не регулярный:}
$L_2 = \Set{a^n b^n}{n \ge 0}$;
его нерегулярность устанавливается леммой о разрастании
(\DMref{п.~1.2.8}{1.2.8.}).

\item \textbf{Контекстно-зависимый, но не контекстно-свободный:}
$L_3 = \Set{a^n b^n c^n}{n \ge 0}$;
этот язык не является контекстно-свободным
(\DMref{п.~1.3.4}{1.3.4.}).

\item Тем самым получаем строгую цепочку включений:
\[
\text{Регулярные языки} \subsetneq \text{КС-языки}
\subsetneq \text{КЗ-языки}.
\]
\end{itemize}

\begin{digress}
Принадлежность языка некоторому классу
и непринадлежность ему
доказываются разными средствами:
характеризационными теоремами,
леммами о разрастании и конструкциями вывода.
\end{digress}

\end{frame}



% СЛАЙД 5: РАСПОЗНАВАТЕЛИ И СЛОЖНОСТЬ
% Текст слайда подготовил Швачко Никита 5030102/20202
% Дата 11.04.2026
\begin{frame}
\frametitle{1.3.9. Классификация контекстно-свободных грамматик (5/6)}

\NB{Классификация и трудоёмкость синтаксического анализа:}

\begin{itemize}
\item Для \odmkey{регулярных грамматик}{Грамматика!регулярная}
и регулярных выражений распознавание осуществляется
конечными автоматами и имеет линейную сложность
(\DMref{п.~1.2}{1.2.}).

\item Для детерминированных контекстно-свободных грамматик
классов \odmkey{$LL(k)$}{LL(k)!Грамматика}
и \odmkey{$LR(k)$}{LR(k)!Грамматика}
синтаксический анализ также может быть организован
за линейное время
(\DMref{п.~1.3.7}{1.3.7.}).

\item Для контекстно-свободных грамматик общего вида
приходится применять универсальные алгоритмы разбора,
например CYK или Эрли
(\DMref{п.~1.4.4}{1.4.4.});
в общем случае они существенно дороже детерминированного анализа.
\end{itemize}

\begin{remark}
Поэтому практически важны не просто КС-грамматики как таковые,
а такие формы их задания,
которые допускают детерминированный синтаксический анализ.
\end{remark}

\end{frame}
% СЛАЙД 6: РЕГУЛЯРНАЯ ФОРМА ЗАПИСИ КСГ
% Текст слайда подготовил Швачко Никита 5030102/20202
% Дата 11.04.2026
\begin{frame}
\frametitle{1.3.9. Классификация контекстно-свободных грамматик (6/6)}

\NB{Практический смысл классификации:}

\begin{enumerate}
\item Регулярные грамматики применяются прежде всего
на этапе лексического анализа.

\item Контекстно-свободные грамматики служат основной
математической моделью синтаксиса языков программирования.

\item Среди КС-грамматик особую роль играют
детерминированные подклассы $LL(k)$ и $LR(k)$,
поскольку именно для них строятся эффективные анализаторы.

\item Поэтому при проектировании языка желательно задавать
синтаксис в регулярной и единообразной форме записи правил,
чтобы грамматика по возможности принадлежала
одному из детерминированных подклассов.
\end{enumerate}

\begin{digress}
Итак, классификация грамматик важна не только сама по себе,
но и как средство выбора формы задания синтаксиса,
методов разбора и средств построения компиляторов.
\end{digress}

\end{frame}



% GA_1_2025.tex, 11564–11564
\subsection{1.5. Методы определения семантики}\label{1.5.}

% GA_1_2025.tex, 11595–11595
\subsubsection{1.5.1. Постановка задачи определения семантики}\label{1.5.1.}

% GA_1_2025.tex, 11612–11647
% СЛАЙД 1: СЕМАНТИКА ЯЗЫКОВ, ОПРЕДЕЛЁННЫХ ГРАММАТИЧЕСКИ
% Текст слайда подготовил Швачко Никита 5030102/20202
% Дата 11.01.2026
\begin{frame}
\frametitle{1.5.1. Постановка задачи определения семантики (2/6)}

\NB{Схема определения семантики языка, заданного грамматикой:}

\medskip
\begin{center}
\begin{tabular}{ccccc}
\textbf{Исходный текст} & $\xrightarrow{\text{Лексер}}$ & \textbf{Токены} & 
$\xrightarrow{\text{Парсер}}$ & \textbf{Дерево разбора} \\
& & & & $\downarrow$ \\
& & & & \textbf{Семантический анализ} \\
& & & & $\downarrow$ \\
& & & & \textbf{Смысл (значение)}
\end{tabular}
\end{center}

\begin{remark}
Синтаксис языка определяет \emph{форму} предложений (что корректно записано).\\
Семантика языка определяет \emph{смысл} предложений (что это означает).
\end{remark}

\begin{example}
Для арифметического выражения $2 + 3 * 4$:\\
~--- синтаксис определяет дерево разбора с учётом приоритета операций,\\
~--- семантика определяет результат вычисления: $14$.
\end{example}

\bigskip
\TODO{Здесь и далее неплохой ПЛАН. \\
Уточнить и реализовать. }

\end{frame}

% GA_2_2025.tex, 1–708
\section{2. Состояния и переходы}
\begin{frame}
\frametitle{\Ok 2. Состояния и переходы}

Алгоритм в неформальном смысле можно мыслить
как {\em последовательность действий} (приводящих к изменению состояния) или 
как {\em переходы состояний} (происходящих в определённой последовательности).
Эти две точки зрения на алгоритмы в известном смысле двойственны и
взаимодополнительны. Последовательности действий хорошо
всем известны в форме машин с неограниченными регистрами, блок-схем или императивных языков программирования.
В этой главе рассматриваются переходы состояний.    

\smallskip
\odmkey{Конечный автомат}{Автомат!конечный} является, видимо, наипростейшей возможной моделью
устройств обработки и преобразования информации, относящейся 
к концепции переходов состояний. Несмотря на свою
простоту, понятие конечного автомата оказывается достаточно
содержательным, чтобы послужить основой большого числа важнейших
приложений.

\smallskip
\begin{remark}
Наряду с термином <<конечный автомат>>
(finite automaton) в англоязычной 
литературе употребляются термины
\odmkey{машина состояний}{Машина!состояний} (state machine)
и 
\odmkey{машина переходов состояний}{Машина!переходов состояний} (state transition machine).
Эти термины обозначают один и тот же класс объектов,
но различаются по контексту употребления.
В формальном математическом контексте 
обычно говорят <<конечный автомат>>, 
а в неформальном программистском контексте
говорят <<машина состояний>>.
\end{remark}

\end{frame}

\subsection{2.1. Конечные автоматы Мили и Мура}\label{2.1}
\begin{frame}
\frametitle{\Qm 2.1. Конечные автоматы Мили и Мура}
Так же, как и в случае с графами, для конечных автоматов имеется
множество определений, различающихся некоторыми деталями. Мы
начнем с классических определений автоматов, принадлежащих
Мили\footnote{George H. Mealy (1927--2010)} и 
Муру\footnote{Edward F. Moore (1925--2003)}.

\end{frame}


\subsubsection{2.1.1. Автомат Мили}\label{2.1.1.}
\begin{frame}
\frametitle{\Ok 2.1.1. Автомат Мили (1/2)}
\odmkey{Конечным автоматом (Мили)}{Автомат!конечный}
\odmkey{}{Автомат!Мили} \odmkey{}{Конечный автомат} называется
система $S=\langle A,Q,B,\delta,\lambda\rangle$, где
\begin{enumerate}
\item
$A=\{a_1,\ldots,a_m\}$~--- конечный алфавит \odmkey{входных
символов}{Символ!входной} \odmkey{(входной
алфавит)}{Алфавит!входной};
\item
$Q=\{q_1,\ldots,q_n\}$~--- конечный алфавит
\odmkey{состояний}{Алфавит!состояний}\odmkey{}{Состояние};
\item
$B=\{b_1,\ldots,b_k\}$~--- конечный алфавит \odmkey{выходных
символов}{Символ!выходной} \odmkey{(выходной
алфавит)}{Алфавит!выходной};
\item
$\delta:Q\times A\rightarrow Q$~--- \odmkey{функция
переходов}{Функция!переходов}, определяющая по состоянию и
входному символу новое состояние;
\item
$\lambda: Q\times A\rightarrow B$~--- \odmkey{функция
выходов}{Функция!выходов}, определяющая по состоянию и входному
символу выходной символ.
\end{enumerate}

Данная математическая модель может быть использована для
абстрактного описания поведения в следующем смысле. Если автомат
находится в состоянии $q_i\in Q$ и на вход поступает символ
$a_j\in A$, то автомат переходит в состояние $q_x=\delta(q_i,a_j)$
и выдает на выход символ $b_y=\lambda(q_i,a_j)$.

\begin{remark}
В базовое определение
автомата не включено понятие 
начального состояния (см. \DMref{п.~2.1.5}{2.1.5.}). Это позволяет задавать одним описанием сразу некоторый класс родственных объектов, но требует дополнительного внимания 
при разборе примеров.  
\end{remark}

\end{frame}

\begin{frame}
\frametitle{\Ok 2.1.1. Автомат Мили (2/2)}
\begin{example}
Автомат $\langle \{1\},
\{0,1\}, \{0,1\},\delta(0,1)\Assign 1,
\delta(1,1)\Assign 0, \lambda(0,1)\Assign 1,\lambda(1,1)\Assign 0\rangle$
с двумя состояниями
по входному символу инвертирует своё состояние
и выдаёт его в качестве выходного символа.
Такой автомат называется \odmkey{триггер}{Триггер}.   
\end{example}
\begin{digress}
Идею конечного автомата можно интерпретировать разными способами.
Можно считать, что  автомат последовательно читает из входного
потока данных входные символы, записывая выходные символы в
выходной поток данных. С другой стороны, можно считать, что
входной символ~--- это событие, которое произошло в системе, а
выходной символ~--- это процедура реакции, которую нужно выполнить
в ответ на возникновение данного события в данном состоянии. С
третьей стороны, можно считать, что автомат~--- это устройство,
которое работает <<по тактам>>, опрашивая на каждом такте
состояния своих входных регистров (входной символ) и меняя
состояния выходных регистров (выходной символ).

\smallskip
Ценность понятия <<конечный автомат>> заключается в том, что оно
отлично подходит для описания всех здесь перечисленных, а также
многих других конкретных интерпретаций.
\end{digress}

\medskip
\begin{example}
Автомат $\langle \{0, 1\},
\{1\}, \{0,1\},\delta(1,0)\Assign 1,
\delta(1,1)\Assign 1, \lambda(1,0)\Assign 1,\lambda(1,1)\Assign 0\rangle$
с одним состоянием
 инвертирует входной символ 
и выдаёт его в качестве выходного символа.
Такой автомат называется \odmkey{инвертор}{Инвертор}.   
\end{example}


\end{frame}

\subsubsection{2.1.2. Задание автомата таблицей}\label{2.1.2.}
\begin{frame}
\frametitle{\Ok 2.1.2. Задание автомата таблицей (1/2)}

Поскольку функции $\delta$ и $\lambda$ определены на конечных
множествах, их можно задать \odmkey{таблицей конечного
автомата}{Таблица!конечного автомата}:

\begin{center}
\begin{tabular}{c|c|c|c|c|c}
    % after \\: \hline or \cline{col1-col2} \cline{col3-col4} ...
   & $a_1$ & $\ldots$ & $a_j$ & $\ldots$ & $a_m$ \\
  \hline
  $q_1$ &  &  &  &  &  \\
  $\ldots$ &  &  &  &  &  \\
  $q_i$ & $\ldots$ & $\ldots$ & $q_x, b_y$ & $\ldots$ & $\ldots$ \\
  $\ldots$ &  &  &  &  &  \\
  $q_n$ &  &  &  &  &  \\
  \end{tabular}
  \end{center}
где $q_x=\delta(q_i,a_j)$, а $b_y=\lambda(q_i,a_j)$.

\begin{columns}[t]
\begin{column}{10cm}
\begin{example}
Пусть $N$~--- натуральное число, а строка $\# a_1\ldots a_s$~---
запись числа в двоичной системе счисления, $a_i\in \{0,1\}$.
Рассмотрим автомат $S=\langle A, Q, B, \delta, \lambda\rangle$,
где $A=\{0,1,\# \}$, $Q=\{1,2,3\}$, $B=\{0, 1, \#\}$, а функции
$\delta$ и $\lambda$ заданы следующей таблицей:
\end{example}
\end{column}
\begin{column}{5cm}
\begin{center}
\begin{tabular}{c|c|c|c}
  & 0 & 1 & $\#$ \\
  \hline
  1 & 3, 1 & 2, 0 & 1, $\#$ \\
  2 & 3, 1 & 2, 0 & 1, 1 \\
  3 & 3, 0 & 3, 1 & 1, $\#$ \\
 \end{tabular}
\end{center}
\end{column}
\end{columns}
\end{frame}

\begin{frame}
\frametitle{\Ok 2.1.2. Задание автомата таблицей (2/2)}

\begin{example} (Продолжение)
Если автомат находится в состоянии 1 и на его вход подана
последовательность символов $a_s, a_{s-1},\ldots,a_1, \#$ (запись
числа $N$, прочитанного справа налево; символ \# является маркером,
ограничивающим запись), то автомат выдаст
последовательность цифр в двоичной записи числа $N+1$ (также в
порядке справа налево).
Например, в следующей таблице
приведен протокол работы автомата
при обработке записи чисел $5$ ($101\#$) слева и $7$ ($111\#$) справа.
\begin{center}
\begin{tabular}{c|c|c|c||c|c|c|c}
    % after \\: \hline or \cline{col1-col2} \cline{col3-col4} ...
  $a$ & $q$ & $\delta$ & $\lambda$ & $a$ & $q$ & $\delta$ & $\lambda$ \\
  \hline
  1 & 1 & 2 & 0 & 1 & 1 & 2 & 0\\
  0 & 2 & 3 & 1 & 1 & 2 & 2 & 0\\
  1 & 3 & 3 & 1 & 1 & 2 & 2 & 0\\
$\#$ & 3 & 1 & $\#$ & $\#$ & 2 & 1 & 1\\
\end{tabular}
\end{center}
\end{example}

\begin{digress}
Конечное число состояний позволяет
<<запомнить и учесть>> 
только {\em конечное} число поступивших символов или событий.  
Поэтому невозможно 
составить конечный автомат, 
который вычисляет $N+1$ для {\em любого}
числа $N$, если цифры числа поступают на вход автомата в привычном порядке~--- слева направо.
В этом проявляется ограниченность конструкции
конечного автомата.   
\end{digress}

\end{frame}


\subsubsection{2.1.3. Задание автомата графом переходов состояний}
\label{2.1.3.}
\begin{frame}
\frametitle{\Ok 2.1.3. Задание автомата графом переходов состояний (1/2)}

Конечный автомат можно задать \odmkey{графом
переходов}{Граф!переходов} (\odmkey{графом переходов
состояний}{граф!переходов!состояний}, \odmkey{диаграммой
состояний}{диаграмма!состояний})~--- ориентированным нагруженным
псевдомультиграфом, узлы которого помечены состояниями автомата
и, если $\delta(q_i,a_j)=q_x$ и $\lambda(q_i,a_j)=b_y$, то в графе
имеется дуга $(q_i,q_x)$, помеченная парой $(a_j,b_y)$.

\begin{example}
Диаграмма графа переходов автомата из предыдущего примера ($N\Assign N+1$)
представлена на рисунке.
\begin{figure}
\begin{center}
\includegraphics[width=120mm]{2_1_3_1.jpg}
\end{center}
\end{figure}
\end{example}

\end{frame}



\begin{frame}
\frametitle{2.1.3. Задание автомата графом переходов состояний (2/2)}
Одно и то же (или очень близкое)
поведение может быть задано существенно
различными автоматами,
подобно тому, как одну и ту же
функцию можно запрограммировать
разными способами.


\begin{example}
На рисунке
приведены диаграмма и таблица автомата,
выполняющего то же преобразование ($N\Assign N+1$),
что и автомат из предыдущих примеров. 
\end{example}

\vspace{-20pt}
\begin{columns}[t]
\begin{column}{7.5cm}

\vspace{10pt}
\begin{center}
\begin{tabular}{c|c|c|c}
   & 0 & 1 & $\#$ \\
  \hline
  1 & 2, 1 & 1, 0 & 1, 1 \\
  2 & 2, 0 & 2, 1 & 1, $\#$ \\
\end{tabular}
\end{center}
\end{column}
\begin{column}{7.5cm}
\begin{center}
\includegraphics[width=60mm]{2_1_3_2.jpg}
\end{center}
\end{column}
\end{columns}

\begin{digress}
Поскольку автомат может служить
средством описания алгоритмов,
составление автоматов является программированием.
В этом контексте рисование диаграмм
является одним из немногих примеров
визуального программирования,
где слово <<визуальный>> употреблено по праву.
\end{digress}

\end{frame}


\subsubsection{2.1.4. Полные и частичные автоматы}\label{2.1.4.}
\begin{frame}
\frametitle{\Ok 2.1.4. Полные и частичные автоматы (1/2)}

Автомат называется \odmkey{полным}{Автомат!полный}, если функции
$\delta$ и $\lambda$ тотальны, то~есть определены для всех пар
значений аргументов. В таблице полного автомата все клетки
заполнены, а граф переходов является регулярным графом степени $m=|A|$ 
(то~есть в любом состоянии исходящие дуги определены для всех
входных символов).

\begin{remark}
Для каждого состояния все исходящие дуги переходов должны быть
помечены {\em разными} входными символами~--- это \odmkey{условие
детерминированности}{Условие!детерминированности} графа переходов
конечного автомата. В некоторых случаях оказывается полезным
рассмотрение графов, подобных графам переходов, но с невыполненным
условием детерминированности. В этом случае из одного состояния
могут выходить несколько дуг, помеченных одним и тем же входным
символом. Такая конструкция называется \odmkey{недетерминированным
автоматом}{Автомат!недетерминированный}. Если не оговорено
противное, далее в этой главе подразумевается, что условие
детерминированности выполнено и рассматриваются детерминированные
автоматы.
\end{remark}

\medskip
Автомат называется \odmkey{частичным}{Автомат!частичный},  если
$\delta$ и/или $\lambda$~---  частичные функции. В таблице
частичного автомата некоторые клетки могут быть не заполнены, а в
графе переходов различные состояния могут иметь разное количество
исходящих дуг.
\end{frame}

\begin{frame}
\frametitle{2.1.4. Полные и частичные автоматы (2/2)}

\begin{digress}
Частичный автомат может быть неработоспособен для некоторых
последовательностей входных символов. Однако это обстоятельство не
мешает практическому применению частичных автоматов. Дело в том,
что не все возможные последовательности входных символов
фактически реализуются. Могут быть такие состояния, в которых
определённые события в принципе невозможны, а потому нет нужды
предусматривать реакцию на них.
\end{digress}

\begin{example}
Автоматы из предыдущих примеров допускают {\em любые}
последовательности символов $0$, $1$, $\#$. Но если на вход автомата
подается {\em правильная} двоичная запись числа
(то есть непустая последовательность 0 и 1,
которая заканчивается $\#$),
то достаточно
следующего частичного автомата (неопределённые значения функций
$\delta$ и $\lambda$ обозначены прочерками):
\end{example}

\vspace{-20pt}
\begin{columns}[t]
\begin{column}{8.4cm}
\begin{center}
\begin{tabular}{c|c|c|c}
  & 0 & 1 & $\#$ \\
  \hline
  1 & 3, 1 & 2, 0 & $-$, $-$ \\
  2 & 3, 1 & 2, 0 & 1, 1 \\
  3 & 3, 0 & 3, 1 & 1, $\#$ \\
 \end{tabular}
\end{center}

\vspace{-10pt}
Для упрощения изложения далее, если не ого\-ворено противное,
автомат считается полным.
\end{column}
\begin{column}{7cm}
\begin{center}
\includegraphics[width=70mm]{2_1_4.jpg}
\end{center}
\end{column}
\end{columns}

\end{frame}

\subsubsection{2.1.5. Начальное и заключительные состояния}\label{2.1.5.}
\begin{frame}
\frametitle{\Ok 2.1.5. Начальное и заключительные состояния (1/2)}

Автомат называется \odmkey{инициальным}{Автомат!инициальный}, если
в нем выделено {\em одно} состояние, которое называется
\odmkey{начальным}{Состояние!начальное}. Обычно начальному
состоянию дают наименьший номер: $1$ (или 0). Инициальный автомат
$S$ с начальным состоянием $q_1$ обозначим $S_{q_1}=\langle A, Q,
B, \delta, \lambda, q_1\rangle$. Если индекс $q_1$
подразумевается, то его можно опустить. Начальное состояние~---
это то состояние, в котором находится автомат, когда на вход
поступает первый символ. На диаграммах графов переходов мы будем
отмечать начальное состояние входной стрелкой, приходящей
извне.

Для данного автомата $S$ существуют $n$ инициальных автоматов, где
$n=|Q|$. Поведение автомата может существенно зависеть от того,
какое состояние является начальным.

\begin{example}
Если в автомате примера \DMref{п.~2.1.2}{2.1.2.} на вход автомата подать
последовательность цифр в записи числа $N$ в двоичной системе
счисления (справа налево) и при этом считать, что начальным
состоянием является состояние 1 или 2, то автомат выдает запись
числа $N+1$, а если начальным состоянием является 3, то автомат
выдает запись числа $N$, то~есть просто копию входной
последовательности.
\end{example}

В автомате может быть выделено некоторое подмножество состояний,
которые называются
\odmkey{заключительными}{Состояние!заключительное}. 
\end{frame}

\begin{frame}
\frametitle{2.1.5. Начальное и заключительные состояния (2/2)}
Иногда удобно
считать, что если автомат переходит в заключительное состояние, то
он прекращает свою работу (то~есть более не реагирует на входные
символы).
На
диаграммах графов переходов мы будем отмечать заключительные
состояния жирной границей.
% Для данного автомата подмножество заключительных
%состояний может быть определено $2^n$ способами, где $n=|Q|$.
\begin{example}
В автомат примера \DMref{п.~2.1.2}{2.1.2.} добавлено заключительное состояние
4, так что автомат завершает свою работу сразу по получении
символа $\#$.
\begin{center}
\begin{tabular}{c|c|c|c}
       & 0 & 1 & $\#$ \\
  \hline
  1 & 3,1 & 2,0 & $4$,$-$ \\
  2 & 3,1 & 2,0 & 4,1 \\
  3 & 3,0 & 3,1 & $4$,$-$ \\
  4 & $-$,$-$ & $-$,$-$ & $-$,$-$ \\
  \end{tabular}
\end{center}
\end{example}

\begin{remark}
Для заключительных состояний функции $\delta$ и $\lambda$ в
таблице автомата задавать не имеет смысла, поэтому строки для
заключительных состояний пусты. Их можно не включать в таблицу,
считая заключительными все состояния с номерами большими, чем
фактическое количество строк в таблице.
В графе переходов узлы
заключительных состояний не имеют исходящих дуг.
\end{remark}

\end{frame}


\subsubsection{2.1.6. Автоматное преобразование}
\begin{frame}
\frametitle{\Ok 2.1.6. Автоматное преобразование (1/2)}

Функции переходов $\delta$ и выходов $\lambda$ можно расширить,
чтобы они воспринимали в качестве второго аргумента {\it
последовательности} входных символов. Определим
\odmkey{расширенную функцию
переходов}{Функция!переходов!расширенная} $\delta^*\colon Q\times
A^* \rightarrow Q$ и \odmkey{расширенную функцию
выходов}{Функция!выходов!расширенная} \\$\lambda^*\colon Q\times A^*
\rightarrow B$  следующим образом:
$$
\A{a\in A}{\A{q\in Q}{\A{\alpha\in
A^*}{\delta^*(q,a)=\delta(q,a)\And\delta^*(q,\alpha a)=
\delta(\delta^*(q,\alpha),a)}}},
$$
$$
\A{a\in A}{\A{q\in Q}{\A{\alpha\in
A^*}{\lambda^*(q,a)=\lambda(q,a)\And\lambda^*(q,\alpha a)=
\lambda(\delta^*(q,\alpha),a)}}}.
$$

Инициальный автомат $S=\langle A, Q, B, \delta, \lambda, q\rangle$
определяет \odmkey{автоматное
преобразование}{Преобразование!автоматное}
$\Map{\mathcal{S}_q}{A^*}{B^*}$ следующим образом:
$$
\A{a \in A}{\mathcal{S}_q(a):= \lambda(q,a)\And \A{\alpha \in
A^*}{\mathcal{S}_q (\alpha a) := \mathcal{S}_q (\alpha) + \lambda
(\delta^*(q,\alpha),a)}},
$$
где $+$ обозначает операцию
\odmkey{конкатенации}{Операция!конкатенации} \odmkey{(сцепления)
строк}{Операция!сцепления}.

\begin{remark}
Если начальное состояние известно
из контекста, то индекс $q$ можно опустить.
\end{remark}

\end{frame}


\begin{frame}
\frametitle{\Ok 2.1.6. Автоматное преобразование (2/2)}


Автоматное преобразование обладает следующими свойствами:
\begin{enumerate}
\item $\A{\alpha\in A^* }{|\alpha|=|\mathcal{S}(\alpha)|}$~---
сохранение длины;

\item $\forall\alpha\in A^*\kern6pt \left(\alpha=\alpha_1\alpha_2
 \And  \beta=\mathcal{S}(\alpha)\And
\beta=\beta_1\beta_2 \And|\alpha_1|=|\beta_1|
\Impl\mathcal{S}(\alpha_1)=\beta_1\right)$~--- образ префикса
есть префикс образа.
\end{enumerate}

\vspace{-2pt}
\begin{remark}
В общем случае свойство 2 не выполняется для постфикса.
\end{remark}

Автоматное преобразование нетрудно реализовать программно,
например, как показано в следующем алгоритме. 
В этом алгоритме
массивы $\delta$ и $\lambda$, задающие соответственно функции переходов и
выходов, считаются глобальными.

\begin{algorithmic}
\REQUIRE входное слово 
\textbf{$\alpha:$ array $[1..s]$ of $A$},
начальное состояние $q\in Q$;
\ENSURE выходное слово 
\textbf{$\beta:$ array $[1..s]$ of $B$}.
\FOR{\textbf{$i$ from $1$ to $s$}}
   \STATE $\beta[i]:=\lambda\left[q,\alpha[i]\right]$
   \COMMENT{ выдать очередной выходной символ }
   \STATE $q:=\delta\left[q,\alpha[i]\right]$
   \COMMENT{ перейти в следующее состояние }
\ENDFOR
\end{algorithmic}


\begin{remark}
Чтобы упростить обозначения,
далее мы обозначаем
автоматное преобразование
той же буквой, что и сам автомат,
если при этом не возникает
недоразумений.
\end{remark}

\end{frame}


\subsubsection{2.1.7. Изоморфизм автоматов}
\label{2.1.7.}
\begin{frame}
\frametitle{\Ok 2.1.7. Изоморфизм автоматов (1/2)}

Говорят, что состояние $q_i$
\odmkey{достижимо}{Состояние!достижимое} из состояния $q_j$, если
$\E{\alpha\in A^*}{q_i = \delta^*(q_j, \alpha)}$. Обозначение:
$\langle q_j,q_i\rangle$. Автомат называется \odmkey{сильно
связным}{Автомат!сильно связный}, если $\A{q_i, q_j}{\langle
q_i,q_j\rangle}$. Сильно связный автомат имеет сильно связный граф
переходов.

В инициальном автомате состояние называется
\odmkey{недостижимым}{Состояние!недостижимое}, если оно
недостижимо из начального состояния. Недостижимые состояния
являются заведомо лишними~--- их присутствие (или отсутствие) не
влияет на работу автомата.

Пусть заданы два автомата $S=\langle A_S, Q_S, B_S, \delta_S, \lambda_S\rangle$ и
$T=\langle A_{T}, Q_{T}, B_{T}, \delta_{T}, \lambda_{T}\rangle$.
Говорят, что существует
\odmkey{гомоморфизм}{Гомоморфизм!автоматов} из $S$ в $T$
(обозначение $S\rightsquigarrow T$), если
существуют функции \\
$$\Map{f}{A_S}{A_T},\quad
   \Map{g}{Q_S}{Q_T},\quad 
   \Map{h}{B_S}{B_T}$$
   такие, что
$$   \A{a\in A_S, q\in Q_S, b\in B_S}
      {\delta_{T}(g(q),f(a))=
g(\delta_{S}(q,a)) \And \lambda_{T}(g(q), f(a))=
h(\lambda_{S}(q,a))}.
$$

Если функции $f$, $g$, $h$ взаимно-однозначны, то автоматы $S$ и $T$
\odmkey{изоморфны}{Изоморфизм!автоматов} (обозначение $S\sim T$).

\end{frame}

\begin{frame}
\frametitle{2.1.7. Изоморфизм автоматов (2/2)}

Изоморфизм автоматов означает, что можно так переименовать буквы в
алфавитах $A$, $Q$ и $B$, что таблицы автоматов совпадут.

Изоморфизм графов переходов (без учета нагрузки на дугах) не
означает изоморфизма автоматов.

\begin{example}
На рисунке автоматы $S$ и $T$ не изоморфны,
 хотя их
графы переходов изоморфны. Существует очевидный гомоморфизм из $T$
в $U$.

\begin{figure}[h]
\includegraphics[width=150mm]{2_1_7.jpg}
\end{figure}

Действительно, $S\not\sim T$,
поскольку легко подобрать функции $f$ и $g$, которые являются биекциями (например, тождественные функции),
но невозможно подобрать 
взаимно-однознач\-ную функцию $h$, так как $|\{0, 1\}|\ne |\{1\}|$.
С другой стороны, $T\rightsquigarrow U$,
поскольку достаточно взять функции $f$ и $h$ 
тождественными, и положить
$g(1)\Assign 1, g(2)\Assign 1$.
\end{example}
\end{frame}

\subsubsection{2.1.8. Автомат Мура}
\begin{frame}
\frametitle{\Ok 2.1.8. Автомат Мура (1/2)}
Конечный автомат называется \odmkey{автоматом Мура}{Автомат!Мура},
 если функция выходов
не зависит от входов, то~есть $\Map{\lambda}{Q}{B}$. Обычно в этом
случае функция обозначается $\mu$  и называется
\odmkey{функцией пометок}{Функция!пометок}, так~как каждому
состоянию сопоставляется пометка~--- выходной символ. В таблице
выходные символы пишут рядом с состояниями, а на диаграмме графа
переходов~--- в узлах, оставляя на дугах только
входные символы.

Автомат Мура работает следующим образом:


\begin{algorithmic}
\REQUIRE {\textbf {$\alpha:$ array $[1..s]$ of $A$, $q\in Q$}};
\ENSURE {\textbf {$\beta:$
 array $[1..s]$ of $B$}}.
 \FOR{\textbf{$i$
 from $1$ to $s$}}
 \STATE $q:=\delta(q,\alpha[i])$ \COMMENT{ переходим в состояние }
 \STATE $\beta[i]:=\mu(q)$ \COMMENT{ выдаем его пометку }
\ENDFOR
\end{algorithmic}

Очевидно, что автомат Мура является автоматом Мили. 
В \DMref{п.~2.3.5}{2.3.5.} показано, что возможности этих двух видов автоматов
совпадают.
\end{frame}


\begin{frame}
\frametitle{2.1.8. Автомат Мура (2/2)}
\begin{example}
Ниже приведены диаграмма и
таблица для автомата Мура, выполняющего то же преобразование
$(N\rightarrow N+1)$, что и автомат Мили из примера \DMref{п.~2.1.2}{2.1.2.}.

\begin{columns}
\begin{column}[T]{4cm}
Переименование пар (состояние, выходной символ) из примера \DMref{п.~2.1.2}{2.1.2.}: \\
    $(1,\# \rightarrow 1/\#)$\\ 
    $(1,1 \rightarrow 2/1)$ \\
    $(2,0 \rightarrow3/0)$\\
    $(3,0 \rightarrow 4/0)$ \\ 
    $(3,1 \rightarrow 5/1)$.\\
\end{column}
\begin{column}[T]{7cm}

\vspace{-10pt}
\begin{figure}[h]
  \includegraphics[width=60mm]{2_1_8.jpg}
\end{figure}
\end{column}
\begin{column}[T]{4cm}
Таблица автомата Мура

    \begin{tabular}{  l | l | l | l  }
     & 0 & 1 & \# \\ \hline
    1/\# & 5 & 3 & 1 \\
    2/1 & 5 & 3 & 1 \\
    3/0 & 5 & 3 & 2 \\
    4/0 & 4 & 5 & 1 \\
    5/1 & 4 & 5 & 1 \\
    \hline
\end{tabular}
\end{column}
\end{columns}
\end{example}

Понятия начальных и заключительных состояний, расширенных функций
переходов и пометок, а также автоматного преобразования очевидным образом
распространяются на случай автомата Мура.


\end{frame}


% GA_2_2025.tex, 739–2276
\subsection{2.2. Распознавание языков конечными автоматами}
\label{2.2.}
\begin{frame}
\frametitle{2.2. Распознавание языков конечными автоматами}
Конечные автоматы --- это один из способов описания алгоритмов.
Данный способ не универсален, но удивительно широко применим. В
этом разделе применимость конечных автоматов иллюстрируется на
примере задачи распознавания языков автоматами~---
одной из самых известных задач программирования.
\end{frame}


\subsubsection{2.2.1. Автоматы-распознаватели}
\label{2.2.1.}
\begin{frame}
\frametitle{2.2.1. Автоматы-распознаватели (1/2)}

Рассмотрим автомат Мура  с выходным алфавитом $B=\{0,1\}$.
Состояния, помеченные 1, будем называть
\odmkey{заключительными}{Состояние!заключительное}. Функцию
пометок $\mu$ можно задать, указав подмножество заключительных состояний $F\subset Q$.
Такой автомат $S=\langle A, Q, \delta, F
\rangle$ называется автоматом \odmkey{без выходов}{Автомат!без
выходов}. Инициальный автомат без выходов $S=\langle A, Q, \delta,
q_0, F\rangle$ называется
\odmkey{автоматом-распознавателем}{Автомат!распознаватель}
(здесь $q_0\in Q$ и
$F\subset Q$).

\begin{remark}
В данном случае попадание в заключительное состояние {\em не}
подразумевает прекращения работы автомата.
\end{remark}

Рассмотрим конечный алфавит $A$ и язык $L\subset A^*$
над алфавитом $A$.
Говорят, что язык $L$ \odmkey{распознается}{Распознавание!языка}
автоматом $S=\langle A, Q, \delta, q_0, F\rangle$, если

\vspace{-10pt}
$$
\A{\alpha\in A^* }{\alpha\in L \Equiv
\delta^*(q_0,\alpha)\in F}.
$$

\vspace{-4pt}

Другими словами, распознавание языка $L$ автоматом $S$ означает
следующее. Если на вход автомата $S$ подать <<правильное>> слово
$\alpha\in L$, то автомат перейдет в заключительное состояние
$f\in F$ (в этом случае говорят, что слово $\alpha$ {\em
допускается} автоматом $S$). Если же на вход автомата $S$ подать
<<неправильное>> слово $\alpha\not\in L$, то автомат перейдет в
состояние $g$, которое не является заключительным, $g\not\in F$
(говорят, что слово $\alpha$ {\em  не допускается} автоматом $S$).
\end{frame}


\begin{frame}
\frametitle{2.2.1. Автоматы-распознаватели (2/2)}
Всякому автомату-распознавателю $S$ соответствует язык $L$,
распознаваемый этим автоматом. А именно, язык $L$~--- это
множество значений всех путей (см. \DMref{п.~1.2.3}{1.2.3.}) из начального состояния в
заключительное.
\begin{remark}
Начальное состояние одно, а заключительных состояний может быть несколько,
но это обстоятельство никак не ограничивает общность модели.
\end{remark}

\smallskip
Таким образом, мы видим, что языки, распознаваемые конечными
автоматами, существуют. Из общей теории множеств следует, что
языки, не распознаваемые конечными автоматами, также существуют.
Действительно, с одной стороны, множество символов $A$ конечно,
множество слов $A^*$  счётно, а множество языков $2^{A^*}$~--- не
счётно. С другой стороны, множество автоматов с $k$ состояниями
конечно, поскольку это число различных орграфов с $k$ узлами,
обладающих свойством автоматной детерминированности. Но счётное
число конечных множеств счётно, поэтому множество автоматов
счётно. Между несчётным и счётным множествами взаимнооднозначного
соответствия быть не может.

\medskip
Возникает вопрос: каково же множество языков, распознаваемых
конечными автоматами?

\end{frame}

% Текст слайда подготовил Тишковец Сергей 5030102/20202 
% Дата 15.05.2026
\begin{frame}
\frametitle{2.2.1. Автоматы-распознаватели (1/3)}

Рассмотрим автомат Мура с выходным алфавитом $B=\{0,1\}$.
Состояния, помеченные 1, будем называть
\odmkey{заключительными}{Состояние!заключительное}. Функцию
пометок $\mu$ можно задать, указав подмножество заключительных состояний $F\subset Q$.
Такой автомат $S=\langle A, Q, \delta, F
\rangle$ называется автоматом \odmkey{без выходов}{Автомат!без
выходов}. Инициальный автомат без выходов $S=\langle A, Q, \delta,
q_0, F\rangle$ называется
\odmkey{автоматом-распознавателем}{Автомат!распознаватель}
(здесь $q_0\in Q$ и $F\subset Q$).

\begin{remark}
В данном случае попадание в заключительное состояние {\em не}
подразумевает прекращения работы автомата.
\end{remark}

\medskip
В теории распознавания языков используются сокращения, основанные на 
\odmkey{условии детерминированности}{Условие!детерминированности} (см.~\DMref{2.1.4}{2.1.4.}):
\begin{itemize}
    \item \textbf{ДКА} --- детерминированный конечный автомат.
    \item \textbf{НКА} --- недетерминированный конечный автомат.
\end{itemize}
Если не оговорено противное, далее в этой главе
под автоматом подразумевается ДКА.
\end{frame}

% Текст слайда подготовил Тишковец Сергей 5030102/20202 
% Дата 15.05.2026
\begin{frame}
\frametitle{2.2.1. Автоматы-распознаватели (2/3)}

Рассмотрим конечный алфавит $A$ и язык $L\subset A^*$ над алфавитом $A$.
Говорят, что язык $L$ \odmkey{распознается}{Распознавание!языка}
автоматом $S=\langle A, Q, \delta, q_0, F\rangle$, если

\vspace{-10pt}
$$
\A{\alpha\in A^* }{\alpha\in L \Equiv
\delta^*(q_0,\alpha)\in F}.
$$

\vspace{-4pt}

Другими словами, распознавание языка $L$ автоматом $S$ означает следующее:
\begin{itemize}
    \item Если на вход подать «правильное» слово $\alpha\in L$, 
    то автомат перейдет в заключительное состояние $f\in F$ 
    (в этом случае говорят, что слово {\em допускается} автоматом $S$).
    \item Если на вход подать «неправильное» слово $\alpha\not\in L$, 
    то автомат перейдет в состояние $g\not\in F$, которое не является заключительным 
    (говорят, что слово {\em не допускается} автоматом $S$).
\end{itemize}

\begin{remark}
Для НКА слово считается допускаемым, если существует хотя бы один 
путь из начального состояния в какое-либо заключительное, 
помеченный данным словом.
\end{remark}


\TODO{30.10.26. ФАН. Здесь недостаточно четко сформулирована главная мысль:
существуют языки, распознаваемые конечными автоматами,
существуют также языке НЕ распознаваемые конечными автоматами,
но при этом, может быть, распознаваемые какими-то другими средствами. 
Нужно починить следующий слайд и добавить примеры. 
}

\end{frame}

% Текст слайда подготовил Тишковец Сергей 5030102/20202 
% Дата 15.05.2026
\begin{frame}
\frametitle{2.2.1. Автоматы-распознаватели (3/3)}
Всякому автомату-распознавателю $S$ соответствует язык $L$,
распознаваемый этим автоматом. А именно, язык $L$~--- это
множество значений всех путей (см. \DMref{п.~1.2.3}{1.2.3.}) из начального состояния в
заключительное.
\begin{remark}
Начальное состояние одно, а заключительных состояний может быть несколько,
но это обстоятельство никак не ограничивает общность модели.
\end{remark}

\smallskip
Из общей теории множеств следует, что языки, 
не распознаваемые конечными автоматами, существуют:
\begin{itemize}
    \item Множество символов $A$ конечно,
    множество слов $A^*$  счётно, а множество языков $2^{A^*}$~--- не
    счётно.
    \item Множество автоматов с $k$ состояниями конечно, 
    поскольку это число различных орграфов с $k$ узлами, 
    обладающих свойством автоматной детерминированности.
\end{itemize}

\smallskip
Но счётное
число конечных множеств счётно, поэтому множество автоматов
счётно. Между несчётным и счётным множествами взаимнооднозначного
соответствия быть не может.

\medskip
Возникает вопрос: каково же семейство языков, распознаваемых конечными автоматами?
\end{frame}


\subsubsection{2.2.2. Распознаватели конечных языков}\label{2.2.2.}
\begin{frame}[shrink=22]
\frametitle{2.2.2. Распознаватели конечных языков (1/4)}


Начнём описание счётного подмножества распознаваемых автоматами
языков с тривиальных наблюдений.

\begin{theorem}
Любой конечный язык распознается конечным автоматом.
\end{theorem}

\begin{proof}
Пусть дан язык $L=\{\alpha_1,\dots,\alpha_n\}$. Сгруппируем слова
$\alpha_1,\dots,\alpha_n$  по первой букве, так что в одну группу
$G_{a_i}$ попадут все слова языка $L$, начинающиеся с одной и той
же буквы $a_i$. Множество всех первых букв обозначим $A_1$.
$$
A_1\subset A, ~L=\bigcup_{a_i\in A_1}G_{a_i},~
G_{a_i}\Assign\Set{\alpha}{\alpha\in L\And\alpha = a_i\beta\And
a_i\in A_1}.
$$

Множество непустых постфиксов (без первой буквы) слов из группы
$G_{a_i}$  обозначим $L_i$ ($L_i\Assign \Set{\beta}{\alpha = a_i
\beta\And \alpha\in G_{a_i} \And \beta\neq \varepsilon}$).

Построим граф переходов автомата $S$ следующим образом. Введем
начальное состояние $q_0$ и состояние недопустимости $\bar{q}$.
Введем состояния $q_i$ для всех $G_{a_i}$. Проведем дуги из $q_0$
в $q_i$ для всех $a_i \in A_1$ и дуги из $q_0$ в $\bar{q}$  для
всех остальных символов. Если в группе $G_{a_i}$ есть
однобуквенное слово $a_i$, то считаем $q_i$ заключительным
состоянием.
\end{proof}
\end{frame}

\begin{frame}
\frametitle{2.2.2. Распознаватели конечных языков (2/4)}
\begin{proof} (Окончание)
Далее рекурсивно применим эту процедуру к каждому языку $L_i$,
используя $q_i$ в качестве начального узла (состояние
недопустимости $\bar{q}$ можно использовать одно для всех языков).
Процесс, очевидно, закончится, так~как язык $L$ конечен. Построенный
граф является графом переходов автомата, распознающего язык $L$.
\end{proof}

\begin{example}
Для конечного языка $L=\{ab, ac, aba\}$ применение алгоритма,
изложенного в доказательстве предыдущей теоремы, дает
автомат-распознаватель, представленный на
рисунке. Здесь
состояние 1~--- начальное, а состояние 6 соответствует
недопустимости слова.

\vspace{-12pt}
\begin{figure}[h]
\begin{center}
\includegraphics[width=100mm]{2_3_2.jpg}
\end{center}
\end{figure}
\end{example}

\end{frame}

% Текст слайда подготовил Тишковец Сергей 5030102/20202 
% Дата 05.01.2026
% Примечание: существующие слайды из раздела 2.2.2 предлагается переименовать следующим образом:
% 2.2.2. Распознаватели конечных языков (1/3) -> 2.2.2. Распознаватели конечных языков (1/4)
% ... (2/3) -> ... (2/4)
% ... (3/3) -> ... (4/4)
\begin{frame}
\frametitle{2.2.2. Распознаватели конечных языков (3/4)}

\begin{example}
Рассмотрим конечный язык ключевых слов Python: $L = \{ \texttt{def}, \texttt{if}, \texttt{for} \}$.
Здесь состояние 1~--- начальное, а состояние 10 соответствует недопустимости слова.
Алфавит $A$ предполагается состоящим из строчных букв английского алфавита $\{a, b, \dots, z\}$.

\begin{figure}[h]
\begin{center}
\includegraphics[width=115mm]{2_2_2.jpg}
\end{center}
\caption{Граф переходов распознавателя ключевых слов Python}
\label{python_keywords_example}
\end{figure}
\end{example}
\end{frame}


\begin{frame}
\frametitle{2.2.2. Распознаватели конечных языков (4/4)}
\begin{remark}
Напомним,  что на рисунках здесь и далее использованы
следующие обозначения. Начальное состояние отмечено входящей извне
дугой. Заключительное состояние обозначено жирным кружком.
Множество символов, написанных на дуге, означает множество
соответствующих мультидуг.
\end{remark}

\medskip
\begin{digress} Если отбросить состояние недопустимости, то граф переходов
автомата-распознава\-теля конечного языка окажется ордеревом. Такое
ордерево является одним из самых эффективных способов
представления конечных словарей (множеств слов) в компьютере.
Время поиска слова в таком словаре пропорционально {\em длине} слова и
{\em не} зависит от общего количества слов в словаре.
Если словарь обладает \odmkey{свойством префикса}{Свойство!префикса},
то есть ни одно слово в словаре не является собственным началом никакого
другого слова, то листья дерева и только они 
оказываются заключительными состояниями.
Однако свойство префикса не является обязательным:
словарь одинаково хорошо работает для {\em любого} конечного множества слов.
Более того, слова можно выровнять по длине, дополнив каким-либо
нейтральным символом (пробелом) до максимальной длины,
и словарь примет форму выровненного дерева, в котором свойство
префикса выполнено автоматически.      
\end{digress}

\end{frame}


\subsubsection{2.2.3. Распознаватели бесконечных языков}
\begin{frame}
\frametitle{2.2.3. Распознаватели бесконечных языков (1/6)}
Множество языков, распознаваемых автоматами, не исчерпывается
конечными языками~--- автоматы могут распознавать и бесконечные
языки.

\begin{example}
Автомат на рисунке распознает бесконечный язык
$$
L=\{\varepsilon, aba, abaaba,\dots,(aba)^n,\dots\}
=\Set{(aba)^n}{n\geq 1}.
$$
Состояние $1$ является
начальным и единственным заключительным, состояние $4$ соответствует
недопустимости. 

\vspace{-10pt}
\begin{figure}[h]
\begin{center}
\includegraphics[width=75mm]{2_2_3_1.jpg}
\end{center}

\end{figure}
\end{example}

\end{frame}

\begin{frame}
\frametitle{2.2.3. Распознаватели бесконечных языков (2/6)}

Язык данного примера эффективно порождается процедурой, описанной
в следующем алгоритме.


\begin{algorithmic}
\STATE{\textbf {$\alpha\Assign\varepsilon$}}
\WHILE{\textbf{true}}
  \STATE{\textbf{yield $\alpha$}}
  \STATE{\textbf{$\alpha\Assign\alpha+  ``aba"$}}
\ENDWHILE
\end{algorithmic}

Однако существуют языки, в том числе эффективно порождаемые,
которые не распознаются никаким конечным автоматом.

Говорят, что автомат \odmkey{распознает}{Распознавание!последовательности }
(бесконечную) последовательность
букв, если он распознает язык, состоящий из всех префиксов данной
последовательности.

\begin{theorem}
Никакая бесконечная непериодическая последовательность не
распознаваема конечным автоматом.
\end{theorem}

\begin{proof}
От противного. Пусть бесконечная непериодическая
последовательность
$
\alpha= a_{i_1} a_{i_2}\dots
$
распознается
автоматом $S=\langle A,Q,\delta,q_0,F\rangle$
 с $n$ состояниями. 
\end{proof}  
\end{frame}

\begin{frame}
\frametitle{2.2.3. Распознаватели бесконечных языков (3/6)}
\begin{proof} (Продолжение)
 Заметим, что для любого начального отрезка
 $a_{i_1}\dots a_{i_j}$ состояние $q_{i_j}=\delta^*(q_0, a_{i_1}\dots
 a_{i_j})$~--- заключительное, то~есть автомат все время проходит
 через заключительные состояния. Возьмем начальный отрезок
 последовательности $\alpha$ длины больше $n$ и рассмотрим
 соответствующую последовательность состояний автомата. Поскольку
 $|Q|=n$, в этой последовательности есть одинаковые состояния
 $q_{i_j}$ и $q_{i_{j+s}}$,
 то~есть $\delta^*(q_{i_j}, a_{i_{j+1}}\dots a_{i_{j+s}})=
 q_{i_j}$. Тогда построим периодическую последовательность
 $\alpha_1\Assign a_{i_1}\dots a_{i_j}(a_{i_{j+1}}\dots
 a_{i_{j+s}})$.
Автомат $S$  для всех начальных отрезков последовательности
$\alpha_1$ будет находиться в заключительных состояниях, то~есть
автомат не отличает $\alpha$ от $\alpha_1$ и, таким образом, не
распознает $\alpha$.
 \end{proof}

\begin{example}
Язык
$$
L=\{a, ab, aba, abab, ababb, ababba, ababbab,ababbabb,
ababbabbb, ababbabbbabbbb, \dots\},$$
который описывается бесконечной
непериодической последовательностью
$$
ababbabbbabbbb \dots,$$
не распознается никаким
конечным автоматом, как видно из предыдущей теоремы.
\end{example}
\end{frame}

\begin{frame}
\frametitle{2.2.3. Распознаватели бесконечных языков (4/6)} 
Однако этот язык эффективно порождается следующим простым алгоритмом. 
\begin{algorithmic}
\STATE{$\alpha\Assign\varepsilon$; $n\Assign1$}
\WHILE{\textbf{true}}
  \STATE{\textbf{$\alpha\Assign\alpha+``a"$}}
  \STATE{\textbf{yield $\alpha$}}
  \FOR{\textbf{$i$ from $1$ to $n$}}
    \STATE{\textbf{$\alpha\Assign\alpha+``b"$}}
    \STATE{\textbf{yield $\alpha$}}
  \ENDFOR
  \STATE{\textbf{$n\Assign n+1$}}
\ENDWHILE
\end{algorithmic}

В двух последних параграфах приведены примеры языков, распознаваемых
и нераспознаваемых конечными автоматами. К сожалению, не
существует общего метода определения того, можно ли распознать
произвольный заданный язык конечным автоматом.

\TODO{Следующие два кадра неудовлетворительны, надо доделать/переделать}
\end{frame}

% Текст слайда подготовил Тишковец Сергей 5030102/20202 
% Дата 17.01.2026
\begin{frame}
\frametitle{2.2.3. Распознаватели бесконечных языков (5/6)}
\begin{theorem}
Свойство языка <<быть распознаваемым конечным автоматом>> алгоритмически неразрешимо.
\end{theorem}
\begin{proof}
От противного.
Предположим, существует алгоритм-анализатор $\Bl{R}$, который для любого генератора определяет, 
является ли порождаемый им язык регулярным.
Пусть $\Bl{M}$~--- произвольная машина Тьюринга (см.~\DMref{2.5.1}{2.5.1.}), 
$\omega$~--- входное слово. 
Построим процедуру-генератор $\Bl{P}$ над алфавитом $A=\{a, b\}$:
\begin{algorithmic}
  \STATE{$k \Assign 1$}
  \WHILE{\textbf{true}}
    \STATE{Симулировать работу $\Bl{M}$ на $\omega$ ровно $k$ шагов (определение шагов МТ см.~\DMref{2.5.1}{2.5.1.}).}
    \IF{$\Bl{M}$ остановилась за эти $k$ шагов}
       \STATE{\textbf{stop}}
    \ENDIF
    \STATE{\textbf{yield} $a^k b^k$}
    \STATE{$k \Assign k+1$}
  \ENDWHILE
\end{algorithmic}

\end{proof}
\end{frame}

% Текст слайда подготовил Тишковец Сергей 5030102/20202 
% Дата 17.01.2026
\begin{frame}
\frametitle{2.2.3. Распознаватели бесконечных языков (6/6)}
\begin{proof} (Продолжение)
Рассмотрим язык $L(\Bl{P})$, порождаемый этой процедурой.
Возможны два случая:
\begin{enumerate}
    \item $\Bl{M}$ зацикливается на $\omega$:
    в этом случае симуляция никогда не обнаруживает остановку, 
    процедура порождает все слова $a^1 b^1, a^2 b^2, \dots$.
    $$L = \{a^k b^k \mid k \ge 1\}$$
    Это бесконечный язык, который нерегулярен (по лемме о разрастании).
    \item $\Bl{M}$ останавливается через $k$ шагов:
 в этом случае порождается конечное множество
    $$L = \{a^1 b^1, \dots, a^{k-1} b^{k-1}\}$$
    Любой конечный язык регулярен (см.~\DMref{2.2.2}{2.2.2.}).
\end{enumerate}

Таким образом: $L(\Bl{P})$ регулярен тогда и только тогда, когда   $\Bl{M}$ останавливается на  $\omega$. Другими словами,
имея $\Bl{R}$, мы могли бы решить проблему остановки машины Тьюринга, которая известна как алгоритмически неразрешимая. 
Получили противоречие.
\end{proof}
\end{frame}

\subsubsection{2.2.4. Построение автомата по грамматике}\label{2.2.4.}
\begin{frame}
\frametitle{2.2.4. Построение автомата по грамматике (1/3)}
Мы видим, что в общем случае 
невозможно узнать, 
можно распознать язык автоматом или нельзя. 

\smallskip
Однако, если для задания языка используется какой-нибудь
конкретный способ, иногда по свойствам 
{\em описания} языка можно
сделать вывод о свойствах самого языка.

\smallskip
В предыдущей главе рассмотрены два средства задания
языков: формальные порождающие грамматики и регулярные выражения.
В этом разделе устанавливается эквивалентная связь между
автоматами-распознавателями,
регулярными выражениями и
автоматными грамматиками.


Первый шаг: покажем, что по автоматной грамматике можно построить
автомат, распознающий язык, порождённый грамматикой.
Входной алфавит автомата --- это терминальный алфавит грамматики $T$, а  множество состояний состоит из нетерминалов $N$,
 начального состояния $q_0$ и состояния недопустимости $\bar{q}$. 
\end{frame}


\begin{frame}
\frametitle{2.2.4. Построение автомата по грамматике (2/3)}

\vspace{-2pt}
{\linespread{0.9}\selectfont
\begin{algorithmic}
\REQUIRE Леволинейная автоматная грамматика $G=\langle N, T, R,I\rangle$.
\ENSURE Автомат-распознаватель $S=\langle T, Q,
\delta, q_0, F\rangle$,
распознающий язык $L(G)$.
\STATE{\textbf{$Q\Assign N\cup \{q_0\}\cup\{\bar{q}\}$}}
\COMMENT{по определению}
\FOR {\textbf{$r\in R$}}
  \IF{\textbf{$r=A\rightarrow a$ }}
    \STATE {\textbf{$\delta(q_0,a)\Assign \delta(q_0,a)\cup A$}}
    \COMMENT{ добавляем переход $q_0\stackrel{a}\longrightarrow A$ }
  \ELSE
    \STATE{\textbf{$\delta(B,a)\Assign \delta(B,a)\cup A$}}
    \COMMENT{ добавляем переход $B \stackrel{a}\longrightarrow A$}
  \ENDIF
\ENDFOR
\FOR{\textbf{$q\in Q$}}
  \FOR{\textbf{$a\in T$}}
    \IF{\textbf{$\neg\exists\delta(q,a)$}}
      \STATE{\textbf{$\delta(q,a)\Assign \bar{q}$}}
      \COMMENT{ недопустимые символы }
    \ENDIF
  \ENDFOR
\ENDFOR
\STATE {\textbf{$F\Assign\{I\}$}}
\end{algorithmic}
}
\end{frame}


\begin{frame}
\frametitle{2.2.4. Построение автомата по грамматике (3/3)}
\begin{ground}
Леволинейная автоматная грамматика имеет правила вида
$A\rightarrow a$ и \\$A\rightarrow Ba$.
Для правила $A\rightarrow a$ в автомат
добавляется переход $q_0\stackrel{a}\longrightarrow A$, а
для правила $A\rightarrow Ba$ в автомат
добавляется переход $B \stackrel{a}\longrightarrow A$.
После этого, для того, чтобы автомат был полным,
добавляются недостающие переходы в состояние недопустимости $\bar{q}$.
Построенный автомат $S$ действительно распознаёт язык,
порождаемый грамматикой $G$

Пусть $I \Rightarrow_{G}^{*} \alpha$ . . . .\TODO{доделать}

\end{ground}

В общем случае такой автомат может быть недетерминированным.
Тогда его нужно преобразовать в детерминированный вид
(аналогично тому, как это делается для недетерминированных
источников, см.~\DMref{п.~1.2.5}{1.2.5.}, исключить недостижимые
состояния и минимизировать.

\begin{example}


\TODO{исправить обозначения\\
Пример с автоматом для чисел . . . . \\
причём тот же автомат подходит к 1.2.4}
\end{example}

\end{frame}


% Текст слайда подготовил Фатахов Тимур Маратович 15.04.2026
\begin{frame}
\frametitle{2.2.4. Построение автомата по грамматике (1/5)}

Мы видим, что в общем случае невозможно заранее узнать,
можно распознать язык автоматом или нельзя.

\smallskip
Однако если язык задан некоторым \emph{конкретным способом},
то по свойствам его описания иногда удаётся сделать вывод
о свойствах самого языка.

\smallskip
В предыдущей главе были рассмотрены два средства задания языков:
формальные порождающие грамматики и регулярные выражения.
В этом разделе устанавливается связь между
автоматами-распознавателями,
регулярными выражениями и
автоматными грамматиками.

\smallskip
Первый шаг: покажем, что по автоматной грамматике можно построить
автомат, распознающий язык, порождённый этой грамматикой.

\smallskip
Пусть дана леволинейная автоматная грамматика
\[
G=\langle N,T,R,I\rangle.
\]
Входным алфавитом автомата служит терминальный алфавит грамматики \(T\),
а множество состояний состоит из нетерминалов \(N\),
начального состояния \(q_0\) и состояния недопустимости \(\bar q\).

\end{frame}

% Текст слайда подготовил Фатахов Тимур Маратович 15.04.2026
\begin{frame}
\frametitle{2.2.4. Построение автомата по грамматике (2/5)}

\vspace{-4pt}
{\linespread{0.92}\selectfont
\begin{algorithmic}
\REQUIRE Леволинейная автоматная грамматика \(G=\langle N,T,R,I\rangle\).
\ENSURE Автомат-распознаватель
\(S=\langle T,Q,\delta,q_0,F\rangle\), распознающий язык \(L(G)\).

\STATE{\textbf{\(Q \Assign N\cup\{q_0,\bar q\}\)}}

\FOR{\textbf{каждого правила \(r\in R\)}}
  \IF{\textbf{\(r=A\to a\)}}
    \STATE{\textbf{\(\delta(q_0,a)\Assign \delta(q_0,a)\cup\{A\}\)}}
  \ELSE
    \STATE{\textbf{\(\delta(B,a)\Assign \delta(B,a)\cup\{A\}\)}}
    \COMMENT{для правила \(A\to Ba\)}
  \ENDIF
\ENDFOR

\FOR{\textbf{каждого \(q\in Q\) и каждого \(a\in T\)}}
  \IF{\textbf{\(\neg\exists \delta(q,a)\)}}
    \STATE{\textbf{\(\delta(q,a)\Assign \{\bar q\}\)}}
  \ENDIF
\ENDFOR

\STATE{\textbf{\(F\Assign\{I\}\)}}
\end{algorithmic}
}

\end{frame}

% Текст слайда подготовил Фатахов Тимур Маратович 15.04.2026
\begin{frame}
\frametitle{2.2.4. Построение автомата по грамматике (3/5)}
\small
Леволинейная автоматная грамматика имеет правила вида
\[
A\to a,
\qquad
A\to Ba.
\]
Для правила \(A\to a\) в автомат добавляется переход \(q_0\xrightarrow{a}A\),
а для правила \(A\to Ba\) --- переход \(B\xrightarrow{a}A\).
После этого все недостающие переходы добавляются в состояние
недопустимости \(\bar q\), чтобы автомат был полным.

\smallskip
Покажем, что построенный автомат \(S\) действительно распознаёт язык,
порождаемый грамматикой \(G\).

Пусть \(\alpha=a_1a_2\dots a_n\) и
\[
I \Rightarrow_G A_{n-1}a_n
\Rightarrow_G A_{n-2}a_{n-1}a_n
\Rightarrow_G \dots
\Rightarrow_G A_1a_2\dots a_n
\Rightarrow_G a_1a_2\dots a_n.
\]
Тогда в автомате существует путь
\[
q_0\xrightarrow{a_1}A_1\xrightarrow{a_2}A_2\xrightarrow{a_3}\dots\xrightarrow{a_n}I.
\]

Обратно, если в автомате имеется путь
\[
q_0\xrightarrow{a_1}A_1\xrightarrow{a_2}A_2\xrightarrow{a_3}\dots\xrightarrow{a_n}I,
\]
то по построению переходов ему соответствуют правила
\[
A_1\to a_1,\quad A_2\to A_1a_2,\quad \dots,\quad I\to A_{n-1}a_n,
\]
и потому \(I\Rightarrow_G^* a_1a_2\dots a_n\).

\remark В общем случае такой автомат может быть недетерминированным. Тогда его нужно
детерминизовать, исключить недостижимые состояния и минимизировать.

\end{frame}

% Текст слайда подготовил Фатахов Тимур Маратович 15.04.2026
\begin{frame}
\frametitle{2.2.4. Построение автомата по грамматике (4/5)}
\small
\example Рассмотрим грамматику
\[
G=\langle N,T,R,I\rangle,
\qquad
N=\{I,A\},
\qquad
T=\{0,1\},
\qquad
R=\{I\to A1,\ A\to A0,\ A\to 1\}.
\]
Она порождает язык
\[
L(G)=\{10^n1\mid n\ge 0\}=\{11,101,1001,\dots\}.
\]

По правилам построения автомата получаем переходы
\[
q_0\xrightarrow{1}A,
\qquad
A\xrightarrow{0}A,
\qquad
A\xrightarrow{1}I.
\]
Все остальные переходы ведут в состояние \(\bar q\), а допускающим является \(I\).

Для цепочки \(1001\) имеем вывод
\[
I\Rightarrow_G A1\Rightarrow_G A01\Rightarrow_G A001\Rightarrow_G 1001
\]
и соответствующий путь
\[
q_0\xrightarrow{1}A\xrightarrow{0}A\xrightarrow{0}A\xrightarrow{1}I.
\]

\remark В этом примере построенный автомат уже детерминирован.

\end{frame}

% Текст слайда подготовил Фатахов Тимур Маратович 15.04.2026
\begin{frame}
\frametitle{2.2.4. Построение автомата по грамматике (5/5)}
\small
\example Рассмотрим грамматику десятичной записи неотрицательных целых чисел
без ведущих нулей:
\[
I\to 0 \mid d \mid Ae,
\qquad
A\to d \mid Ae,
\]
где \(d\in\{1,\dots,9\}\), \(e\in\{0,1,\dots,9\}\).

По правилам построения автомата получаем переходы
\[
q_0\xrightarrow{0}I,
\quad
q_0\xrightarrow{d}I,
\quad
q_0\xrightarrow{d}A,
\quad
A\xrightarrow{e}I,
\quad
A\xrightarrow{e}A.
\]

Например,
\[
I\Rightarrow_G A4\Rightarrow_G A04\Rightarrow_G 204,
\qquad
q_0\xrightarrow{2}A\xrightarrow{0}A\xrightarrow{4}I.
\]

Этот автомат недетерминирован: из \(q_0\) по цифрам \(1,\dots,9\) и из \(A\)
по любой цифре имеются по два перехода.
После детерминизации получаем автомат
\[
s\xrightarrow{0}z,
\qquad
s\xrightarrow{d}n,
\qquad
n\xrightarrow{e}n,
\]
где \(z\) и \(n\) --- допускающие, а остальные переходы ведут в \(\bar q\).
Это тот же автомат для записи чисел, который рассматривался в п.~1.2.4.

\end{frame}


% Текст слайда подготовил Фатахов Тимур Маратович 15.04.2026
\begin{frame}
\frametitle{2.2.4. Построение автомата по грамматике (1/5)}

Мы видим, что в общем случае невозможно заранее узнать,
можно распознать язык автоматом или нельзя.

\smallskip
Однако если язык задан одним из {\em конкретных} способов описания,
то по свойствам такого описания иногда удаётся сделать вывод
о свойствах самого языка.

\smallskip
В предыдущей главе были рассмотрены два средства задания языков:
формальные порождающие грамматики и регулярные выражения.
В этом разделе устанавливается связь между автоматами-распознавателями,
регулярными выражениями и автоматными грамматиками.

\smallskip
Сначала покажем, что по автоматной грамматике можно построить
автомат, распознающий язык, порождённый этой грамматикой.

\smallskip
Пусть дана леволинейная автоматная грамматика
\[
G=\langle \mathrm{M},\Sigma,\Pi,S\rangle.
\]
Входным алфавитом автомата служит терминальный алфавит \(\Sigma\),
а множество состояний состоит из нетерминалов \(\mathrm{M}\),
начального состояния \(q_0\) и состояния недопустимости \(\bar q\).

\end{frame}

% Текст слайда подготовил Фатахов Тимур Маратович 15.04.2026
\begin{frame}
\frametitle{2.2.4. Построение автомата по грамматике (2/5)}

Построим автомат-распознаватель
\[
\langle \Sigma,Q,\delta,q_0,F\rangle.
\]
Положим
\[
Q\Assign \mathrm{M}\cup\{q_0,\bar q\},
\qquad
F\Assign\{S\}.
\]

\smallskip
Далее для каждого правила из \(\Pi\) добавляются переходы:
если \(A\to a\in\Pi\), то
\[
A\in\delta(q_0,a),
\]
а если \(A\to Ba\in\Pi\), то
\[
A\in\delta(B,a).
\]

\smallskip
После этого для каждой пары \((q,a)\in Q\times\Sigma\),
для которой переход ещё не определён, полагаем
\[
\delta(q,a)\Assign\{\bar q\}.
\]

\smallskip
Построенный автомат полон; в общем случае он может быть
недетерминированным.

\end{frame}

% Текст слайда подготовил Фатахов Тимур Маратович 15.04.2026
\begin{frame}
\frametitle{2.2.4. Построение автомата по грамматике (3/5)}

\begin{ground}
Леволинейная автоматная грамматика имеет правила вида
\(A\to a\) и \(A\to Ba\).
Правилу \(A\to a\) соответствует переход \(q_0\xrightarrow{a}A\),
а правилу \(A\to Ba\) --- переход \(B\xrightarrow{a}A\).
Все недостающие переходы добавляются в состояние недопустимости \(\bar q\).

Покажем, что построенный автомат распознаёт язык \(L(G)\).
Пусть \(\alpha=a_1a_2\dots a_n\) и в грамматике \(G\) имеется вывод
\[
S \Rightarrow_G A_{n-1}a_n \Rightarrow_G \dots
\Rightarrow_G A_1a_2\dots a_n \Rightarrow_G a_1a_2\dots a_n.
\]
Это и есть последовательность секвенциальных форм вывода цепочки \(\alpha\).
Ей соответствует путь
\[
q_0\xrightarrow{a_1}A_1\xrightarrow{a_2}\dots\xrightarrow{a_n}S.
\]

Обратно, всякому такому пути соответствуют правила
\(A_1\to a_1,\ A_2\to A_1a_2,\ \dots,\ S\to A_{n-1}a_n\),
поэтому \(S\Rightarrow_G^* a_1a_2\dots a_n\).
Следовательно, построенный автомат распознаёт язык \(L(G)\).
\end{ground}

\end{frame}

% Текст слайда подготовил Фатахов Тимур Маратович 15.04.2026
\begin{frame}
\frametitle{2.2.4. Построение автомата по грамматике (4/5)}
\small

\begin{remark}
В общем случае построенный автомат может быть недетерминированным;
тогда его детерминизуют, исключают недостижимые состояния и минимизируют.
\end{remark}

\begin{example}
Рассмотрим грамматику
\[
G=\langle \mathrm{M},\Sigma,\Pi,S\rangle,\quad
\mathrm{M}=\{S,A\},\quad
\Sigma=\{0,1\},
\]
\[
\Pi=\{S\to A1,\ A\to A0,\ A\to 1\}.
\]
Она порождает язык \(L(G)=\{10^n1\mid n\ge 0\}\).

По правилам построения получаем переходы
\[
q_0\xrightarrow{1}A,\qquad
A\xrightarrow{0}A,\qquad
A\xrightarrow{1}S,
\]
причём допускающим состоянием служит \(S\), а все остальные переходы ведут в \(\bar q\).

Например,
\[
S\Rightarrow_G A1\Rightarrow_G A01\Rightarrow_G A001\Rightarrow_G 1001,
\qquad
q_0\xrightarrow{1}A\xrightarrow{0}A\xrightarrow{0}A\xrightarrow{1}S.
\]
\end{example}

\begin{remark}
В данном примере построенный автомат уже детерминирован.
\end{remark}

\end{frame}

% Текст слайда подготовил Фатахов Тимур Маратович 15.04.2026
\begin{frame}
\frametitle{2.2.4. Построение автомата по грамматике (5/5)}
\small

\begin{example}
Рассмотрим грамматику записи неотрицательных десятичных чисел
без ведущих нулей. Пусть
\[
\mathrm{M}=\{S,A\},
\qquad
\Sigma=\{0,1,\dots,9\},
\]
\[
\Pi=\{S\to 0\}\cup
\{S\to i,\ A\to i\mid 1\le i\le 9\}\cup
\{S\to Aj,\ A\to Aj\mid 0\le j\le 9\}.
\]

По правилам построения получаем переходы
\[
q_0\xrightarrow{0}S,\qquad
q_0\xrightarrow{i}S,\qquad
q_0\xrightarrow{i}A
\quad (1\le i\le 9),
\]
\[
A\xrightarrow{j}S,\qquad
A\xrightarrow{j}A
\quad (0\le j\le 9).
\]

Например,
\[
S\Rightarrow_G A4\Rightarrow_G A04\Rightarrow_G 204,
\qquad
q_0\xrightarrow{2}A\xrightarrow{0}A\xrightarrow{4}S.
\]

Этот автомат недетерминирован:
из \(q_0\) по цифрам \(1,\dots,9\) имеются переходы в \(S\) и в \(A\);
из \(A\) по любой цифре также имеются переходы в \(S\) и в \(A\).
\end{example}

\begin{remark}
После детерминизации получается тот же автомат для записи чисел,
который рассматривался в \DMref{п.~1.2.4}{1.2.4.}.
\end{remark}

\end{frame}

% Текст слайда подготовил Фатахов Тимур Маратович 15.04.2026
\begin{frame}
\frametitle{2.2.4. Построение автомата по грамматике (1/5)}

Мы видим, что в общем случае невозможно заранее узнать,
можно распознать язык автоматом или нельзя.

\smallskip
Однако если язык задан одним из {\em конкретных} способов описания,
то по свойствам такого описания иногда удаётся сделать вывод
о свойствах самого языка.

\smallskip
В предыдущей главе были рассмотрены два средства задания языков:
формальные порождающие грамматики и регулярные выражения.
В этом разделе устанавливается связь между автоматами-распознавателями,
регулярными выражениями и автоматными грамматиками.

\smallskip
Сначала покажем, что по автоматной грамматике можно построить
автомат, распознающий язык, порождённый этой грамматикой.

\smallskip
Пусть дана леволинейная автоматная грамматика
\[
G=\langle \mathrm{M},\Sigma,\Pi,S\rangle.
\]
Входным алфавитом автомата служит терминальный алфавит \(\Sigma\),
а множество состояний состоит из нетерминалов \(\mathrm{M}\),
начального состояния \(q_0\) и состояния недопустимости \(\bar q\).

\end{frame}

% Текст слайда подготовил Фатахов Тимур Маратович 15.04.2026
\begin{frame}
\frametitle{2.2.4. Построение автомата по грамматике (2/5)}

Построим автомат-распознаватель
\[
\langle \Sigma,Q,\delta,q_0,F\rangle.
\]
Положим
\[
Q\Assign \mathrm{M}\cup\{q_0,\bar q\},
\qquad
F\Assign\{S\}.
\]

\smallskip
Далее для каждого правила из \(\Pi\) добавляются переходы:
если \(A\to a\in\Pi\), то
\[
A\in\delta(q_0,a),
\]
а если \(A\to Ba\in\Pi\), то
\[
A\in\delta(B,a).
\]

\smallskip
После этого для каждой пары \((q,a)\in Q\times\Sigma\),
для которой переход ещё не определён, полагаем
\[
\delta(q,a)\Assign\{\bar q\}.
\]

\smallskip
Построенный автомат полон; в общем случае он может быть
недетерминированным.

\end{frame}

% Текст слайда подготовил Фатахов Тимур Маратович 15.04.2026
\begin{frame}
\frametitle{2.2.4. Построение автомата по грамматике (3/5)}

\begin{ground}
Леволинейная автоматная грамматика имеет правила вида
\(A\to a\) и \(A\to Ba\).
Правилу \(A\to a\) соответствует переход \(q_0\xrightarrow{a}A\),
а правилу \(A\to Ba\) --- переход \(B\xrightarrow{a}A\).
Все недостающие переходы добавляются в состояние недопустимости \(\bar q\).

Покажем, что построенный автомат распознаёт язык \(L(G)\).
Пусть \(\alpha=a_1a_2\dots a_n\) и в грамматике \(G\) имеется вывод
\[
S \Rightarrow_G A_{n-1}a_n \Rightarrow_G \dots
\Rightarrow_G A_1a_2\dots a_n \Rightarrow_G a_1a_2\dots a_n.
\]
Это и есть последовательность секвенциальных форм вывода цепочки \(\alpha\).
Ей соответствует путь
\[
q_0\xrightarrow{a_1}A_1\xrightarrow{a_2}\dots\xrightarrow{a_n}S.
\]

Обратно, всякому такому пути соответствуют правила
\(A_1\to a_1,\ A_2\to A_1a_2,\ \dots,\ S\to A_{n-1}a_n\),
поэтому \(S\Rightarrow_G^* a_1a_2\dots a_n\).
Следовательно, построенный автомат распознаёт язык \(L(G)\).
\end{ground}

\end{frame}

% Текст слайда подготовил Фатахов Тимур Маратович 15.04.2026
\begin{frame}
\frametitle{2.2.4. Построение автомата по грамматике (4/5)}
\small

\begin{remark}
В общем случае построенный автомат может быть недетерминированным;
тогда его детерминизуют, исключают недостижимые состояния и минимизируют.
\end{remark}

\begin{example}
Рассмотрим грамматику
\[
G=\langle \mathrm{M},\Sigma,\Pi,S\rangle,\quad
\mathrm{M}=\{S,A\},\quad
\Sigma=\{0,1\},
\]
\[
\Pi=\{S\to A1,\ A\to A0,\ A\to 1\}.
\]
Она порождает язык \(L(G)=\{10^n1\mid n\ge 0\}\).

По правилам построения получаем переходы
\[
q_0\xrightarrow{1}A,\qquad
A\xrightarrow{0}A,\qquad
A\xrightarrow{1}S,
\]
причём допускающим состоянием служит \(S\), а все остальные переходы ведут в \(\bar q\).

Например,
\[
S\Rightarrow_G A1\Rightarrow_G A01\Rightarrow_G A001\Rightarrow_G 1001,
\qquad
q_0\xrightarrow{1}A\xrightarrow{0}A\xrightarrow{0}A\xrightarrow{1}S.
\]
\end{example}

\begin{remark}
В данном примере построенный автомат уже детерминирован.
\end{remark}

\end{frame}

% Текст слайда подготовил Фатахов Тимур Маратович 15.04.2026
\begin{frame}
\frametitle{2.2.4. Построение автомата по грамматике (5/5)}
\small

\begin{example}
Рассмотрим грамматику записи неотрицательных десятичных чисел
без ведущих нулей. Пусть
\[
\mathrm{M}=\{S,A\},
\qquad
\Sigma=\{0,1,\dots,9\},
\]
\[
\Pi=\{S\to 0\}\cup
\{S\to i,\ A\to i\mid 1\le i\le 9\}\cup
\{S\to Aj,\ A\to Aj\mid 0\le j\le 9\}.
\]

По правилам построения получаем переходы
\[
q_0\xrightarrow{0}S,\qquad
q_0\xrightarrow{i}S,\qquad
q_0\xrightarrow{i}A
\quad (1\le i\le 9),
\]
\[
A\xrightarrow{j}S,\qquad
A\xrightarrow{j}A
\quad (0\le j\le 9).
\]

Например,
\[
S\Rightarrow_G A4\Rightarrow_G A04\Rightarrow_G 204,
\qquad
q_0\xrightarrow{2}A\xrightarrow{0}A\xrightarrow{4}S.
\]

Этот автомат недетерминирован:
из \(q_0\) по цифрам \(1,\dots,9\) имеются переходы в \(S\) и в \(A\);
из \(A\) по любой цифре также имеются переходы в \(S\) и в \(A\).
\end{example}

\begin{remark}
После детерминизации получается тот же автомат для записи чисел,
который рассматривался в \DMref{п.~1.2.4}{1.2.4.}.
\end{remark}

\end{frame}

\subsubsection{2.2.5. Синтез автоматов}
\begin{frame}
\frametitle{2.2.5. Синтез автоматов (1/3)}
Задача синтеза автомата~--- это задача построения автомата для заданного
языка.

Прежде всего, заметим,
что детерминированный источник и автомат-распознаватель~--- это
одно и то же~--- нагруженный символами орграф с выделенным начальным
узлом и множеством заключительных узлов.

\begin{theorem} {\NB{[Клини о синтезе автоматов]}}
Для любого регулярного языка существует конечный автомат $H$,
распознающий этот язык.
\end{theorem}

\begin{proof}
Возьмем регулярный язык $L$. По определению он описывается
регулярным выражением $E$. Согласно теореме~\DMref{1.2.4}{1.2.4.}, существует
двухполюсник $H$, представляющий $E$ (может быть, с большим
количеством лишних пустых дуг). По теореме
детерминизации существует
 детерминированный источник $H'$, представляющий $E$.

Детерминированный источник является 
автоматом-распознавателем.
\end{proof}

\begin{remark}
На практике в $H'$ может быть меньше $2^n$  узлов, так~как можно не
включать те подмножества, которые не достижимы из начальных
состояний. Однако в худшем случае число узлов равно всё-таки $2^n$.
\end{remark}

\TODO{В подходящем месте ввести сокращения НКА и ДКА.}
\end{frame}

% Текст слайда подготовил Тишковец Сергей 5030102/20202 
% Дата 17.01.2026
\begin{frame}
\frametitle{2.2.5. Синтез автоматов (2/3)}
\begin{theorem} \NB{[Оценка детерминизации]}
Пусть $S$~--- недетерминированный конечный автомат (НКА) с $n$ состояниями над алфавитом $A$.
Тогда эквивалентный ему детерминированный автомат (ДКА) $S'$ имеет:
\begin{itemize}
    \item в лучшем случае: $n$ состояний и $n \cdot |A|$ переходов (если $S$ уже детерминирован);
    \item в худшем случае: $2^n$ состояний и $2^n \cdot |A|$ переходов.
\end{itemize}
\end{theorem}
Верхний предел $2^n$ достижим (см.~\DMref{п.~1.2.4}{1.2.4.} для доказательства теоремы детерминизации 
и~\DMref{п.~2.2.6}{2.2.6.} для анализа).

\begin{example}
Рассмотрим язык $L=\Set{\omega \in A^*}{\text{символ на } k \text{-й позиции с конца равен } a}$, где $A=\{a,b\}$.
\begin{itemize}
    \item НКА для этого языка имеет $\cal{O}(k)$ состояний и $\cal{O}(k)$ переходов.
     
    \item Минимальный ДКА требует $2^k$ состояний и $2^k \cdot |A|$ переходов.
\end{itemize}
\end{example}
\begin{remark}
При синтезе из регулярного выражения длины $m$, промежуточный НКА имеет $\cal{O}(m)$ состояний и переходов, 
но итоговый ДКА может иметь $2^{\cal{O}(m)}$ состояний и $2^{\cal{O}(m)} \cdot |A|$ переходов.
\end{remark}
\end{frame}

% Текст слайда подготовил Тишковец Сергей 5030102/20202 
% Дата 17.01.2026
\begin{frame}
\frametitle{2.2.5. Синтез автоматов (3/3)}
Пусть даны два детерминированных конечных автомата с $n$ и $m$ состояниями соответственно. 
При построении автомата для операций над языками число состояний и переходов может возрастать мультипликативно в худших случаях.
\begin{table}[h]
\centering
\scriptsize
\setlength{\tabcolsep}{4pt}
\renewcommand{\arraystretch}{1.05}
\resizebox{\textwidth}{!}{
\begin{tabular}{|l|c|c|c|c|}
\hline
\textbf{Операция} & \multicolumn{2}{c|}{\textbf{Состояния}} & \multicolumn{2}{c|}{\textbf{Переходы}} \\ \cline{2-5}
 & \textbf{лучшее} & \textbf{худшее} & \textbf{лучшее} & \textbf{худшее} \\ \hline
Объединение ($L_1 \cup L_2$) & $\cal{O}$$(n + m)$ & $n \cdot m$ & $\cal{O}$$((n+m) \cdot |A|)$ & $n m \cdot |A|$ \\ \hline
Пересечение ($L_1 \cap L_2$) & $\cal{O}$$(n + m)$ & $n \cdot m$ & $\cal{O}$$((n+m) \cdot |A|)$ & $n m \cdot |A|$ \\ \hline
Дополнение ($A^* \setminus L_1$) & $n$ & $n$ & $n \cdot |A|$ & $n \cdot |A|$ \\ \hline
Конкатенация ($L_1 L_2$) & $\cal{O}$$(n + m)$ & $(n+m) \cdot 2^{n+m}$ & $\cal{O}$$((n+m) \cdot |A|)$ & $(n+m) 2^{n+m} \cdot |A|$ \\ \hline
\end{tabular}
}
\normalsize
\end{table}
\begin{remark}
При анализе детерминированного конечного автомата с $n$ состояниями по методу удаления состояний (см.~\DMref{п.~2.2.6}{2.2.6.}) 
размер полученного регулярного выражения в лучшем случае составляет $\cal{O}$$(n^2)$, а в худшем — $\cal{O}$$(\frac{2^n}{\sqrt{n}})$. 
На практике применяют минимизацию (алгоритм Хопкрофта) с временной сложностью
 $\cal{O}(n \log n)$.
\end{remark}
\TODO{ФАН 20 января. Хопкрофт?}
\end{frame}


% Текст слайда подготовил Тишковец Сергей 5030102/20202 
% Дата 15.05.2026
% Примечание: существующие слайды из раздела 2.2.5 предлагается переименовать следующим образом:
% 2.2.5. Синтез автоматов (1/3) -> 2.2.5. Синтез автоматов (1/5)
% ... (2/3) -> ... (2/5)
% ... (3/3) -> ... (3/5)
\begin{frame}
\frametitle{2.2.5. Синтез автоматов (4/5)}

На практике для минимизации ДКА применяется алгоритм Хопкрофта 
со сложностью $\mathcal{O}(n \log n)$, где $n = |Q|$.

\medskip
\NB{Идея:}
Вместо полной проверки всех состояний на каждом шаге, 
алгоритм Хопкрофта использует очередь активных разделителей $W$:

\begin{enumerate}
    \item Все состояния делятся на финальные ($F$) и нефинальные ($Q \setminus F$).
    \item Выбирается блок-разделитель $C$ из очереди и символ $a$. 
    \item Каждый существующий блок разбивается на две части в 
    зависимости от того, попадает ли переход по символу $a$ в блок $C$.
\end{enumerate}

\medskip
\begin{remark}
В очередь добавляется только \emph{меньшая} часть разбитого блока. 
Это гарантирует, что каждое состояние участвует в роли разделителя  
не более $\log_2 n$ раз.
\end{remark}
\end{frame}

% Текст слайда подготовил Тишковец Сергей 5030102/20202 
% Дата 15.05.2026
\begin{frame}
\frametitle{2.2.5. Синтез автоматов (5/5)}

\begin{algorithmic}
\REQUIRE ДКА $S=\langle A, Q, \delta, q_0, F\rangle$
\ENSURE Минимальный эквивалентный ДКА
\STATE $\mathcal{P} \Assign \{F,\;Q\setminus F\}$ \COMMENT{начальное разбиение}
\STATE $W \Assign \{F,\;Q\setminus F\}$ \COMMENT{очередь разделителей}
\WHILE{$W \neq \emptyset$}
    \STATE Извлечь блок $C$ из $W$
    \FORALL{$a \in A$}
        \STATE $X \Assign \{q \in Q \mid \delta(q, a) \in C\}$ \COMMENT{предшественники $C$}
        \FORALL{$D \in \mathcal{P} : (D \cap X \neq \emptyset) \wedge (D \setminus X \neq \emptyset)$}
            \STATE Разбить $D$ на $D_1 = D \cap X$ и $D_2 = D \setminus X$
            \STATE Заменить $D$ в $\mathcal{P}$ на $D_1$ и $D_2$
            \STATE Добавить в $W$ меньший из блоков $D_1, D_2$ 
        \ENDFOR
    \ENDFOR
\ENDWHILE
\STATE Построить ДКА по финальному разбиению $\mathcal{P}$
\end{algorithmic}

\end{frame}

\subsubsection{2.2.6. Анализ автоматов}
\label{2.2.6.}
\begin{frame}
\frametitle{2.2.6. Анализ автоматов (1/3)}
Анализ автоматов~--- это
построение языка для заданного автомата. Пусть задан автомат
$S=\langle A, Q, \delta, q, F\rangle$. Определим индуктивно язык
$E^k_{ij}$.
\begin{enumerate}
  \item $E^0_{ij}\Def \{a\in A~|~ \delta(q_i, a)= q_j\}$.
  \item $E^k_{ij}\Def E^{k-1}_{ij}\cup E^{k-1}_{ik} (E^{k-1}_{kk})^* E^{k-1}_{kj}$.
\end{enumerate}

Обозначим через $M_{ij}^k$ множество всех входных слов $\alpha\in A^*$,
ведущих из
$q_i$ в $q_j$ \\($q_j=\delta^{*}(q_i,\alpha)$)
и проходящих только через вершины $1..k$.

\begin{lemma}
$E^k_{ij}= M^k_{ij}$.
\end{lemma}
\begin{proof}
Индукция по $k$. База: $k=0$, $E^0_{ij}=\{a\in A~|~ \delta(q_i, a)=
q_j\}= M^0_{ij}$~--- множество дуг, не проходящих через
промежуточные вершины. Индукционный переход: пусть $E^{k-1}_{ij}= M^{k-1}_{ij}$.
Рассмотрим $M^k_{ij}$. Это множество можно разбить на две части:\\
 $M_{ij}^k = \bar{M}_{ij}^k
\cup\tilde{M}_{ij}^k$, где $\bar{M}_{ij}^k\Assign \Set{\langle q_i,
q_j\rangle\in M_{ij}^k}{q_k \not\in \langle q_i, q_j\rangle}$,\\
а $\tilde{M}_{ij}^k\Assign \Set{\langle q_i, q_j\rangle \in M_{ij}^k}
{q_k\in \langle q_i, q_j\rangle}$. Имеем $\bar{M}_{ij}^k =
M_{ij}^{k-1} = E_{ij}^{k-1}$. Пусть $\alpha\in \tilde{M}_{ij}^k$ и
проходит через $q_k$ $s$ раз (то~есть образует $s-1$ петлю через
узлы с номерами $1..k-1$). Тогда $\alpha= \alpha_0
\alpha_1\dots \alpha_{s-1}\alpha_s$, где $\alpha_0 \in
M^{k-1}_{ik}$, $\alpha_s \in M_{kj}^{k-1}$, $\alpha_1, \dots,
\alpha_{s-1}\in M^{k-1}_{kk}$.  
\end{proof}
\end{frame}

\begin{frame}
\frametitle{2.2.6. Анализ автоматов (2/3)}
\begin{proof} (Окончание)
Таким образом, \\$\alpha\in
M^{k-1}_{ik}(M^{k-1}_{kk})^* M_{kj}^{k-1} =
E_{ik}^{k-1}(E^{k-1}_{kk})^* E^{k-1}_{kj} = E^k_{ij}$. \\Пусть теперь
$\alpha \in E^{k-1}_{ik}(E^{k-1}_{kk})^* E_{kj}^{k-1}$. Тогда
$\alpha\in M^{k-1}_{ik}(M^{k-1}_{kk})^* M^{k-1}_{kj}$,  и значит,
$\alpha\in \tilde{M}^k_{ij}$. 
Наконец, $\tilde{M}_{ij}^k =
E^{k-1}_{ik}(E^{k-1}_{kk})^* E^{k-1}_{kj}$ и, следовательно,\\
$M_{ij}^k = \bar{M}^k_{ij}\cup \tilde{M}^k_{ij} = E^{k-1}_{ij}\cup
E^{k-1}_{ij}(E^{k-1}_{kk})^* E_{kj}^{k-1}$.
\end{proof}
\begin{theorem}
\NB{[Клини об анализе автоматов]}
Всякий язык, представимый конечным автоматом, регулярен.
\end{theorem}

\begin{proof}
Пусть задан автомат-распознаватель $S=\langle A, Q, \delta, q,
F\rangle$,  $|Q|=n$. Язык $L$, определяемый автоматом $S$~--- это
множество всех путей из начального состояния в заключительное. В
обозначениях предыдущей леммы
$$
L=\bigcup_{z\in F}M^n_{qz},
$$
но $M_{qz}^n= E^n_{qz}$ и потому регулярно, а значит, $L$
регулярен.
\end{proof}
\end{frame}

% Текст слайда подготовил Черняховский Лев
% Дата 09.01.2026
\begin{frame}
\frametitle{2.2.6. Анализ автоматов (3/3)}

\begin{example}
Пусть задан автомат-распознаватель $S=\langle A,Q,\delta,q,F\rangle$, \\ где
$A=\{a,b\}$, $Q=\{q_1,q_2\}$, $q=q_1$, $F=\{q_2\}$ (см. рисунок),
\begin{columns}[T,onlytextwidth]
  \begin{column}{0.55\textwidth}
  а функция переходов $\delta$ задана таблицей
  
  \begin{center}
  $\displaystyle
    \begin{array}{c|cc}
      \delta & a & b\\\hline
      q_1 & q_1 & q_2\\
      q_2 & q_2 & q_2
    \end{array}
    $
  \end{center}
  \end{column}
  \begin{column}{0.45\textwidth}
 \includegraphics[width=50mm]{2_2_6_3.png}
  \end{column}
\end{columns}
\small
По доказательству предыдущего слайда имеем $\displaystyle L=\bigcup_{z\in F}E^{|Q|}_{qz}=E^{2}_{12}$.

\[
\begin{aligned}
k=0:\quad
&E^{0}_{11}=\{a\},\ E^{0}_{12}=\{b\},\ E^{0}_{21}=\varnothing,\ E^{0}_{22}=\{a,b\}.\\
k=1:\quad
&E^{1}_{12}
=E^{0}_{12}\cup E^{0}_{11}(E^{0}_{11})^{*}E^{0}_{12}
=\{b\}\cup \{a\}\{a\}^{*}\{b\}
=\{a\}^{*}\{b\}.\\
&\text{При этом }E^{1}_{22}=\{a,b\}.\\
k=2:\quad
&E^{2}_{12}
=E^{1}_{12}\cup E^{1}_{12}(E^{1}_{22})^{*}E^{1}_{22}
=\{a\}^{*}\{b\}(\{a,b\})^{*}
=\{a\}^{*}\{b\}A^{*}.
\end{aligned}
\]

Следовательно, $L=\{a\}^{*}\{b\}A^{*}$ --- язык всех слов над $A=\{a,b\}$,
содержащих хотя бы один символ $b$.
\end{example}
\end{frame}

\subsubsection{2.2.7. Множество языков, распознаваемых автоматами}
\label{2.2.7.}
% Текст слайда подготовил Тишковец Сергей 5030102/20202 
% Дата 17.01.2026
\begin{frame}
\frametitle{2.2.7. Множество языков, распознаваемых автоматами (1/2)}

Конечные автоматы-распознаватели,
регулярные выражения,
регулярные порождающие грамматики и
источники (графы языков) оказываются равнообъёмными
понятиями~--- все они позволяют описать
одно и то же множество языков
(\odmkey{регулярных}{Язык!регулярный} или
\odmkey{автоматных}{Язык!автоматный} языков).

Наличие нескольких альтернативных способов
описания не является избыточным: на практике
для описания регулярных языков
применяют разные способы в зависимости от целей описания.
\begin{itemize}
\item
Конечные автоматы удобнее
всего применять, если нужно запрограммировать
распознаватель языка.
\item
Регулярные выражения являются самым компактным описанием.
\item
Порождающая грамматика
часто применяется для
формального описания в стандартизующих документах.
\item
Граф языка в виде диаграммы является, видимо,
наиболее наглядным.
\end{itemize}

\end{frame}

% Текст слайда подготовил Тишковец Сергей 5030102/20202 
% Дата 17.01.2026
\begin{frame}
\frametitle{2.2.7. Множество языков, распознаваемых автоматами (2/2)}
На следующем рисунке просуммированы связи между этими понятиями.

\begin{figure}[h]
\begin{center}
\includegraphics[width=130mm]{2_2_7.jpg}
\end{center}
\end{figure}

\end{frame}


\begin{frame}
\frametitle{Вычисления, совмещённые с разбором}
% GA_2_2025.tex, 2287–2288
При лексическом анализе можно и нужно вычислять значения 
некоторых лексем прямо во время анализа. Например, это нетрудно сделать для чисел, несколько сложнее для дат. При этом применяются различные трюки.
\end{frame}

% GA_2_2025.tex, 2307–2838
\subsection{2.3. Эквивалентность автоматов}\label{2.3.}
\begin{frame}
\frametitle{2.3. Эквивалентность автоматов}

Конечный автомат реализует определённое поведение, реагируя на
цепочки входных символов. Можно поставить следующий важный вопрос:
в каких случаях разные автоматы реализуют одно и то же поведение?
Другими словами, в каких случаях автоматы эквивалентны? Конечный
автомат является, в сущности, одним из способов описания алгоритма. В общем
случае (для произвольных алгоритмов) задача определения
эквивалентности алгоритмически неразрешима. Однако для конечных
автоматов это не так. В данном разделе устанавливаются
необходимое и достаточное условия эквивалентности конечных
автоматов и рассматриваются некоторые важные следствия.
\end{frame}


\subsubsection{2.3.1. Неотличимость состояний и автоматов}\label{2.3.1.}
\begin{frame}
\frametitle{\Ok 2.3.1. Неотличимость состояний и автоматов (1/2)}

Пусть даны два инициальных автомата над общими алфавитами входных
и выходных символов, $S=\langle A, Q_S, B, \delta_S, \lambda_S,
q\rangle$ и $T=\langle A, Q_{T}, B, \delta_{T}, \lambda_{T},
r\rangle$. Состояния $q\in Q_S$ и $r\in Q_{T}$ называются
\odmkey{неотличимыми}{состояние!неотличимое} (обозначение $q\equiv
r$), если $\A{\alpha\in A^*}{S(\alpha)=T(\alpha)}$, то~есть если
совпадают автоматные преобразования (\DMref{п.~2.1.6}{2.1.6.}), определяемые инициальными
автоматами.

\begin{example}
Состояния автомата, диаграмма которого представлена на
рисунке, неотличимы. $A=\{0,1\}$, $B=\{0,1\}$, $Q=\{1,2\}$.
\begin{figure}[h]
\includegraphics[width=80mm]{2_3_1_1.jpg}

\end{figure}
\end{example}

Неотличимость состояний, очевидно, является эквивалентностью.

\begin{remark}
Неотличимые состояния обычно называют
\odmkey{эквивалентными}{состояние!эквивалентное}. Термин
<<неотличимость состояний>> введен, чтобы избежать таких
словосочетаний, как <<отношение эквивалентности состояний является
отношением эквивалентности>>.
\end{remark}

\end{frame}

\begin{frame}
\frametitle{\Ok 2.3.1. Неотличимость состояний и автоматов (2/2)}

Инициальные конечные автоматы называются
\odmkey{эквивалентными}{автомат!эквивалентный},  если их
входные алфавиты совпадают, а начальные состояния неотличимы. Обозначение:
$S\equiv T$.

Конечные автоматы $S=\langle A, Q_S, B, \delta_S,
\lambda_S\rangle$ и $T=\langle A, Q_{T}, B, \delta_{T},
\lambda_{T}\rangle$ называются эквивалентными, если
$$
\A{q\in Q_S}{\E{r \in Q_{T}}{\langle A, Q_S, B, \delta_S,
\lambda_S, q\rangle \equiv \langle A, Q_{T}, B, \delta_{T},
\lambda_{T}, r\rangle }},
$$
$$
\A{r\in Q_{T}}{\E{q \in Q_S}{ \langle A, Q_{T}, B, \delta_{T},
\lambda_{T}, r\rangle \equiv \langle A, Q_S, B, \delta_S, \lambda_S,
q\rangle}}.
$$

Изоморфные автоматы с общими алфавитами
входных и выходных символов, очевидно, эквивалентны.
В то же время эквивалентные автоматы не
обязательно изоморфны.

%\vspace{-10pt}
%\begin{columns}[t]
%\begin{column}{6cm}
\begin{example}
На рисунке представлены
диаграммы эквивалентных, но неизоморфных автоматов.
\end{example}
%\end{column}
%\begin{column}{9cm}
%\vspace{-10pt}
\begin{figure}[h]
\includegraphics[width=140mm]{{2_3_1_2.jpg}}
\end{figure}
%\end{column}
%\end{columns}
\end{frame}


\subsubsection{2.3.2. Прямое произведение автоматов}\label{2.3.2.}
\begin{frame}
\frametitle{\Ok  2.3.2. Прямое произведение автоматов (1/3)}

Необходимое и достаточное условие эквивалентности инициальных
конечных автоматов устанавливается путём введения следующей
простой операции.

\odmkey{Прямым произведением}{Произведение!автоматов!прямое}
инициальных конечных автоматов $S=\langle A, Q_S, B_S, \delta_S,
\lambda_S, q\rangle$ и $T=\langle A, Q_{T}, B_{T}, \delta_{T},
\lambda_{T}, r\rangle$ (обозначается $S\times T$) называется
инициальный автомат \\$S\times {T} = \langle A,
Q_{S\times{T}}=Q_S\times{Q}_{T}, B_{S\times {T}}= B_S\times B_{T},
\delta_{S\times{T}}, \lambda_{S\times{T}}, (q,r)\rangle$, где
$$
\A{q'\in Q_S}{\A{r'\in Q_{T}}{\A{a\in A}{
\delta_{S\times{T}}\left((q',r'),a\right)=\left(\delta_S(q',a),
\delta_{T}(r',a)\right)}}},
$$

\vspace{-12pt}
$$
\A{q'\in Q_S}{\A{r'\in Q_{T}}{\A{a\in A}{
\lambda_{S\times{T}}\left((q',r'),a\right)=\left(\lambda_S(q',a),
\lambda_{T}(r',a)\right)}}}.
$$
Другими словами, прямое произведение автоматов~--- это просто два
невзаимодействующих автомата, синхронно работающих с общим
входом. Состояниями прямого произведения являются пары состояний
сомножителей, а выходными символами~---
пары выходных символов сомножителей.

\begin{theorem} {\NB{[Мура]}}
Два инициальных конечных автомата эквивалентны тогда и только
тогда, когда для любого достижимого состояния в их прямом
произведении функции выходов совпадают.
\end{theorem}
\end{frame}


\begin{frame}
\frametitle{\Ok 2.3.2. Прямое произведение автоматов (2/3)}
\begin{remark} Если
$S=\langle A, Q_S, B_S,
\delta_S, \lambda_S, q\rangle$ и $T=\langle A, Q_{T}, B_{T},
\delta_{T}, \lambda_{T}, r\rangle$,
то
$S\equiv T \Equiv$ \\$\Equiv\! \A{(q', r')\!\in\! Q_{S\!\times\! T}}{\E{\alpha\!\in\! A^*}{
\delta^*_{S\!\times\! T}\left((q,r),\alpha\right)\! =\! (q',r')\! \Impl\!\A{a\!\in\! A}{\lambda_S(q', a)\!= \!\lambda_{T}(r',a)}}}.
$

\end{remark}
\begin{proof}
\proofsection{Необходимость.}
От противного. Пусть
$$
\neg\A{(q', r')\!\in\! Q_{S\!\times\! T}}{\E{\alpha\!\in\! A^*}{
\delta^*_{S\!\times\! T}\left((q,r),\alpha\right)\! =\! (q',r')\! \Impl\!\A{a\!\in\! A}{\lambda_S(q', a)\!= \!\lambda_{T}(r',a)}}}.
$$
Тогда
$$
\E{(q', r')\!\in\! Q_{S\!\times\! T}}{\Not\E{\alpha\!\in\! A^*}{
\delta^*_{S\!\times\! T}\left((q,r),\alpha\right)\! =\! (q',r')\! \Impl\!\A{a\!\in\! A}{\lambda_S(q', a)\!= \!\lambda_{T}(r',a)}}},
$$
что можно преобразовать в
$$
\E{(q', r')\!\in\! Q_{S\!\times\! T}}{\A{\alpha\!\in\! A^*}{
\delta^*_{S\!\times\! T}\left((q,r),\alpha\right)\! =\! (q',r')\! \And\!\E{a\!\in\! A}{\lambda_S(q', a)\!\neq \!\lambda_{T}(r',a)}}}.
$$
Положим $\beta\Assign\alpha a$. Имеем
\begin{align*}
S(\beta)&=S(\alpha a)= S(\alpha)+\lambda_S^*(\delta_S^*(q,\alpha),a)= 
T(\alpha)+\lambda_S^*(\delta_S^*(q,\alpha),a) =T(\alpha)+\lambda_S^*(q',a)\neq \\
&\neq T(\alpha)+\lambda^*_{T}(r',a)= 
T(\alpha)+\lambda_{T}^*(\delta^*_{T}(r,\alpha), a) = T(\alpha a)=T(\beta).
\end{align*}
Получили противоречие: $\E {\beta\in A^*}{S(\beta)\neq T(\beta)}$.
\end{proof}
\end{frame}


\begin{frame}
\frametitle{\Ok 2.3.2. Прямое произведение автоматов (3/3)}
\begin{proof} (Продолжение)
\proofsection{Достаточность.} Индукция по $|\alpha|$, $\alpha \in
A^*$. База: $|\alpha|=1$, $\alpha = a \in A$. Имеем $(q,r)$
достижимо по определению, следовательно, $\lambda_S(q,a)=
\lambda_{T}(r,a)$ по условию, и значит, $S(a) = T(a)$. Индукционный
переход: пусть $\A{\alpha\in A^k}{\lambda^*_S(q,\alpha)=
\lambda^*_{T}(r,\alpha)}$, 
рассмотрим $\alpha\in A^{k+1}$,
$|\alpha|=k+1$. Для любой пары $(q', r')\in 
Q_{S \times T}$ имеем
$$
\E{\alpha\in A^*}
{\delta^*_{S\times T} \left((q,r),\alpha\right)=(q',r')\Impl
\A{ a\in A }{\lambda_S(q',a)=\lambda_{T}(r',a)}},
$$
следовательно,
$$
\A{a\in A}{\lambda_S\left(\delta^*_S (q,\alpha), a\right) =
\lambda_{T}\left(\delta^*_{T}(r,\alpha), a\right)},
$$
откуда
$$
\A{a\in A}{\A{\alpha\in A^k }{\lambda^*_S (q,\alpha a) =
\lambda^*_{T}(r,\alpha a)}},
$$
и значит,
$
\A{\alpha\in A^{k+1}}{\lambda^*_S(q,\alpha) = \lambda^*_T(r,\alpha)}.$
\end{proof}

\end{frame}


\subsubsection{2.3.3. Минимальный конечный автомат}\label{2.3.3.}
\begin{frame}
\frametitle{\Ok 2.3.3. Минимальный конечный автомат (1/2)}
Автоматы могут функционировать одинаково, даже если у них разное
число состояний.

\begin{example}
Автоматы на рисунке эквивалентны независимо от того,
какое состояние взять в качестве начального для левого автомата. Оба автомата инвертируют входную двоичную последовательность.

\vspace{-6pt}
\begin{figure}[h]
\includegraphics[width=120mm]{{2_3_1_2.jpg}}
\end{figure}
\end{example}

%\vspace{-6pt}
Таким образом, возникает задача
\odmkey{минимизации}{Минимизация!автоматов} конечных автоматов:
для данного автомата найти эквивалентный автомат, имеющий
минимальное число состояний. Такой автомат называется
\odmkey{минимальным}{Автомат!минимальный}.

\begin{theorem}
Для любого автомата $S$ существует минимальный эквивалентный
автомат $T$, причем количество состояний в минимальном автомате
равно мощности фактормножества состояний по отношению
неотличимости,
$$
\A{S}{\E{T}{S\equiv T \And |Q_{T}| = \min_{U\equiv S} |Q_U|}} \And
|Q_{T}| = |Q_S/\equiv|.
$$
\end{theorem}
\end{frame}


\begin{frame}
\frametitle{\Ok 2.3.3. Минимальный конечный автомат (2/2)}
\begin{proof}
Пусть $\mathcal{C}=\{C_1,\dots, C_t\}$~--- разбиение множества
состояний $Q_S$ на классы эквивалентности по отношению
неотличимости состояний $\equiv$:
$$
\bigcup_{i=1}^t C_i = Q_S,\quad
\A{i,j \in 1..t}{ i\neq j \Impl C_i\cap C_j = \emptyset}, \quad
\mathcal{C}=Q_S/\equiv.
$$
Заметим, что
$q_1\in C_i \And q_2 \in C_i \Impl \A{a \in A}{\delta(q_1, a)\in
C_j \And \delta(q_2, a)\in C_j}$. Действитель\-но, от противного,
если $\delta(q_1,a)\not\equiv\delta(q_2,a)$ (обозначим
$q_1':=\delta(q_1,a)$, $q_2':= \delta(q_2,a)$), то
$\E{\alpha\in A^*}{S_{q_1}(\alpha)\neq S_{q_2}(\alpha)}$, но
тогда $S_{q_1'}(a\alpha)\neq S_{q_2'}(a\alpha)$, то есть $q_1
\not\equiv q_2$, что противоречит тому, что
$q_1 \in C_i \And q_2\in C_i$.

Определим автомат $T =\langle A, Q_{T}, B, \delta_{T},
\lambda_{T}\rangle$ следующим образом: $Q_{T}:= \{C_1, \dots,
C_t\}$, $\delta_{T}(C_i, a): = C_j$, такое что $\delta_S(q,a)\in
C_j$, а $\lambda_{T}(C_i, a)\Assign \lambda_S(q,a)$, где $q$~---
любое состояние из класса $C_i, q\in C_i$. Заметим, что $T\equiv
S$ по построению и $T$ не имеет неотличимых состояний. Покажем от
противного, что $T$~--- минимальный.

Пусть $\E{T'}{T'\equiv S \And |Q_{T'}|< |Q_{T}|}$. По
транзитивности $T'\equiv T$. Но $|Q_{T'}|< |Q_{T}|$, следовательно,
$\E{t'\in Q_{T'}}{\E{t_1, t_2\in Q_{T}}{t_1\equiv t'\And t_2\equiv
t'}}$, откуда $t_1\equiv t_2$, что противоречит построению.
\end{proof}

\end{frame}


\subsubsection{2.3.4. Алгоритм минимизации Мили}\label{2.3.4.}
\begin{frame}
\frametitle{\Ok 2.3.4. Алгоритм минимизации Мили (1/5)}
Приведённый ниже алгоритм Мили строит минимальный эквивалентный автомат
путем последовательного разбиения множества состояний исходного
автомата на блоки, пока эти блоки не станут классами эквивалентности
по неотличимости.

\begin{algorithmic}
\REQUIRE {Автомат $S=\langle A, Q, B, \delta, \lambda \rangle$};
\ENSURE {Разбиение $\mathcal{G}$ множества $Q$}.
\STATE{\textbf{$\mathcal{C}_0 := Q/\equiv_\lambda$}}
\COMMENT{ начальное разбиение }
\FOR{\textbf{$i$ from $1$ to $|Q|$}}
 \STATE $\mathcal{C}_i\Assign\emptyset$
 \COMMENT{ инициализация очередного разбиения }
 \FOR{\textbf{$K\in \mathcal{C}_{i-1}$}}
  \STATE $\mathcal{C}_i\Assign\mathcal{C}_i 
  \cup K/ \equiv_\delta^i$
  \COMMENT{ измельчение блока разбиения }
 \ENDFOR
 \IF{$\mathcal{C}_i = \mathcal{C}_{i-1}$}
   \STATE \textbf{ exit }\COMMENT{ прекращение процесса }
 \ENDIF
\ENDFOR
\end{algorithmic}
\end{frame}

\begin{frame}
\frametitle{\Ok 2.3.4. Алгоритм минимизации Мили (2/5)}

В  этом алгоритме операция вычисления фактормножества
считается заданной, а отношения $\equiv_\lambda$ и $\equiv_\delta^i$
определяются следующим образом:
\begin{align*}
q'\equiv_\lambda q''& \Def \A{a\in A}{\lambda(q',a)=\lambda(q'', a)},\\
q'\equiv_\delta^i q''&\Def 
\A{a \in A}{\E{ K\in C_{i-1}}{
{\delta(q', a)\in K \And  \delta(q'',a)\in K}}},
\quad q',q'' \in C_i.
\end{align*}

\begin{remark}
При реализации алгоритма Мили хранить все разбиения не
обязательно~--- достаточно двух последних.
Явная индексация разбиений введена в алгоритме
только для удобства обоснования.
\end{remark}

\begin{ground}
Пусть алгоритм остановился на шаге с номером $k\leq|Q|$. Покажем,
что полученное разбиение $C_k$ множества $Q$ является
фактормножеством по неотличимости, то~есть
$$
q'\in K \And  q''
\in K \in C_k \Impl q'\equiv q''
$$
и
$$
q'\in K'\in C_k \And  q''
\in K'' \in C_k \And  K'\neq K'' \Impl q'\not\equiv q''.
$$

\vspace{-\baselineskip}
\end{ground}
\end{frame}

\begin{frame}
\frametitle{\Ok 2.3.4. Алгоритм минимизации Мили (3/5)}
\begin{ground} (Продолжение)
Действительно, пусть $q', q'' \in K\in C_k$.
 По условию остановки алгоритма $q'\equiv_\delta^k q''$, то~есть $\A{a\in A}
 {\E{K'\in C_{k-1}}{\delta(q', a)\in K' \And \delta(q'', a)\in K'}}$
 (иначе блок $K$ можно было бы
 измельчить на меньшие блоки).
 Таким образом, \\$\delta(q', a) \equiv_\delta^{k-1} \delta(q'', a)$,
 аналогично, $\A{b\in A}{\delta(\delta(q',a),b) \equiv_\delta^{k-2}
 \delta(\delta(q'',a),b)}$ и т.~д.
 Следо\-ва\-тельно,
$
\E{ K''}{ \A{ \alpha\in A^*}{\delta^*(q',\alpha)\in K''\And  \delta^*(q'',
 \alpha) \in K''}}.
 $
 Но если $q'\equiv_\delta^i q''$, \\то $q'\equiv_\lambda q''$
 (потому что последующие разбиения получаются измельчением блоков
 исходного разбиения), и значит, $\lambda(q',a)= \lambda(q'',a)$,
 откуда
 $S_{q'}(\alpha)= S_{q''}(\alpha)$,
 то~есть $q'\equiv q''$.

 Пусть теперь $q'\in K'\in C_k$, $q''\in K'' \in C_k$ и $K'\neq
 K''$.
Тогда
$$
\E{a'\in A}{\delta(q',a')\in K_{k-1}' \And  \delta
(q'',a')\in K_{k-1}''},
$$
где $K_{k-1}'\neq K_{k-1}''$, $K_{k-1}'$,
$K_{k-1}'' \in C_{k-1}$. Другими словами, можно найти такое $a'$,
что $\delta(q', a')$ и $\delta(q'', a')$ обязательно принадлежат
разным классам разбиения на предыдущем уровне (иначе $q'$ и
$q''$ попали бы в один класс на данном уровне).
 Но тогда\\
$
\E{a''\in A}{ \delta\left(\delta(q', a'), a''\right)\in K_{k-2}'\And
              \delta\left(\delta(q'',a'),a''\right)\in K_{k-2}''}.
$
\end{ground}
\end{frame}

\begin{frame}
\frametitle{\Ok 2.3.4. Алгоритм минимизации Мили (4/5)}
\begin{ground} (Окончание)
Заметим, что $ K_{k-2}'\neq  K_{k-2}''$, $K_{k-2}', K_{k-2}''\in C_{k-2}$
и так~далее,
то есть можно
построить слово $\alpha$, $|\alpha|< k$, такое что
$\delta^*(q',\alpha)\not\equiv_\lambda \delta^*(q'',\alpha)$.
Это означает, что $\delta^*(q', \alpha)$ и
 $\delta^*(q'',\alpha)$ принадлежат разным классам начального
 разбиения $Q/\equiv_\lambda$,  откуда $\lambda^*(q',
 \alpha)\neq\lambda^*(q'',\alpha)$, и значит, 
 $q'\not\equiv q''$.
\end{ground}

\medskip
\begin{example}
Применение алгоритма Мили к автомату, диаграмма которого приведена
на следующем рисунке слева, дает автомат, диаграмма которого
приведена на том же рисунке справа. При этом последовательность
разбиений следующая:

\begin{enumerate}
\item $Q=\{1,2,3,4,5\}$~--- множество состояний;

\item $C_0 = \{1,4,5\}$, $\{2\}$, $\{3\}$~--- начальное разбиение по $\lambda$;

\item $C_1 = \{1\}, \{4, 5\}$, $\{2\}$, 
$\{3\}$~--- измельчение блока $\{1,4,5\}$;

\item $C_2=C_1$~--- остановка алгоритма.
\end{enumerate}
\end{example}
\end{frame}

\begin{frame}
\frametitle{\Ok 2.3.4. Алгоритм минимизации Мили (5/5)}


\begin{example} (Продолжение -- рисунок)
\begin{figure}[h]
\begin{center}
\includegraphics[width=120mm]{2_2_4_1.jpg}
\end{center}
\end{figure}
\end{example}

\vspace{-10pt}
Автомат на рисунке выполняет следующее преобразование:
если в последовательности битов встречаются два одинаковых подряд,
то второй из них инвертируется. Остальные биты не изменяются.

\begin{digress}
Минимизация автоматов, построенных вручную,
мало что даёт, поскольку эти автоматы 
небольшие и, как правило, близки к оптимальным.
Другое дело,
когда речь идёт об автоматах,
сгенерированных
автоматически по формальным правилам.
В таких автоматах бывают тысячи и десятки тысяч
состояний, и минимизация даёт значительный эффект. 
\end{digress}

\end{frame}


\subsubsection{2.3.5. Эквивалентность автоматов Мура и Мили}
\begin{frame}
\frametitle{\Ok 2.3.5. Эквивалентность автоматов Мура и Мили (1/2)}
Очевидно, что автомат Мура является частным случаем автомата Мили.
Действительно, достаточно перенести отметку вершины на входящие
ребра, и из автомата Мура получится автомат Мили (с тем же числом
состояний).

Следующая теорема устанавливает тот факт, что выразительные
возможности двух классов автоматов совпадают.

\begin{theorem}
Для любого автомата Мили существует эквивалентный автомат Мура.
\end{theorem}

\begin{proof}
Пусть $S=\langle A, Q, B, \delta, \lambda\rangle$~--- автомат
Мили. Определим автомат Мура
$$
M= \langle A_M, Q_M, B_M,
\delta_M,\mu\rangle
$$
следующим образом:
$A_M\Assign A$;
$B_M\Assign B$;
$Q_M\Assign Q\times A \cup Q =
\left\{q^M_{ij}\right\}^{k,m}_{i=1, j=1} \cup
\left\{q^M_i\right\}^k_{i=1}$, где $\delta_M(q_i^M, a_j)\Assign
q^M_{ij}$, а $\delta_M(q^M_{ij}, a_s)\Assign q_{ts}$ такое, что $q_t
= \delta(q_i, a_s)$;
$\mu(q^M_i)\!\Assign\! \varepsilon$,
$\mu(q^M_{ij})\!\Assign\!\lambda(q_i, a_j)$.

Покажем индукцией по $|\alpha|$, что $S(q_i, \alpha)=
M(q_i^M,\alpha)$. База: $|\alpha|=1\Impl \alpha =a$. По построению
$\A{a}{\A{i}{ S(q_i, a)= M(q_i^M, a)}}$. Индукционный переход: пусть \\$S(q_i, \alpha)=
M(q_i^M,\alpha)$. Рассмотрим $\alpha a$. Имеем $\delta^*(q_i,\alpha)
= q_j \Impl \delta_M^*(q_i^M,\alpha)= q^M_{ij}$,\\
$\lambda(\delta^*(a_i,\alpha), a) = \mu(\delta_M^*(q_i^M, \alpha))$, и
значит, $S(q_i, \alpha a)= M(q_i^M, \alpha a)$.
\end{proof}
\end{frame}


\begin{frame}
\frametitle{\Ok 2.3.5. Эквивалентность автоматов Мура и Мили (2/2)}
\begin{example}
На рисунке слева приведен автомат Мили с двумя
состояниями, а справа~--- эквивалентный автомат Мура с шестью
состояниями. Эти автоматы инвертируют биты через один.

\begin{figure}[h]
\begin{center}
\includegraphics[width=120mm]{2_2_5.jpg}
\end{center}
%\caption{Автомат Мили (слева) и эквивалентный автомат Мура
%(справа)} \label{odm33-6}
\end{figure}
\end{example}

\begin{remark}
После построения автомата Мура его можно минимизировать. Но в
большинстве случаев минимальный автомат Мили содержит меньше
состояний, чем эквивалентный ему минимальный автомат Мура.
\end{remark}

\end{frame}


% GA_2_2025.tex, 3714–4026
\subsection{2.5. Автоматы с памятью}\label{2.5}
% 2.5. Автоматы с памятью. Г. А. Гусев, 24.10.2026.
% Основание (отправлено ФАН 04.10.2026, правки приняты):
%   АТОМ_слайд570_2.5_память_автомата_n_состояний;
%   АТОМ_слайд571_2.5_свойство_префикса_ключевые_слова;
%   АТОМ_слайд571_2.5_ссылка_на_теоремы_Клини.
% Что сделано:
%   1) кадр 1/3: автомат с n состояниями «помнит» не n событий, а различает
%      не более n историй; для k последних символов нужно |A|^k состояний,
%      то есть автомат хранит не более log_2 n бит (п. 2.4.3);
%   2) кадр 2/3, замечание: эффективность дерева решений не зависит от
%      свойства префикса (как и сказано в отступлении п. 2.2.2 (4/4)); снято
%      фактически неверное «именно так устроены ключевые слова» (do/double);
%   3) кадр 2/3, последний абзац: ссылка на п. 1.2.5 (расширения регулярных
%      выражений) заменена ссылками на обе теоремы Клини — п. 2.2.5 и 2.2.6 —
%      и п. 2.2.7; добавлено «и только их».
% Вторая редакция (по замечаниям к первой редакции): восстановлены выделения
% оригинала — «дерева решений», «свойство префикса» (определяемые понятия),
% T_3 записано готическим, как в оригинале; промежутки между абзацами — как
% в оригинале.
% Кадр 3/3 оставлен без изменений.
%%
% 16.11.26. Немного поправил и добавил шпаргалку. ФАН

\begin{frame}
\frametitle{2.5. Автоматы с памятью (1/4)}

В самом названии <<конечные автоматы>> слово <<конечные>> очень существенно, и
во многом характеризует свойства этого способа описания алгоритмов.

\smallskip
Во всех формальных определениях алгоритмов, наряду со способностью абстрактных
машин совершать некоторые действия, в той или иной форме присутствует
{\em память}, то есть способность хранить и использовать результаты ранее
совершенных действий. Заметим, что в универсальных моделях вычислимости память
потенциально бесконечна: машина Тьюринга имеет бесконечную ленту (или ленты),
регистры в МНР неограничены, слова в нормальных алгорифмах Маркова могут иметь
любую длину.

\smallskip
В рассматриваемых моделях конечных автоматов в явном виде памяти нет. В неявном
виде памятью автомата является множество состояний, но эта память ограничена:
автомат <<знает>>, в каком он пребывает состоянии, но не более того.
Например, 
чтобы различать, какая именно последовательность $k$ входных символов 
алфавита $A$ была прочитана,
автомату нужно не менее $|A|^k$ состояний;
в противном случае по принципу Дирихле  найдутся такие различные слова,
которые переводят автомат в одно и то же состояние.
Иначе говоря, автомат с $n$
состояниями хранит не более $\log_2 n$ бит информации
(см.~п.~\DMref{2.4.3}{2.4.3.}).
\end{frame}

\begin{frame}
\frametitle{2.5. Автоматы с памятью (2/4)}

В частном случае комбинационных автоматов (см.~п.~\DMref{2.4.1}{2.4.1.}),
автомат вовсе не имеет памяти и является просто функцией над конечным
алфавитом. Применительно к задаче распознавания языков, автомат без памяти может
распознавать элементарные языки (то есть одиночные символы,
см.~п.~\DMref{1.2.2}{1.2.2.}), и не более того.
\smallskip

Конечный язык, то есть множество слов в заданном алфавите, проще всего
распознаётся с помощью \odmkey{дерева решений}{Дерево!решений}, которое получится, если действовать, как
указано в п.~\DMref{2.2.2}{2.2.2.}, но не добавлять состояние недопустимости.

\begin{remark}
Дерево решений весьма эффективно, если алфавит не слишком велик и слова не очень
длинные: время распознавания пропорционально длине слова и не зависит от
количества слов (см.~п.~\DMref{2.2.2}{2.2.2.}). Так обычно распознаются
ключевые слова языков программирования. \odmkey{Свойство префикса}{Свойство!префикса} (никакое слово
языка не является собственным началом другого слова) упрощает дерево~---
заключительными оказываются только листья,~--- но наличие свойства префикса не обязательно для распознавания конечных языков.
\end{remark}

Конечные автоматы, то есть автоматы с конечной памятью состояний, могут
распознавать регулярные языки, и только их (см.~\DMref{п.~2.2.5}{2.2.5.},
\DMref{п.~2.2.6}{2.2.6.} и \DMref{п.~2.2.7}{2.2.7.}), класс которых совпадает
с классом языков, порождаемых грамматиками типа $\frak{T}_3$.
\end{frame}

\begin{frame}
\frametitle{2.5. Автоматы с памятью (3/4)}

Как показано в \DMref{п.~2.5.3}{2.5.3.}, если допустить использование
простейшей неадресуемой, но потенциально неограниченной памяти (стека), то конечные автоматы
способны эффективно распознавать контекстно-свободные языки, и это
обстоятельство является основой современных языков программирования.

\smallskip
Для того, чтобы распознавать контекстно-зависимые языки достаточно допустить
наличие адресуемой памяти, ограниченной длиной входного слова 
(\DMref{п.~2.5.2}{2.5.2.}). 

\smallskip
Наконец, если снабдить автомат неограниченной памятью, то получается машина
Тьюрин\-га, способная распознавать любой перечислимый язык (\DMref{п.~1.1.3}{1.1.3.}).

\smallskip
В этом разделе  в начале рассматриваются различные варианты добавления к конечному
автомату памяти в том или ином виде, и рассматривается вопрос о том, как
расширения модели конечных автоматов сказываются на их свойствах и возможностях
в части распознавания языков в порядке от общего к частному.
В конце раздела для сравнения рассматриватся сети Петри~--- модель описания 
поведения, также основанная на графах переходов состояний, но содержащая память 
{\em внутри} модели, а не снаружи, как в случае конечных автоматов.  
\end{frame}

\begin{frame}
\frametitle{2.5. Автоматы с памятью (4/4)}
\begin{tabular}{p{10mm}|p{42mm}|p{91mm}}
  № & Название параграфа
  & Ключевые термины и обозначения \\
  \hline
 \text{\DMref{2.5.1.}{2.5.1.}} &  Машины Тьюринга
  & {  
 }
   \\
\text{\DMref{2.5.2.}{2.5.2.}} & Линейно-ограниченные автоматы 
  & 
{ 
}
  \\
\text{\DMref{2.5.3.}{2.5.3.}}  & Автоматы с магазинной памятью &
{
}
\\
\text{\DMref{2.5.4.}{2.5.4.}}  & Сети Петри &
{
 } 
\\
\text{\DMref{2.5.5.}{2.5.5.}}  & Языки, распознаваемые сетями Петри &
{
 } 
\end{tabular}


\end{frame}

\subsubsection{2.5.1. Машины Тьюринга}\label{2.5.1.}
\begin{frame}
\frametitle{2.5.1. Машины Тьюринга}

Машина Тьюринга~--- это конечный автомат,
к которому добавлена адресуемая память,
в виде потенциально бесконечной ленты.
Автомат наделяется способностью 
читать с ленты  и писать на ленту один символ,
находящийся по текущему адресу, а также
уменьшать и увеличивать адрес на единицу, то есть сдвигаться по ленте влево и вправо. 
\begin{remark}
Можно рассматривать 
\odmkey{многоленточные}{Машина!Тьюринга!многоленточная}
машины Тьюринга,
в которых используется 
{\em несколько} (конечное число) лент.
Для любой многоленточной машины Тьюринга
можно построить эквивалентную одноленточную машину
ценой некоторого усложнения кодирования алгоритмов.
Зато использование многоленточной модели
иногда упрощает доказательства свойств алгоритмов. 
\end{remark}

\medskip
Говорят, что машина Тьюринга 
{\em принимает} некоторое входное слово, записанное на ленте,
если она останавливается в заключительном состоянии. 
Говорят, что машина Тьюринга {\em распознаёт} язык,
если она принимает все слова этого языка и только их.

\medskip
\begin{theorem}
Если язык $L$ порождается некоторой 
грамматикой $G$ типа 0, то существует машина 
Тьюринга $M$, которая распознаёт язык $L$.
\end{theorem}
\TODO{Эссе + недетерминированные машины Тьюринга.} 
\end{frame}


% СЛАЙД 1: НЕДЕТЕРМИНИРОВАННЫЕ МАШИНЫ ТЬЮРИНГА
% Текст слайда подготовил Швачко Никита 5030102/20202
% Дата 11.01.2026
\begin{frame}
\frametitle{2.5.1. Машины Тьюринга (2/5)}

\odmkey{Недетерминированная машина Тьюринга}{Машина!Тьюринга!недетерминированная} (НМТ)~--- 
это обобщение машины Тьюринга, в котором функция переходов 
возвращает \emph{множество} возможных конфигураций:
\[
\delta: Q \times \Gamma \to 2^{Q \times \Gamma \times \{L, R\}},
\]
где $Q$~--- множество состояний, $\Gamma$~--- алфавит ленты, $L, R$~--- сдвиг влево/вправо.

\begin{remark}
НМТ \emph{принимает} слово $\omega$, если \emph{существует} хотя бы одна 
последовательность переходов, приводящая к принимающему состоянию.
Вычисление НМТ образует \odmkey{дерево вычислений}{Дерево!вычислений}: 
каждая ветвь соответствует одному возможному пути выполнения.
\end{remark}

\begin{theorem}
Для любой недетерминированной машины Тьюринга можно построить 
эквивалентную детерминированную машину Тьюринга.
\end{theorem}

\begin{proof}
ДМТ моделирует НМТ \emph{обходом дерева вычислений в ширину}:
перебирает все возможные последовательности переходов длины $1, 2, 3, \ldots$,
пока не найдёт принимающую ветвь (если она существует).
\end{proof}

\end{frame}


% СЛАЙД 2: ПРИМЕР 1 — Машина Тьюринга для языка a^n b^n
% Текст слайда подготовил Швачко Никита 5030102/20202
% Дата 11.01.2026
\begin{frame}
\frametitle{2.5.1. Пример: МТ для языка $\{a^n b^n\}$ (3/5)}

\begin{example}
Построим ДМТ $M$, распознающую язык $L = \{a^n b^n \mid n \ge 0\}$.

\NB{Идея:} На каждом проходе помечаем по одному $a$ и одному $b$, 
проверяя, что они исчерпываются одновременно.

\NB{Алгоритм:}
\begin{enumerate}
\item Если лента пуста (только $\sqcup$)~--- \textbf{принять}.
\item Заменить первый непомеченный $a$ на $X$, сдвинуться вправо.
\item Пропустить все $a$ и $Y$, найти первый непомеченный $b$.
\item Заменить этот $b$ на $Y$, вернуться влево до первого $X$.
\item Сдвинуться вправо и повторить с шага 2.
\item Когда все $a$ помечены: проверить, что все $b$ тоже помечены.
\item Если остались непомеченные $a$ или $b$~--- \textbf{отвергнуть}.
\end{enumerate}

\NB{Сложность:} $O(n^2)$ шагов для слова длины $2n$.
\end{example}

\end{frame}


% СЛАЙД 3: ПРИМЕР 2 — Недетерминированная МТ
% Текст слайда подготовил Швачко Никита 5030102/20202
% Дата 11.01.2026
\begin{frame}
\frametitle{2.5.1. Пример: НМТ для проверки составности числа (4/5)}

\begin{example}
Построим НМТ для проверки, является ли число $n$ (в двоичной записи) \emph{составным}.

\NB{Идея:} Недетерминированно <<угадываем>> делитель $d$, затем 
детерминированно проверяем делимость.

\NB{Алгоритм:}
\begin{enumerate}
\item \textbf{Недетерминированный выбор:} записать на вспомогательную ленту 
      двоичное число $d$, где $1 < d < n$ (выбор каждого бита~--- недетерминированный).
\item \textbf{Детерминированная проверка:} вычислить $n \mod d$.
\item Если $n \mod d = 0$~--- \textbf{принять}, иначе~--- отвергнуть эту ветвь.
\end{enumerate}

Если $n$ составное, существует ветвь вычисления с делителем $d \mid n$, и НМТ принимает.

Если $n$ простое, все ветви отвергают, и НМТ не принимает.
\end{example}

\begin{remark}
Эквивалентная ДМТ перебирает все $d$ от $2$ до $n-1$ последовательно.
Поскольку число $n$ записано в двоичной системе (размер входа $\approx \log_2 n$),
количество делителей $n - 2$ экспоненциально по размеру входа.
\end{remark}

\end{frame}


% СЛАЙД 4: ПРИМЕР 3 — Конструкция моделирования НМТ
% Текст слайда подготовил Швачко Никита 5030102/20202
% Дата 11.01.2026
\begin{frame}
\frametitle{2.5.1. Моделирование НМТ детерминированной МТ (5/5)}

\begin{example}
Пусть НМТ $N$ имеет не более $r$ вариантов перехода из каждой конфигурации.
Построим эквивалентную ДМТ $M$ (используя 3 ленты).

\NB{Структура лент:}
\begin{itemize}
\item \textbf{Лента 1:} хранит входное слово $\omega$ (только для чтения).
\item \textbf{Лента 2:} рабочая лента для симуляции одной ветви $N$.
\item \textbf{Лента 3:} <<адрес>> текущей ветви~--- последовательность $(c_1, \ldots, c_\ell)$, $c_i \in \{1, \ldots, r\}$.
\end{itemize}

\NB{Алгоритм} (обход в ширину):
\begin{enumerate}
\item Для $\ell = 0, 1, 2, \ldots$:
\item \quad Для каждой последовательности $(c_1, \ldots, c_\ell)$ длины $\ell$:
\item \quad \quad Скопировать $\omega$ на ленту 2, симулировать $N$, выбирая $c_i$-й переход на шаге $i$.
\item \quad \quad Если $N$ приняла~--- \textbf{принять}.
\end{enumerate}
\end{example}

\bigskip
\TODO{К сожалению, не сделано главного = не доказана основная теорема, которая 
необходима для связывания грамматик и автоматов и используется во многих других местах.}

\end{frame}



% GA_2_2025.tex, 5346–5828
\subsection{2.6. Программирование автоматов}
\begin{frame}
\frametitle{2.6. Программирование автоматов}

Конечные автоматы являются одним из самых распространенных приёмов
описания алгоритмов и, следовательно, способ реализации конечного
автомата в программе практически очень важен. В данном курсе
рассматриваются только некоторые основные модели, связанные с понятием
конечного автомата. На практике используется гораздо больше
специфических конструкций, в основе которых лежит идея конечного
графа переходов состояний. Буквально для каждой автоматной модели предложены свои
(иногда многочисленные) способы представления автоматов в
программе. Обсудить их все не представляется возможным~--- мы
приводим только несколько характерных примеров.

\medskip
\TODO{Различные реализации автоматов. Доделать. }
\end{frame}


\subsubsection{2.6.1. Интерпретаторы таблиц}\label{2.6.1.}
\begin{frame}
\frametitle{2.6.1. Интерпретаторы таблиц (1/3)}


Для начала рассмотрим способы представления,
основанные на интерпретации хранимых данных.

Пусть задан инициальный конечный автомат-преобразователь $S=\langle A, Q, B, \delta,\lambda, q_1
\rangle$, см.~\DMref{п.~2.1.1}{2.1.1.}, \DMref{п.~2.1.2}{2.1.2.}, \DMref{п.~2.1.5}{2.1.5.}.

\medskip
Алфавиты, участвующие в описании конечного автомата, можно
представить массивами:

\textbf{$A:$ array $[1..m]$ of $\ldots$}~--- входной алфавит,

\textbf{$B:$ array $[1..k]$ of $\ldots$}~--- выходной алфавит,

\textbf{$Q:$ array $[1..n]$ of $\ldots$}~--- алфавит состояний.

\begin{remark}
Здесь намеренно не указан тип элементов массивов~--- он может быть
любым. В дальнейших построениях используются только индексы в этих
массивах, то~есть целые числа. 
\end{remark}

\medskip
Функции переходов и выходов
также можно представить массивами:

$  \delta: {\textbf { array }} [1.. n, 1.. m] \textbf{ of } 1..n $,

$ \lambda: {\textbf { array }} [1.. n, 1.. m] \textbf{ of } 1..k$.
\end{frame}

\begin{frame}
\frametitle{2.6.1. Интерпретаторы таблиц (2/3)}
В таком случае автоматное отображение, определяемое инициальным
автоматом-преоб\-разователем $S=\langle A, Q, B, \delta,\lambda, q_1
\rangle$, реализуется следующим алгоритмом.
\begin{algorithmic}
\REQUIRE входное слово {\textbf {$\alpha:$ array $[1..s]$ of $1..m$}}
\ENSURE выходное слово {\textbf {$\beta:$ array $[1..s]$ of $1..k$}}
\STATE $q:=q_1$ \COMMENT { начальное состояние }
\FOR{\textbf{$i$
 from $1$ to $s$}}
 \STATE $a:=\alpha[i]$ \COMMENT{ очередной входной символ }
 \STATE $\beta[i]:=\lambda[q,a]$ \COMMENT{ очередной выходной символ }
\STATE $q:=\delta[q,a]$ \COMMENT{ новое состояние }
\ENDFOR
\end{algorithmic}

Другие типы детерминированных автоматов также допускают подобное
представление.

\medskip
Важным обстоятельством является то, что алгоритм работы зависит не
от конкретных таблиц функций переходов и выходов, а только
от типа автомата. Таким образом, можно реализовать интерпретатор
автомата нужного типа один раз, передавая ему необходимую информацию
в виде аргументов.

\end{frame}

\begin{frame}
\frametitle{2.6.1. Интерпретаторы таблиц (3/3)}

Рассмотрим, например, универсальный автомат-распознаватель,
представленный в следующем алгоритме.

\begin{algorithmic}
\REQUIRE {{Язык $L$, заданный таблицей переходов $\delta:$
\textbf{array} $[1..k, 1..m]$ \textbf{of} $1..n$;\\
начальное состояние $q_1\in 1..k$;\\ подмножество заключительных
состояний $F$: \textbf{set of} $1..n$; \\входное слово $\alpha:$
\textbf{array} $[1..s]$ \textbf{of} $1..m$ }}
\ENSURE
{\textbf{true}} если $\alpha\in L$, {\textbf{false}} в противном
случае.
\STATE $q:= q_1$  \COMMENT { начальное состояние }
\FOR{\textbf{$i$
 from $1$ to $s$}}
 \STATE $q:=\delta[q, \alpha[i]]$ \COMMENT{ новое состояние }
 \ENDFOR
 \STATE {\textbf{return $q\in F$}} \COMMENT{ ответ распознавателя }
\end{algorithmic}


Табличные интерпретаторы очень просты. 
Если массивы $\delta$ и
$\lambda$ не слишком велики, то такой способ представления
автомата весьма эффективен.

\medskip
\TODO{Верифицировать, вставить ссылки, насытить примерами.}
\end{frame}


\subsubsection{2.6.2. Представление автомата в программном коде}\label{2.5.2.}
\begin{frame}
\frametitle{2.6.2. Представление автомата в программном коде (1/3)}

Альтернативой интерпретаторам, рассмотренным в предыдущем параграфе,
является \\ 
<<компиляция>> переходов автомата в программный код. Если
используемый язык программирования допускает ветвление по значению
переменной (оператор \textbf{case}), 
то код выполнения одного шага автомата 
$S= \langle A, Q, B, \delta,
\lambda\rangle$ может иметь структуру, 
представленную в следующей схеме.

\begin{algorithmic}
\STATE \textbf{case $q$}
 \COMMENT{ ветвление по текущему состоянию }
 \STATE \textbf{$\quad q_1:$ case $a$}
  \COMMENT{ ветвление по входному символу }
    \STATE \textbf{$\quad\phantom{q_1:} \quad a_1: 
    \quad b\Assign b_{i_1}$}
     \COMMENT{ выходной символ }
    \STATE {$\quad \phantom{q_1:}
    \quad \phantom{a_1:} \quad\ q:= q_{j_1}$}
     \COMMENT{ новое состояние }
    \STATE $\quad\phantom{q_1:} \quad a_2:
    \quad b:= b_{i_2}$
    \STATE {$\quad \phantom{q_1:}
    \quad \phantom{a_1:} \quad\ q:= b_{j_2}$}
\STATE $\quad \phantom{q_1:}
    \quad \phantom{a_1:} \quad \dots$
\STATE\textbf{$\quad\phantom{q_1:}$ end case}
\STATE \textbf{$\quad q_2:$ case $a$}
 \STATE $\dots$
\STATE\textbf{end case} 
\end{algorithmic}

\end{frame}


\begin{frame}
\frametitle{2.6.2. Представление автомата в программном коде (2/3)}

\begin{remark} 
\begin{enumerate}
\item В правых частях операторов присваивания $q\Assign\dots$ и
$b\Assign\dots$, реализующих функции переходов $\delta$ и функцию
выходов $\lambda$, соответственно, \\стоят 
{\em константы} $b_{i_1},
b_{i_2},\dots, q_{i_1}, q_{j_2}, \dots$, не требующие вычислений.
\item Если нужно предусмотреть какие-то дополнительные действия
на переходе, которые автомат должен выполнять помимо определения
выходного символа и смены состояния, то их очень удобно
запрограммировать прямо в соответствующей ветви оператора
\textbf{case}.
\item В зависимости от типа автомата схема 
должна быть дополнена средствами задания начального состояния и получения очередного символа (или события).
\item Общий внешний цикл, обрамляющий выполнение одного шага, определяется назначением автомата.
Если это автомат-распознаватель или преобразователь,
то применяют цикл \textbf{for} по входному слову.
Если же это управляющий автомат, то применяют цикл
\textbf{while}, управляемый достижением заключительного состояния или иным условием. 
\end{enumerate}
\end{remark}

\TODO{Нужны простые, но убедительные примеры.}
\end{frame}


\begin{frame}
\frametitle{2.6.2. Представление автомата в программном коде (3/3)}

Сравнительные характеристики двух подходов к реализации автоматов
собраны в следующую таблицу.

\begin{tabular}{p{4cm}|p{5cm}|p{5cm}}
  % after \\: \hline or \cline{col1-col2} \cline{col3-col4} ...
  \hline
  Характеристика & Интерпретация & Компиляция \\
  \hline
  Объем памяти для хранения данных &
  Значительный. Линейно зависит от размера автомата & Незначительный.
  Не зависит от размера автомата \\
    \hline
Объем исполняемого кода & Незначительный.
  Не зависит от размера автомата & Значительный. Линейно зависит от размера автомата \\
    \hline
Скорость работы & Замедленная. Выборка из массивов & Максимальная. Присваивание констант \\
  \hline
  Встраивание дополнительных процедур & Затруднено & Очень просто \\ \hline
\end{tabular}

\end{frame}

\subsubsection{2.6.3. Задача поиска вхождений}
\label{2.6.3.}
\begin{frame}
\frametitle{2.6.3. Задача поиска вхождений (1/2)}


Рассмотрим применение техники конечных
автоматов для решения одной распространённой задачи программирования, тесно связанной с синтаксическим анализом~---
поиск вхождений образца в текст.
Задан текст
\textbf{$T:$ array $[1..n]$ of $A$}
и образец
\textbf{$P:$ array $[1..m]$ of $A$},
$m\leq n$. Требуется найти
все вхождения $P$ в $T$, то~есть определить все такие
значения $s$, что $\A{i \in 1..m}{T[s+i]= P[i]}$. Если $P$ не
входит в $T$, то множество ответов пусто.
Следующий простой алгоритм, очевидно, решает
данную задачу.

\begin{algorithmic}
\REQUIRE Текст \textbf{$T:$ array $[1..n]$ of char}, образец
\textbf{$P:$ array $[1..m]$ of char}.
\ENSURE Множество
позиций, с которых начинаются вхождения $P$ в $T$. 
\FOR {\textbf{$s$ from $1$ to $n-m+1$}}
  \FOR {\textbf{$i$ from $1$ to $m$}}
    \IF {$P[i]\neq T[s+i]$ }
      \STATE \textbf{next for $s$}
      \COMMENT{ нет совпадения~--- переходим к следующей позиции }
    \ENDIF
 \ENDFOR
 \STATE {\textbf{yield $s$}} \COMMENT{ нашли вхождение }
\ENDFOR
\end{algorithmic}

\end{frame}


\begin{frame}
\frametitle{2.6.3. Задача поиска вхождений (2/2)}
\begin{remark}
Если $P$ не
входит в $T$, то алгоритм ничего не возвращает.
\end{remark}

\medskip
\begin{digress}
Эта задача кажется очень простой, однако это не так. Приведённый алгоритм
 поиска вхождений <<в лоб>> имеет неприемлемо высокую
трудоемкость $O(mn)$.
Дело в том, что задача поиска вхождений находится в центре многих самых насущных приложений.
В частности, в задаче 
\odmkey{полнотекстового поиска}{Задача!полнотекстового поиска}
в Интернете объёмы текстов, в которых надобно искать, измеряются терабайтами,
а количество образцов (поисковых запросов) 
исчисляется миллионами.
Не проще и задачи в \odmkey{биоинформатике}{Биоинформатика}, 
где постоянно приходится искать вхождения достаточно
длинных фрагментов в очень длинные последовательности,
описывающие белки и другие органические полимеры.
В символическом искусственном интеллекте поиск вхождений
является основой 
\odmkey{унификации формул}{Унификация!формул}.
Хотя формулы в логическом программировании 
и не так длинны, как тексты в Интернете и
молекулы в биоинформатике,
но зато из-за {\em комбинаторного взрыва}
при полном переборе искать вхождения 
приходится миллионы и миллиарды раз.
\end{digress}


Известно несколько весьма остроумных способов более эффективного
решения этой задачи. В последующих параграфах рассматривается способ,
основанный на применении понятия конечного автомата.

\end{frame}

\subsubsection{2.6.4. Функция постфикса}\label{2.6.4.}
\begin{frame}
\frametitle{2.6.4. Функция постфикса (1/3)}
Определим \odmkey{функцию постфикса}{Функция!постфикса}
$\Map{\sigma_P}{A^*}{1..m}$ для данного образца $P$ как длину
максимального постфикса слова $\alpha$, который является одновременно
префиксом образца $P$ (см. рисунок):
$$\sigma_P(\alpha):= \max k, \text{ где } \A{i\in 1..k}{P[i]=\alpha[|\alpha|-k+i]}.$$

\begin{figure}[h]
\begin{center}
\includegraphics[width=100mm]{2_6_4_1.jpg}
\end{center}
\end{figure}

\begin{example}
Если $P=ab$, то $\sigma_P(\varepsilon)=0$; $\sigma_P(ccaca)=1$;
$\sigma_P(ccab)=2$.
\end{example}

Очевидно, что $\sigma_P(\alpha)=m$ тогда и только тогда, когда
$P$ является постфиксом $\alpha$. Если $y$ является постфиксом
$x$, то $\sigma_P(x)\geq\sigma_P(y)$. Если образец $P$
подразумевается, то индекс  $P$ в обозначении $\sigma_P$ можно
опустить.
\end{frame}

\begin{frame}
\frametitle{2.6.4. Функция постфикса (2/3)}
\begin{lemma} {\NB [Неравенство для функции постфикса]}
$$\A{\alpha}{\A{a}{\sigma(\alpha a)\leq \sigma(\alpha)
+1}}.$$
\end{lemma}

\begin{proof}
От противного. Пусть $\sigma(\alpha a)> \sigma(\alpha) +1$, тогда
отбросим последний символ $a$ от наибольшего постфикса $\alpha a$,
который является одновременно префиксом $P$. Получаем, что
$\sigma(\alpha) > \sigma(\alpha)$~--- противоречие
(см. рисунок).
\end{proof}

\begin{figure}[h]
\begin{center}
\includegraphics[width=100mm]{2_6_4_2.jpg}
\end{center}
\end{figure}

\end{frame}

\begin{frame}
\frametitle{2.6.4. Функция постфикса (3/3)}
\begin{lemma} {\NB [Реккурентная формула для функции постфикса]}
$$\A{\alpha}{\A{a}{\sigma(\alpha a)= \sigma (P[1..
\sigma(\alpha)] a)}}.$$
\end{lemma}


\begin{proof}
Пусть $k:=\sigma(\alpha)$. По предыдущей лемме
$\sigma(\alpha a)\leq k+1$,
поэтому $\sigma(\alpha a) = \sigma (\alpha
[|\alpha|-k+1..|\alpha|] a)$. Но $\alpha[|\alpha|- k +1 ..
|\alpha|] = P[1..k]$ по определению $\sigma$.
\end{proof}

\begin{remark}
Запись $X[a..b]$ означает последовательность
элементов массива $X$ с индесами от $a$ до $b$.
Такая операция называется \odmkey{вырезкой}{Вырезка!из массива}
из массива.
\end{remark}

\end{frame}



\subsubsection{2.6.5. Автомат поиска вхождений}\label{2.6.5}
\begin{frame}
\frametitle{2.6.5. Автомат поиска вхождений }
Построим конечный автомат-распознаватель $S=\langle A, Q, \delta,
q_0, F\rangle$, соответствующий образцу $P$ следующим образом: $Q:
0..m$; $q_0:= 0$; $F: \{m\}$; $\delta(q,a): = \sigma_P
(P[1..q] a)$. Другими словами, если автомат находится в
состоянии $q$ и входным символом является $a$,  то автомат
переходит  в состояние с номером, который определяется значением
функции постфикса для строки, составленной из первых $q$ букв
образца $P$  и буквы $a$.

Автомат-распознаватель работает следующим образом.

\begin{algorithmic}
\REQUIRE { Текст  \textbf { $T:$ array $[1..n]$ of char},
 образец
\textbf{$P:$ array $[1..m]$ of char}}
\ENSURE {Множество
позиций, с которых начинаются вхождения $P$ в $T$. Если $P$ не
входит в $T$, то алгоритм ничего не возвращает}
 \STATE $q\Assign 0$
 \COMMENT{ начальное состояние }
 \FOR {\textbf{$i$ from 1 to $n$}}
   \STATE $q\Assign \delta(q, T[i])$
   \COMMENT{ следующее состояние }
   \IF {$q=m$}
     \STATE {\textbf{yield $i-m$}}
     \COMMENT{ нашли вхождение }
   \ENDIF
 \ENDFOR
\end{algorithmic}

\begin{ground}
Основной цикл алгоритма имеет следующий инвариант: $\A{i\geq
0}{\delta^* (0, T[1..i])= \sigma (T[1..i])}$. Действительно,
по индукции, для $i=0$ это очевидно: $q=0$ и
$\sigma(\varepsilon)=0$. Если инвариант выполнен на шаге $i$, то
по рекуррентной формуле для функции постфикса
он выполнен и на шаге $i+1$. Отсюда получается, что
если $q=m$, то автомат только что прочел в тексте $T$ последний
символ образца $P$ и следует вернуть позицию $i-m$ в качестве
ответа.
\end{ground}

Время работы алгоритма, очевидно, пропорционально
длине текста, $O(n)$,  что существенно лучше, чем $O(nm)$.

\TODO{Примеры.}
\end{frame}



\subsubsection{2.6.6. Построение функции переходов}\label{2.6.6.}
\begin{frame}
\frametitle{2.6.6. Построение функции переходов (1/3)}
Функцию переходов $\delta$, соответствующую образцу $P$, нетрудно
вычислить, например, с помощью следующего алгоритма.

\begin{algorithmic}
\REQUIRE Образец \textbf{$P:$ array $[1..m]$ of $A$}
\ENSURE Функция переходов \textbf{$\delta:$ array $[0..m,A]$ of $0..m$}
\FOR {\textbf{$q$ from $0$ to $m$}}
  \FOR {$a\in A$} \STATE
    $k:= \min(m+1, q+2)$
    \COMMENT{ наибольшее возможное $k$ }
    \REPEAT
      \STATE $k:= k-1$
      \COMMENT{ пытаемся уменьшить $k$ }
    \UNTIL{$P[1..k]=P[1..q]+a \Or k = 0$}
    \STATE {$\delta[q,a]:=k$}
  \ENDFOR
\ENDFOR
\end{algorithmic}
\end{frame}

\begin{frame}
\frametitle{2.6.6. Построение функции переходов (2/3)}

\begin{ground}
В двух вложенных циклах перебираются все пары $q,a$
и вычисляется число
$k= \sigma(P[1..q]a)$,
как это требуется для построения автомата. Данное вычисление
выполняется <<в лоб>>: начинаем с максимально возможного $k$ и
уменьшаем его, пока  префикс $P[1..k]$  не совпадет с
постфиксом слова $P[1..q] a$ (см. рисунок).

\begin{figure}[h]
\begin{center}
\includegraphics[width=100mm]{odm33-29.eps}
\end{center}
\end{figure}
\end{ground}

Внутренний цикл выполняется не более $m+1$ раз и на
каждом шаге требует не более  $m+1$ сравнений
для проверки условия окончания. Таким образом, общая
трудоемкость данного алгоритма построения функции переходов составляет
$O(m^3|A|)$.
\end{frame}

\begin{frame}
\frametitle{2.6.6. Построение функции переходов (3/3)}
Существуют значительно более эффективные алгоритмы, которые строят
функцию переходов за $O(m|A|)$.

\TODO{Развить эту тему. 
 См., например, Кормен.
 Сократить число параграфов. 
 {\color{blue} Тимофеев Александр}}

\end{frame}


% GA_3_2025.tex, 1–1
\section{3. Языки предметных областей}

% GA_3_2025.tex, 660–833
\subsection{3.2. Методы описания языков}\label{3.2.}
\begin{frame}
\frametitle{3.2. Методы описания языков}
Мы рассматриваем исключительно
описание {\em языков предметной области}, применяемых в компьютерных
технологиях. 

\smallskip
Прежде всего, заметим, что языки программирования, в
отличии от естественных языков, создаются конкретными людьми с
определённой целью, и авторы языков могут придать им нужные
свойства. Поэтому для описания {\em синтаксиса} языков
программирования не применяются грамматики типов $\frak{T}_0$ и $\frak{T}_1$
(все алгоритмы работы с такими грамматиками
 слишком неэффективны), а применяются более простые грамматики
типов $\frak{T}_2$ и $\frak{T}_3$. Но и в этих простых случаях
выделяются ещё более простые подклассы, чтобы облегчить
компьютерную обработку описаний таких языков.


\smallskip
Далее, заметим, что хотя языки предметных областей
являются формальными языками, но сами предметные области
редко имеют готовую формализацию. Поэтому формальные 
методы определения {\em семантики} часто оказываются неприменимы
и тогда приходится использовать описания ad hoc.

\medskip
Наконец, {\em прагматика}, будучи важнейшей частью 
описания языка с практической точки зрения,
вовсе не имеет развитых формальных методов описания.


\end{frame}


\subsubsection{3.2.1. Структура описания языка программирования}
\begin{frame}
\frametitle{3.2.1. Структура описания языка программирования (1/4)}
В лингвистике для естественных языков используется довольно сложная
иерархия описаний. Даже если ограничиться только письменным языком
(текстом), оставив в стороне устную речь, мы видим,
что правила для образования слов (морфология),
совсем не похожи на правила составления предложений
(синтаксис). Если же принять во внимание иероглифическую письменность,
или правила произношения, интонации и прочее, 
то видимая способность лингвистов справиться с задачей
вызывает уважение.   

\smallskip
В программировании при описании различных языков используются самые 
разнообразные методы. Здесь принято выделять
шесть частей описания языка: алфавит, лексика, синтаксис, контекстные условия, семантика и прагматика. Это деление субъективно, но оно
хорошо работает в простых случаях и оставляет место для
применения других методов в более сложных случаях.

\smallskip
{\color{blue}\bf [1.]} \odmkey{Алфавит}{Алфавит}~--- описание множества допустимых символов языка. Как
правило, это просто конечное множество символов, которые можно
явно перечислить, но иногда встречаются и особенности. В некоторых
языках программирования и разметки текста встречаются конструкции,
в которых могут участвовать {\em любые} символы, а не только те,
которые перечислены как основные.

\end{frame}


\begin{frame}
\frametitle{3.2.1. Структура описания языка программирования (2/4)}
\begin{example} Таковы
 комментарии и строковые константы в некоторых языках или
конструкция CDATA в XML.
\end{example}
\begin{digress}
В современном языкотворчестве наблюдается тенденция
к усложнению алфавитов и повышению их роли в описании языка.
В частности, кодировки семества Unicode позволяют
использовать намного более представительные множества символов,
нежели коды ASCII.
Другое направление~--- \odmkey{атрибутированные символы}{Атрибут!символа}.
Действительно, такие атрибуты, как цвет, размер шрифта, начертание, 
гарнитура и т.~д. несут определённую информацию и могут использоваться как 
части языка.  
\end{digress}

\smallskip
{\color{blue}\bf [2.]} \odmkey{Лексика}{Лексика}~--- описание множества
элементарных конструкций, называемых \odmkey{лексическими
единицами}{Лексическая единица} (или \odmkey{лексемами}{Лексема}).
Отличительной особенностью лексем является то, что они не содержат
в своем составе других лексем, т.~е. являются <<атомами>>,
то есть терминалами с точки зрения синтаксиса.
  Для
описания лексики в настоящее время обычно используются автоматные
грамматики (см. \DMref{п.~1.1.6}{1.1.6.}) и регулярные выражения (что
эквивалентно, см. \DMref{п.~1.2.2}{1.2.2.}).
\begin{example} Изображения чисел,
строковые константы, знаки операций, имена и ключевые слова~--- типичные лексемы
современных языков программирования.
\end{example}

\end{frame}


\begin{frame}
\frametitle{3.2.1. Структура описания языка программирования (3/4)}
{\color{blue}\bf [3.]} \odmkey{Синтаксис}{Синтаксис}~--- описание правил,
по которым из лексем строятся предложения (или
\odmkey{операторы}{Оператор!языка программирования}) языка, а из
операторов строятся программы. Как правило, для описания
синтаксиса языков программирования используются
контекстно-свободные порождающие грамматики в одной из описанных
выше форм. При этом некоторые языки устроены так, что даже
отдельные операторы  допускают немедленное выполнение, такие языки
называются \odmkey{интерпретируемыми}{Язык!интерпретируемый} (или
\odmkey{скриптовыми}{Язык!скриптовый}). Другие же языки требуют
наличия определенной конструкции для того, чтобы программа могла
считаться законченной.
\begin{example}
Текст программы на языке VBScript может состоять из одного
оператора присваивания или вывода сообщения. Программа на языке Cи
должна содержать описание функции с именем main для того, чтобы
быть выполнимой.
\end{example}

\smallskip
{\color{blue}\bf [4.]}  \odmkey{Контекстные условия}{Условие!контекстное}~---
описание синтаксических правил,
которые невозможно или неудобно выразить средствами
контекстно-свободной грамматики. Некоторые способы описания
контекстных условий обсуждаются в  
\DMref{п.~1.5.4}{1.5.4.} и \DMref{п.~3.2.5}{3.2.5.}.
\begin{example}
Фактические аргументы в операторе вызова процедуры должны
соответствовать по количеству, порядку и типам формальным
параметрам в описании этой процедуры.
\end{example}
\end{frame}


\begin{frame}
\frametitle{3.2.1. Структура описания языка программирования (4/4)}

{\color{blue}\bf [5.]} \odmkey{Семантика}{Семантика}~--- описание правил
приписывания смысла синтаксически правильным конструкциям языка.
Для языков программирования смыслом программ является их
выполнение. В настоящее время формальные методы описания семантики
языков программирования разработаны недостаточно, хотя постоянные
усилия в этом направлении предпринимаются. Как правило, семантика
описывается либо на естественном языке, либо в терминах работы
какой-то абстрактной вычислительной машины.

\smallskip
\begin{example}
Семантика языка Java описана в терминах абстрактной машины~--- JVM
(Java Virtual Machine).
\end{example}

{\color{blue}\bf [6.]} \odmkey{Прагматика}{Прагматика}~--- описание 
практически значимых особенностей языка, не выраженных
средствами описания синтаксиса и семантики. 
Например: производительность,
лаконичность, наличие сообщества носителей языка, наличие стандартов и тому подобное.
Для разных языков специфичны различные аспекты прагматики.

\begin{example}
Для языка Forth объём памяти для хранения программы во 
внутреннем представлении (\odmkey{шитый код}{Шитый код}) меньше,
чем объём самой компактной эквивалентной программы на машинном языке.
Это прагматическое свойство очень ценно для встроенных и бортовых систем.
\end{example}  

\end{frame}


\end{document}
